\documentclass[a4paper,11pt]{book}
\usepackage[T1]{fontenc}
\usepackage[utf8]{inputenc}
\usepackage{lmodern}
\usepackage[english]{babel}
\usepackage{amsmath,amssymb,amsthm}
\usepackage{mathtools}
\usepackage{booktabs}
\usepackage{enumitem}
\usepackage{float}
\usepackage{algorithm}
\usepackage{algpseudocode}
\usepackage[margin=3.1cm]{geometry}
\usepackage{tikz}
\usetikzlibrary{positioning,arrows.meta}
\usepackage[hidelinks]{hyperref}

\theoremstyle{plain}
\newtheorem{theorem}{Theorem}[chapter]
\newtheorem{proposition}[theorem]{Proposition}
\newtheorem{lemma}[theorem]{Lemma}
\newtheorem{corollary}[theorem]{Corollary}
\theoremstyle{definition}
\newtheorem{definition}[theorem]{Definition}
\newtheorem{example}[theorem]{Example}
\theoremstyle{remark}
\newtheorem{remark}[theorem]{Remark}

\numberwithin{equation}{chapter}
\numberwithin{algorithm}{chapter}
\renewcommand{\algorithmicrequire}{\textbf{Input:}}
\renewcommand{\algorithmicensure}{\textbf{Output:}}

\newcommand{\R}{\mathbb{R}}
\newcommand{\N}{\mathbb{N}}
\newcommand{\Sp}{\mathcal{S}}
\newcommand{\proj}[2]{\Pi_{#2}\!\left(#1\right)}
\newcommand{\norm}[1]{\left\|#1\right\|}
\newcommand{\abs}[1]{\left|#1\right|}
\newcommand{\inner}[2]{\left\langle #1, #2 \right\rangle}
\newcommand{\diag}{\operatorname{diag}}
\newcommand{\dom}{\operatorname{dom}}
\newcommand{\intr}{\operatorname{int}}
\newcommand{\relint}{\operatorname{relint}}
\newcommand{\relbd}{\operatorname{rbd}}
\newcommand{\aff}{\operatorname{aff}}
\newcommand{\conv}{\operatorname{conv}}
\newcommand{\epi}{\operatorname{epi}}
\newcommand{\gph}{\operatorname{gph}}
\newcommand{\range}{\operatorname{range}}
\newcommand{\rank}{\operatorname{rank}}
\newcommand{\tr}{\operatorname{tr}}
\newcommand{\Aut}{\operatorname{Aut}}
\newcommand{\prox}{\operatorname{prox}}
\newcommand{\svec}{\operatorname{svec}}
\newcommand{\smat}{\operatorname{smat}}
\newcommand{\Kexp}{\mathcal{K}_{\exp}}
\newcommand{\Kpow}{\mathcal{K}_{\alpha}}
\DeclareMathOperator{\Fix}{Fix}
\newcommand{\Lor}{\mathcal{L}}
\DeclareMathOperator{\dist}{dist}
\DeclareMathOperator*{\argmin}{argmin}
\DeclareMathOperator*{\argmax}{argmax}

\title{\bfseries ProxGQP\\[3mm]\large Reference}
\author{}
\date{}

\begin{document}
	\maketitle
	\tableofcontents
	
	\chapter{Introduction}
	
	This document serves as a reference for the solver proxgqp. It contains the
	theory needed to maintain the codebase, together with the implementation
	details.
	
	It is separated into four parts: the first introduces the theoretical tools
	used to develop the two methods, the second develops the two methods on that
	basis, the third give some details on the implementatin and the last give some benchmark against other solvers. A closing chapter give the way the solver was developped and possible direction for the futur. 
	
	
	\part{Background}
	
	\chapter{Numerical linear algebra recalled}
	
	The objective of this part is to introduce the notions of numerical linear
	algebra used in the solver.
	\section{Conditioning}
	
	We first recall what a submultiplicative norm is.
	
	\begin{definition}[Submultiplicative matrix norm]
		A norm $\norm{\cdot}$ on square matrices is \emph{submultiplicative} if
		\begin{equation}
			\norm{AB}\leq\norm{A}\,\norm{B}
			\qquad\text{for all }A,B .
		\end{equation}
		The spectral norm and the Frobenius norm are submultiplicative; the entrywise
		maximum is not.
	\end{definition}
	
	\begin{definition}[Conditioning]
		Let $A\in\R^{n\times n}$ be invertible and $\norm{\cdot}$ a submultiplicative
		matrix norm. The \emph{condition number} of $A$ is
		\begin{equation}
			\kappa(A)=\norm{A}\,\norm{A^{-1}} .
		\end{equation}
	\end{definition}
	
	Generally we have in any submultiplicative norm:
	$$\kappa(A^TA) \leq \kappa(A)^2$$
	
	
	
	\begin{definition}[Singular values and spectral norm]
		Let $A\in\R^{m\times n}$. The \emph{singular values} of $A$ are the
		nonnegative square roots of the eigenvalues of $A^\top A$, written
		\[
		\sigma_1(A)\geq\cdots\geq\sigma_n(A)\geq0 .
		\]
		The largest singular value is denoted by $\sigma_{\max}(A)$.
		When $A$ has full column rank, all these singular values are positive,
		and the smallest one is denoted by $\sigma_{\min}(A)$.
		
		The \emph{spectral norm}, or matrix $2$-norm, is
		\[
		\norm{A}_2
		=\max_{\norm{x}_2=1}\norm{Ax}_2
		=\sigma_{\max}(A).
		\]
		For a rectangular matrix of full column rank, its spectral condition
		number is defined by
		\[
		\kappa_2(A)
		=\frac{\sigma_{\max}(A)}{\sigma_{\min}(A)}.
		\]
		And if $A$ symetric definie positive:
		\[
		\kappa_2(A)
		=\frac{\lambda_{\max}(A)}{\lambda_{\min}(A)}.
		\]
		
	\end{definition}
	
	
	\begin{proposition}[Conditioning of a normal product]
		Let $A\in\R^{m\times n}$ have full column rank. In the spectral norm
		\begin{equation}
			\kappa_2\bigl(A^\top A\bigr)=\kappa_2(A)^2 .
		\end{equation}
	\end{proposition}
	\begin{proof}
		The eigenvalues of $A^\top A$ are the squares of the singular values of $A$,
		which are positive by the rank assumption, so
		$\kappa(A^\top A)=\sigma_{\max}(A)^2/\sigma_{\min}(A)^2=\kappa(A)^2$.
	\end{proof}
	
	We now define backward stability. An algorithm is backward stable when its
	output for a problem $(A,b)$ is the exact solution of a nearby problem in the
	data space, whence the name.
	\begin{definition}[Backward error]
		Let $\hat x$ be a computed solution from an algorithm of $Ax=b$. Its \emph{backward error} is
		\begin{equation}
			\eta(\hat x)=\min\bigl\{\varepsilon\geq0:\ (A+\Delta A)\hat x=b+\Delta b,\
			\norm{\Delta A}\leq\varepsilon\norm{A},\
			\norm{\Delta b}\leq\varepsilon\norm{b}\bigr\} .
		\end{equation}
		An algorithm is \emph{backward stable} when $\eta(\hat x)$ is of the order of
		the unit roundoff 	 for every input.
	\end{definition}
	
	\begin{definition}[Forward error]
		With $x$ the exact solution of $Ax=b$ and $\hat x$ a computed one, the
		\emph{forward error} is $\norm{\hat x-x}/\norm{x}$.
	\end{definition}
	
	The forward error is bounded by a combination of the stability of the
	algorithm, through the backward error, and the intrinsic difficulty of the
	problem, through the conditioning.
	\begin{proposition}[Forward error from the condition number, \cite{Higham2002}]
		If $\kappa(A)\,\eta(\hat x)<1$ then
		\begin{equation}
			\frac{\norm{\hat x-x}}{\norm{x}}
			\ \leq\ \frac{2\,\kappa(A)\,\eta(\hat x)}{1-\kappa(A)\,\eta(\hat x)} ,
		\end{equation}
		which to first order is $\kappa(A)\,\eta(\hat x)$.
	\end{proposition}
	
	\section{Factorisation}
	
	In numerical linear algebra one does not generally form $A^{-1}$ to solve
	$Ax=b$, mainly because the associated algorithms are not stable enough or are
	too costly. A factorisation is used instead, of the general form
	$A=R_1\times\cdots\times R_n$, where each $R_i$ makes $R_i^{-1}x$ easy to
	obtain, as for a diagonal or a triangular matrix. Only symmetric $A$ concerns
	us here, since those are the ones appearing in the program. But first we define
	the unit in which the cost of an algorithm is measured.
	
	\begin{definition}[Flop]
		A \emph{flop} is one floating point arithmetic operation. The cost of an
		algorithm is its flop count.
	\end{definition}
	
	
	\begin{definition}[Dense kernel]
		A \emph{dense kernel} is an operation on matrices stored entry by entry, in
		contrast with the sparse matrix(defined later), so that its cost depends on the dimensions
		alone.
	\end{definition}
	
	
	\begin{definition}[BLAS levels]\label{def:blaslevels}
		Dense kernels are graded by the ratio of arithmetic to memory traffic.
		\emph{Level 1} kernels act on vectors, performing $O(n)$ flops on $O(n)$ data;
		\emph{level 2} kernels are matrix--vector, performing $O(n^2)$ flops on
		$O(n^2)$ data; \emph{level 3} kernels are matrix--matrix, performing $O(n^3)$
		flops on $O(n^2)$ data. Only level 3 kernels perform substantially more
		computation per data element fetched, and hence, since fetching data is often
		the bottleneck, it is better to use level 3 operations whenever possible.
	\end{definition}
	
	
	\begin{definition}[Blocked algorithm]
		An algorithm is \emph{blocked} when it is organised so that its inner loop is
		a level 3 kernel: the matrix is partitioned into submatrices, and each step
		updates a submatrix by a matrix--matrix product rather than by a sequence of
		scalar or vector operations. 
	\end{definition}
	\begin{definition}[Cholesky factorisation]
		For $A\in\Sp^n$ a \emph{Cholesky factorisation} is $A=LL^\top$ with $L$ lower
		triangular of positive diagonal. It costs $\tfrac13n^3+O(n^2)$ flops and is
		blocked, its inner kernel being a level 3 update.
	\end{definition}
	
	\begin{proposition}[Existence of the Cholesky factorisation]
		$A\in\Sp^n$ admits a Cholesky factorisation if and only if $A\succ0$, in which
		case the factor is unique. The factorisation is then backward stable and needs
		no pivoting.
	\end{proposition}
	
	\begin{definition}[$LDL^\top$ factorisation]
		For $A\in\Sp^n$ an \emph{$LDL^\top$ factorisation} is $PAP^\top=LDL^\top$ with
		$P$ a permutation, $L$ unit lower triangular and $D$ block diagonal with blocks
		of size $1$ and $2$. It costs $\tfrac13n^3+O(n^2)$ flops and is blocked.
	\end{definition}
	
	\begin{proposition}[Existence of the $LDL^\top$ factorisation, \cite{BunchKaufman1977}]
		Every $A\in\Sp^n$ admits an $LDL^\top$ factorisation with $D$ block diagonal of
		blocks of size at most $2$, and under the Bunch--Kaufman pivoting rule the
		factorisation is backward stable with a growth factor bounded independently of
		$A$. 
	\end{proposition}
	
	Both factorisations are the same procedure, and it is worth writing out once,
	because every later statement about cost, fill and ordering is a statement
	about one line of it. The form below is the \emph{right looking} one: at each
	step a column of $L$ is produced and the whole remaining submatrix is updated
	at once.

	\begin{algorithm}[H]
		\caption{Cholesky, right looking}
		\begin{algorithmic}[1]
			\Require $A\in\Sp^n$ with $A\succ0$
			\Ensure $L$ lower triangular with $A=LL^{\top}$
			\For{$k=1,\dots,n$}
			\State $L_{kk}\gets\sqrt{A_{kk}}$
			\State $L_{k+1:n,\,k}\gets A_{k+1:n,\,k}/L_{kk}$
			\Comment{the $k$th column of the factor}
			\State $A_{k+1:n,\,k+1:n}\gets A_{k+1:n,\,k+1:n}
			-L_{k+1:n,\,k}L_{k+1:n,\,k}^{\top}$
			\Comment{the rank one update}
			\EndFor
		\end{algorithmic}
	\end{algorithm}

	\begin{algorithm}[H]
		\caption{$LDL^{\top}$, right looking, with $D$ diagonal}
		\begin{algorithmic}[1]
			\Require $A\in\Sp^n$ for which the factorisation exists without pivoting
			\Ensure $L$ unit lower triangular and $D$ diagonal with $A=LDL^{\top}$
			\For{$k=1,\dots,n$}
			\State $D_{kk}\gets A_{kk}$
			\State $L_{kk}\gets1$,\quad
			$L_{k+1:n,\,k}\gets A_{k+1:n,\,k}/D_{kk}$
			\State $A_{k+1:n,\,k+1:n}\gets A_{k+1:n,\,k+1:n}
			-D_{kk}\,L_{k+1:n,\,k}L_{k+1:n,\,k}^{\top}$
			\EndFor
		\end{algorithmic}
	\end{algorithm}

	The two differ in one place only. Cholesky puts $\sqrt{A_{kk}}$ on the
	diagonal and so requires $A_{kk}>0$ at every step, which is why it exists
	exactly for $A\succ0$. The $LDL^{\top}$ form keeps $A_{kk}$ itself in $D$ and
	takes no square root, so a negative pivot is admissible and the factorisation
	can proceed on an indefinite matrix; what it cannot tolerate is a pivot that
	is zero or, in floating point, small relative to the entries it must divide.
	That is what pivoting exists to avoid, and what the class of the next
	definitions avoids by construction. The loop order $1,\dots,n$ is written
	that way only for readability; any symmetric permutation may be applied
	first.

	\begin{definition}[Inertia]\label{def:inertia}
		The \emph{inertia} of $A\in\Sp^n$ is the triple
		$\bigl(n_+(A),n_-(A),n_0(A)\bigr)$ counting its positive, negative and zero
		eigenvalues.
	\end{definition}

	\begin{proposition}[Sylvester's law of inertia]\label{prop:sylvester}
		For any invertible $X$, the matrices $A$ and $XAX^{\top}$ have the same
		inertia. In particular a symmetric permutation does not change it, and if
		$A=LDL^{\top}$ then the inertia of $A$ is that of $D$, which for a diagonal
		$D$ is read off the signs of its entries.
	\end{proposition}

	\begin{definition}[Quasi-definite matrix]\label{def:quasidefinite}
		A symmetric matrix is \emph{quasi-definite} when it can be written in blocks
		as
		\begin{equation}
			\begin{pmatrix}
				H & B^{\top}\\
				B & -C
			\end{pmatrix},
			\qquad
			H\succ0,
			\qquad
			C\succ0 ,
		\end{equation}
		with $B$ arbitrary. Such a matrix is indefinite, of inertia
		$(n_H,n_C,0)$ where $n_H$ and $n_C$ are the sizes of the two diagonal
		blocks.
	\end{definition}

	\begin{proposition}[Vanderbei]\label{prop:quasidef}
		A quasi-definite matrix is invertible, and for \emph{every} symmetric
		permutation $P$ the matrix $PAP^{\top}$ admits a factorisation
		$LDL^{\top}$ with $D$ diagonal and no pivoting required. The signs of $D$
		are those of the inertia, by Proposition~\ref{prop:sylvester}.
	\end{proposition}

	This is the property that separates the two systems of the solver from a
	general indefinite one. Because the factorisation exists for every symmetric
	permutation, the permutation may be chosen for sparsity alone, by
	Definition~\ref{def:mindegree} and its refinements, without any of it being
	revisited for stability once numbers arrive. A general indefinite matrix
	requires pivoting, which is decided from the values and therefore changes the
	pattern between iterations, defeating the split between the symbolic and the
	numeric factorisation which will be defined later in this chapter.

	\begin{definition}[Low-rank update of a factorisation]
		Let $A\in\Sp^n$ be factorised, $U\in\R^{n\times p}$ and $C\in\Sp^p$ with
		$p\ll n$. A \emph{rank-$p$ update} is the passage from a factorisation of $A$
		to one of
		\begin{equation}
			\tilde A=A+UCU^\top .
		\end{equation}
	\end{definition}
	
	\begin{proposition}[Woodbury formula]
		If $A$, $C$ and $C^{-1}+U^\top A^{-1}U$ are invertible then
		\begin{equation}
			\bigl(A+UCU^\top\bigr)^{-1}
			=A^{-1}-A^{-1}U\bigl(C^{-1}+U^\top A^{-1}U\bigr)^{-1}U^\top A^{-1} .
		\end{equation}
	\end{proposition}
	\begin{proposition}[Cost of a dense low-rank update, \cite{GillGolubMurraySaunders1974}]
		A Cholesky or $LDL^\top$ factorisation of a dense $A\in\Sp^n$ can be updated to
		one of $A+UCU^\top$ with $\rank U=p$ in $O(pn^2)$ flops, against $O(n^3)$ for
		recomputation.
	\end{proposition}
	
	\begin{proposition}[Symmetric eigendecomposition]
		Every $A\in\Sp^n$ admits $A=Q\Lambda Q^\top$ with $Q$ orthogonal and $\Lambda$
		real diagonal. It is computed backward stably in $O(n^3)$ flops, with a
		constant several times that of a Cholesky factorisation.
	\end{proposition}
	
	\section{Sparse matrices}
	
	When the matrix has many zero entries, as happens for specifically structured
	data, storing and computing with every entry of $A$ is wasteful. This leads to
	the notion of a sparse matrix.
	

	\begin{definition}[Sparse matrix]
		A matrix is \emph{sparse} when its number of nonzero entries is $O(n)$ rather
		than $O(n^2)$, and is represented so that the cost of an operation is
		proportional to that number rather than to $n^2$.
	\end{definition}
	
	\begin{definition}[Pattern]
		The \emph{pattern} of $A$ is $\mathcal{P}(A)=\{(i,j):A_{ij}\neq0\}$. Two
		matrices have the same pattern when these sets agree. The \emph{fill} of a
		factorisation is $\mathcal{P}(L)\setminus\mathcal{P}(A)$, the entries that
		vanish in $A$ and not in its factor.
	\end{definition}
	
	
	The sparsity pattern determines which entries of a matrix are nonzero, while a
	sparse storage format determines how those entries are represented in memory.
	The three standard formats are coordinate (COO), compressed sparse row (CSR),
	and compressed sparse column (CSC).
	
	\begin{definition}[Coordinate format]
		For a matrix $A\in\mathbb{R}^{m\times n}$ with $\operatorname{nnz}(A)$
		nonzero entries, the \emph{coordinate format} stores three arrays:
		\[
		\texttt{row}[k],\qquad \texttt{col}[k],\qquad \texttt{val}[k],
		\qquad k=1,\dots,\operatorname{nnz}(A),
		\]
		where $\texttt{val}[k]=A_{\texttt{row}[k],\texttt{col}[k]}$. Thus, each
		nonzero is stored together with its row and column indices.
	\end{definition}
	
	The trouble with the previous representation is that it is not minimal in
	size. If the values are laid out in row order, left to right and top to bottom,
	it is enough to record where each row starts in \texttt{val} to recover the row
	index.
	
	\begin{definition}[Compressed sparse row format]
		The \emph{compressed sparse row} (CSR) format stores three arrays:
		\[
		\texttt{val},\qquad \texttt{col\_ind},\qquad \texttt{row\_ptr}.
		\]
		The arrays $\texttt{val}$ and $\texttt{col\_ind}$ have length
		$\operatorname{nnz}(A)$, while $\texttt{row\_ptr}$ has length $m+1$. The
		nonzeros in row $i$ are stored in positions
		\[
		\texttt{row\_ptr}[i],\dots,\texttt{row\_ptr}[i+1]-1
		\]
		of $\texttt{val}$ and $\texttt{col\_ind}$. Hence,
		\[
		A_{ij}=\texttt{val}[k]
		\quad\text{when}\quad
		k\in[\texttt{row\_ptr}[i],\texttt{row\_ptr}[i+1])
		\ \text{and}\ 
		\texttt{col\_ind}[k]=j.
		\]
	\end{definition}
	
	CSC is the same idea, with the entries of \texttt{val} laid out in column
	order, top to bottom and left to right.
	\begin{definition}[Compressed sparse column format]
		The \emph{compressed sparse column} (CSC) format is the column-oriented
		analogue of CSR. It stores three arrays:
		\[
		\texttt{val},\qquad \texttt{row\_ind},\qquad \texttt{col\_ptr}.
		\]
		The arrays $\texttt{val}$ and $\texttt{row\_ind}$ have length
		$\operatorname{nnz}(A)$, while $\texttt{col\_ptr}$ has length $n+1$. The
		nonzeros in column $j$ are stored in positions
		\[
		\texttt{col\_ptr}[j],\dots,\texttt{col\_ptr}[j+1]-1
		\]
		of $\texttt{val}$ and $\texttt{row\_ind}$.
	\end{definition}
	
	For a symmetric matrix only the upper or the lower triangular part need be
	kept, which halves the memory required and leads to the upper and lower CSR and
	CSC representations.
	
	
	
	For a sparse matrix the factorisation has two phases, the first being the
	symbolic factorisation, which is the costlier.
	
	\begin{definition}[Fill-reducing permutation]
		A \emph{fill-reducing permutation} is a permutation $P$ chosen so that the
		factor of $PAP^\top$ has fewer nonzeros than that of $A$. It is computed from
		$\mathcal{P}(A)$ alone by a heuristic, minimising the fill exactly being
		NP-hard.
	\end{definition}
	
	The factorisations above are computed one column at a time, and the cost of
	each step is what decides the cost of the whole.

	Read the third line of either factorisation algorithm again, this time on a
	sparse matrix. The column $L_{k+1:n,k}$ has nonzeros only in the rows where
	$A_{k+1:n,k}$ did, so the outer product touches the submatrix indexed by
	exactly those rows and nothing else. How many there are is the quantity that
	governs everything.

	\begin{definition}[Elimination step and degree]\label{def:elimdegree}
		Let $A\in\Sp^n$ be factored in the order $1,\dots,n$. The \emph{elimination}
		of column $k$ computes that column of $L$ and then updates the remaining
		submatrix by the rank one term $-\ell_kd_k\ell_k^{\top}$. The
		\emph{elimination degree} $d_k$ is the number of nonzeros below the diagonal
		in $\ell_k$, that is the number of rows still to be eliminated that column
		$k$ touches.

		Three consequences follow from the update being an outer product on those
		$d_k$ rows. The step costs about $d_k^2/2$ flops, the update being
		symmetric, so the whole factorisation costs
		\begin{equation}\label{eq:flopcount}
			\text{flops}\ \approx\ \sum_{k=1}^{n}\frac{d_k^2}{2} .
		\end{equation}
		The step makes those $d_k$ rows pairwise adjacent, so any pair among them
		that was not already an entry of $A$ becomes fill. And a column with
		$d_k=1$ updates a single diagonal entry, so it costs one flop and creates
		no fill at all.
	\end{definition}

	A dense matrix has $d_k=n-k$, and \eqref{eq:flopcount} then gives $n^3/6$,
	which is the familiar count. Everything a fill-reducing permutation does is an
	attempt to keep the $d_k$ small.

	\begin{definition}[Minimum degree ordering]\label{def:mindegree}
		The \emph{minimum degree} heuristic chooses, at each step, a column of least
		current degree among those not yet eliminated, updating the degrees after
		each elimination. It is the heuristic behind the orderings used in practice,
		of which \textsc{amd} is the standard refinement.
	\end{definition}

	A column of degree one is therefore eliminated first by such an ordering, and
	by Definition~\ref{def:elimdegree} that costs one flop and creates no fill.

	\begin{definition}[Elimination tree]
		Let $A\in\Sp^n$ and suppose that, after applying a symmetric permutation, it
		has a factorisation
		\[
		PAP^\top=LL^\top
		\qquad\text{or}\qquad
		PAP^\top=LDL^\top,
		\]
		where $L$ is lower triangular.
		
		The \emph{elimination tree} describes the dependencies between the columns of
		the factor. For each column $j$, its parent is
		\[
		\operatorname{parent}(j)
		=
		\min\{i>j:L_{ij}\neq0\},
		\]
		
	\end{definition}
	
	
	\begin{definition}[Symbolic factorisation]\label{def:symbolic}
		The \emph{symbolic factorisation} of $A$ is the computation, from
		$\mathcal{P}(A)$ alone, of a fill-reducing permutation, of the elimination
		tree, of $\mathcal{P}(L)$ and of the storage for $L$. It reads no numerical
		value, and is therefore common to all matrices of the same pattern.
	\end{definition}
	
	Once this is done for concrete value a numeric factorisation can be done much faster.
	\begin{definition}[Numeric factorisation]\label{def:numeric}
		The \emph{numeric factorisation} is the computation of the values of $L$ and
		$D$ into the structure produced by the symbolic factorisation.
	\end{definition}


Consequently, if the sparsity pattern of $A$ is fixed even as its values
change, the symbolic factorisation need be performed only once, which is an
important gain.

	Two numeric factorisations of the same pattern can differ greatly in speed,
	according to whether they treat columns one at a time or in groups.

	\begin{definition}[Simplicial and supernodal factorisation]
		\label{def:supernodal}
		A \emph{simplicial} factorisation eliminates one column at a time, and its
		inner loop is a sparse vector operation. A \emph{supernodal} factorisation
		first groups columns with identical or nearly identical patterns into
		\emph{supernodes}, and eliminates each group as a unit; its inner loop is
		then a dense kernel on a small dense block, and it can call the level three
		routines of Definition~\ref{def:blaslevels}.
	\end{definition}

	The choice is not free. Supernodes only exist when many columns share a
	pattern, which is to say when the factor is locally dense, so a supernodal
	algorithm gains on a dense factor and loses on a very sparse one, where the
	grouping finds nothing and only adds overhead. A backend that offers both
	usually decides between them from the predicted density of $L$.

	
	
	
	\section{Schur complement}
	
	Given a block matrix, it is possible to reduce its size and obtain an
	equivalent smaller system. It is obtained by expressing one variable directly
	in terms of the others and substituting it into the remaining relations. What
	results is generally worse conditioned, but smaller, and sometimes gains a
	property such as definiteness.
	
	\begin{proposition}[Schur complement]
		Let $K=\begin{pmatrix}A&B\\ B^\top&D\end{pmatrix}$ with $D$ invertible, and put
		$S=A-BD^{-1}B^\top$. Then $K$ is invertible if and only if $S$ is, the system
		$K\binom{u}{v}=\binom{a}{b}$ is equivalent to
		\begin{equation}
			S\,u=a-BD^{-1}b,
			\qquad
			v=D^{-1}\bigl(b-B^\top u\bigr) ,
		\end{equation}
		
	\end{proposition}
	
	\section{Iterative refinement}
	
	Given a factorisation, it is sometimes possible to reduce the error by a
	process named iterative refinement.
	
	\begin{definition}[Iterative refinement]
		Given a factorisation of $A$ and an approximate solution $\hat x^{(0)}$ of
		$Ax=b$, \emph{iterative refinement} computes, for $k=0,1,\dots$,
		\begin{equation}
			\rho^{(k)}=b-A\hat x^{(k)},
			\qquad
			A\,\delta^{(k)}=\rho^{(k)},
			\qquad
			\hat x^{(k+1)}=\hat x^{(k)}+\delta^{(k)} ,
		\end{equation}
		the correction being solved with the existing factorisation.
	\end{definition}
	
	\chapter{Convex analytic background}
	
	\section{Convex sets and cones}
	
	We recall some basic notions connected with convexity.
	\begin{definition}[Convex set]
		$C\subseteq\R^m$ is \emph{convex} if $\lambda v+(1-\lambda)w\in C$ for all
		$v,w\in C$ and all $\lambda\in[0,1]$.
	\end{definition}
	
	\begin{definition}[Convex hull]
		A \emph{convex combination} of $v_1,\dots,v_r$ is $\sum_i\lambda_iv_i$ with
		$\lambda_i\geq0$ and $\sum_i\lambda_i=1$. The \emph{convex hull} $\conv C$ is
		the set of convex combinations of finitely many points of $C$, equivalently the
		smallest convex set containing $C$.
	\end{definition}
	
	We first define what a cone is. Cones appear naturally in optimisation
	problems, mainly because most constraints posed in practice are invariant under
	scaling by a positive number.
	
	\begin{definition}[Cone]
		$K\subseteq\R^m$ is a \emph{cone} if $v\in K$ and $t\geq0$ imply $tv\in K$. It
		is a \emph{convex cone} if it is in addition convex, equivalently if
		$v,w\in K$ and $s,t\geq0$ imply $sv+tw\in K$.
	\end{definition}
	
	
	Sometimes we want to say that an element lies strictly inside a cone, but the
	cone is often contained in an affine space of strictly lower dimension, because
	of equality constraints. We therefore prefer the notion of relative interior.
	
	\begin{definition}[Affine hull]
		The \emph{affine hull} $\aff C$ is the smallest affine subspace of $\R^m$
		containing $C$.
	\end{definition}
	
	\begin{definition}[Relative interior]
		\begin{equation}
			\relint C=\bigl\{v\in C:\ \exists\varepsilon>0,\
			B(v,\varepsilon)\cap\aff C\subseteq C\bigr\} .
		\end{equation}
	\end{definition}
	
	\begin{definition}[Relative boundary]
		$\relbd C=\operatorname{cl}C\setminus\relint C$.
	\end{definition}
	
	\begin{definition}[Pointed cone]
		A convex cone $K$ is \emph{pointed} if $K\cap(-K)=\{0\}$.
	\end{definition}
	
	\begin{definition}[Solid cone]
		A convex cone $K\subseteq\R^m$ is \emph{solid} if $\intr K\neq\emptyset$.
	\end{definition}
	
	\begin{definition}[Proper cone]
		A convex cone is \emph{proper} if it is closed, pointed and solid.
	\end{definition}
	
	\section{Generalised inequalities}
	
	The simplest constraints after the scalar inequalities are memberships of a
	half space, $x\leq0$ for instance. For a general cone, membership can be read
	in the same fashion once the notion of a generalised inequality is defined.
	
	\begin{definition}[Generalised inequality]
		For $K\subseteq\R^m$ a closed convex cone,
		\begin{equation}
			x\preceq_K y\iff y-x\in K,
			\qquad
			x\prec_K y\iff y-x\in\intr K .
		\end{equation}
	\end{definition}
	
	\begin{proposition}[Order properties of $\preceq_K$]
		$\preceq_K$ is reflexive and transitive, and if $x\preceq_Ky$ and
		$x'\preceq_Ky'$ then $x+x'\preceq_Ky+y'$ and $tx\preceq_Kty$ for every
		$t\geq0$.
	\end{proposition}
	
	
	\begin{proposition}[Antisymmetry]
		$\preceq_K$ is antisymmetric, and hence a partial order, if and only if $K$ is
		pointed.
	\end{proposition}
	
	\begin{proposition}[Emptiness of the strict relation]
		There exist $x,y$ with $x\prec_Ky$ if and only if $K$ is solid.
	\end{proposition}
	
	\begin{definition}[The two classical instances]
		For $K=\R^n_+$, $x\preceq_Ky$ is $x_i\leq y_i$ for every $i$ and $x\prec_Ky$ is
		$x_i<y_i$ for every $i$. For $K=\Sp^n_+$, $X\preceq_KY$ is $Y-X$ positive
		semidefinite and $X\prec_KY$ is $Y-X$ positive definite; both are written
		without the subscript, as $X\preceq Y$ and $X\prec Y$.
	\end{definition}
	
	\section{Dual and polar, $\Aut(K)$ and symmetric cones}
	
	Two notions are natural here. The first is the dual cone, the set of elements
	of the space forming an angle of at most $\frac{\pi}{2}$ with every element of
	the cone.
	
	\begin{definition}[Dual cone]
		$K^*=\{z\in\R^m:\inner{z}{v}\geq0\ \text{ for all }v\in K\}$.
	\end{definition}
	\label{key}
	The second is its opposite, the set of all elements forming an angle of at
	least $\frac{\pi}{2}$ with every element of the cone.
	
	\begin{definition}[Polar cone]
		$K^\circ=\{z\in\R^m:\inner{z}{v}\leq0\ \text{ for all }v\in K\}=-K^*$.
	\end{definition}
	
	No cone equals its polar: $K\subseteq K^\circ$ forces $\inner{v}{v}\leq0$ on
	$K$, hence $K=\{0\}$, and $\{0\}^\circ=\R^m$. The question is worth asking of
	the dual instead, where it happens often enough to be worth a name, and leads
	to the properties discussed below.
	\begin{definition}[Self-dual cone]
		$K$ is \emph{self-dual} if $K=K^*$.
	\end{definition}
	
	As the cone get bigger the dual get smaller self-duality is hence a first idea of symmetry for a cone.
	
	Topologically the dual is always the better behaved of the two, even when the
	cone itself is not, being the set of its supporting hyperplanes.
	\begin{proposition}[The dual is a closed convex cone]
		For any $K\subseteq\R^m$, $K^*$ is a closed convex cone.
	\end{proposition}
	
	\begin{proof}
		$K^*=\bigcap_{v\in K}\{z:\inner{z}{v}\geq0\}$ is an intersection of closed
		convex cone. And
		all three properties are preserved by arbitrary intersections.
	\end{proof}
	
	\begin{proposition}[Inverse monotonicity of duality]
		$K_1\subseteq K_2$ implies $K_2^*\subseteq K_1^*$.
	\end{proposition}
	
	
	\begin{proposition}[Bipolar theorem, \cite{Rockafellar1970}]
		$K$ is a closed convex cone if and only if $K^{**}=K$.
	\end{proposition}
	
	
	When constraints of different kinds are present, the cone is written as a
	product of smaller cones, one per constraint. The following proposition states
	that the dual is compatible with that product.
	
	\begin{proposition}[Duality and Cartesian products]
		If $K=K_1\times\cdots\times K_p$ then
		$K^*=K_1^*\times\cdots\times K_p^*$.
	\end{proposition}
	
	\begin{proof}
		Write $z=(z_1,\dots,z_p)$ and $v=(v_1,\dots,v_p)$, so that
		$\inner{z}{v}=\sum_j\inner{z_j}{v_j}$. If every $z_j\in K_j^*$ then each term
		is nonnegative for $v\in K$, so $z\in K^*$. Conversely let $z\in K^*$ and fix
		$j$ and $v_j\in K_j$. Taking $v$ with that $j$th block and zero elsewhere,
		which lies in $K$, gives $\inner{z_j}{v_j}\geq0$, so $z_j\in K_j^*$.
	\end{proof}
	
	Once we have a cone a natural question what change of coordinate that preserve it :
	\begin{definition}[Automorphism group]
		For $K\subseteq\R^m$ a proper cone,
		\begin{equation}
			\Aut(K)=\{A\in GL(m,\R):\ AK=K\} .
		\end{equation}
	\end{definition}
	
	Another natural question is whether some points are more special than others,
	in the sense that no automorphism carries one to another while keeping the cone
	intact. An $A\in\Aut(K)$ is a linear homeomorphism of $\R^m$ carrying $K$ onto
	$K$, so it carries $\intr K$ onto $\intr K$ and hence $\partial K$ onto
	$\partial K$: the boundary is preserved, and it is enough to ask the question on
	the interior.
	
	\begin{definition}[Homogeneous cone]
		A proper cone $K$ is \emph{homogeneous} if $\Aut(K)$ acts transitively on
		$\intr K$: for all $u,v\in\intr K$ there is $A\in\Aut(K)$ with $Au=v$.
	\end{definition}
	
	This gives us a second symmetry property which together with self-duality give the following definition: 
	\begin{definition}[Symmetric cone]
		A proper cone is \emph{symmetric} if it is homogeneous and self-dual.
	\end{definition}
	
	Symmetric cones have far better properties than non symmetric ones, and this
	is reflected in the methods.
	
	\section{Cone catalogue}
	
	This section gives the definition of the six cones implemented in the solver,
	together with some of their properties.
	\begin{definition}[Zero cone]
		$\{0\}\subseteq\R^d$ is a closed convex cone with $\{0\}^*=\R^d$. It is not
		self-dual and not solid.
	\end{definition}
	
	\begin{definition}[Nonnegative orthant]
		$\R^d_+=\{x\in\R^d:x_i\geq0\}$ is a symmetric cone, and $(\R^d_+)^*=\R^d_+$.
	\end{definition}
	
	\begin{definition}[Lorentz cone]
		$\Lor^d=\{(t,a)\in\R\times\R^{d-1}:\norm{a}\leq t\}$ is a symmetric cone, and
		$(\Lor^d)^*=\Lor^d$.
	\end{definition}
	
	\begin{definition}[Semidefinite cone]
		$\Sp^k_+=\{X\in\Sp^k:X\succeq0\}$ is a symmetric cone for the trace inner
		product $\inner{X}{Y}=\tr(XY)$, and $(\Sp^k_+)^*=\Sp^k_+$.
	\end{definition}
	
	\begin{definition}[Exponential cone]
		\begin{equation}
			\Kexp=\operatorname{cl}\bigl\{(x_1,x_2,x_3)\in\R^3:\
			x_2\,e^{x_1/x_2}\leq x_3,\ x_2>0\bigr\} ,
		\end{equation}
		a proper cone that is not self-dual, with
		\begin{equation}
			\Kexp^*=\operatorname{cl}\bigl\{(z_1,z_2,z_3)\in\R^3:\
			-z_1e^{z_2/z_1}\leq e\,z_3,\ z_1<0\bigr\} .
		\end{equation}
	\end{definition}
	
	\begin{definition}[Power cone]
		For $\alpha\in(0,1)$,
		\begin{equation}
			\Kpow=\bigl\{(x_1,x_2,x_3)\in\R^3:\
			x_1^{\alpha}x_2^{1-\alpha}\geq|x_3|,\ x_1\geq0,\ x_2\geq0\bigr\} ,
		\end{equation}
		a proper cone that is not self-dual, with
		\begin{equation}
			\Kpow^*=\Bigl\{(z_1,z_2,z_3):\
			\Bigl(\frac{z_1}{\alpha}\Bigr)^{\alpha}
			\Bigl(\frac{z_2}{1-\alpha}\Bigr)^{1-\alpha}\geq|z_3|,\
			z_1\geq0,\ z_2\geq0\Bigr\} .
		\end{equation}
	\end{definition}
	
	\section{Jordan algebra and symmetric cones}
	
	Symmetric cones enjoy a great many properties, owing to their links with
	Jordan algebras.
	\begin{definition}[Euclidean Jordan algebra]
		A \emph{Euclidean Jordan algebra} (EJA) is a finite dimensional real vector space $V$
		with an inner product $\inner{\cdot}{\cdot}$ and a bilinear product $\circ$
		such that, for all $u,v,w\in V$ and with $u^2=u\circ u$,
		\begin{equation}
			u\circ v=v\circ u,
			\qquad
			u\circ(u^2\circ v)=u^2\circ(u\circ v),
			\qquad
			\inner{u\circ v}{w}=\inner{v}{u\circ w} .
		\end{equation}
		It is with an unit $e$ satisfying $e\circ u=u$.
	\end{definition}
	
	Fixing a basis and writing $x\circ y$ in it shows that $\circ$ is $C^{\infty}$,
	being a bilinear map.
	
	\begin{example}[Three Jordan algebras]
		$V=\R^d$ with $u\circ v=u\odot v$ the componentwise product and $e=(1,\dots,1)$.
		$V=\R\times\R^{d-1}$ with
		$(t,a)\circ(s,b)=\bigl(ts+\inner{a}{b},\,tb+sa\bigr)$ and $e=(1,0)$.
		$V=\Sp^k$ with $U\circ W=\tfrac12(UW+WU)$ and $e=I$.
	\end{example}
	
	The following theorems link symmetric with EJA:
	\begin{definition}[Cone of squares]
		The \emph{cone of squares} of $V$ is $K_V=\{u\circ u:u\in V\}$.
	\end{definition}
	
	\begin{proposition}[Jordan, von Neumann and Wigner, \cite{FarautKoranyi1994}]
		A proper cone is symmetric if and only if it is the cone of squares of a
		Euclidean Jordan algebra.
	\end{proposition}
	
	We have the fundamental theorem:
	
	\begin{proposition}[Spectral decomposition, \cite{FarautKoranyi1994}]
		Let $V$ be a Euclidean Jordan algebra(EJA). There is an unique $r\in \mathbb{N}$ named the rank of $V$ such that every $x\in V$ admits
		\begin{equation}
			x=\sum_{i=1}^{r}\lambda_i c_i ,
		\end{equation}
		with $\lambda_1,\dots,\lambda_r\in\R$ the \emph{eigenvalues} of $x$ and
		$c_1,\dots,c_r$ a Jordan frame, that is idempotents with $c_i\circ c_j=0$ for
		$i\neq j$, $c_i\circ c_i=c_i$ and $\sum_ic_i=e$. The eigenvalues are unique up
		to their order, the frame need not be. 
	\end{proposition}
	
	It is then immediate that:
	
	\begin{corollary}[Trace, determinant, inverse, and powers]
		Let $V$ be a Euclidean Jordan algebra of rank $r$, and let
		\[
		x=\sum_{i=1}^{r}\lambda_i c_i
		\]
		be a spectral decomposition of $x$. The \emph{trace} and \emph{determinant} of
		$x$ are defined by
		\[
		\operatorname{tr}(x)=\sum_{i=1}^{r}\lambda_i,
		\qquad
		\det(x)=\prod_{i=1}^{r}\lambda_i.
		\]
		These quantities are independent of the choice of Jordan frame because the
		eigenvalues are unique up to permutation.
		
		Moreover,
		\[
		x\in K_V
		\quad\Longleftrightarrow\quad
		\lambda_i\geq0
		\quad\text{for every }i,
		\]
		and
		\[
		x\in\operatorname{int}K_V
		\quad\Longleftrightarrow\quad
		\lambda_i>0
		\quad\text{for every }i.
		\]
		
		The element $x$ is invertible if and only if
		\[
		\det(x)\neq0,
		\]
		or equivalently, if and only if $\lambda_i\neq0$ for every $i$. In this case,
		\[
		x^{-1}
		=
		\sum_{i=1}^{r}\lambda_i^{-1}c_i.
		\]
		
		More generally, whenever the scalar powers $\lambda_i^p$ are defined,
		\[
		x^p
		=
		\sum_{i=1}^{r}\lambda_i^p c_i.
		\]
		
		Finally, every primitive idempotent $c_i$ in a Jordan frame satisfies
		\[
		\operatorname{tr}(c_i)=1.
		\]
		If $r>1$, then
		\[
		\det(c_i)=0.
		\]
	\end{corollary}
	
	\begin{definition}[Ideal of a Jordan algebra]
		Let $(J,\circ)$ be a Jordan algebra. A linear subspace $I\subseteq J$ is
		called a \emph{Jordan ideal} if
		\[
		x\circ y\in I
		\qquad\text{for all }x\in J\text{ and }y\in I.
		\]
		Equivalently,
		\[
		J\circ I\subseteq I.
		\]
	\end{definition}
	
	\begin{definition}[Simple Jordan algebra]
		A Jordan algebra $J$ is called \emph{simple} if its only Jordan ideals are
		the trivial ideals
		\[
		\{0\}
		\qquad\text{and}\qquad
		J.
		\]
		Equivalently, if $I\subseteq J$ is a linear subspace satisfying
		\[
		x\circ y\in I
		\qquad\text{for all }x\in J\text{ and }y\in I,
		\]
		then either $I=\{0\}$ or $I=J$.
	\end{definition}
	
	
	Simple Euclidean Jordan algebras are the building blocks of all the others.
	
	
	\begin{theorem}[Decomposition of Euclidean Jordan algebras]
		Every finite-dimensional Euclidean Jordan algebra $V$ is an orthogonal direct
		sum of simple Euclidean Jordan algebras. More precisely, there exist simple
		Jordan ideals $V_1,\dots,V_p$ such that
		\[
		V=V_1\oplus\cdots\oplus V_p,
		\]
		\[
		V_i\circ V_j=\{0\}\qquad\text{for }i\neq j,
		\]
		and
		\[
		\langle V_i,V_j\rangle=0\qquad\text{for }i\neq j.
		\]
		
	\end{theorem}
	
	
	\begin{proposition}[Form of the inner product]
		Let $V$ be a finite-dimensional Euclidean Jordan algebra with orthogonal
		decomposition into simple ideals
		\[
		V=V_1\oplus\cdots\oplus V_r.
		\]
		Then :
		\[
		\langle x,y\rangle
		=
		\sum_{i=1}^r c_i\,
		\operatorname{tr}_i(x_i\circ y_i),
		\qquad c_i>0,
		\]
		where
		\[
		x=(x_1,\dots,x_r),\qquad
		y=(y_1,\dots,y_r),
		\]
		and $\operatorname{tr}_i$ denotes the Jordan trace on the simple component
		$V_i$. In particular, the canonical trace inner product is obtained by
		taking $c_i=1$ for every $i$:
		\[
		\langle x,y\rangle_{\mathrm{can}}
		=
		\operatorname{tr}(x\circ y)
		=
		\sum_{i=1}^r\operatorname{tr}_i(x_i\circ y_i).
		\]
	\end{proposition}
	
	Some property that we will use:
	
	\begin{theorem}[Orthogonality in the cone of squares]
		Let $V$ be a Euclidean Jordan algebra, let $K$ denote its cone of squares,
		and equip $V$ with the trace inner product
		\[
		\langle x,y\rangle=\operatorname{tr}(x\circ y).
		\]
		Then, for all $x,y\in K$,
		\[
		x\circ y=0
		\qquad\Longleftrightarrow\qquad
		\langle x,y\rangle=0.
		\]
	\end{theorem}
	
	\begin{definition}[Left multiplication operator]
		Let $V$ be a Jordan algebra. For each $x\in V$, the
		\emph{left multiplication operator} associated with $x$ is the linear map
		\[
		L(x):V\to V,
		\qquad
		L(x)y=x\circ y.
		\]
	\end{definition}
	\begin{definition}[Quadratic representation]
		For $x\in V$, the \emph{quadratic representation} of $x$ is the linear map
		\[
		Q(x):V\to V
		\]
		defined by
		\[
		Q(x)=2L(x)^2-L(x^2).
		\]
		Equivalently, for $y\in V$,
		\[
		Q(x)y
		=
		2x\circ(x\circ y)-x^2\circ y.
		\]
	\end{definition}
	
	
	\begin{theorem}[Derivatives of the log-determinant]
		Let $V$ be a finite-dimensional Euclidean Jordan algebra with cone of squares
		$K$, and equip $V$ with the trace inner product
		\[
		\langle x,y\rangle=\operatorname{tr}(x\circ y).
		\]
		Define
		\[
		f:\operatorname{int}K\to\mathbb{R},
		\qquad
		f(x)=\log\det(x).
		\]
		Then
		\[
		\nabla f(x)=x^{-1}
		\]
		and
		\[
		\nabla^2 f(x)=-Q(x)^{-1}=-Q(x^{-1}).
		\]
		
	\end{theorem}
	
	\begin{theorem}
		Let \(V\) be a Euclidean Jordan algebra with unit element \(e\), and let
		\(K\subset V\) be its symmetric cone. Then
		\[
		\left\{
		A\in\operatorname{Aut}(K):
		A=A^{\mathsf T},\ A\succ0
		\right\}
		=
		\left\{
		Q(w):w\in\operatorname{int}K
		\right\}.
		\]
		Equivalently, an automorphism \(A\) of \(K\) is self-adjoint and positive
		definite if and only if there exists \(w\in\operatorname{int}K\) such that
		\[
		A=Q(w).
		\]
	\end{theorem}
	\begin{theorem}
		Let \(V\) be a Euclidean Jordan algebra with unit element \(e\), and let
		\(K\subset V\) be its symmetric cone. For every
		\(s,x\in\operatorname{int}K\), there exists a unique operator
		\[
		W\in\operatorname{Aut}(K)
		\]
		such that for the canonical inner product:
		\[
		W=W^{\mathsf T},\qquad W\succ0,\qquad Ws=x.
		\]
		Moreover, \(W\) is a quadratic representation:
		\[
		W=Q(w),
		\]
		where the unique element \(w\in\operatorname{int}K\) is given by
		\[
		w=
		Q\!\left(s^{-1/2}\right)
		\left(Q\!\left(s^{1/2}\right)x\right)^{1/2}
		\]
		
	\end{theorem}
	
	For each symmetric cone considered here, the canonical trace inner product is a positive scalar multiple of the standard inner product:
	
	$$
	\langle x,y\rangle_{\mathrm{tr}}
	=
	\gamma_K\langle x,y\rangle_{\mathrm{std}},
	$$
	
	where
	
	$$
	\gamma_K=
	\begin{cases}
		1, & K=\mathbb R_+^n,\\
		2, & K=\mathcal L^n,\\
		1, & K=\mathbb S_+^n.
	\end{cases}
	$$
	
	For the semidefinite cone, the standard inner product is the Frobenius inner product
	
	$$
	\langle X,Y\rangle_F=\operatorname{Tr}(XY).
	$$
	
	\section{Convex functions}
	
	In this section we recall the definition of convex functions and some of their property.
	
	\begin{definition}[Extended real valued convex function]
		$f:\R^N\to\R\cup\{+\infty\}$ is \emph{convex} if
		$f(\lambda v+(1-\lambda)w)\leq\lambda f(v)+(1-\lambda)f(w)$ for all $v,w$ and
		all $\lambda\in(0,1)$, with the convention that the inequality holds whenever
		the right side is $+\infty$. The function $f$ is \emph{concave} if $-f$ is
		convex.
	\end{definition}
	
	
	Convexity is stable under several operations, among them addition and
	multiplication by a positive scalar. Two further facts are used .
	
	\begin{proposition}[Supremum and partial minimization]
		Let $(f_\theta)_\theta$ be a family of convex functions. Then
		$h(v)=\sup_\theta f_\theta(v)$ is convex. If $g$ is convex jointly in its two
		arguments, then the partial minimum $e(v)=\inf_u g(v,u)$ is convex. A pointwise
		infimum $v\mapsto\inf_\theta f_\theta(v)$ is not convex in general.
	\end{proposition}
	
	
	
	
	\begin{definition}[Domain and properness]
		$\dom f=\{v:f(v)<+\infty\}$, and $f$ is \emph{proper} if
		$\dom f\neq\emptyset$.
	\end{definition}
	
	\begin{definition}[Epigraph and lower semicontinuity]
		$\epi f=\{(v,t)\in\R^N\times\R:f(v)\leq t\}$, and $f$ is \emph{lower
			semicontinuous} at $v$ if $\liminf_{w\to v}f(w)\geq f(v)$.
	\end{definition}
	
	\begin{definition}[Closed function]
		$f$ is \emph{closed} if $\epi f$ is closed, equivalently if $f$ is lower
		semicontinuous everywhere.
	\end{definition}
	
	\begin{definition}[Coercivity]
		$f$ is \emph{coercive} if $f(v)\to+\infty$ whenever $\norm{v}\to+\infty$.
	\end{definition}
	
	
	\begin{proposition}[Attainment under coercivity]
		A closed proper coercive $f:\R^N\to\R\cup\{+\infty\}$ attains its infimum.
	\end{proposition}
	
	
	\begin{definition}[Strong convexity]
		$f$ is \emph{strongly convex with modulus $\sigma>0$} in the metric
		$\mathcal{M}\succ0$ if $f-\tfrac{\sigma}{2}\norm{\cdot}^2_{\mathcal{M}}$ is
		convex, where $\norm{v}^2_{\mathcal{M}}=\inner{v}{\mathcal{M}v}$. For
		$\mathcal{M}=I$ the metric is omitted.
	\end{definition}
	
	\begin{proposition}[Attainment under strong convexity]
		A closed proper strongly convex $f$ attains its infimum at exactly one point.
	\end{proposition}
	
	\begin{definition}[Indicator]
		$\delta_C(v)=0$ for $v\in C$ and $\delta_C(v)=+\infty$ otherwise.
	\end{definition}
	
	Given $v\in C$, the normal cone at $v$ collects the directions making an
	obtuse angle with every direction pointing from $v$ into $C$. 
	\begin{definition}[Normal cone]
		For $C$ convex and $v\in C$,
		$N_C(v)=\{z:\inner{z}{w-v}\leq0\ \text{ for all }w\in C\}$.
	\end{definition}
	
	\begin{definition}[Subdifferential]
		For $f$ convex,
		\begin{equation}
			\partial f(v)=\bigl\{z:f(w)\geq f(v)+\inner{z}{w-v}\ \text{ for all }w
			\bigr\} .
		\end{equation}
	\end{definition}
	
	\begin{proposition}[Normal for a convex cone]
		For $v\in K$, $K$ convex cone, the normal cone at $v$ is given by:
		\begin{equation}
			N_K(v) = \{z \in K^{\circ} : \langle z,v \rangle = 0\}
		\end{equation}
	\end{proposition}
	\begin{proof}
		Let $z\in N_K(v)$. For every $u\in K$, the conic property gives $v+u\in K$.
		Therefore,
		\[
		\langle z,u\rangle
		=
		\langle z,(v+u)-v\rangle
		\leq0.
		\]
		Since this holds for every $u\in K$, we have $z\in K^\circ$.
		
		Taking $w=0$ in the definition of $N_K(v)$ gives
		\[
		\langle z,-v\rangle\leq0,
		\qquad\text{and hence}\qquad
		\langle z,v\rangle\geq0.
		\]
		On the other hand, $z\in K^\circ$ and $v\in K$ imply
		\[
		\langle z,v\rangle\leq0.
		\]
		Consequently,
		\[
		\langle z,v\rangle=0.
		\]
		
		Conversely, suppose that
		\[
		z\in K^\circ
		\qquad\text{and}\qquad
		\langle z,v\rangle=0.
		\]
		Then, for every $w\in K$,
		\[
		\langle z,w-v\rangle
		=
		\langle z,w\rangle-\langle z,v\rangle
		=
		\langle z,w\rangle
		\leq0.
		\]
		Thus, $z\in N_K(v)$.
	\end{proof}
	\begin{proposition}[Singleton subdifferential]
		For $f$ convex, if $\partial f(v)$ is a singleton then $f$ is differentiable
		at $v$.
	\end{proposition}
	
	\begin{proposition}[Subdifferential of the indicator]
		For $C$ closed convex and $v\in C$, $\partial\delta_C(v)=N_C(v)$, and
		$\partial\delta_C(v)=\emptyset$ for $v\notin C$. 
		
	\end{proposition}
	
	From the definition of the subdifferential it is immediate that
	\begin{proposition}[Fermat]
		$p$ minimises a convex $f$ over $\R^N$ if and only if $0\in\partial f(p)$.
	\end{proposition}
	
	\section{Differentiability}
	
	
	\begin{definition}[Lipschitz]
		$\mathcal{R}:\R^m\to\R^p$ is \emph{Lipschitz with constant $L$} if
		\begin{equation}
			\norm{\mathcal{R}(v)-\mathcal{R}(w)}\leq L\norm{v-w}
			\qquad\text{for all }v,w ,
		\end{equation}
		and \emph{locally Lipschitz} if every point has a neighbourhood on which it
		is Lipschitz with some constant.
	\end{definition}
	
	\begin{proposition}[Rademacher, \cite{EvansGariepy2015}]
		A locally Lipschitz $\mathcal{R}:\R^m\to\R^p$ is differentiable outside a set
		of Lebesgue measure zero.
	\end{proposition}
	
	It is then natural to extend the definition of differentiability in the following way:
	\begin{definition}[Bouligand and Clarke Jacobians]
		Let \(\mathcal R:\R^m\to\R^m\) be locally Lipschitz, and let
		\(D_{\mathcal R}\) denote the set of points at which \(\mathcal R\) is
		differentiable. The \emph{Bouligand Jacobian} of \(\mathcal R\) at \(v\) is
		\begin{equation}
			\partial_B\mathcal R(v)
			=
			\left\{
			\lim_{k\to\infty}\mathcal R'(v_k):
			v_k\to v,\quad v_k\in D_{\mathcal R}
			\right\}.
		\end{equation}
		The \emph{Clarke Jacobian} of \(\mathcal R\) at \(v\) is the convex hull of
		its Bouligand Jacobian:
		\begin{equation}
			\partial_C\mathcal R(v)
			=
			\conv\bigl(\partial_B\mathcal R(v)\bigr).
		\end{equation}
		We write
		\(\partial\mathcal R(v)=\partial_C\mathcal R(v)\).
	\end{definition}
	
	
	\section{Projection}
	
	
	
	\begin{definition}[Projection]
		For $C\subseteq\R^m$ nonempty closed convex and $v\in\R^m$,
		$\proj{v}{C}$ is a $w\in C$ minimising $\norm{v-w}$ over $C$.
	\end{definition}
	
	\begin{proposition}[Existence, uniqueness and characterisation]
		For $C$ nonempty closed convex the minimiser exists, is unique, and $p=\proj{v}{C}$
		is characterised by
		\begin{equation}
			p\in C,
			\qquad
			\inner{v-p}{w-p}\leq0\ \text{ for all }w\in C .
		\end{equation}
	\end{proposition}
	
	\begin{proof}
		Consider the function
		\[
		w\longmapsto
		\frac12\norm{v-w}^2+\delta_C(w).
		\]
		Since $C$ is nonempty, closed and convex, this function is closed, proper and
		strongly convex with modulus $1$. The preceding attainment result therefore
		shows that it has a unique minimiser. Thus, $\proj{v}{C}$ exists and is unique.
		
		Let $p=\proj{v}{C}$, and fix $w\in C$. By convexity of $C$,
		\[
		p+t(w-p)\in C
		\qquad\text{for every }t\in[0,1].
		\]
		Consequently, the function
		\[
		\varphi(t)
		=
		\frac12\norm{v-p-t(w-p)}^2
		\]
		has a minimum over $[0,1]$ at $t=0$. Its right derivative at zero therefore
		satisfies
		\[
		0\leq\varphi'(0)
		=
		-\inner{v-p}{w-p}.
		\]
		Hence,
		\[
		\inner{v-p}{w-p}\leq0
		\qquad\text{for every }w\in C.
		\]
		
		Conversely, suppose that $p\in C$ satisfies
		\[
		\inner{v-p}{w-p}\leq0
		\qquad\text{for every }w\in C.
		\]
		Then, for every $w\in C$,
		\[
		\begin{aligned}
			\norm{v-w}^2
			&=\norm{(v-p)-(w-p)}^2\\
			&=\norm{v-p}^2
			-2\inner{v-p}{w-p}
			+\norm{w-p}^2\\
			&\geq\norm{v-p}^2.
		\end{aligned}
		\]
		Thus, $p$ minimises the distance from $v$ to $C$, and therefore
		$p=\proj{v}{C}$.
	\end{proof}
	
	\begin{definition}[Firm nonexpansiveness]
		$T:\R^m\to\R^m$ is \emph{firmly nonexpansive} if
		\begin{equation}
			\norm{T(v)-T(w)}^2\leq\inner{T(v)-T(w)}{v-w}
			\qquad\text{for all }v,w .
		\end{equation}
	\end{definition}
	\begin{definition}[Monotone single-valued map]
		$T:\R^m\to\R^m$ is \emph{monotone} if
		\begin{equation}
			\inner{T(v)-T(w)}{v-w}\geq0
			\qquad\text{for all }v,w .
		\end{equation}
	\end{definition}
	
	\begin{proposition}[Firm nonexpansiveness implies Lipschitz and monotone]
		A firmly nonexpansive $T$ is Lipschitz with constant $1$ and monotone.
	\end{proposition}
	
	\begin{proof}
		Monotonicity is immediate, the right side of the defining inequality being at
		least the left, which is nonnegative. For the Lipschitz bound, Cauchy--Schwarz
		gives $\norm{T(v)-T(w)}^2\leq\norm{T(v)-T(w)}\norm{v-w}$, and dividing by
		$\norm{T(v)-T(w)}$ when it is nonzero gives the claim; when it is zero the
		claim is trivial.
	\end{proof}
	
	\begin{proposition}[Firm nonexpansiveness of the projection]
		$\proj{\cdot}{C}$ is firmly nonexpansive for every nonempty closed convex $C$.
	\end{proposition}
	
	\begin{proof}
		Write $p=\proj{v}{C}$ and $q=\proj{w}{C}$. The characterisation at $v$ with the
		test point $q$ gives $\inner{v-p}{q-p}\leq0$, and at $w$ with the test point
		$p$ gives $\inner{w-q}{p-q}\leq0$. Adding them,
		$\inner{(v-p)-(w-q)}{q-p}\leq0$, that is
		$\inner{v-w}{q-p}\leq\inner{p-q}{q-p}=-\norm{p-q}^2$, which on negating is the
		required inequality.
	\end{proof}
	
	
	Rademacher's theorem therefore applies and the Clarke Jacobian is well
	defined. The following theorem links the projection on $K$ with the projection
	on its polar, generalising the orthogonal decomposition along a subspace. 
	
	
	\begin{proposition}[Moreau]
		Let $K$ be a closed convex cone. Every $v\in\R^m$ decomposes uniquely as
		\begin{equation}
			v=\proj{v}{K}+\proj{v}{K^\circ},
			\qquad
			\inner{\proj{v}{K}}{\proj{v}{K^\circ}}=0 .
		\end{equation}
	\end{proposition}
	
	\begin{proof}
		Set
		\[
		p=\proj{v}{K},
		\qquad
		q=v-p.
		\]
		By the characterisation of the projection,
		\[
		q\in N_K(p)=K^\circ\cap p^\perp.
		\]
		Consequently,
		\[
		q\in K^\circ
		\qquad\text{and}\qquad
		\langle p,q\rangle=0.
		\]
		
		Since $K$ is a closed convex cone, the bipolar theorem gives
		$(K^\circ)^\circ=K$. Therefore,
		\[
		v-q=p
		\in K\cap q^\perp
		=(K^\circ)^\circ\cap q^\perp
		=N_{K^\circ}(q).
		\]
		The characterisation of the projection now gives
		\[
		q=\proj{v}{K^\circ}.
		\]
		Hence,
		\[
		v=\proj{v}{K}+\proj{v}{K^\circ},
		\qquad
		\left\langle\proj{v}{K},\proj{v}{K^\circ}\right\rangle=0.
		\]
		
		For uniqueness, suppose that
		\[
		v=p'+q',
		\qquad
		p'\in K,
		\qquad
		q'\in K^\circ,
		\qquad
		\langle p',q'\rangle=0.
		\]
		Then
		\[
		v-p'=q'\in K^\circ\cap(p')^\perp=N_K(p').
		\]
		Thus,
		\[
		p'=\proj{v}{K},
		\qquad
		q'=v-p'=\proj{v}{K^\circ}.
		\]
	\end{proof}
	\begin{proposition}
		Let $K$ be a nonempty closed convex cone, and define
		\[
		f(x)=\frac12\|\Pi_K(x)\|^2.
		\]
		Then $f$ is continuously differentiable and
		\[
		\nabla f(x)=\Pi_K(x).
		\]
		
	\end{proposition}
	
	\begin{proof}
		Set $p=\Pi_K(x)$. The Moreau decomposition gives
		$\langle x,p\rangle=\|p\|^2$. Moreover, by completing the square, $p$ is the
		unique maximiser over $y\in K$ of
		$\langle x,y\rangle-\frac12\|y\|^2$. Consequently,
		$f(x)=\max_{y\in K}\{\langle x,y\rangle-\frac12\|y\|^2\}$.
		
		Now let $p_h=\Pi_K(x+h)$. Using $p$ as a candidate at $x+h$ gives
		$f(x+h)\geq f(x)+\langle p,h\rangle$. On the other hand, using the optimality
		of $p$ at $x$ gives
		$f(x+h)\leq f(x)+\langle p_h,h\rangle$. Therefore,
		\[
		0
		\leq
		f(x+h)-f(x)-\langle p,h\rangle
		\leq
		\langle p_h-p,h\rangle
		\leq
		\|p_h-p\|\,\|h\|
		\leq
		\|h\|^2,
		\]
		where the last inequality follows from the $1$-Lipschitz continuity of the
		projection. Hence
		$f(x+h)-f(x)-\langle p,h\rangle=o(\|h\|)$, which proves that
		$\nabla f(x)=p=\Pi_K(x)$. 
	\end{proof}
	\begin{proposition}[Distance in terms of a projection]
		For $K$ a closed convex cone,
		\begin{equation}
			\dist(v,K)=\norm{\proj{v}{K^\circ}}=\norm{\proj{-v}{K^*}} ,
		\end{equation}
		so $v\in K$ if and only if $\proj{-v}{K^*}=0$.
	\end{proposition}
	
	\begin{proposition}[The Clarke Jacobian of a projection]
		Let $K$ be a closed convex cone and $F\in\partial\proj{\cdot}{K^*}(\sigma)$.
		Then $F=F^\top$ and $0\preceq F\preceq I$.
	\end{proposition}
	\begin{proof}
		Set $P=\proj{\cdot}{K^*}$. From the preceding proposition,
		\[
		P=\nabla\phi,
		\qquad
		\phi(x)=\frac12\|P(x)\|^2,
		\]
		and $\phi$ is convex. Moreover, the Moreau decomposition gives
		\[
		I-P=\proj{\cdot}{(K^*)^\circ}=\proj{\cdot}{-K}.
		\]
		Therefore,
		\[
		I-P=\nabla\psi,
		\qquad
		\psi(x)=\frac12\|\proj{x}{-K}\|^2,
		\]
		where $\psi$ is also convex.
		
		Consider a point $x$ at which $P$ is differentiable, and write
		$J=DP(x)$. Since $P=\nabla\phi$, the matrix $J=\nabla^2\phi(x)$ is symmetric
		and positive semidefinite. Hence
		\[
		J=J^\top,
		\qquad
		J\succeq0.
		\]
		Similarly, $I-P=\nabla\psi$ is differentiable at $x$, with derivative $I-J$.
		Since $\psi$ is convex,
		\[
		I-J\succeq0.
		\]
		Thus, at every point where $P$ is differentiable,
		\[
		J=J^\top,
		\qquad
		0\preceq J\preceq I.
		\]
		
		The projection $P$ is Lipschitz and is therefore differentiable almost
		everywhere. By definition, every matrix in the Clarke generalized Jacobian
		$\partial P(\sigma)$ is a convex combination of limits of matrices
		$DP(x_k)$, where $x_k\to\sigma$ and $P$ is differentiable at every $x_k$.
		Since the set
		\[
		\left\{A:A=A^\top,\ 0\preceq A\preceq I\right\}
		\]
		is closed and convex, these properties are preserved under limits and convex
		combinations. Consequently, every $F\in\partial P(\sigma)$ satisfies
		\[
		F=F^\top,
		\qquad
		0\preceq F\preceq I.
		\]
	\end{proof}
	
	
	
	\begin{proposition}[Action on the argument]
		$F\sigma=\proj{\sigma}{K}$ for every
		$F\in\partial\proj{\cdot}{K}(\sigma)$ and every $\sigma$.
	\end{proposition}
	\begin{proof}
		Set $P=\proj{\cdot}{K}$. Since $K$ is a cone, its projection is positively
		homogeneous:
		\[
		P(tx)=tP(x)
		\qquad\text{for every }t\geq0.
		\]
		At any point $x$ where $P$ is differentiable, differentiating this identity
		with respect to $t$ at $t=1$ gives
		\[
		DP(x)x=P(x).
		\]
		
		Now let $J$ be a limiting Jacobian of $P$ at $\sigma$. There is a sequence
		$x_k\to\sigma$ such that $P$ is differentiable at every $x_k$ and
		$DP(x_k)\to J$. Since $P$ is $1$-Lipschitz,
		$\|DP(x_k)\|\leq1$. Therefore,
		\[
		\begin{aligned}
			\|DP(x_k)\sigma-P(\sigma)\|
			&\leq
			\|DP(x_k)(\sigma-x_k)\|
			+\|DP(x_k)x_k-P(\sigma)\| \\
			&=
			\|DP(x_k)(\sigma-x_k)\|
			+\|P(x_k)-P(\sigma)\| \\
			&\leq
			2\|x_k-\sigma\|.
		\end{aligned}
		\]
		Letting $k\to\infty$ gives
		\[
		J\sigma=P(\sigma).
		\]
		
		Finally, the Clarke generalized Jacobian is the convex hull of such limiting
		Jacobians. Since the identity $J\sigma=P(\sigma)$ is preserved under limits
		and convex combinations, every $F\in\partial P(\sigma)$ satisfies
		\[
		F\sigma=P(\sigma)=\proj{\sigma}{K}.
		\]
	\end{proof}
	Replacing $\sigma$ by his projection on $K$ and using the fact that $\Pi_{K}\circ \Pi_{K} = \Pi_{K}$ , we deduce that: 
	\begin{proposition}[Fixed point formulation]
		$p = \Pi_{K}(\sigma)$ then: $$ Fp = p$$
		
	\end{proposition}
	
	
	A natural question is which metrics leave the Euclidean projection onto $K$
	unchanged.
	
	\begin{definition}[Admissible metric]
		Let $K\subseteq\mathbb R^n$ be a nonempty closed convex cone, and let
		$M\in\mathbb S_{++}^n$. Define
		\[
		\langle x,y\rangle_M=\langle x,My\rangle,
		\qquad
		\norm{x}_M^2=\langle x,Mx\rangle.
		\]
		The projection onto $K$ in the metric $M$ is
		\[
		\proj{v}{K}^{M}
		=
		\argmin_{p\in K}\frac12\norm{v-p}_M^2.
		\]
		The metric $M$ is said to be \emph{admissible for $K$} if
		\[
		\proj{v}{K}^{M}=\proj{v}{K}
		\qquad\text{for every }v\in\mathbb R^n.
		\]
	\end{definition}
	
	\begin{definition}[Polar cone in the metric $M$]
		The polar cone of $K$ in the metric $M$ is
		\[
		K^{\circ_M}
		=
		\left\{
		z\in\mathbb R^n:
		\langle z,w\rangle_M\leq0
		\text{ for every }w\in K
		\right\}.
		\]
		Equivalently,
		\[
		K^{\circ_M}=M^{-1}K^\circ.
		\]
	\end{definition}
	
	\begin{proposition}[Characterisation of the metric projection]
		Let $p_M\in\mathbb R^n$. Then
		\[
		p_M=\proj{v}{K}^{M}
		\]
		if and only if
		\[
		p_M\in K,
		\qquad
		v-p_M\in K^{\circ_M},
		\qquad
		\langle p_M,v-p_M\rangle_M=0.
		\]
		Equivalently,
		\[
		p_M\in K,
		\qquad
		M(v-p_M)\in K^\circ,
		\qquad
		\langle p_M,M(v-p_M)\rangle=0.
		\]
	\end{proposition}
	
	There is a simple characterisation of admissibility.
	\begin{theorem}[Commutation characterisation of admissibility]
		Let $K\subseteq\mathbb R^n$ be a nonempty closed convex cone and let
		$M\in\mathbb S_{++}^n$. Then $M$ is admissible for $K$ if and only if
		\[
		M\proj{v}{K}=\proj{Mv}{K}
		\qquad\text{for every }v\in\mathbb R^n.
		\]
	\end{theorem}
	
	\begin{proof}
		Write $P=\proj{\cdot}{K}$. Recall that the metric projection is characterised
		by
		\[
		p=\proj{v}{K}^{M}
		\quad\Longleftrightarrow\quad
		p\in K,\quad M(v-p)\in K^\circ,\quad
		\langle p,M(v-p)\rangle=0.
		\]
		
		Suppose first that $M$ is admissible. Let $p\in K$ and
		$r\in N_K(p)$. Since
		\[
		p=P(p+r)=\proj{p+r}{K}^{M},
		\]
		the characterisation of the metric projection gives
		\[
		Mr\in N_K(p).
		\]
		Thus,
		\[
		M N_K(p)\subseteq N_K(p)
		\qquad\text{for every }p\in K.
		\]
		In particular, taking $p=0$ gives
		\[
		MK^\circ\subseteq K^\circ.
		\]
		Since $M$ is symmetric, this also implies $MK\subseteq K$. Indeed, for
		$x\in K$ and $z\in K^\circ$,
		\[
		\langle Mx,z\rangle=\langle x,Mz\rangle\leq0.
		\]
		
		Now let
		\[
		p=P(v),
		\qquad
		q=v-p.
		\]
		Then $q\in N_K(p)$. Consequently,
		\[
		Mq\in N_K(p)
		\qquad\text{and}\qquad
		M^2q\in N_K(p).
		\]
		It follows that $Mq\in K^\circ$ and
		\[
		\langle Mp,Mq\rangle
		=
		\langle p,M^2q\rangle
		=
		0.
		\]
		Since $Mp\in K$ and
		\[
		Mv=Mp+Mq,
		\]
		the Moreau characterisation gives
		\[
		P(Mv)=Mp=MP(v).
		\]
		
		Conversely, suppose that
		\[
		P(Mv)=MP(v)
		\qquad\text{for every }v.
		\]
		Since $\proj{v}{K^\circ}=v-P(v)$, we also have
		\[
		\proj{Mv}{K^\circ}
		=
		M\proj{v}{K^\circ}.
		\]
		
		Fix $v$, and set
		\[
		p=P(v),
		\qquad
		q=\proj{v}{K^\circ}=v-p.
		\]
		By repeated application of the commutation identity, for every integer
		$k\geq0$,
		\[
		P(M^kv)=M^kp,
		\qquad
		\proj{M^kv}{K^\circ}=M^kq.
		\]
		The Moreau decomposition therefore gives
		\[
		0
		=
		\langle M^kp,M^kq\rangle
		=
		\langle p,M^{2k}q\rangle.
		\]
		
		Let $\mu_1,\dots,\mu_s$ be the distinct eigenvalues of $M$, and let
		$E_1,\dots,E_s$ be the corresponding orthogonal spectral projectors. Writing
		\[
		a_j=\langle p,E_jq\rangle,
		\]
		the preceding identities for $k=0,\dots,s-1$ become
		\[
		\sum_{j=1}^s\mu_j^{2k}a_j=0.
		\]
		Since the numbers $\mu_j^2$ are distinct, the corresponding Vandermonde matrix
		is invertible. Hence $a_j=0$ for every $j$, and therefore
		\[
		\langle p,Mq\rangle
		=
		\sum_{j=1}^s\mu_j a_j
		=
		0.
		\]
		
		Moreover,
		\[
		Mq=\proj{Mv}{K^\circ}\in K^\circ.
		\]
		Thus,
		\[
		p\in K,\qquad M(v-p)=Mq\in K^\circ,\qquad
		\langle p,M(v-p)\rangle=0.
		\]
		The characterisation of the metric projection gives
		\[
		p=\proj{v}{K}^{M}.
		\]
		Since $p=P(v)$, the metric $M$ is admissible.
	\end{proof}
	
	
	By Moreau we have immediately that:
	\begin{proposition}[The three forms of admissibility]
		For $M\succ0$ the three conditions
		$M\proj{\cdot}{K}=\proj{M\,\cdot}{K}$,
		$M\proj{\cdot}{K^*}=\proj{M\,\cdot}{K^*}$ and
		$M\proj{\cdot}{K^\circ}=\proj{M\,\cdot}{K^\circ}$ are equivalent.
	\end{proposition}
	
	Below we give some property of admissible metric we will use.
	\begin{proposition}[Structure and closure of admissible metrics]
		Let $K$ be a closed convex cone and let $M\in\mathcal A(K)$. Then:
		
		\begin{enumerate}
			\item The cones are invariant under $M$:
			\[
			MK=K
			\]
			In particular, $M\in\operatorname{Aut}(K)$.
			
			\item For every $\sigma$ and every
			$F\in\partial\proj{\cdot}{K^*}(\sigma)$,
			\[
			MF=FM.
			\]
			Consequently,
			\[
			M^{-1}F=(M^{-1}F)^\top,
			\qquad
			0\preceq M^{-1}F\preceq M^{-1}.
			\]
			
			\item The set $\mathcal A(K)$ is closed under positive linear combinations,
			inversion and positive-definite limits. If $M,N\in\mathcal A(K)$ and $MN=NM$,
			then
			\[
			MN\in\mathcal A(K).
			\]
			Thus, $\mathcal A(K)\cup\{0\}$ is a convex cone.
			
			\item For every $\alpha\in\mathbb R$,
			\[
			M^\alpha\in\mathcal A(K).
			\]
			
		\end{enumerate}
	\end{proposition}
	
	\begin{proof}
		We first record a useful characterisation. For every $p\in K$,
		\[
		M\in\mathcal A(K)
		\quad\Longleftrightarrow\quad
		MN_K(p)\subseteq N_K(p).
		\]
		Indeed, if $q\in N_K(p)$, then
		\[
		p=\proj{p+q}{K}.
		\]
		Admissibility and the characterisation of the metric projection imply
		\[
		Mq\in N_K(p).
		\]
		Conversely, if $MN_K(p)\subseteq N_K(p)$ for every $p\in K$, then, for
		$p=\proj{v}{K}$ and $q=v-p\in N_K(p)$, one has $Mq\in N_K(p)$. Hence
		$p=\proj{v}{K}^{M}$, so $M$ is admissible.
		
		\paragraph{Convexity.}
		Let $M,N\in\mathcal A(K)$ and let $\alpha,\beta>0$. For every $p\in K$ and
		$q\in N_K(p)$,
		\[
		Mq\in N_K(p),
		\qquad
		Nq\in N_K(p).
		\]
		Since $N_K(p)$ is a convex cone,
		\[
		(\alpha M+\beta N)q
		=
		\alpha Mq+\beta Nq
		\in N_K(p).
		\]
		Therefore,
		\[
		\alpha M+\beta N\in\mathcal A(K).
		\]
		Thus, $\mathcal A(K)$ is closed under positive linear combinations.
		
		\paragraph{Inversion and invariance of the cones.}
		Write $P=\proj{\cdot}{K}$. By the commutation characterisation,
		\[
		MP(v)=P(Mv).
		\]
		Replacing $v$ by $M^{-1}v$ gives
		\[
		P(M^{-1}v)=M^{-1}P(v),
		\]
		so
		\[
		M^{-1}\in\mathcal A(K).
		\]
		
		If $v\in K$, then $P(v)=v$, and therefore
		\[
		P(Mv)=MP(v)=Mv.
		\]
		Thus, $Mv\in K$, which gives $MK\subseteq K$. Applying the same argument to
		$M^{-1}$ gives the reverse inclusion, so
		\[
		MK=K.
		\]
		Consequently, $M\in\operatorname{Aut}(K)$.
		
		\paragraph{Products and limits.}
		Let $M,N\in\mathcal A(K)$ and suppose that $MN=NM$. Then $MN$ is symmetric
		positive definite, and
		\[
		MNP(v)
		=
		MP(Nv)
		=
		P(MNv).
		\]
		Hence,
		\[
		MN\in\mathcal A(K).
		\]
		
		Now suppose that $M_j\in\mathcal A(K)$ and
		\[
		M_j\longrightarrow M\succ0.
		\]
		Using the continuity of the projection,
		\[
		MP(v)
		=
		\lim_{j\to\infty}M_jP(v)
		=
		\lim_{j\to\infty}P(M_jv)
		=
		P(Mv).
		\]
		Therefore, $M\in\mathcal A(K)$.
		
		\paragraph{Commutation with generalized Jacobians.}
		Set
		\[
		P_*=\proj{\cdot}{K^*}.
		\]
		Let $x$ be a point at which $P_*$ is differentiable, and write
		\[
		J=DP_*(x).
		\]
		Since
		\[
		P_*(y)=MP_*(M^{-1}y),
		\]
		the projection is differentiable at $Mx$, with
		\[
		J'=DP_*(Mx)=MJM^{-1}.
		\]
		Both $J$ and $J'$ are symmetric. Therefore,
		\[
		MJM^{-1}
		=
		J'
		=
		(J')^\top
		=
		M^{-1}JM,
		\]
		and hence
		\[
		M^2J=JM^2.
		\]
		Since $M\succ0$, the matrices $M$ and $M^2$ have the same eigenspaces. It
		follows that
		\[
		MJ=JM.
		\]
		
		This relation is preserved under limits and convex combinations. Therefore,
		for every $\sigma$ and every
		$F\in\partial P_*(\sigma)$,
		\[
		MF=FM.
		\]
		Since $F=F^\top$ and $0\preceq F\preceq I$, simultaneous diagonalisation gives
		\[
		M^{-1}F=(M^{-1}F)^\top,
		\qquad
		0\preceq M^{-1}F\preceq M^{-1}.
		\]
		
		\paragraph{Square roots and real powers.}
		We first prove that $M^{1/2}$ is admissible. Consider the Newton iteration
		\[
		X_0=I,
		\qquad
		X_{j+1}
		=
		\frac12\left(X_j+MX_j^{-1}\right).
		\]
		Every $X_j$ is a positive-definite rational function of $M$ and therefore
		commutes with $M$. Suppose that $X_j$ is admissible. Then $X_j^{-1}$ is
		admissible, and the commuting-product property gives
		\[
		MX_j^{-1}\in\mathcal A(K).
		\]
		Closure under positive linear combinations then gives
		\[
		X_{j+1}\in\mathcal A(K).
		\]
		Since $X_0=I$ is admissible, every $X_j$ is admissible.
		
		Diagonalising $M$ reduces this iteration to the scalar Newton iteration for
		the positive square root. Hence,
		\[
		X_j\longrightarrow M^{1/2}.
		\]
		Closure under positive-definite limits gives
		\[
		M^{1/2}\in\mathcal A(K).
		\]
		
		Repeated square roots, commuting products and inversion now give
		\[
		M^{m/2^r}\in\mathcal A(K)
		\]
		for every $m\in\mathbb Z$ and every nonnegative integer $r$.
		
		Let $\alpha\in\mathbb R$, and choose dyadic rational numbers $\alpha_j$ such
		that $\alpha_j\to\alpha$. We then have,
		\[
		M^{\alpha_j}\longrightarrow M^\alpha.
		\]
		Therefore,
		\[
		M^\alpha\in\mathcal A(K).
		\]
		
	\end{proof}
	
	\begin{proposition}[Weighted squared norm of the projection]
		Let $K\subseteq\mathbb R^n$ be a closed convex cone and let
		$M\in\mathcal A(K)$. Define
		\[
		\phi_M(x)
		=
		\frac12\norm{\proj{x}{K}}_M^2.
		\]
		Then $\phi_M$ is convex and continuously differentiable, with
		\[
		\nabla\phi_M(x)=M\proj{x}{K}.
		\]
		Equivalently,
		\[
		\partial\phi_M(x)
		=
		\left\{M\proj{x}{K}\right\}.
		\]
	\end{proposition}
	
	\begin{proof}
		Since $M$ is admissible,
		\[
		\proj{x}{K}^{M}=\proj{x}{K}.
		\]
		Completing the square gives
		\[
		\langle Mx,y\rangle-\frac12\norm{y}_M^2
		=
		\frac12\norm{x}_M^2-\frac12\norm{x-y}_M^2.
		\]
		Consequently,
		\[
		\proj{x}{K}
		=
		\argmax_{y\in K}
		\left\{
		\langle Mx,y\rangle-\frac12\norm{y}_M^2
		\right\}.
		\]
		If $p=\proj{x}{K}$, metric complementarity gives
		\[
		\langle Mp,x-p\rangle=0,
		\]
		and hence
		\[
		\phi_M(x)
		=
		\max_{y\in K}
		\left\{
		\langle Mx,y\rangle-\frac12\norm{y}_M^2
		\right\}.
		\]
		This representation also shows that $\phi_M$ is convex.
		
		Now let
		\[
		p=\proj{x}{K},
		\qquad
		p_h=\proj{x+h}{K}.
		\]
		Using $p$ as a candidate at $x+h$ gives
		\[
		\phi_M(x+h)
		\geq
		\phi_M(x)+\langle Mp,h\rangle.
		\]
		Using the optimality of $p$ at $x$ and the maximiser $p_h$ at $x+h$ gives
		\[
		\phi_M(x+h)
		\leq
		\phi_M(x)+\langle Mp_h,h\rangle.
		\]
		Therefore,
		\[
		\begin{aligned}
			0
			&\leq
			\phi_M(x+h)-\phi_M(x)-\langle Mp,h\rangle\\
			&\leq
			\langle M(p_h-p),h\rangle\\
			&\leq
			\norm{p_h-p}_M\norm{h}_M\\
			&\leq
			\norm{h}_M^2,
		\end{aligned}
		\]
		where the last inequality follows from the nonexpansiveness of the metric
		projection. Thus,
		\[
		\phi_M(x+h)-\phi_M(x)-\langle Mp,h\rangle
		=
		o(\norm{h}),
		\]
		which proves
		\[
		\nabla\phi_M(x)=Mp=M\proj{x}{K}.
		\]
		Continuity of the gradient follows from continuity of the projection.
	\end{proof}
	\begin{proposition}[contraction formula]
		Let \(P=\Pi_{K^*}\), let \(M\in\mathcal A(K^*)\), and suppose that \(P\) is
		twice differentiable at \(\sigma\). Set
		\[
		p=P(\sigma)
		\qquad\text{and}\qquad
		F=DP(\sigma).
		\]
		Then
		\begin{equation}
			D^2P(\sigma)[h,Mp]
			=
			MF(I-F)h
			\qquad
			\text{for every }h\in\R^m.
		\end{equation}
		Equivalently,
		\begin{equation}
			D^2P(\sigma)[\,\cdot,Mp]
			=
			MF(I-F).
		\end{equation}
	\end{proposition}
	
	\begin{proof}
		At every point \(x\) at which \(P\) is differentiable, the derivative of
		the projection satisfies
		\begin{equation}
			DP(x)P(x)=P(x).
		\end{equation}
		Differentiating this identity at \(x=\sigma\) in the direction \(h\) gives
		\[
		D^2P(\sigma)[h,p]+F^2h=Fh.
		\]
		It follows that
		\begin{equation}
			D^2P(\sigma)[h,p]
			=
			F(I-F)h.
		\end{equation}
		
		Since \(M\) is admissible, it commutes with the derivative of \(P\):
		\[
		DP(x)M=MDP(x)
		\]
		at every point where \(P\) is differentiable. Differentiating this identity
		at \(\sigma\) in the direction \(h\), and then applying the resulting
		identity to \(p\), gives
		\[
		D^2P(\sigma)[h,Mp]
		=
		M D^2P(\sigma)[h,p].
		\]
		Using the preceding formula, we obtain
		\[
		D^2P(\sigma)[h,Mp]
		=
		MF(I-F)h.
		\]
	\end{proof}
	
	
	
	
	\begin{proposition}[Diagonal admissible metrics on a product]
		Let
		\[
		K=K_1\times\cdots\times K_p
		\]
		be a product of cones from the catalogue, and let
		\[
		M=\diag(D_1,\ldots,D_p)\succ0,
		\]
		where each \(D_j\) is diagonal in the coordinates of the corresponding
		block. Then \(M\in\mathcal A(K)\) if and only if \(D_j\) is a positive
		multiple of the identity on every block, except on zero-cone and
		nonnegative-orthant blocks, where any positive diagonal matrix is allowed.
	\end{proposition}
	
	\begin{proof}
		Since
		\[
		K^*=K_1^*\times\cdots\times K_p^*
		\]
		and the projection onto \(K^*\) acts blockwise, we have
		\[
		\Pi_{K^*}(v_1,\ldots,v_p)
		=
		\bigl(\Pi_{K_1^*}(v_1),\ldots,\Pi_{K_p^*}(v_p)\bigr).
		\]
		Consequently, \(M\) is admissible for \(K\) if and only if \(D_j\) is
		admissible for \(K_j\) for every \(j\). It is therefore enough to consider
		one block.
		
		For a zero-cone block, the projection is constant, while the projection
		onto its dual is the identity. Hence every positive diagonal matrix is
		admissible.
		
		For a nonnegative-orthant block,
		\[
		\Pi_{\R_+^d}(v)=\max(v,0)
		\]
		componentwise. If \(D\succ0\) is diagonal, then
		\[
		\Pi_{\R_+^d}(Dv)
		=
		\max(Dv,0)
		=
		D\max(v,0)
		=
		D\Pi_{\R_+^d}(v).
		\]
		Thus every positive diagonal matrix is admissible.
		
		Now consider any other cone \(C\) from the catalogue. Such a cone admits a
		boundary point \(x\) whose normal cone is a ray
		\[
		N_C(x)=\R_+n,
		\]
		where every coordinate of \(n\) is nonzero. If \(D\) is admissible, the
		normal-cone characterization of admissibility gives
		\[
		DN_C(x)\subseteq N_C(x).
		\]
		In particular, \(Dn\in\R_+n\), so there exists \(\lambda>0\) such that
		\[
		Dn=\lambda n.
		\]
		Writing \(D=\diag(d_1,\ldots,d_d)\), we obtain
		\[
		d_i n_i=\lambda n_i
		\qquad\text{for every }i.
		\]
		Since every \(n_i\) is nonzero,
		\[
		d_i=\lambda
		\qquad\text{for every }i.
		\]
		Therefore,
		\[
		D=\lambda I.
		\]
		
		Such a full-support normal can be chosen on every Lorentz, semidefinite,
		exponential, and power-cone block. For example, on a Lorentz block one may
		take
		\[
		x=(1,u),
		\qquad
		\|u\|=1,
		\]
		with every coordinate of \(u\) nonzero. Then
		\[
		N_C(x)=\R_+(-1,u),
		\]
		whose generator has no zero coordinate. The other blocks are handled by
		choosing a generic smooth boundary point, whose normal likewise has full
		support.
		
		Conversely, every positive multiple of the identity commutes with a
		positively homogeneous projection:
		\[
		\Pi_{C^*}(\lambda v)
		=
		\lambda\Pi_{C^*}(v),
		\qquad \lambda>0.
		\]
		Hence every block matrix of the stated form is admissible, and therefore so
		is their block-diagonal product.
	\end{proof}
	
	
	\chapter{Convex program, Lagrangian and KKT}
	
	This chapter introduces what are convex program and the KKT conditions: a set of relations equivalent to optimality
	that guides the design of every algorithm that follows.
	\section{Convex program and Lagrangian}
	\begin{definition}[Convex program]
		Let $f:\R^n\to\R$ be convex, let $c:\R^n\to\R^m$ be affine, written
		$c(x)=Cx+d$, and let $K\subseteq\R^m$ be a closed convex cone. The problem
		is
		\begin{equation}
			\min_{x\in\R^n}\ f(x)
			\qquad\text{subject to}\qquad
			c(x)\preceq_K0 .
		\end{equation}
		Its optimal value is written $p^\star$, and a point satisfying the constraint
		is called \emph{feasible}.
	\end{definition}
	
	The problem the solver addresses is a conic quadratic program.
	\begin{definition}[Conic quadratic program]
		The \emph{conic quadratic program} is the particular case in which $f$ is the
		convex quadratic
		\begin{equation}
			f(x)=\tfrac12\inner{x}{Px}+\inner{q}{x},
			\qquad P\in\Sp^n,\ P\succeq0,\ q\in\R^n ,
		\end{equation}
		so that $f$ is differentiable with $\nabla f(x)=Px+q$. This is the problem
		carried by the rest of the document.
	\end{definition}
	
	\begin{definition}[Lagrangian]
		\begin{equation}
			L(x,z)=f(x)+\inner{z}{c(x)},
			\qquad z\in K^* .
		\end{equation}
	\end{definition}
	
	\begin{proposition}[Dual function]
		Let
		\[
		g(z)=\inf_{x}L(x,z),
		\qquad z\in K^*.
		\]
		Then \(g\) is concave and upper semicontinuous. In the conic quadratic case,
		\(g(z)>-\infty\) if and only if
		\begin{equation}
			Px+q+C^\top z=0
			\label{eq:dual-stationarity}
		\end{equation}
		admits a solution \(x_z\). In that case,
		\begin{equation}
			g(z)
			=
			-\frac12\inner{x_z}{Px_z}
			+
			\inner{z}{d},
		\end{equation}
		and this value is independent of the choice of \(x_z\).
	\end{proposition}
	
	\begin{proof}
		For every fixed \(x\), the mapping
		\[
		z\longmapsto L(x,z)
		=
		f(x)+\inner{z}{Cx+d}
		\]
		is affine and continuous. Since a pointwise infimum of affine functions is
		concave, \(g\) is concave. Moreover, for every \(\alpha\in\R\),
		\[
		\{z:g(z)<\alpha\}
		=
		\bigcup_x\{z:L(x,z)<\alpha\},
		\]
		which is open. Hence \(g\) is upper semicontinuous.
		
		In the conic quadratic case,
		\[
		L(x,z)
		=
		\frac12\inner{x}{Px}
		+
		\inner{q+C^\top z}{x}
		+
		\inner{z}{d},
		\]
		where \(P\succeq0\). Set
		\[
		b_z=q+C^\top z.
		\]
		If there exists \(x_z\) satisfying
		\[
		Px_z+b_z=0,
		\]
		then, for every \(h\),
		\[
		L(x_z+h,z)
		=
		L(x_z,z)
		+
		\frac12\inner{h}{Ph}.
		\]
		Since \(P\succeq0\), this shows that \(x_z\) is a global minimizer and that
		\(g(z)=L(x_z,z)>-\infty\).
		
		Conversely, suppose that \eqref{eq:dual-stationarity} has no solution. Since
		\(P\) is symmetric,
		\[
		\operatorname{range}(P)=\ker(P)^\perp,
		\]
		so \(b_z\notin\operatorname{range}(P)\). There therefore exists
		\(h\in\ker(P)\) such that
		\[
		\inner{b_z}{h}\neq0.
		\]
		Replacing \(h\) by \(-h\), if necessary, we may assume that
		\(\inner{b_z}{h}<0\). Then
		\[
		L(th,z)
		=
		t\inner{b_z}{h}
		+
		\inner{z}{d}
		\longrightarrow-\infty
		\qquad\text{as }t\to+\infty.
		\]
		Thus \(g(z)=-\infty\).
		
		Finally, if \(x_z\) satisfies \(Px_z+b_z=0\), then
		\begin{align}
			g(z)
			&=L(x_z,z)\\
			&=\frac12\inner{x_z}{Px_z}
			+\inner{b_z}{x_z}
			+\inner{z}{d}\\
			&=-\frac12\inner{x_z}{Px_z}
			+\inner{z}{d}.
		\end{align}
		If \(x_z'\) is another solution, then
		\[
		P(x_z-x_z')=0.
		\]
		Consequently,
		\[
		Px_z=Px_z'
		\]
		and
		\[
		\inner{x_z}{Px_z}
		=
		\inner{x_z'}{Px_z'}.
		\]
		Therefore, the displayed value of \(g(z)\) does not depend on the chosen
		solution.
	\end{proof}
	
	\begin{definition}[Dual problem]
		\begin{equation}
			\max_{z\in K^*}\ g(z) ,
		\end{equation}
		with optimal value $g^\star$. A $z\in K^*$ at which $g(z)$ is finite is
		\emph{dual feasible}.
	\end{definition}
	
	\begin{proposition}[Weak duality]
		For every primal feasible point \(x\) and every dual feasible point \(z\),
		\[
		g(z)\leq f(x).
		\]
		Consequently,
		\[
		g^\star\leq p^\star.
		\]
	\end{proposition}
	
	\begin{proof}
		Let \(x\) be primal feasible and \(z\) be dual feasible. Then
		\[
		-c(x)\in K
		\qquad\text{and}\qquad
		z\in K^*.
		\]
		By the definition of the dual cone,
		\[
		\inner{z}{-c(x)}\geq0,
		\]
		and hence
		\[
		\inner{z}{c(x)}\leq0.
		\]
		Using the definition of the dual function, we obtain
		\begin{align}
			g(z)
			&=\inf_u L(u,z)\\
			&\leq L(x,z)\\
			&=f(x)+\inner{z}{c(x)}\\
			&\leq f(x).
		\end{align}
		Taking the supremum over all dual feasible \(z\), and then the infimum over
		all primal feasible \(x\), gives
		\[
		g^\star
		=
		\sup_{z\in K^*}g(z)
		\leq
		\inf_{\,-c(x)\in K}f(x)
		=
		p^\star.
		\]
	\end{proof}
	
	\begin{definition}[Duality gap]
		The \emph{duality gap} of a feasible $x$ and a dual feasible $z$ is
		$f(x)-g(z)\geq0$. The \emph{optimal duality gap} is $p^\star-g^\star\geq0$.
	\end{definition}
	
	\begin{definition}[Strong duality]
		\emph{Strong duality} holds when $p^\star=g^\star$.
	\end{definition}
	
	Strong duality does not hold unconditionally; it requires a qualification
	condition. A classical one follows.
	\begin{definition}[Slater condition]
		The problem satisfies the \emph{Slater condition} if there is $x$ with
		$-c(x)\in\relint K$.
	\end{definition}
	
	\begin{proposition}[Slater implies strong duality]
		If the problem is bounded below and satisfies the Slater condition, then
		$p^\star=g^\star$ and the supremum defining $g^\star$ is attained.
	\end{proposition}
	
	\section{KKT conditions}
	
	In what follows the problem is solved through a system of relations equivalent
	to optimality whenever strong duality holds.
	
	\begin{proposition}[Karush, Kuhn and Tucker conditions]
		Assume that \(p^\star\) is finite, that strong duality holds,
		\[
		p^\star=g^\star,
		\]
		and that the dual supremum is attained. Then a feasible point \(x^\star\) is
		optimal if and only if there exists \(z^\star\) satisfying
		\begin{align}
			0&\in\partial f(x^\star)+C^\top z^\star,
			&&\text{stationarity},\\
			-c(x^\star)&\in K,
			&&\text{primal feasibility},\\
			z^\star&\in K^*,
			&&\text{dual feasibility},\\
			\inner{z^\star}{c(x^\star)}&=0,
			&&\text{complementarity}.
		\end{align}
		For \(z^\star\in K^*\), primal feasibility and complementarity are
		equivalent to
		\begin{equation}
			c(x^\star)\in N_{K^*}(z^\star).
		\end{equation}
		Thus the KKT conditions can equivalently be written as
		\begin{equation}
			0\in\partial f(x^\star)+C^\top z^\star,
			\qquad
			c(x^\star)\in N_{K^*}(z^\star).
		\end{equation}
	\end{proposition}
	
	\begin{proof}
		Suppose first that \(x^\star\) is primal optimal, and let \(z^\star\) be a
		dual optimal point. Strong duality gives
		\[
		f(x^\star)=p^\star=g^\star=g(z^\star).
		\]
		By the definition of the dual function,
		\[
		g(z^\star)\leq L(x^\star,z^\star).
		\]
		We have:
		\[
		f(x^\star)
		=
		g(z^\star)
		\leq
		L(x^\star,z^\star)
		\leq
		f(x^\star).
		\]
		Hence all the inequalities are equalities. In particular,
		\[
		\inner{z^\star}{c(x^\star)}=0.
		\]
		
		The equality
		\[
		g(z^\star)=L(x^\star,z^\star)
		\]
		also shows that \(x^\star\) minimizes the convex function
		\(L(\cdot,z^\star)\). Fermat's rule therefore gives
		\[
		0\in\partial_xL(x^\star,z^\star).
		\]
		Since
		\[
		L(x,z^\star)
		=
		f(x)+\inner{z^\star}{Cx+d},
		\]
		the subdifferential sum rule with an affine function gives
		\[
		\partial_xL(x^\star,z^\star)
		=
		\partial f(x^\star)+C^\top z^\star.
		\]
		Consequently,
		\[
		0\in\partial f(x^\star)+C^\top z^\star.
		\]
		The other 2 conditions are simply admissibility. Thus all four KKT conditions hold.
		
		Conversely, suppose that \(x^\star\) and \(z^\star\) satisfy the four KKT
		conditions. The stationarity condition and the affine subdifferential rule
		give
		\[
		0\in\partial_xL(x^\star,z^\star).
		\]
		Since \(L(\cdot,z^\star)\) is convex, Fermat's rule implies that
		\(x^\star\) minimizes \(L(\cdot,z^\star)\). Therefore,
		\[
		g(z^\star)=L(x^\star,z^\star).
		\]
		Using complementarity, we obtain
		\[
		g(z^\star)
		=
		f(x^\star)+\inner{z^\star}{c(x^\star)}
		=
		f(x^\star).
		\]
		
		Since \(x^\star\) is primal feasible and \(z^\star\) is dual feasible, weak
		duality gives
		\[
		g(z^\star)
		\leq
		g^\star
		\leq
		p^\star
		\leq
		f(x^\star).
		\]
		The two endpoints are equal, so equality holds throughout:
		\[
		g(z^\star)=g^\star=p^\star=f(x^\star).
		\]
		Therefore, \(x^\star\) is primal optimal and \(z^\star\) is dual optimal.
		
		It remains to verify the normal-cone formulation. Since \(K\) is a closed
		convex cone,
		\[
		(K^*)^\circ=-K.
		\]
		For \(z^\star\in K^*\), the normal cone to \(K^*\) is therefore
		\[
		N_{K^*}(z^\star)
		=
		(K^*)^\circ\cap(z^\star)^\perp
		=
		-K\cap(z^\star)^\perp.
		\]
		Consequently,
		\[
		c(x^\star)\in N_{K^*}(z^\star)
		\]
		is equivalent to
		\[
		-c(x^\star)\in K
		\qquad\text{and}\qquad
		\inner{z^\star}{c(x^\star)}=0.
		\]
		This proves the equivalent normal-cone formulation.
	\end{proof}
	
	\begin{corollary}[The differentiable and the quadratic cases]
		If $f$ is differentiable at $x^\star$ then $\partial f(x^\star)
		=\{\nabla f(x^\star)\}$ and the first condition reads
		\begin{equation}
			\nabla f(x^\star)+C^\top z^\star=0 ,
		\end{equation}
		which for the conic quadratic program is
		\begin{equation}
			Px^\star+q+C^\top z^\star=0 .
		\end{equation}
	\end{corollary}
	
	\chapter{Monotone operator and proximal representation}
	
	This chapter introduces monotone operators and proves some of their
	properties. They are useful because they supply a theory from which the
	convergence of a great many algorithms follows directly, by reading each as a
	realisation of the abstract proximal point algorithm applied to finding their
	zeros.
	
	\section{Set-valued operators}
	
	\begin{definition}[Set-valued operator]
		A \emph{set-valued operator} on $\R^N$ is a subset $T\subseteq\R^N\times\R^N$.
		One writes
		\begin{equation}
			T(u)=\{v:(u,v)\in T\},
			\qquad
			\gph T=T,
			\qquad
			T^{-1}=\{(v,u):(u,v)\in T\},
		\end{equation}
		and $\dom T=\{u:T(u)\neq\emptyset\}$.
	\end{definition}
	
	For such an operator, monotonicity is the natural generalisation of a
	nondecreasing map on the line, and of the gradient of a convex function.
	\begin{definition}[Monotone operator]
		$T$ is \emph{monotone} if
		\begin{equation}
			\inner{v-v'}{u-u'}\geq0
			\qquad\text{whenever }(u,v)\in T\text{ and }(u',v')\in T .
		\end{equation}
	\end{definition}
	
	\begin{definition}[Order on operators]
		$T_1\subseteq T_2$ means containment of the graphs.
	\end{definition}
	
	\begin{definition}[Maximal monotone operator]
		A monotone $T$ is \emph{maximal monotone} if the only monotone $T'$ with
		$T\subseteq T'$ is $T$ itself.
	\end{definition}
	
	\begin{proposition}[Closed graph]
		Let \(T:\R^n\rightrightarrows\R^n\) be maximal monotone. Then
		\(\operatorname{gra}T\) is closed.
	\end{proposition}
	
	\begin{proof}
		Let \((u_k,v_k)\) be a sequence in \(\operatorname{gra}T\) such that
		\[
		(u_k,v_k)\longrightarrow(u,v).
		\]
		For every \((u',v')\in\operatorname{gra}T\), the monotonicity of \(T\) gives
		\[
		\inner{v_k-v'}{u_k-u'}\geq0.
		\]
		Passing to the limit as \(k\to\infty\), we obtain
		\[
		\inner{v-v'}{u-u'}\geq0
		\qquad
		\text{for every }(u',v')\in\operatorname{gra}T.
		\]
		Thus the pair \((u,v)\) is monotonically related to every point of
		\(\operatorname{gra}T\). Consequently,
		\[
		\operatorname{gra}T\cup\{(u,v)\}
		\]
		is the graph of a monotone extension of \(T\). Since \(T\) is maximal
		monotone, it admits no proper monotone extension. Therefore,
		\[
		(u,v)\in\operatorname{gra}T.
		\]
		Hence \(\operatorname{gra}T\) is closed.
	\end{proof}
	\begin{proposition}[Affine operators with positive semidefinite symmetric part]
		Let \(Q\in\R^{N\times N}\) satisfy
		\[
		\inner{Qu}{u}\geq0
		\qquad
		\text{for every }u\in\R^N,
		\]
		and let \(b\in\R^N\). Then the affine operator
		\[
		T(u)=Qu+b
		\]
		is maximal monotone and has domain \(\R^N\).
	\end{proposition}
	
	\begin{proof}
		For any \(u,u'\in\R^N\),
		\begin{align}
			\inner{T(u)-T(u')}{u-u'}
			&=
			\inner{Q(u-u')}{u-u'}\\
			&\geq0.
		\end{align}
		Thus \(T\) is monotone.
		
		To prove maximality, suppose that \((\bar u,\bar v)\) is monotonically
		related to \(\operatorname{gra}T\). This means that
		\[
		\inner{\bar v-T(u)}{\bar u-u}\geq0
		\qquad
		\text{for every }u\in\R^N.
		\]
		Fix \(h\in\R^N\) and \(t>0\), and choose
		\[
		u=\bar u-th.
		\]
		The preceding inequality becomes
		\[
		\inner{\bar v-Q(\bar u-th)-b}{th}\geq0,
		\]
		or equivalently,
		\[
		\inner{\bar v-Q\bar u-b+tQh}{h}\geq0.
		\]
		Letting \(t\downarrow0\), we obtain
		\[
		\inner{\bar v-Q\bar u-b}{h}\geq0
		\qquad
		\text{for every }h\in\R^N.
		\]
		Applying the same inequality to \(-h\) shows that
		\[
		\inner{\bar v-Q\bar u-b}{h}=0
		\qquad
		\text{for every }h\in\R^N.
		\]
		Therefore,
		\[
		\bar v=Q\bar u+b,
		\]
		so \((\bar u,\bar v)\in\operatorname{gra}T\). Hence no point can be added to
		the graph while preserving monotonicity, and \(T\) is maximal monotone.
	\end{proof}
	\begin{proposition}[Rockafellar, \cite{Rockafellar1970}]
		If $\phi:\R^N\to\R\cup\{+\infty\}$ is closed proper convex then
		$\partial\phi$ is maximal monotone.
	\end{proposition}
	
	\begin{proposition}[Sum of maximal monotone operators, \cite{Rockafellar1970}]
		If $T_1$ and $T_2$ are maximal monotone and
		$\relint\dom T_1\cap\relint\dom T_2\neq\emptyset$, then $T_1+T_2$ is maximal
		monotone. In particular this holds when one of them has domain $\R^N$.
	\end{proposition}
	\section{The resolvent}
	
	\begin{definition}[Metric resolvent]
		Let \(T:\R^N\rightrightarrows\R^N\) be set-valued and let
		\(\mathcal M\succ0\). The \emph{metric resolvent} of \(T\) is
		\begin{equation}
			J_T^{\mathcal M}
			=
			(\mathcal M+T)^{-1}\mathcal M.
		\end{equation}
		Equivalently,
		\begin{equation}
			p\in J_T^{\mathcal M}(v)
			\quad\Longleftrightarrow\quad
			\mathcal M(v-p)\in T(p).
		\end{equation}
	\end{definition}
	
	\begin{proposition}[Uniqueness on the domain]
		If \(T\) is monotone and \(\mathcal M\succ0\), then
		\(J_T^{\mathcal M}\) is single-valued on its domain.
	\end{proposition}
	
	\begin{proof}
		Suppose that \(p,p'\in J_T^{\mathcal M}(v)\). Then
		\[
		\mathcal M(v-p)\in T(p),
		\qquad
		\mathcal M(v-p')\in T(p').
		\]
		The monotonicity of \(T\) gives
		\[
		\inner{
			\mathcal M(v-p)-\mathcal M(v-p')
		}{p-p'}
		\geq0.
		\]
		The left-hand side is
		\[
		-\inner{\mathcal M(p-p')}{p-p'}
		=
		-\norm{p-p'}_{\mathcal M}^2.
		\]
		Hence
		\[
		\norm{p-p'}_{\mathcal M}^2\leq0,
		\]
		and therefore \(p=p'\).
	\end{proof}
	
	\begin{proposition}[Firm nonexpansiveness on the domain]
		Let \(T\) be monotone and let \(\mathcal M\succ0\). For every
		\(v,v'\in\dom J_T^{\mathcal M}\),
		\begin{equation}
			\norm{
				J_T^{\mathcal M}(v)-J_T^{\mathcal M}(v')
			}_{\mathcal M}^2
			\leq
			\inner{
				J_T^{\mathcal M}(v)-J_T^{\mathcal M}(v')
			}{v-v'}_{\mathcal M}.
			\label{eq:resolvent-firm}
		\end{equation}
		In particular,
		\begin{equation}
			\norm{
				J_T^{\mathcal M}(v)-J_T^{\mathcal M}(v')
			}_{\mathcal M}
			\leq
			\norm{v-v'}_{\mathcal M}.
		\end{equation}
	\end{proposition}
	
	\begin{proof}
		Set
		\[
		p=J_T^{\mathcal M}(v),
		\qquad
		p'=J_T^{\mathcal M}(v').
		\]
		By the definition of the resolvent,
		\[
		\mathcal M(v-p)\in T(p),
		\qquad
		\mathcal M(v'-p')\in T(p').
		\]
		Monotonicity gives
		\[
		\inner{
			\mathcal M(v-p)-\mathcal M(v'-p')
		}{p-p'}
		\geq0.
		\]
		Rearranging, we obtain
		\[
		\norm{p-p'}_{\mathcal M}^2
		\leq
		\inner{p-p'}{v-v'}_{\mathcal M},
		\]
		which proves firm nonexpansiveness. The Cauchy-Schwarz inequality then gives
		\[
		\norm{p-p'}_{\mathcal M}^2
		\leq
		\norm{p-p'}_{\mathcal M}\norm{v-v'}_{\mathcal M},
		\]
		and hence
		\[
		\norm{p-p'}_{\mathcal M}
		\leq
		\norm{v-v'}_{\mathcal M}.
		\]
	\end{proof}
	
	\begin{proposition}[Minty's theorem, \cite{Minty1962}]
		A monotone operator \(T:\R^N\rightrightarrows\R^N\) is maximal monotone if
		and only if
		\[
		\range(I+T)=\R^N.
		\]
	\end{proposition}
	
	\begin{proposition}[Minty's theorem in a metric]
		Let \(T:\R^N\rightrightarrows\R^N\) be monotone. Then \(T\) is maximal
		monotone if and only if
		\[
		\range(\mathcal M+T)=\R^N
		\]
		for one, and consequently for every, \(\mathcal M\succ0\).
	\end{proposition}
	
	\begin{proof}
		Let \(R=\mathcal M^{1/2}\), and define
		\[
		\widetilde T(y)
		=
		R^{-1}T(R^{-1}y).
		\]
		If
		\[
		u\in T(x),
		\qquad
		u'\in T(x'),
		\]
		and
		\[
		y=Rx,
		\qquad
		y'=Rx',
		\qquad
		w=R^{-1}u,
		\qquad
		w'=R^{-1}u',
		\]
		then
		\[
		\inner{w-w'}{y-y'}
		=
		\inner{u-u'}{x-x'}.
		\]
		Thus the change of variables preserves monotonicity and maximal
		monotonicity. Consequently, \(T\) is maximal monotone if and only if
		\(\widetilde T\) is maximal monotone.
		
		Moreover,
		\[
		u\in\mathcal Mx+T(x)
		\]
		is equivalent, after multiplication by \(R^{-1}\), to
		\[
		R^{-1}u\in y+\widetilde T(y).
		\]
		Since \(R^{-1}\) is bijective,
		\[
		\range(\mathcal M+T)=\R^N
		\quad\Longleftrightarrow\quad
		\range(I+\widetilde T)=\R^N.
		\]
		The result follows from Minty's theorem applied to \(\widetilde T\).
	\end{proof}
	
	\begin{proposition}[Resolvent of a maximal monotone operator]
		Let \(T\) be maximal monotone and let \(\mathcal M\succ0\). Then
		\(J_T^{\mathcal M}\) is single-valued and defined on all of \(\R^N\).
		Moreover,
		\begin{equation}
			\Fix(J_T^{\mathcal M})=T^{-1}(0).
		\end{equation}
	\end{proposition}
	
	\begin{proof}
		Single-valuedness follows from monotonicity. The metric form of Minty's
		theorem gives
		\[
		\range(\mathcal M+T)=\R^N.
		\]
		Since \(\mathcal M\) is bijective, \(J_T^{\mathcal M}\) is defined on all
		of \(\R^N\).
		
		Finally,
		\begin{align}
			p=J_T^{\mathcal M}(p)
			&\iff
			\mathcal M(p-p)\in T(p)\\
			&\iff
			0\in T(p).
		\end{align}
		Thus
		\[
		\Fix(J_T^{\mathcal M})=T^{-1}(0).
		\]
	\end{proof}
	
	\begin{proposition}[One-step estimate]
		Let \(T\) be maximal monotone, let \(\mathcal M\succ0\), and let
		\(v^\star\in T^{-1}(0)\). If
		\[
		p=J_T^{\mathcal M}(v),
		\]
		then
		\begin{equation}
			\norm{p-v^\star}_{\mathcal M}^2
			\leq
			\norm{v-v^\star}_{\mathcal M}^2
			-
			\norm{p-v}_{\mathcal M}^2.
			\label{eq:resolvent-one-step}
		\end{equation}
	\end{proposition}
	
	\begin{proof}
		Since \(v^\star\in T^{-1}(0)\), it is a fixed point of the resolvent:
		\[
		J_T^{\mathcal M}(v^\star)=v^\star.
		\]
		Firm nonexpansiveness therefore gives
		\[
		\norm{p-v^\star}_{\mathcal M}^2
		\leq
		\inner{p-v^\star}{v-v^\star}_{\mathcal M}.
		\]
		Using
		\[
		v-v^\star=(v-p)+(p-v^\star),
		\]
		we obtain
		\[
		\inner{p-v^\star}{v-p}_{\mathcal M}\geq0.
		\]
		Expanding the squared distance gives
		\begin{align}
			\norm{v-v^\star}_{\mathcal M}^2
			&=
			\norm{v-p}_{\mathcal M}^2
			+
			\norm{p-v^\star}_{\mathcal M}^2
			+
			2\inner{v-p}{p-v^\star}_{\mathcal M}\\
			&\geq
			\norm{v-p}_{\mathcal M}^2
			+
			\norm{p-v^\star}_{\mathcal M}^2.
		\end{align}
		Rearranging proves the result.
	\end{proof}
	
	\section{Variable metric proximal point algorithm}
	
	\begin{definition}[Variable-metric quasi-Fejér sequence]
		Let \(S\subseteq\R^N\) be nonempty, and let
		\((\mathcal M_k)\) be a sequence of positive definite matrices. A sequence
		\((v^k)\) is \emph{variable-metric quasi-Fejér monotone} with respect to
		\(S\) if, for every \(v^\star\in S\), there exists a summable sequence
		\((\delta_k)\) of nonnegative numbers such that
		\begin{equation}
			\norm{v^{k+1}-v^\star}_{\mathcal M_{k+1}}^2
			\leq
			(1+\eta_k)
			\norm{v^k-v^\star}_{\mathcal M_k}^2
			+
			\delta_k,
			\label{eq:variable-quasi-fejer}
		\end{equation}
		where \(\eta_k\geq0\) and
		\[
		\sum_{k=0}^{\infty}\eta_k<\infty.
		\]
		When \(\mathcal M_k=I\), \(\eta_k=0\), and \(\delta_k=0\), this reduces to
		the usual notion of Fejér monotonicity.
	\end{definition}
	
	\begin{lemma}[Variable-metric quasi-Fejér convergence]
		Let \(S\subseteq\R^N\) be nonempty, and suppose that
		\[
		\mu I\preceq\mathcal M_k\preceq\nu I
		\]
		for some \(0<\mu\leq\nu<\infty\). Suppose that, for every
		\(v^\star\in S\),
		\begin{equation}
			d_{k+1}(v^\star)
			\leq
			(1+\eta_k)d_k(v^\star)-s_k+\delta_k,
			\label{eq:quasi-fejer-descent}
		\end{equation}
		where
		\[
		d_k(v^\star)
		=
		\norm{v^k-v^\star}_{\mathcal M_k}^2,
		\]
		the sequences \((\eta_k)\), \((\delta_k)\), and \((s_k)\) are nonnegative,
		and
		\[
		\sum_{k=0}^{\infty}\eta_k<\infty,
		\qquad
		\sum_{k=0}^{\infty}\delta_k<\infty.
		\]
		Then:
		
		\begin{enumerate}
			\item \((v^k)\) is bounded;
			\item \(d_k(v^\star)\) converges for every \(v^\star\in S\);
			\item \(\sum_{k=0}^{\infty}s_k<\infty\);
			\item if every cluster point of \((v^k)\) belongs to \(S\), then
			\((v^k)\) converges to a point of \(S\).
		\end{enumerate}
	\end{lemma}
	
	\begin{proof}
		Define
		\[
		\Lambda_0=1,
		\qquad
		\Lambda_k=\prod_{j=0}^{k-1}(1+\eta_j).
		\]
		Since
		\[
		\log\Lambda_k
		\leq
		\sum_{j=0}^{k-1}\eta_j,
		\]
		the sequence \((\Lambda_k)\) converges to a finite positive limit.
		
		Fix \(v^\star\in S\), write \(d_k=d_k(v^\star)\), and divide
		\eqref{eq:quasi-fejer-descent} by
		\(\Lambda_{k+1}=(1+\eta_k)\Lambda_k\). Dropping the nonnegative term \(s_k\)
		gives
		\[
		\frac{d_{k+1}}{\Lambda_{k+1}}
		\leq
		\frac{d_k}{\Lambda_k}
		+
		\frac{\delta_k}{\Lambda_{k+1}}.
		\]
		Set
		\[
		b_k=\frac{d_k}{\Lambda_k},
		\qquad
		r_k=\frac{\delta_k}{\Lambda_{k+1}}.
		\]
		Then
		\[
		b_{k+1}\leq b_k+r_k,
		\qquad
		\sum_{k=0}^{\infty}r_k<\infty.
		\]
		The sequence
		\[
		\widetilde b_k
		=
		b_k+\sum_{j=k}^{\infty}r_j
		\]
		is nonnegative and nonincreasing. It therefore converges. Since the tail
		sum tends to zero, \((b_k)\) converges. As \((\Lambda_k)\) also converges,
		\((d_k)\) converges.
		
		In particular, \((d_k)\) is bounded. Since
		\[
		\mu\norm{v^k-v^\star}^2\leq d_k,
		\]
		the sequence \((v^k)\) is bounded.
		
		Rearranging \eqref{eq:quasi-fejer-descent} yields
		\[
		s_k
		\leq
		d_k-d_{k+1}+\eta_kd_k+\delta_k.
		\]
		Summing from \(k=0\) to \(n\), we obtain
		\[
		\sum_{k=0}^{n}s_k
		\leq
		d_0-d_{n+1}
		+
		\sum_{k=0}^{n}\eta_kd_k
		+
		\sum_{k=0}^{n}\delta_k.
		\]
		The right-hand side is bounded uniformly in \(n\), so
		\[
		\sum_{k=0}^{\infty}s_k<\infty.
		\]
		
		Finally, suppose that every cluster point of \((v^k)\) belongs to \(S\).
		Since \((v^k)\) is bounded, it has a cluster point \(\bar v\in S\). Choose
		a subsequence such that
		\[
		v^{k_i}\longrightarrow\bar v.
		\]
		Applying the preceding argument with \(v^\star=\bar v\) shows that
		\[
		d_k(\bar v)
		=
		\norm{v^k-\bar v}_{\mathcal M_k}^2
		\]
		has a limit. Along the subsequence,
		\[
		d_{k_i}(\bar v)
		\leq
		\nu\norm{v^{k_i}-\bar v}^2
		\longrightarrow0.
		\]
		Hence \(d_k(\bar v)\to0\). Since
		\[
		\mu\norm{v^k-\bar v}^2
		\leq
		d_k(\bar v),
		\]
		we conclude that \(v^k\to\bar v\in S\).
	\end{proof}
	
	\begin{proposition}[Convergence in the exact case]
		Let \(T\) be maximal monotone with \(T^{-1}(0)\neq\emptyset\). Suppose that
		\((\mathcal M_k)\) satisfies
		\begin{equation}
			\mu I\preceq\mathcal M_k\preceq\nu I,
			\qquad
			\mathcal M_{k+1}\preceq(1+\eta_k)\mathcal M_k,
			\label{eq:variable-metric-assumptions}
		\end{equation}
		where \(0<\mu\leq\nu<\infty\), \(\eta_k\geq0\), and
		\[
		\sum_{k=0}^{\infty}\eta_k<\infty.
		\]
		Then the sequence
		\begin{equation}
			v^{k+1}=J_T^{\mathcal M_k}(v^k)
		\end{equation}
		converges, for every initial point \(v^0\), to an element of \(T^{-1}(0)\).
	\end{proposition}
	
	\begin{proof}
		Set
		\[
		S=T^{-1}(0),
		\]
		and fix \(v^\star\in S\). The one-step estimate gives
		\[
		\norm{v^{k+1}-v^\star}_{\mathcal M_k}^2
		\leq
		\norm{v^k-v^\star}_{\mathcal M_k}^2
		-
		\norm{v^{k+1}-v^k}_{\mathcal M_k}^2.
		\]
		Using the metric variation condition, we obtain
		\[
		\norm{v^{k+1}-v^\star}_{\mathcal M_{k+1}}^2
		\leq
		(1+\eta_k)
		\norm{v^k-v^\star}_{\mathcal M_k}^2
		-
		\norm{v^{k+1}-v^k}_{\mathcal M_k}^2.
		\]
		Thus the quasi-Fejér lemma applies with
		\[
		s_k=\norm{v^{k+1}-v^k}_{\mathcal M_k}^2,
		\qquad
		\delta_k=0.
		\]
		It follows that \((v^k)\) is bounded and
		\[
		\sum_{k=0}^{\infty}
		\norm{v^{k+1}-v^k}_{\mathcal M_k}^2
		<\infty.
		\]
		Since \(\mathcal M_k\succeq\mu I\),
		\[
		v^{k+1}-v^k\longrightarrow0.
		\]
		
		Let \(\bar v\) be a cluster point, and choose a subsequence
		\[
		v^{k_i}\longrightarrow\bar v.
		\]
		The resolvent relation gives
		\[
		u^{k_i}
		:=
		\mathcal M_{k_i-1}
		\bigl(v^{k_i-1}-v^{k_i}\bigr)
		\in T(v^{k_i}).
		\]
		Moreover,
		\[
		\norm{u^{k_i}}
		\leq
		\nu\norm{v^{k_i-1}-v^{k_i}}
		\longrightarrow0.
		\]
		Since the graph of \(T\) is closed,
		\[
		0\in T(\bar v).
		\]
		Thus every cluster point belongs to \(S\). The quasi-Fejér lemma now gives
		\[
		v^k\longrightarrow\bar v\in T^{-1}(0).
		\]
	\end{proof}
	
	\begin{proposition}[Convergence in the inexact case, \cite{Rockafellar1976a}]
		Under the assumptions of the preceding proposition, suppose instead that
		\(v^{k+1}\) satisfies
		\begin{equation}
			\norm{
				v^{k+1}-J_T^{\mathcal M_k}(v^k)
			}_{\mathcal M_k}
			\leq\varepsilon_k,
			\qquad
			\sum_{k=0}^{\infty}\varepsilon_k<\infty.
			\label{eq:inexact-resolvent-condition}
		\end{equation}
		Then \((v^k)\) converges to a point of \(T^{-1}(0)\).
	\end{proposition}
	
	\begin{proof}
		Set
		\[
		\widehat v^{k+1}=J_T^{\mathcal M_k}(v^k),
		\qquad
		e_k=v^{k+1}-\widehat v^{k+1}.
		\]
		Then
		\[
		\norm{e_k}_{\mathcal M_k}\leq\varepsilon_k.
		\]
		Fix \(v^\star\in T^{-1}(0)\), and write
		\[
		a_k=\norm{v^k-v^\star}_{\mathcal M_k},
		\qquad
		t_k=\norm{\widehat v^{k+1}-v^k}_{\mathcal M_k}.
		\]
		The one-step estimate gives
		\[
		\norm{\widehat v^{k+1}-v^\star}_{\mathcal M_k}^2
		\leq
		a_k^2-t_k^2.
		\]
		In particular,
		\[
		\norm{\widehat v^{k+1}-v^\star}_{\mathcal M_k}
		\leq a_k.
		\]
		The triangle inequality and the metric variation condition imply
		\[
		a_{k+1}
		\leq
		(1+\eta_k)^{1/2}(a_k+\varepsilon_k).
		\]
		Iterating this inequality and using
		\[
		\sum_k\eta_k<\infty,
		\qquad
		\sum_k\varepsilon_k<\infty,
		\]
		shows that \((a_k)\) is bounded. Choose \(A>0\) such that
		\[
		a_k\leq A
		\qquad
		\text{for every }k.
		\]
		
		We also have
		\begin{align}
			\norm{v^{k+1}-v^\star}_{\mathcal M_k}^2
			&=
			\norm{\widehat v^{k+1}-v^\star+e_k}_{\mathcal M_k}^2\\
			&\leq
			\norm{\widehat v^{k+1}-v^\star}_{\mathcal M_k}^2
			+
			2\norm{\widehat v^{k+1}-v^\star}_{\mathcal M_k}
			\norm{e_k}_{\mathcal M_k}
			+
			\norm{e_k}_{\mathcal M_k}^2\\
			&\leq
			a_k^2-t_k^2+2A\varepsilon_k+\varepsilon_k^2.
		\end{align}
		Using the metric variation condition gives
		\[
		a_{k+1}^2
		\leq
		(1+\eta_k)a_k^2
		-
		(1+\eta_k)t_k^2
		+
		(1+\eta_k)
		\bigl(2A\varepsilon_k+\varepsilon_k^2\bigr).
		\]
		The quasi-Fejér lemma applies with
		\[
		s_k=(1+\eta_k)t_k^2
		\]
		and
		\[
		\delta_k
		=
		(1+\eta_k)
		\bigl(2A\varepsilon_k+\varepsilon_k^2\bigr).
		\]
		The sequence \((\delta_k)\) is summable. Consequently,
		\[
		\sum_{k=0}^{\infty}t_k^2<\infty,
		\]
		and hence
		\[
		\widehat v^{k+1}-v^k\longrightarrow0.
		\]
		Moreover, \(\varepsilon_k\to0\) and
		\(\mathcal M_k\succeq\mu I\) imply
		\[
		\norm{e_k}
		\leq
		\frac{\varepsilon_k}{\sqrt{\mu}}
		\longrightarrow0.
		\]
		Thus
		\[
		v^{k+1}-v^k
		=
		\widehat v^{k+1}-v^k+e_k
		\longrightarrow0.
		\]
		
		The resolvent relation gives
		\[
		u^{k+1}
		=
		\mathcal M_k
		\bigl(v^k-\widehat v^{k+1}\bigr)
		\in T(\widehat v^{k+1}).
		\]
		Since \(\mathcal M_k\preceq\nu I\),
		\[
		\norm{u^{k+1}}
		\leq
		\nu\norm{v^k-\widehat v^{k+1}}
		\longrightarrow0.
		\]
		
		Let \(\bar v\) be a cluster point of \((v^k)\), and choose a subsequence
		\(v^{k_i}\to\bar v\). Since
		\[
		\widehat v^{k_i+1}-v^{k_i}\longrightarrow0,
		\]
		we also have
		\[
		\widehat v^{k_i+1}\longrightarrow\bar v.
		\]
		The graph of \(T\) is closed, so
		\[
		0\in T(\bar v).
		\]
		Thus every cluster point of \((v^k)\) belongs to \(T^{-1}(0)\). The
		quasi-Fejér lemma now implies that the entire sequence converges to a point
		of \(T^{-1}(0)\).
	\end{proof}
	
	\begin{algorithm}[H]
		\caption{Variable metric proximal point algorithm}
		\begin{algorithmic}[1]
			\Require $T$ maximal monotone, $v^0\in\R^N$, metrics $(\mathcal{M}_k)$,
			tolerances $(\varepsilon_k)$ summable
			\Ensure a point of $T^{-1}(0)$
			\For{$k=0,1,2,\dots$}
			\State compute $v^{k+1}$ with
			$\norm{v^{k+1}-J^{\mathcal{M}_k}_T(v^k)}_{\mathcal{M}_k}\leq\varepsilon_k$,
			that is solve $\mathcal{M}_k(v^k-v^{k+1})\in T(v^{k+1})$ approximately
			\State choose $\mathcal{M}_{k+1}$ satisfying the bounds of the convergence
			hypotheses
			\EndFor
		\end{algorithmic}
	\end{algorithm}
	
	
	\section{Proximal operator and Moreau envelope}
	
	\begin{definition}[Proximal operator and Moreau envelope]
		For $\phi$ closed proper convex and $\mathcal{M}\succ0$,
		\begin{equation}
			\prox^{\mathcal{M}}_{\phi}(v)
			=\argmin_{w}\Bigl\{\phi(w)+\tfrac12\norm{w-v}^2_{\mathcal{M}}\Bigr\},
			\qquad
			e^{\mathcal{M}}_{\phi}(v)
			=\min_{w}\Bigl\{\phi(w)+\tfrac12\norm{w-v}^2_{\mathcal{M}}\Bigr\} .
		\end{equation}
	\end{definition}
	
	The Moreau envelope selects, for each $v$, the point striking the best trade
	off between minimising $\phi$ and staying close to $v$. One advantage is the
	regularity it gains, as the properties below show; and as $\mathcal{M}$ grows
	the envelope increases to $\phi$ itself.
	
	\begin{proposition}[Well-posedness]
		Let \(\phi:\R^N\to\R\cup\{+\infty\}\) be closed, proper, and convex, and let
		\(\mathcal M\succ0\). Then, for every \(v\in\R^N\), the problem
		\[
		\min_{w\in\R^N}
		\left\{
		\phi(w)+\frac12\norm{w-v}_{\mathcal M}^2
		\right\}
		\]
		has a unique solution. Consequently,
		\(\prox_{\phi}^{\mathcal M}\) is single-valued and defined on all of
		\(\R^N\), and \(e_{\phi}^{\mathcal M}\) is finite everywhere.
	\end{proposition}
	
	\begin{proof}
		For fixed \(v\), the function
		\[
		w\longmapsto
		\phi(w)+\frac12\norm{w-v}_{\mathcal M}^2
		\]
		is closed and proper. Since \(\mathcal M\succ0\), its quadratic term is
		strongly convex and coercive. The entire function is therefore coercive and
		strongly convex. Coercivity and closedness give existence of a minimizer,
		while strong convexity gives uniqueness.
	\end{proof}
	
	\begin{proposition}[Resolvent form]
		For \(v,p\in\R^N\),
		\[
		p=\prox_{\phi}^{\mathcal M}(v)
		\]
		if and only if
		\[
		\mathcal M(v-p)\in\partial\phi(p).
		\]
		Equivalently,
		\[
		\prox_{\phi}^{\mathcal M}
		=
		\bigl(I+\mathcal M^{-1}\partial\phi\bigr)^{-1}.
		\]
	\end{proposition}
	
	\begin{proof}
		By Fermat's rule,
		\begin{align}
			p=\prox_{\phi}^{\mathcal M}(v)
			&\iff
			0\in
			\partial\phi(p)+\mathcal M(p-v)\\
			&\iff
			\mathcal M(v-p)\in\partial\phi(p).
		\end{align}
		Rearranging the last inclusion gives
		\[
		v\in p+\mathcal M^{-1}\partial\phi(p),
		\]
		which is precisely the resolvent formulation.
	\end{proof}
	As a consequence of the general monotone operator case:
	\begin{proposition}[Firm nonexpansiveness]
		The mapping \(\prox_{\phi}^{\mathcal M}\) is firmly nonexpansive in the
		metric induced by \(\mathcal M\). More precisely, for every \(v,v'\in\R^N\),
		\begin{equation}
			\norm{
				\prox_{\phi}^{\mathcal M}(v)
				-
				\prox_{\phi}^{\mathcal M}(v')
			}_{\mathcal M}^2
			\leq
			\inner{
				\prox_{\phi}^{\mathcal M}(v)
				-
				\prox_{\phi}^{\mathcal M}(v')
			}{v-v'}_{\mathcal M}.
		\end{equation}
		In particular,
		\[
		\norm{
			\prox_{\phi}^{\mathcal M}(v)
			-
			\prox_{\phi}^{\mathcal M}(v')
		}_{\mathcal M}
		\leq
		\norm{v-v'}_{\mathcal M}.
		\]
	\end{proposition}
	
	
	\begin{proposition}[Regularity of the envelope]
		The Moreau envelope \(e_{\phi}^{\mathcal M}\) is convex, finite everywhere,
		and continuously differentiable. Its gradient is
		\begin{equation}
			\nabla e_{\phi}^{\mathcal M}(v)
			=
			\mathcal M
			\bigl(
			v-\prox_{\phi}^{\mathcal M}(v)
			\bigr).
		\end{equation}
	\end{proposition}
	
	\begin{proof}
		Finiteness follows from the well-posedness proposition. The function
		\[
		(v,w)\longmapsto
		\phi(w)+\frac12\norm{w-v}_{\mathcal M}^2
		\]
		is jointly convex in \((v,w)\). Its partial minimization over \(w\) is
		therefore convex, so \(e_{\phi}^{\mathcal M}\) is convex.
		
		Fix \(v\), and write
		\[
		p=\prox_{\phi}^{\mathcal M}(v),
		\qquad
		r=\mathcal M(v-p).
		\]
		For \(h\in\R^N\), evaluating the minimization problem defining
		\(e_{\phi}^{\mathcal M}(v+h)\) at \(p\) gives
		\begin{align}
			e_{\phi}^{\mathcal M}(v+h)
			&\leq
			\phi(p)+\frac12\norm{p-v-h}_{\mathcal M}^2\\
			&=
			e_{\phi}^{\mathcal M}(v)
			+\inner{r}{h}
			+\frac12\norm{h}_{\mathcal M}^2.
		\end{align}
		Thus
		\begin{equation}
			e_{\phi}^{\mathcal M}(v+h)-e_{\phi}^{\mathcal M}(v)
			\leq
			\inner{r}{h}
			+\frac12\norm{h}_{\mathcal M}^2.
			\label{eq:envelope-upper-bound}
		\end{equation}
		
		Now set
		\[
		p_h=\prox_{\phi}^{\mathcal M}(v+h),
		\qquad
		r_h=\mathcal M(v+h-p_h).
		\]
		Applying \eqref{eq:envelope-upper-bound} with \(v+h\) in place of \(v\) and
		with increment \(-h\) gives
		\[
		e_{\phi}^{\mathcal M}(v)-e_{\phi}^{\mathcal M}(v+h)
		\leq
		-\inner{r_h}{h}
		+\frac12\norm{h}_{\mathcal M}^2.
		\]
		Equivalently,
		\[
		e_{\phi}^{\mathcal M}(v+h)-e_{\phi}^{\mathcal M}(v)
		\geq
		\inner{r_h}{h}
		-\frac12\norm{h}_{\mathcal M}^2.
		\]
		
		The proximal mapping is Lipschitz, so \(p_h\to p\) as \(h\to0\). Hence
		\[
		r_h=\mathcal M(v+h-p_h)\longrightarrow\mathcal M(v-p)=r.
		\]
		Combining the upper and lower estimates yields
		\[
		e_{\phi}^{\mathcal M}(v+h)
		=
		e_{\phi}^{\mathcal M}(v)
		+
		\inner{r}{h}
		+
		o(\norm{h}).
		\]
		Therefore,
		\[
		\nabla e_{\phi}^{\mathcal M}(v)
		=
		r
		=
		\mathcal M
		\bigl(
		v-\prox_{\phi}^{\mathcal M}(v)
		\bigr).
		\]
		
		Finally, the proximal mapping is continuous, so the displayed gradient is
		continuous. Thus \(e_{\phi}^{\mathcal M}\) is continuously differentiable
		on \(\R^N\).
	\end{proof}
	
	\begin{proposition}[Envelope of an indicator]
		Let \(C\subseteq\R^N\) be nonempty, closed, and convex, and let
		\(\mathcal M\succ0\). Then
		\begin{equation}
			\prox_{\delta_C}^{\mathcal M}(v)
			=
			\proj{v}{C}^{\mathcal M},
			\qquad
			e_{\delta_C}^{\mathcal M}(v)
			=
			\frac12\dist_{\mathcal M}(v,C)^2,
		\end{equation}
		where \(\proj{\cdot}{C}^{\mathcal M}\) denotes the projection onto \(C\) in
		the metric induced by \(\mathcal M\). In particular, if \(C=K^*\) and
		\(\mathcal M\in\mathcal A(K)\), then
		\[
		\proj{v}{K^*}^{\mathcal M}
		=
		\proj{v}{K^*},
		\]
		so
		\[
		\prox_{\delta_{K^*}}^{\mathcal M}(v)
		=
		\proj{v}{K^*}.
		\]
	\end{proposition}
	
	\begin{proof}
		By the definition of the proximal mapping,
		\[
		\prox_{\delta_C}^{\mathcal M}(v)
		=
		\argmin_{w\in\R^N}
		\left\{
		\delta_C(w)+\frac12\norm{w-v}_{\mathcal M}^2
		\right\}.
		\]
		Since \(\delta_C(w)=0\) for \(w\in C\) and \(+\infty\) otherwise, this reduces
		to
		\[
		\prox_{\delta_C}^{\mathcal M}(v)
		=
		\argmin_{w\in C}
		\frac12\norm{w-v}_{\mathcal M}^2
		=
		\proj{v}{C}^{\mathcal M}.
		\]
		
		Similarly, the definition of the Moreau envelope gives
		\begin{align}
			e_{\delta_C}^{\mathcal M}(v)
			&=
			\inf_{w\in\R^N}
			\left\{
			\delta_C(w)+\frac12\norm{w-v}_{\mathcal M}^2
			\right\}\\
			&=
			\inf_{w\in C}\frac12\norm{w-v}_{\mathcal M}^2\\
			&=
			\frac12\dist_{\mathcal M}(v,C)^2.
		\end{align}
		
		Finally, if \(C=K^*\) and \(\mathcal M\in\mathcal A(K)\), admissibility
		ensures that the metric projection onto \(K^*\) agrees with the Euclidean
		projection. Therefore,
		\[
		\proj{v}{K^*}^{\mathcal M}
		=
		\proj{v}{K^*}.
		\]
	\end{proof}
	
	\chapter{Newton method}
	
	\section{Classic Newton method}
	
	\begin{definition}[Equality constrained minimisation]
		Given $f:\R^n\to\R$ twice continuously differentiable and convex,
		$A\in\R^{q\times n}$ and $b\in\R^{q}$, the problem is
		\begin{equation}
			\min_{x\in\R^n}\ f(x)
			\qquad\text{subject to}\qquad
			Ax=b .
		\end{equation}
	\end{definition}
	
	\begin{definition}[Zero finding]
		Given $\mathcal{R}:\R^N\to\R^N$, the problem is to find $v$ with
		$\mathcal{R}(v)=0$.
	\end{definition}
	
	\begin{algorithm}[H]
		\caption{Newton's method for $\mathcal{R}(v)=0$}
		\begin{algorithmic}[1]
			\Require $\mathcal{R}$ differentiable, $v^0\in\R^N$, tolerance $\varepsilon>0$
			\Ensure $v$ with $\norm{\mathcal{R}(v)}\leq\varepsilon$
			\For{$k=0,1,2,\dots$}
			\If{$\norm{\mathcal{R}(v^k)}\leq\varepsilon$} \State \Return $v^k$ \EndIf
			\State solve $\mathcal{R}'(v^k)\,\Delta v^k=-\mathcal{R}(v^k)$
			\State $v^{k+1}\gets v^k+\Delta v^k$
			\EndFor
		\end{algorithmic}
	\end{algorithm}
	
	\begin{algorithm}[H]
		\caption{Newton's method on the optimality conditions of the equality
			constrained problem}
		\begin{algorithmic}[1]
			\Require $f$, $A$, $b$, a starting pair $(x^0,y^0)$, tolerance $\varepsilon>0$
			\Ensure $(x,y)$ with $\norm{\mathcal{R}(x,y)}\leq\varepsilon$, where
			$\mathcal{R}(x,y)=\bigl(\nabla f(x)+A^\top y,\ Ax-b\bigr)$
			\For{$k=0,1,2,\dots$}
			\If{$\norm{\mathcal{R}(x^k,y^k)}\leq\varepsilon$}
			\State \Return $(x^k,y^k)$ \EndIf
			\State solve
			$\begin{pmatrix}\nabla^2f(x^k)&A^\top\\ A&0\end{pmatrix}
			\begin{pmatrix}\Delta x^k\\ \Delta y^k\end{pmatrix}
			=-\begin{pmatrix}\nabla f(x^k)+A^\top y^k\\ Ax^k-b\end{pmatrix}$
			\State $(x^{k+1},y^{k+1})\gets(x^k+\Delta x^k,\ y^k+\Delta y^k)$
			\EndFor
		\end{algorithmic}
	\end{algorithm}
	
	\begin{proposition}[Local quadratic convergence]
		Let \(\mathcal R(v^\star)=0\). Suppose that \(\mathcal R\) is continuously
		differentiable in a neighborhood of \(v^\star\), that \(\mathcal R'\) is
		Lipschitz continuous there with constant \(L\), and that
		\(\mathcal R'(v^\star)\) is invertible with
		\[
		\norm{\mathcal R'(v^\star)^{-1}}\leq\beta.
		\]
		Then there exists \(\varepsilon>0\) such that, for every \(v^0\) satisfying
		\[
		\norm{v^0-v^\star}\leq\varepsilon,
		\]
		Newton's iteration
		\[
		v^{k+1}
		=
		v^k-\mathcal R'(v^k)^{-1}\mathcal R(v^k)
		\]
		is well defined, remains in the ball
		\(\overline B(v^\star,\varepsilon)\), and satisfies
		\begin{equation}
			\norm{v^{k+1}-v^\star}
			\leq
			\beta L\norm{v^k-v^\star}^2.
		\end{equation}
		In particular, \(v^k\to v^\star\) quadratically.
	\end{proposition}
	
	\begin{proof}
		Since \(\mathcal R'(v^\star)\) is invertible and \(\mathcal R'\) is
		continuous, there exists \(\varepsilon>0\) such that
		\(\mathcal R'(v)\) is invertible and
		\begin{equation}
			\norm{\mathcal R'(v)^{-1}}\leq2\beta
			\label{eq:local-inverse-bound}
		\end{equation}
		whenever
		\[
		\norm{v-v^\star}\leq\varepsilon.
		\]
		Shrinking \(\varepsilon\), if necessary, we may also assume that
		\(\overline B(v^\star,\varepsilon)\) lies in the neighborhood where
		\(\mathcal R'\) is \(L\)-Lipschitz and that
		\[
		\beta L\varepsilon<1.
		\]
		
		Let \(v\in\overline B(v^\star,\varepsilon)\), and set
		\[
		e=v-v^\star.
		\]
		Since \(\mathcal R(v^\star)=0\):
		\[
		\mathcal R(v)
		=
		\int_0^1
		\mathcal R'(v^\star+te)e\,dt.
		\]
		Consequently,
		\[
		\mathcal R'(v)e-\mathcal R(v)
		=
		\int_0^1
		\bigl(
		\mathcal R'(v)-\mathcal R'(v^\star+te)
		\bigr)e\,dt.
		\]
		The Lipschitz continuity of \(\mathcal R'\) yields
		\begin{align}
			\norm{\mathcal R'(v)e-\mathcal R(v)}
			&\leq
			\int_0^1
			L\norm{v-(v^\star+te)}\norm{e}\,dt\\
			&=
			L\norm{e}^2\int_0^1(1-t)\,dt\\
			&=
			\frac{L}{2}\norm{e}^2.
		\end{align}
		
		Let
		\[
		v^+
		=
		v-\mathcal R'(v)^{-1}\mathcal R(v)
		\]
		be the Newton iterate. Then
		\[
		v^+-v^\star
		=
		\mathcal R'(v)^{-1}
		\bigl(
		\mathcal R'(v)e-\mathcal R(v)
		\bigr).
		\]
		Using \eqref{eq:local-inverse-bound}, we obtain
		\[
		\norm{v^+-v^\star}
		\leq
		2\beta\frac{L}{2}\norm{e}^2
		=
		\beta L\norm{v-v^\star}^2.
		\]
		Since \(\norm{v-v^\star}\leq\varepsilon\),
		\[
		\norm{v^+-v^\star}
		\leq
		\beta L\varepsilon\norm{v-v^\star}
		<
		\norm{v-v^\star}
		\leq
		\varepsilon.
		\]
		Thus the Newton iterate remains in
		\(\overline B(v^\star,\varepsilon)\). The argument can therefore be applied
		inductively to every iterate, giving
		\[
		\norm{v^{k+1}-v^\star}
		\leq
		\beta L\norm{v^k-v^\star}^2
		\]
		for every \(k\). Hence \(v^k\to v^\star\) quadratically.
	\end{proof}
	\section{Semismooth Newton method}
	\begin{definition}[Semismoothness]
		Let \(\mathcal R:\R^N\to\R^N\) be locally Lipschitz. The mapping
		\(\mathcal R\) is \emph{semismooth at \(v\)} if it is directionally
		differentiable at \(v\) and
		\begin{equation}
			\sup_{V\in\partial_C\mathcal R(v+h)}
			\norm{
				\mathcal R(v+h)-\mathcal R(v)-Vh
			}
			=
			o(\norm{h})
			\qquad
			\text{as }h\to0.
			\label{eq:semismoothness}
		\end{equation}
		It is \emph{semismooth of order \(\gamma>0\)} at \(v\) if
		\begin{equation}
			\sup_{V\in\partial_C\mathcal R(v+h)}
			\norm{
				\mathcal R(v+h)-\mathcal R(v)-Vh
			}
			=
			O(\norm{h}^{1+\gamma}).
			\label{eq:semismoothness-order}
		\end{equation}
		The case \(\gamma=1\) is called \emph{strong semismoothness}.
	\end{definition}
	
	\begin{definition}[CD-regularity]
		A point \(v^\star\) is \emph{CD-regular} for \(\mathcal R\) if every matrix
		in the Clarke Jacobian
		\[
		\partial\mathcal R(v^\star)
		\]
		is nonsingular.
	\end{definition}
	
	\begin{algorithm}[H]
		\caption{Semismooth Newton method}
		\begin{algorithmic}[1]
			\Require A locally Lipschitz and semismooth mapping
			\(\mathcal R:\R^N\to\R^N\), an initial point \(v^0\), and a tolerance
			\(\varepsilon>0\)
			\Ensure A point \(v\) satisfying \(\norm{\mathcal R(v)}\leq\varepsilon\)
			\For{\(k=0,1,2,\ldots\)}
			\If{\(\norm{\mathcal R(v^k)}\leq\varepsilon\)}
			\State \Return \(v^k\)
			\EndIf
			\State Choose \(V_k\in\partial_C\mathcal R(v^k)\)
			\State Solve
			\[
			V_k\Delta v^k=-\mathcal R(v^k)
			\]
			\State Set \(v^{k+1}\gets v^k+\Delta v^k\)
			\EndFor
		\end{algorithmic}
	\end{algorithm}
	
	\begin{definition}[Q-convergence]
		Let \(v^k\to v^\star\).
		
		\begin{enumerate}
			\item The convergence is \emph{Q-linear} with rate
			\(\rho\in(0,1)\) if
			\[
			\norm{v^{k+1}-v^\star}
			\leq
			\rho\norm{v^k-v^\star}
			\]
			for all sufficiently large \(k\).
			
			\item The convergence is \emph{Q-superlinear} if
			\[
			\norm{v^{k+1}-v^\star}
			=
			o\bigl(\norm{v^k-v^\star}\bigr).
			\]
			
			\item The convergence has \emph{Q-order at least \(1+\gamma\)} if
			\[
			\norm{v^{k+1}-v^\star}
			=
			O\bigl(\norm{v^k-v^\star}^{1+\gamma}\bigr).
			\]
		\end{enumerate}
	\end{definition}
	
	\begin{proposition}[Closed graph of the Clarke Jacobian]
		Let \(\mathcal R:\R^N\to\R^N\) be locally Lipschitz. Then the set-valued
		mapping
		\[
		v\longmapsto\partial\mathcal R(v)
		\]
		has a closed graph. More precisely, if
		\[
		v_k\to v,
		\qquad
		V_k\to V,
		\qquad
		V_k\in\partial\mathcal R(v_k),
		\]
		then
		\[
		V\in\partial\mathcal R(v).
		\]
	\end{proposition}
	
	
	\begin{lemma}[Uniform invertibility near a CD-regular point]
		Suppose that \(\mathcal R\) is locally Lipschitz near \(v^\star\) and that
		every matrix in \(\partial\mathcal R(v^\star)\) is nonsingular. Then there
		exist a neighborhood \(U\) of \(v^\star\) and a constant \(\beta>0\) such
		that every
		\[
		V\in\partial\mathcal R(v),
		\qquad v\in U,
		\]
		is nonsingular and satisfies
		\[
		\norm{V^{-1}}\leq\beta.
		\]
	\end{lemma}
	
	\begin{proof}
		Suppose that the conclusion were false. Then there would exist sequences
		\[
		v_j\to v^\star,
		\qquad
		V_j\in\partial\mathcal R(v_j),
		\]
		such that \(V_j\) is singular or
		\[
		\norm{V_j^{-1}}\longrightarrow+\infty.
		\]
		
		
		Since \(\mathcal R\) is locally Lipschitz, its Clarke Jacobians are locally
		bounded. After passing to a subsequence, we may assume that
		\[
		V_j\to V.
		\]
		In either case, we can choose \(x_j\in\R^N\) such that
		\[
		\norm{x_j}=1
		\qquad\text{and}\qquad
		\norm{V_jx_j}\longrightarrow0.
		\]
		
		The unit sphere is compact, so, after passing to another subsequence,
		\[
		x_j\to x
		\qquad\text{with}\qquad
		\norm{x}=1.
		\]
		Since the Clarke Jacobian has a closed graph,
		\[
		V\in\partial\mathcal R(v^\star).
		\]
		Moreover,
		\[
		Vx
		=
		\lim_{j\to\infty}V_jx_j
		=
		0.
		\]
		Because \(x\neq0\), the matrix \(V\) is singular. This contradicts the
		CD-regularity of \(v^\star\). Therefore, the asserted neighborhood and
		uniform inverse bound exist.
	\end{proof}
	\begin{proposition}[Local convergence of the semismooth Newton method,
		\cite{QiSun1993}]
		Suppose that
		\[
		\mathcal R(v^\star)=0,
		\]
		that \(\mathcal R\) is semismooth at \(v^\star\), and that \(v^\star\) is
		CD-regular. Then there exists \(\varepsilon>0\) such that the semismooth
		Newton method is well defined for every initial point satisfying
		\[
		\norm{v^0-v^\star}\leq\varepsilon.
		\]
		The iterates remain in \(\overline B(v^\star,\varepsilon)\) and converge to
		\(v^\star\) Q-superlinearly.
		
		If \(\mathcal R\) is semismooth of order \(\gamma>0\) at \(v^\star\), then
		the convergence has Q-order at least \(1+\gamma\):
		\[
		\norm{v^{k+1}-v^\star}
		=
		O\bigl(\norm{v^k-v^\star}^{1+\gamma}\bigr).
		\]
		In particular, strong semismoothness gives local Q-quadratic convergence.
	\end{proposition}
	
	\begin{proof}
		By the uniform invertibility lemma, there exist a neighborhood \(U\) of
		\(v^\star\) and a constant \(\beta>0\) such that
		\[
		\norm{V^{-1}}\leq\beta
		\]
		for every \(v\in U\) and every
		\(V\in\partial_C\mathcal R(v)\).
		
		For \(v\) near \(v^\star\), set
		\[
		e=v-v^\star
		\]
		and choose any \(V\in\partial\mathcal R(v)\). Since
		\(\mathcal R(v^\star)=0\), semismoothness gives
		\[
		\mathcal R(v)-Ve=o(\norm{e}),
		\]
		uniformly over all \(V\in\partial\mathcal R(v)\).
		
		The corresponding semismooth Newton iterate is
		\[
		v^+=v-V^{-1}\mathcal R(v).
		\]
		Its error satisfies
		\begin{align}
			v^+-v^\star
			&=
			e-V^{-1}\mathcal R(v)\\
			&=
			-V^{-1}\bigl(\mathcal R(v)-Ve\bigr).
		\end{align}
		Consequently,
		\begin{align}
			\norm{v^+-v^\star}
			&\leq
			\norm{V^{-1}}
			\norm{\mathcal R(v)-Ve}\\
			&\leq
			\beta
			\norm{\mathcal R(v)-Ve}\\
			&=
			o(\norm{v-v^\star}).
		\end{align}
		
		Therefore, there exists \(\varepsilon>0\) such that
		\(\overline B(v^\star,\varepsilon)\subseteq U\) and
		\[
		\norm{v^+-v^\star}
		\leq
		\frac12\norm{v-v^\star}
		\]
		whenever
		\[
		0<\norm{v-v^\star}\leq\varepsilon.
		\]
		It follows inductively that every matrix \(V_k\) is nonsingular, every
		Newton step is well defined, and all iterates remain in
		\(\overline B(v^\star,\varepsilon)\). Moreover,
		\[
		\norm{v^{k+1}-v^\star}
		\leq
		\frac12\norm{v^k-v^\star},
		\]
		so \(v^k\to v^\star\). The sharper estimate
		\[
		\norm{v^{k+1}-v^\star}
		=
		o\bigl(\norm{v^k-v^\star}\bigr)
		\]
		shows that the convergence is Q-superlinear.
		
		If \(\mathcal R\) is semismooth of order \(\gamma\), then there exists
		\(C>0\) such that
		\[
		\norm{\mathcal R(v)-Ve}
		\leq
		C\norm{e}^{1+\gamma}
		\]
		for all \(v\) sufficiently close to \(v^\star\) and every
		\(V\in\partial_C\mathcal R(v)\). Hence
		\[
		\norm{v^+-v^\star}
		\leq
		\beta C\norm{v-v^\star}^{1+\gamma}.
		\]
		Applying this estimate to the iterates gives
		\[
		\norm{v^{k+1}-v^\star}
		=
		O\bigl(\norm{v^k-v^\star}^{1+\gamma}\bigr).
		\]
		For \(\gamma=1\), this is Q-quadratic convergence.
	\end{proof}
	
	\section{Merit functions and line searches}
	
	\begin{definition}[Merit function]
		Let $S\subseteq\R^N$ be the solution set of a problem. A function
		$\psi:\R^N\to\R\cup\{+\infty\}$ is a \emph{merit function} for it if
		$\argmin\psi=S$. 
	\end{definition}
	
	\begin{definition}[Descent direction]
		$d$ is a \emph{descent direction} for $\psi$ at $w$ if $\psi$ is differentiable
		at $w$ and $\inner{\nabla\psi(w)}{d}<0$.
	\end{definition}
	
	\begin{definition}[Line search]
		Given $w$, a direction $d$ and a merit $\psi$, a \emph{line search} is a rule
		returning $\alpha>0$ from the restriction
		$\alpha\mapsto\psi(w+\alpha d)$ and its slope
		$\varsigma(\alpha)=\inner{\nabla\psi(w+\alpha d)}{d}$.
	\end{definition}
	
	\begin{definition}[Exact line search]
		$\alpha\in\argmin_{\alpha>0}\psi(w+\alpha d)$.
	\end{definition}
	
	\begin{definition}[Secant search]
		Given $0\leq\alpha_{\mathrm{lo}}<\alpha_{\mathrm{hi}}$ with
		$\varsigma(\alpha_{\mathrm{lo}})\leq0\leq\varsigma(\alpha_{\mathrm{hi}})$, the
		\emph{secant search} performs at most $N$ steps of
		\begin{equation}
			\alpha=\alpha_{\mathrm{lo}}-\varsigma(\alpha_{\mathrm{lo}})\,
			\frac{\alpha_{\mathrm{hi}}-\alpha_{\mathrm{lo}}}
			{\varsigma(\alpha_{\mathrm{hi}})-\varsigma(\alpha_{\mathrm{lo}})},
		\end{equation}
		replacing $\alpha_{\mathrm{lo}}$ or $\alpha_{\mathrm{hi}}$ by $\alpha$
		according to the sign of $\varsigma(\alpha)$, falling back to the midpoint when
		$\alpha$ is not interior to the bracket, and stopping when the bracket is
		narrower than a tolerance. It returns $\alpha_{\mathrm{hi}}$.
	\end{definition}
	
	\begin{definition}[Armijo line search]
		Given $\gamma\in(0,1)$, $\beta\in(0,1)$ and a descent direction $d$, the
		\emph{backtracking Armijo} rule accepts the first
		$\alpha\in\{\beta^j\}_{j\geq0}$ with
		\begin{equation}
			\psi(w+\alpha d)\leq\psi(w)+\gamma\,\alpha\,\varsigma(0) .
		\end{equation}
	\end{definition}
	
	\begin{definition}[Backtracking decrease line search]
		Given $\beta\in(0,1)$, the rule accepts the first
		$\alpha\in\{\beta^j\}_{j\geq0}$ with $\psi(w+\alpha d)<\psi(w)$.
	\end{definition}
	
	\begin{algorithm}[H]
		\caption{Semismooth Newton method with a line search on a merit function}
		\begin{algorithmic}[1]
			\Require $\mathcal{R}$, a merit $\psi$ for $\{\mathcal{R}=0\}$, a line search
			rule, $v^0$, tolerance $\varepsilon>0$
			\Ensure $v$ with $\norm{\mathcal{R}(v)}\leq\varepsilon$
			\For{$k=0,1,2,\dots$}
			\If{$\norm{\mathcal{R}(v^k)}\leq\varepsilon$} \State \Return $v^k$ \EndIf
			\State choose $V_k\in\partial\mathcal{R}(v^k)$
			\State solve $V_k\,\Delta v^k=-\mathcal{R}(v^k)$
			\State obtain $\alpha_k$ from the line search rule applied to $\psi$ at $v^k$
			along $\Delta v^k$
			\State $v^{k+1}\gets v^k+\alpha_k\Delta v^k$
			\EndFor
		\end{algorithmic}
	\end{algorithm}
	
	
	
	\begin{proposition}[Semismoothness of the catalogue projections]
		The projections onto the following cones, and onto their duals, are
		semismooth at every point:
		
		\begin{enumerate}
			\item the zero cone;
			\item the nonnegative orthant;
			\item the Lorentz cone;
			\item the semidefinite cone;
			\item the exponential cone;
			\item the power cone.
		\end{enumerate}
	\end{proposition}
	
	\section{Self-concordant functions and damped Newton methods}
	
	\begin{definition}[Self-concordant function]
		Let \(D\subseteq\R^N\) be open and convex, and let
		\(f:D\to\R\) be convex and three times differentiable. The \emph{local
			seminorm} induced by \(f\) at \(x\in D\) is
		\begin{equation}
			\norm{h}_x
			=
			\bigl(D^2f(x)[h,h]\bigr)^{1/2}.
		\end{equation}
		The function \(f\) is \emph{self-concordant} if
		\begin{equation}
			\left|D^3f(x)[h,h,h]\right|
			\leq
			2\norm{h}_x^3
			\qquad
			\text{for every }x\in D\text{ and }h\in\R^N.
		\end{equation}
	\end{definition}
	
	\begin{proposition}[Stability under linear maps]
		Let \(f\) be self-concordant on \(D\), and let
		\(A\in\R^{N\times M}\). Then
		\[
		g=f\circ A
		\]
		is self-concordant on
		\[
		A^{-1}(D)=\{y\in\R^M:Ay\in D\}.
		\]
	\end{proposition}
	
	\begin{proof}
		The chain rule gives
		\[
		D^2g(y)[h,h]
		=
		D^2f(Ay)[Ah,Ah]
		\]
		and
		\[
		D^3g(y)[h,h,h]
		=
		D^3f(Ay)[Ah,Ah,Ah].
		\]
		Applying the self-concordance inequality for \(f\) at \(Ay\) in the
		direction \(Ah\), we obtain
		\begin{align}
			\left|D^3g(y)[h,h,h]\right|
			&=
			\left|D^3f(Ay)[Ah,Ah,Ah]\right|\\
			&\leq
			2\bigl(D^2f(Ay)[Ah,Ah]\bigr)^{3/2}\\
			&=
			2\bigl(D^2g(y)[h,h]\bigr)^{3/2}.
		\end{align}
		Thus \(g\) is self-concordant. No injectivity or surjectivity assumption on
		\(A\) is required.
	\end{proof}
	
	\begin{proposition}[Stability under sums and positive scaling]
		Let \(f_1,\ldots,f_r\) be self-concordant on the open convex sets
		\(D_1,\ldots,D_r\), and suppose that
		\[
		D=\bigcap_{i=1}^rD_i\neq\emptyset.
		\]
		Then
		\[
		f=\sum_{i=1}^rf_i
		\]
		is self-concordant on \(D\).
		
		Moreover, if \(f\) is self-concordant and \(t\geq1\), then \(tf\) is
		self-concordant.
	\end{proposition}
	
	
	
	
	\begin{definition}[Feasible points and directions]
		Consider the equality-constrained problem
		\[
		\min_{x\in D}f(x)
		\qquad
		\text{subject to}
		\qquad
		Ax=b.
		\]
		A point \(x\) is \emph{feasible} if
		\[
		x\in D
		\qquad\text{and}\qquad
		Ax=b.
		\]
		A direction \(h\) is a \emph{feasible direction} at \(x\) if
		\[
		Ah=0
		\]
		and \(x+th\in D\) for all sufficiently small \(t>0\).
	\end{definition}
	
	\begin{definition}[Dikin ellipsoid]
		Let \(f\) be twice differentiable on \(D\), with
		\(\nabla^2f(x)\succ0\). For \(x\in D\) and \(r>0\), the \emph{Dikin
			ellipsoid of radius \(r\)} centered at \(x\) is
		\begin{equation}
			\mathcal E_x(r)
			=
			\left\{
			x+h:
			\norm{h}_x<r
			\right\}.
		\end{equation}
		Equivalently,
		\[
		\mathcal E_x(r)
		=
		\left\{
		x+h:
		\inner{h}{\nabla^2f(x)h}<r^2
		\right\}.
		\]
		The set \(\mathcal E_x(1)\) is called the \emph{unit Dikin ellipsoid}.
	\end{definition}
	
	\begin{proposition}[Dikin ellipsoid property]
		Let
		\[
		f:\R^N\to\R\cup\{+\infty\}
		\]
		be closed, proper, and convex, with open domain
		\[
		D=\dom f.
		\]
		Suppose that \(f\) is self-concordant on \(D\) and that
		\[
		\nabla^2f(x)\succ0
		\qquad
		\text{for every }x\in D.
		\]
		Then
		\begin{equation}
			\mathcal E_x(1)
			=
			\left\{
			x+h:\norm{h}_x<1
			\right\}
			\subseteq D
			\qquad
			\text{for every }x\in D.
		\end{equation}
	\end{proposition}
	
	\begin{definition}[Newton step and Newton decrement]
		Let \(f\) be a closed self-concordant function with open domain
		\(D=\dom f\), and suppose that
		\[
		\nabla^2f(x)\succ0
		\qquad
		\text{for every }x\in D.
		\]
		Consider the equality-constrained problem
		\[
		\min_{x\in D}f(x)
		\qquad
		\text{subject to}
		\qquad
		Ax=b,
		\]
		where \(A\) has full row rank.
		
		At a feasible point \(x\), the \emph{Newton step} \(\Delta x\) and the
		associated multiplier \(w\) are defined by
		\begin{equation}
			\begin{pmatrix}
				\nabla^2f(x) & A^\top\\
				A            & 0
			\end{pmatrix}
			\begin{pmatrix}
				\Delta x\\
				w
			\end{pmatrix}
			=
			\begin{pmatrix}
				-\nabla f(x)\\
				0
			\end{pmatrix}.
			\label{eq:constrained-newton-system}
		\end{equation}
		The \emph{Newton decrement} is
		\begin{equation}
			\lambda(x)
			=
			\norm{\Delta x}_x
			=
			\bigl(
			D^2f(x)[\Delta x,\Delta x]
			\bigr)^{1/2}.
		\end{equation}
	\end{definition}
	
	\begin{proposition}[Properties of the Newton decrement]
		Let \(x\) be feasible. Then
		\begin{equation}
			\lambda(x)^2
			=
			-\inner{\nabla f(x)}{\Delta x}.
			\label{eq:newton-decrement-descent}
		\end{equation}
		If \(\lambda(x)>0\), then \(\Delta x\) is a feasible descent direction.
		Moreover,
		\[
		\lambda(x)=0
		\]
		if and only if \(x\) minimizes \(f\) subject to \(Ax=b\).
	\end{proposition}
	
	
	\begin{algorithm}[H]
		\caption{Damped Newton method for a self-concordant function}
		\begin{algorithmic}[1]
			\Require A closed self-concordant function \(f\) with open domain
			\(D=\dom f\) and \(\nabla^2f\succ0\), a full-row-rank matrix \(A\), a
			feasible point \(x^0\), a threshold
			\(\eta\in(0,\tfrac14]\), and a tolerance
			\(\varepsilon\in(0,\tfrac1{16}]\)
			\Ensure A feasible point \(x\) satisfying
			\(f(x)-p^\star\leq\varepsilon\)
			\For{\(k=0,1,2,\ldots\)}
			\State Compute \(\Delta x^k\) and
			\(\lambda_k=\lambda(x^k)\)
			\If{\(\lambda_k^2\leq\varepsilon\)}
			\State \Return \(x^k\)
			\EndIf
			\If{\(\lambda_k<\eta\)}
			\State \(t_k\gets1\)
			\Else
			\State \(t_k\gets\dfrac{1}{1+\lambda_k}\)
			\EndIf
			\State \(x^{k+1}\gets x^k+t_k\Delta x^k\)
			\EndFor
		\end{algorithmic}
	\end{algorithm}
	
	\begin{theorem}[Convergence of damped Newton,
		\cite{NesterovNemirovskii1994}]
		Let \(f\) be a closed self-concordant function with open domain
		\(D=\dom f\), and suppose that
		\[
		\nabla^2f(x)\succ0
		\qquad
		\text{for every }x\in D.
		\]
		Suppose that \(f\) attains its finite minimum \(p^\star\) on
		\[
		\{x\in D:Ax=b\}.
		\]
		Define
		\[
		\varrho(t)=t-\log(1+t).
		\]
		Then the damped Newton method has the following properties:
		
		\begin{enumerate}
			\item Every iterate belongs to \(D\) and satisfies \(Ax^k=b\).
			
			\item During the damped phase, when \(\lambda_k\geq\eta\),
			\[
			f(x^k)-f(x^{k+1})
			\geq
			\varrho(\eta).
			\]
			
			\item Once \(\lambda_k<\eta\leq\tfrac14\), full Newton steps are taken and
			\[
			2\lambda_{k+1}
			\leq
			(2\lambda_k)^2.
			\]
			
			\item Whenever \(\lambda_k\leq\tfrac14\),
			\[
			f(x^k)-p^\star
			\leq
			\lambda_k^2.
			\]
		\end{enumerate}
		
		Consequently, the number of iterations required to obtain
		\[
		f(x^k)-p^\star\leq\varepsilon
		\]
		is at most
		\begin{equation}
			\frac{f(x^0)-p^\star}{\varrho(\eta)}
			+
			O\bigl(\log\log(1/\varepsilon)\bigr).
		\end{equation}
		For fixed \(\eta\), the dependence on the problem data enters only through
		the initial objective gap \(f(x^0)-p^\star\).
	\end{theorem}
	
	
	
	\section{Smoothing families}
	
	A smoothing family replaces the cone indicator by a sequence of convex
	functions that are differentiable on suitable open domains. Each resulting
	subproblem is an equality-constrained smooth convex problem.
	
	\begin{definition}[Smoothing of an indicator]
		Let \(K\subseteq\R^m\) be a closed convex cone. A \emph{smoothing of
			\(\delta_K\)} is a family
		\[
		\left\{
		(\phi_\theta,A_\theta,b_\theta)
		\right\}_{\theta\in\Theta}
		\]
		with the following properties for every \(\theta\in\Theta\):
		
		\begin{enumerate}
			\item The function
			\[
			\phi_\theta:\R^m\to\R\cup\{+\infty\}
			\]
			is closed, proper, and convex.
			
			\item Its interior domain
			\[
			D_\theta=\intr(\dom\phi_\theta)
			\]
			is nonempty.
			
			\item The function \(\phi_\theta\) is differentiable on \(D_\theta\), and
			\(\nabla\phi_\theta\) is semismooth there.
			
			\item For every \(s\in D_\theta\) and every
			\[
			H\in\partial(\nabla\phi_\theta)(s),
			\]
			the matrix \(H\) is symmetric positive semidefinite.
			
			\item The matrix \(A_\theta\) and vector \(b_\theta\) specify the affine
			constraint carried by the corresponding smoothed problem.
		\end{enumerate}
	\end{definition}
	
	\begin{definition}[Interior and exterior smoothing]
		A smoothing family is \emph{interior} if
		\[
		D_\theta\subseteq K
		\qquad
		\text{for every }\theta\in\Theta.
		\]
		It is \emph{exterior} if
		\[
		K\subseteq D_\theta
		\qquad
		\text{for every }\theta\in\Theta.
		\]
	\end{definition}
	
	\begin{algorithm}[H]
		\caption{Penalty method based on a smoothing family}
		\begin{algorithmic}[1]
			\Require A differentiable convex function \(f\), an affine mapping
			\(c(x)=Cx+d\), a cone \(K\), a smoothing family
			\(\{(\phi_\theta,A_\theta,b_\theta)\}_{\theta\in\Theta}\), parameters
			\((\theta_k)\), and tolerances \((\varepsilon_k)\)
			\Ensure An approximate solution of
			\[
			\min_x f(x)
			\qquad
			\text{subject to}
			\qquad
			-c(x)\in K
			\]
			\For{\(k=0,1,2,\ldots\)}
			\State Compute \(x^{k+1}\), to accuracy \(\varepsilon_k\), by solving
			\[
			\min_x
			\left\{
			f(x)+\phi_{\theta_k}\bigl(-c(x)\bigr)
			\right\}
			\qquad
			\text{subject to}
			\qquad
			A_{\theta_k}x=b_{\theta_k}
			\]
			\State Choose the next parameter \(\theta_{k+1}\)
			\EndFor
		\end{algorithmic}
	\end{algorithm}
	
	\begin{proposition}[Optimality conditions for a smoothed problem]
		Let
		\[
		c(x)=Cx+d,
		\]
		fix \(\theta\in\Theta\), and suppose that
		\[
		-c(x)\in D_\theta,
		\qquad
		A_\theta x=b_\theta.
		\]
		Then \(x\) minimizes
		\[
		f(x)+\phi_\theta\bigl(-c(x)\bigr)
		\]
		subject to \(A_\theta x=b_\theta\) if and only if there exist vectors
		\(y\) and \(z\) such that
		\begin{equation}
			\begin{aligned}
				\nabla f(x)+C^\top z+A_\theta^\top y&=0,\\
				z&=-\nabla\phi_\theta\bigl(-c(x)\bigr),\\
				A_\theta x&=b_\theta.
			\end{aligned}
			\label{eq:smoothed-optimality-system}
		\end{equation}
	\end{proposition}
	
	\begin{proof}
		Define
		\[
		F_\theta(x)
		=
		f(x)+\phi_\theta\bigl(-c(x)\bigr).
		\]
		This function is convex and differentiable at \(x\). Since
		\[
		D(-c)(x)=-C,
		\]
		the chain rule gives
		\[
		\nabla F_\theta(x)
		=
		\nabla f(x)
		-
		C^\top\nabla\phi_\theta\bigl(-c(x)\bigr).
		\]
		
		A feasible point \(x\) minimizes \(F_\theta\) subject to
		\(A_\theta x=b_\theta\) if and only if its gradient is orthogonal to every
		feasible direction:
		\[
		\inner{\nabla F_\theta(x)}{h}=0
		\qquad
		\text{for every }h\in\ker A_\theta.
		\]
		Since
		\[
		(\ker A_\theta)^\perp=\range(A_\theta^\top),
		\]
		this is equivalent to the existence of \(y\) such that
		\[
		\nabla F_\theta(x)+A_\theta^\top y=0.
		\]
		Setting
		\[
		z=-\nabla\phi_\theta\bigl(-c(x)\bigr)
		\]
		gives
		\[
		\nabla f(x)+C^\top z+A_\theta^\top y=0.
		\]
		Together with \(A_\theta x=b_\theta\), this is precisely
		\eqref{eq:smoothed-optimality-system}. Since the problem is convex, these
		first-order conditions are both necessary and sufficient.
	\end{proof}
	
	\chapter{Interior-point method basics}
	
	\section{Barrier smoothing}
	
	\begin{definition}[Conic quadratic problem]
		Let
		\[
		f(x)=\frac12\inner{x}{Px}+\inner{q}{x},
		\qquad
		c(x)=Cx+d,
		\]
		where \(P\succeq0\), and let \(K\subseteq\R^m\) be a closed convex cone. We
		consider
		\begin{equation}
			\min_{x\in\R^n} f(x)
			\qquad
			\text{subject to}
			\qquad
			-c(x)\in K.
			\label{eq:conic-quadratic-problem}
		\end{equation}
	\end{definition}
	
	\begin{definition}[Slack formulation]
		Introducing the slack variable
		\[
		s=-c(x),
		\]
		problem \eqref{eq:conic-quadratic-problem} becomes
		\begin{equation}
			\min_{x,s} f(x)
			\qquad
			\text{subject to}
			\qquad
			s+c(x)=0,
			\quad
			s\in K.
			\label{eq:slack-formulation}
		\end{equation}
	\end{definition}
	
	
	\begin{definition}[Self-concordant barrier]
		Let \(K\) be a closed convex cone . A function
		\[
		B:\R^m\to\R\cup\{+\infty\}
		\]
		is a \emph{self-concordant barrier for \(K\)} if:
		
		\begin{enumerate}
			\item \(B\) is closed, proper, and convex, with
			\[
			\dom B=\intr K;
			\]
			
			\item \(B\) is three times differentiable and self-concordant on
			\(\intr K\);
			
			\item
			\[
			B(s_k)\longrightarrow+\infty
			\]
			whenever \(s_k\in\relint K\) converges to a point of \(\relbd K\).
		\end{enumerate}
	\end{definition}
	
	\begin{definition}[Barrier parameter]
		A self-concordant barrier \(B\) has \emph{parameter \(\nu>0\)} if
		\begin{equation}
			\norm{\nabla B(s)}_{s,B}^*
			\leq
			\sqrt{\nu}
			\qquad
			\text{for every }s\in\intr K,
			\label{eq:barrier-parameter}
		\end{equation}
		where
		\[
		\norm{h}_{s,B}
		=
		\bigl(
		h^\top\nabla^2B(s)h
		\bigr)^{1/2}
		\]
		and
		\[
		\norm{u}_{s,B}^*
		=
		\bigl(
		u^\top[\nabla^2B(s)]^{-1}u
		\bigr)^{1/2}.
		\]
		Equivalently,
		\[
		|DB(s)[h]|
		\leq
		\sqrt{\nu}\,\norm{h}_{s,B}
		\qquad
		\text{for every }h\in\R^m.
		\]
	\end{definition}
	
	\begin{definition}[Barrier problem]
		Let \(B\) be a self-concordant barrier for \(K\). For \(\tau>0\), the
		\emph{barrier problem} is
		\begin{equation}
			\min_{x,s}
			\left\{
			f(x)+\tau B(s)
			\right\}
			\qquad
			\text{subject to}
			\qquad
			s+c(x)=0.
			\label{eq:barrier-problem}
		\end{equation}
		Since \(\tau>0\), it has the same minimizers as
		\begin{equation}
			\min_{x,s}
			\left\{
			\frac{f(x)}{\tau}+B(s)
			\right\}
			\qquad
			\text{subject to}
			\qquad
			s+c(x)=0.
			\label{eq:rescaled-barrier-problem}
		\end{equation}
	\end{definition}
	
	The rescaled formulation \eqref{eq:rescaled-barrier-problem} is important.
	Although \(B\) is self-concordant, the function \(\tau B\) is
	self-concordant with the standard normalization only when \(\tau\geq1\).
	By contrast, \(f/\tau+B\) is self-concordant for every \(\tau>0\), because
	\(f/\tau\) is a convex quadratic and has zero third derivative.
	
	\begin{proposition}[Optimality conditions for the barrier problem]
		Let \(s\in\intr K\) and \(s+c(x)=0\). Then \((x,s)\) solves
		\eqref{eq:barrier-problem} if and only if there exist \(y,z\in\R^m\)
		satisfying
		\begin{align}
			\nabla f(x)+C^\top y&=0,
			\label{eq:barrier-kkt-1}\\
			y-z&=0,
			\label{eq:barrier-kkt-2}\\
			s+c(x)&=0,
			\label{eq:barrier-kkt-3}\\
			z+\tau\nabla B(s)&=0.
			\label{eq:barrier-kkt-4}
		\end{align}
	\end{proposition}
	\begin{proof}
		Consequence of the general smoothing optimality condition.
	\end{proof}
	
	\begin{definition}[Central path]
		Suppose that the barrier problem has a unique solution for every
		\(\tau>0\). The \emph{central path} is
		\begin{equation}
			\left\{
			\bigl(x(\tau),s(\tau),y(\tau),z(\tau)\bigr):
			\tau>0
			\right\},
		\end{equation}
		where the point at parameter \(\tau\) satisfies
		\eqref{eq:barrier-kkt-1} to \eqref{eq:barrier-kkt-4}.
	\end{definition}
	
	\begin{algorithm}[H]
		\caption{Barrier method}
		\begin{algorithmic}[1]
			\Require A self-concordant barrier \(B\) with parameter \(\nu\), a strictly
			feasible point, an initial parameter \(\tau_0>0\), a reduction factor
			\(\rho\in(0,1)\), a centering tolerance
			\(\eta\in(0,\tfrac14]\), and a final tolerance \(\varepsilon>0\)
			\Ensure An approximate primal-dual solution
			\For{\(k=0,1,2,\ldots\)}
			\If{\(\nu\tau_k\leq\varepsilon\)}
			\State \Return the current iterate
			\EndIf
			\State Approximately minimize
			\[
			\frac{f(x)}{\tau_k}+B(s)
			\qquad
			\text{subject to}
			\qquad
			s+c(x)=0
			\]
			by damped Newton steps until the Newton decrement is at most \(\eta\)
			\State Set \(z\gets-\tau_k\nabla B(s)\)
			\State Set \(\tau_{k+1}\gets\rho\tau_k\)
			\EndFor
		\end{algorithmic}
	\end{algorithm}
	
	\section{Short-step path following}
	
	\begin{definition}[Local norm and dual local norm]
		Let \(g\) be self-concordant and suppose that
		\(\nabla^2g(v)\succ0\). Define
		\begin{equation}
			\norm{h}_{v,g}
			=
			\bigl(
			h^\top\nabla^2g(v)h
			\bigr)^{1/2},
			\qquad
			\norm{u}_{v,g}^*
			=
			\bigl(
			u^\top[\nabla^2g(v)]^{-1}u
			\bigr)^{1/2}.
		\end{equation}
		The Newton decrement is
		\[
		\lambda_g(v)
		=
		\norm{\nabla g(v)}_{v,g}^*.
		\]
		For an equality-constrained problem, the decrement is instead computed from
		the constrained Newton direction.
	\end{definition}
	
	For the barrier problem, write
	\[
	g_\tau(x,s)
	=
	\frac{f(x)}{\tau}+B(s)
	\]
	on the affine set \(s+c(x)=0\), and denote its constrained Newton decrement
	by \(\lambda_\tau(x,s)\).
	
	\begin{proposition}[Effect of decreasing the barrier parameter]
		Suppose that
		\[
		\lambda_\tau(v)\leq\eta
		\]
		and set
		\[
		\tau^+
		=
		\left(
		1-\frac{\gamma}{\sqrt{\nu}}
		\right)\tau,
		\qquad
		0<\gamma<\sqrt{\nu}.
		\]
		Then
		\begin{equation}
			\lambda_{\tau^+}(v)
			\leq
			\frac{\eta+\gamma}
			{1-\gamma/\sqrt{\nu}}.
			\label{eq:one-step-cut-bound}
		\end{equation}
	\end{proposition}
	
	\begin{algorithm}[H]
		\caption{Short-step path-following method}
		\begin{algorithmic}[1]
			\Require A barrier \(B\) with parameter \(\nu\), constants
			\(\eta,\gamma>0\) satisfying
			\[
			\eta'
			:=
			\frac{\eta+\gamma}{1-\gamma/\sqrt{\nu}}
			\leq\frac14,
			\qquad
			2(\eta')^2\leq\eta,
			\]
			an initial parameter \(\tau_0>0\), a tolerance \(\varepsilon>0\), and a
			point \(v^0\) satisfying \(\lambda_{\tau_0}(v^0)\leq\eta\)
			\Ensure A point with duality gap at most \(\varepsilon\)
			\For{\(k=0,1,2,\ldots\)}
			\If{\(\nu\tau_k\leq\varepsilon\)}
			\State \Return \(v^k\)
			\EndIf
			\State Set
			\[
			\tau_{k+1}
			\gets
			\left(
			1-\frac{\gamma}{\sqrt{\nu}}
			\right)\tau_k
			\]
			\State Take one full Newton step for \(g_{\tau_{k+1}}\) at \(v^k\)
			\State Denote the resulting point by \(v^{k+1}\)
			\EndFor
		\end{algorithmic}
	\end{algorithm}
	
	The Newton system for \(g_\tau\) is
	\begin{equation}
		\begin{pmatrix}
			\nabla^2f(x)/\tau & 0                  & C^\top\\
			0                 & \nabla^2B(s)       & I\\
			C                 & I                  & 0
		\end{pmatrix}
		\begin{pmatrix}
			\Delta x\\
			\Delta s\\
			\Delta y
		\end{pmatrix}
		=
		-
		\begin{pmatrix}
			\nabla f(x)/\tau+C^\top y\\
			\nabla B(s)+y\\
			s+c(x)
		\end{pmatrix}.
		\label{eq:short-step-newton-system}
	\end{equation}
	
	\begin{proposition}[Iteration count of the short-step method,
		\cite{NesterovNemirovskii1994}]
		Every iterate of the short-step method satisfies
		\[
		\lambda_{\tau_k}(v^k)\leq\eta.
		\]
		Moreover, the method reaches
		\[
		\nu\tau_k\leq\varepsilon
		\]
		after at most
		\begin{equation}
			\left\lceil
			\frac{
				\log(\nu\tau_0/\varepsilon)
			}{
				-\log(1-\gamma/\sqrt{\nu})
			}
			\right\rceil
		\end{equation}
		iterations. In particular, for fixed \(\gamma\),
		\begin{equation}
			k
			=
			O\left(
			\sqrt{\nu}\log\frac{\nu\tau_0}{\varepsilon}
			\right).
		\end{equation}
	\end{proposition}
	
	
	The interest of the short step method is at each step you don't have an inner damped newton only one step per $\tau$.
	\section{Logarithmically homogeneous barriers}
	
	\begin{definition}[Logarithmic homogeneity]
		A barrier \(B\) for \(K\) is \emph{logarithmically homogeneous with
			parameter \(\nu\)} if
		\begin{equation}
			B(ts)
			=
			B(s)-\nu\log t
			\qquad
			\text{for every }s\in\intr K\text{ and }t>0.
			\label{eq:logarithmic-homogeneity}
		\end{equation}
	\end{definition}
	
	
	
	\section{Barriers for the catalogue}
	
	\begin{proposition}[Barrier of a symmetric cone]
		Let \(K\) be a symmetric cone associated with a Euclidean Jordan algebra
		\(V\) of rank \(r\). With respect to the canonical trace inner product,
		\[
		B(s)=-\log\det(s)
		\]
		is a logarithmically homogeneous self-concordant barrier with parameter
		\[
		\nu=r.
		\]
		Its derivatives are in the canonical inner product:
		\begin{equation}
			\nabla B(s)=-s^{-1},
			\qquad
			\nabla^2B(s)=Q(s)^{-1}=Q(s^{-1}).
		\end{equation}
	\end{proposition}
	
	The standard barriers for the classical symmetric cones are:
	
	\begin{enumerate}
		\item For the nonnegative orthant,
		\[
		K=\R_+^n,
		\qquad
		B(x)=-\sum_{i=1}^n\log x_i,
		\qquad
		\nu=n.
		\]
		
		\item For the Lorentz cone,
		\[
		K=\Lor^{n+1}
		=
		\left\{
		(t,u)\in\R\times\R^n:
		t\geq\norm{u}
		\right\},
		\]
		the standard Euclidean barrier is
		\[
		B(t,u)
		=
		-\log\bigl(t^2-\norm{u}^2\bigr),
		\qquad
		\nu=2.
		\]
		
		\item For the positive semidefinite cone,
		\[
		K=\Sp_+^n,
		\qquad
		B(X)=-\log\det X,
		\qquad
		\nu=n.
		\]
	\end{enumerate}
	
	
	\begin{remark}[Lorentz cone specific case]
		The Lorentz cone requires particular care because its canonical trace inner
		product is twice the usual Euclidean inner product:
		\[
		\inner{x}{y}_{\mathrm{tr}}
		=
		2\inner{x}{y}_{\mathrm{Euc}}.
		\]
		
		Hence in the classic inner product we have :
		
		$$ \nabla B(s) =  -2s^{-1}$$
	\end{remark}
	
	\begin{proposition}[Barriers for the exponential and power cones,
		\cite{NesterovNemirovskii1994,Chares2009}]
		The function
		\begin{equation}
			B_{\exp}(s)
			=
			-\log\bigl(
			s_2\log(s_3/s_2)-s_1
			\bigr)
			-\log s_2
			-\log s_3
		\end{equation}
		is a logarithmically homogeneous self-concordant barrier for the exponential
		cone with parameter \(3\).
		
		For \(\alpha\in(0,1)\), the function
		\begin{equation}
			B_\alpha(s)
			=
			-\log\bigl(
			s_1^{2\alpha}s_2^{2(1-\alpha)}-s_3^2
			\bigr)
			-(1-\alpha)\log s_1
			-\alpha\log s_2
		\end{equation}
		is a logarithmically homogeneous self-concordant barrier for the power cone
		with parameter \(3\).
	\end{proposition}
	
	\begin{proposition}[Euler's identity for a logarithmically homogeneous barrier]
		\label{prop:euler}
		Let $K\subseteq\R^m$ be a cone and let $B$ be differentiable and
		$\nu$-logarithmically homogeneous, that is
		\begin{equation}\label{eq:loghom}
			B(\lambda s)=B(s)-\nu\log\lambda
			\qquad\text{for all }s\in\intr K,\ \lambda>0 .
		\end{equation}
		Then
		\begin{equation}\label{eq:euler}
			\inner{s}{-\nabla B(s)}=\nu
			\qquad\text{for all }s\in\intr K .
		\end{equation}
	\end{proposition}
	
	\begin{proof}
		Fix $s\in\intr K$ and let $$\varphi(\lambda)=B(\lambda s)$$ for $\lambda>0$, which
		is well defined since $\intr K$ is a cone. 
		
		
		Differentiating the left side by the
		chain rule gives $$\varphi'(\lambda)=\inner{\nabla B(\lambda s)}{s}$$, and
		differentiating the right side of \eqref{eq:loghom} gives
		$\varphi'(\lambda)=-\nu/\lambda$. Evaluating at $\lambda=1$ yields
		$\inner{\nabla B(s)}{s}=-\nu$, which is \eqref{eq:euler}.
	\end{proof}
	
	Euler identity is useful cause it permit you to know where you are on the central path.
	
	\begin{corollary}[The pairing on the central path]\label{cor:pairing}
		If $z=-\tau\nabla B(s)$ for some $\tau>0$, then
		\begin{equation}
			\inner{s}{z}=\nu\tau .
		\end{equation}
	\end{corollary}
	
	\begin{proof}
		$\inner{s}{z}=\inner{s}{-\tau \nabla B(s)}=\nu\tau$
		by Proposition~\ref{prop:euler}.
	\end{proof}
	
	\begin{proposition}[Barrier of a product]
		Let
		\[
		K=K_1\times\cdots\times K_p,
		\]
		and let \(B_j\) be a logarithmically homogeneous self-concordant barrier for
		\(K_j\) with parameter \(\nu_j\). Then
		\begin{equation}
			B(s)
			=
			\sum_{j=1}^pB_j(s_j)
		\end{equation}
		is a logarithmically homogeneous self-concordant barrier for \(K\) with
		parameter
		\begin{equation}
			\nu=\sum_{j=1}^p\nu_j.
		\end{equation}
	\end{proposition}
	
	\begin{remark}[Zero-cone blocks]
		Let
		\[
		K_j=\{0\}.
		\]
		This cone has empty interior in its ambient space, but its relative interior
		is
		\[
		\relint K_j=\{0\},
		\]
		and its relative boundary is empty. Under the relative-interior convention
		adopted here, any constant function on \(\{0\}\) is therefore a barrier by
		vacuity. We choose
		\[
		B_j\equiv0,
		\]
		with barrier parameter \(\nu_j=0\).
		
		Traditionally, zero-cone blocks are excluded from the definition of a
		barrier and treated directly as equality constraints. We retain them in the
		present framework to obtain a uniform treatment of product cones.
		
		Since the Hessian of a constant barrier is zero, it provides no curvature
		to the Newton system. In the implementation, the zero block is therefore
		replaced by a fixed positive regularization block of the form
		\[
		H_j=\Lambda_j I,
		\qquad
		\Lambda_j>0.
		\]
		This block comes from the augmented-Lagrangian or proximal regularization( that we will view later on );
		it is not the Hessian of the zero-cone barrier.
	\end{remark}
	
	\section{Primal-dual Newton systems}
	
	Let
	\[
	v=(x,s,y,z),
	\]
	where \(y\) is the multiplier of \(s+c(x)=0\) and \(z\) is the cone dual
	variable. The residuals common to all the systems considered below are
	\begin{align}
		r_1&=Px+q+C^\top y,\\
		r_2&=y-z,\\
		r_3&=s+c(x).
	\end{align}
	The final residual \(r_4\) describes complementarity.
	
	A general linearized system has the form
	\begin{equation}
		\begin{pmatrix}
			P & 0 & C^\top & 0\\
			0 & 0 & I      & -I\\
			C & I & 0      & 0\\
			0 & H & 0      & I
		\end{pmatrix}
		\begin{pmatrix}
			\Delta x\\
			\Delta s\\
			\Delta y\\
			\Delta z
		\end{pmatrix}
		=
		-
		\begin{pmatrix}
			r_1\\
			r_2\\
			r_3\\
			r_4
		\end{pmatrix}.
		\label{eq:general-primal-dual-system}
	\end{equation}
	
	Eliminating
	\[
	\Delta z=-H\Delta s-r_4
	\]
	gives
	\begin{equation}
		\begin{pmatrix}
			P & 0 & C^\top\\
			0 & H & I\\
			C & I & 0
		\end{pmatrix}
		\begin{pmatrix}
			\Delta x\\
			\Delta s\\
			\Delta y
		\end{pmatrix}
		=
		-
		\begin{pmatrix}
			r_1\\
			r_2+r_4\\
			r_3
		\end{pmatrix}.
		\label{eq:three-by-three-primal-dual-system}
	\end{equation}
	This matrix is symmetric whenever \(H\) is symmetric.
	
	Further elimination gives
	\begin{equation}
		\bigl(
		P+C^\top HC
		\bigr)\Delta x
		=
		-r_1
		+
		C^\top(r_2+r_4)
		-
		C^\top Hr_3.
		\label{eq:condensed-primal-dual-system}
	\end{equation}
	The remaining directions are recovered from
	\begin{align}
		\Delta s&=-r_3-C\Delta x,\\
		\Delta z&=-H\Delta s-r_4,\\
		\Delta y&=\Delta z-r_2.
	\end{align}
	
	\begin{definition}[Admissible Newton system]
		Let \(G_\tau\) denote the nonlinear residual whose linearization produces
		the Newton system. The system is \emph{admissible} if
		\[
		G_\tau(v)=0
		\]
		holds exactly at the central-path points with parameter \(\tau\).
	\end{definition}
	
	\begin{definition}[Tangent Newton system]
		An admissible Newton system is \emph{tangent} if, at every central-path
		point,
		\begin{equation}
			H=\tau\nabla^2B(s).
		\end{equation}
		Thus its linearization agrees with the exact barrier linearization on the
		central path. Away from the path, \(H\) may be an approximation.
	\end{definition}
	
	\begin{remark}
		Admissibility concerns the zero set of the nonlinear residual. Tangency
		concerns its first-order behavior at that zero set. These are different
		properties.
		
		At a central-path point, the residual and the Newton step are both zero, so
		the value of the step cannot be used to test tangency. Tangency is a
		statement about the derivative of the residual.
	\end{remark}
	
	\subsection{Exact barrier Newton system}
	
	For the barrier residual
	\[
	r_4=z+\tau\nabla B(s),
	\]
	the exact derivative with respect to \(s\) is
	\[
	H=\tau\nabla^2B(s).
	\]
	The resulting system is both admissible and tangent.
	
	After eliminating \(\Delta z\), it coincides with the Newton system
	\eqref{eq:short-step-newton-system}, up to the rescaling of the equality
	multiplier caused by dividing the objective by \(\tau\).
	
	\subsection{Averaged Newton system}
	
	\begin{definition}[Averaged Newton system]
		Let \(B\) be logarithmically homogeneous with parameter \(\nu\), and define
		the current average complementarity by
		\begin{equation}
			\mu=\frac{\inner{s}{z}}{\nu}.
		\end{equation}
		The \emph{averaged Newton system} uses
		\[
		r_4=z+\tau\nabla B(s)
		\]
		and
		\[
		H=\mu\nabla^2B(s).
		\]
		Thus
		\begin{equation}
			\begin{pmatrix}
				P & 0 & C^\top & 0\\
				0 & 0 & I      & -I\\
				C & I & 0      & 0\\
				0 & \mu\nabla^2B(s) & 0 & I
			\end{pmatrix}
			\begin{pmatrix}
				\Delta x\\
				\Delta s\\
				\Delta y\\
				\Delta z
			\end{pmatrix}
			=
			-
			\begin{pmatrix}
				Px+q+C^\top y\\
				y-z\\
				s+c(x)\\
				z+\tau\nabla B(s)
			\end{pmatrix}.
		\end{equation}
		After eliminating \(\Delta z\),
		\begin{equation}
			\begin{pmatrix}
				P & 0 & C^\top\\
				0 & \mu\nabla^2B(s) & I\\
				C & I & 0
			\end{pmatrix}
			\begin{pmatrix}
				\Delta x\\
				\Delta s\\
				\Delta y
			\end{pmatrix}
			=
			-
			\begin{pmatrix}
				Px+q+C^\top y\\
				y+\tau\nabla B(s)\\
				s+c(x)
			\end{pmatrix}.
		\end{equation}
	\end{definition}
	
	\begin{proposition}[Admissibility and tangency of the averaged system]
		The averaged Newton system is admissible. It is also tangent on the central
		path.
	\end{proposition}
	
	\begin{proof}
		The nonlinear residuals are
		\[
		r_1,\qquad r_2,\qquad r_3,\qquad
		r_4=z+\tau\nabla B(s).
		\]
		Their simultaneous vanishing is exactly the optimality system of the barrier
		problem. The choice of \(H\) changes only the linearized direction and does
		not change the zero set. Hence the system is admissible.
		
		On the central path by Euler identity:
		\[
		\inner{s}{z}
		=
		-\tau\inner{s}{\nabla B(s)}
		=
		\nu\tau.
		\]
		Thus
		\[
		\mu=\frac{\inner{s}{z}}{\nu}=\tau,
		\]
		and therefore
		\[
		H
		=
		\mu\nabla^2B(s)
		=
		\tau\nabla^2B(s).
		\]
		Hence the system is tangent.
	\end{proof}
	
	\subsection{Jordan Newton system}
	
	In this subsection, \(K\)  is symmetric and solid and the Jordan variables are written
	using the canonical trace inner product.
	
	\begin{proposition}[Jordan form of centrality]
		Let
		\[
		B(s)=-\log\det(s).
		\]
		For \(s,z\in\int K\),
		\[
		z+\tau\nabla B(s)=0
		\]
		supposing that $\inner{\cdot}{\cdot}_{\mathrm{can}}=C\inner{\cdot}{\cdot}$ to: 
		\begin{equation}
			s\circ z=C\tau e.
			\label{eq:jordan-centrality}
		\end{equation}
	\end{proposition}
	
	
	\begin{proof}
		With respect to the  used inner product:
		\[
		\nabla B(s)=-Cs^{-1}.
		\]
		Therefore,
		\[
		z+\tau\nabla B(s)=0
		\quad\Longleftrightarrow\quad
		z=C\tau s^{-1}.
		\]
		Taking the Jordan product with \(s\) gives
		\[
		s\circ z
		=
		C\tau(s\circ s^{-1})
		=
		C\tau e.
		\]
	\end{proof}
	
	\begin{definition}[Jordan Newton system]
		Linearizing
		\[
		s\circ z-C\tau e=0
		\]
		gives
		\[
		L(z)\Delta s+L(s)\Delta z
		=
		-\bigl(s\circ z-C\tau e\bigr).
		\]
		Since \(s\in\intr K\), the operator \(L(s)\) is invertible. Multiplying by
		\(L(s)^{-1}\) gives
		\[
		\Delta z+L(s)^{-1}L(z)\Delta s
		=
		-L(s)^{-1}(s\circ z-C\tau e).
		\]
		Thus the Jordan system uses
		\begin{equation}
			H=L(s)^{-1}L(z)
		\end{equation}
		and
		\begin{equation}
			r_4
			=
			L(s)^{-1}(s\circ z-C\tau e)
			=
			z+\tau\nabla B(s).
		\end{equation}
	\end{definition}
	
	The operator \(L(s)^{-1}L(z)\) is generally not symmetric because
	\(L(s)\) and \(L(z)\) need not commute. This is the reason for introducing
	Nesterov-Todd scaling.
	
	\begin{proposition}[Admissibility of the Jordan system]
		The Jordan Newton system is admissible.
	\end{proposition}
	
	\begin{proof}
		For \(s,z\in\intr K\),
		\[
		s\circ z=C\tau e
		\]
		is equivalent to
		\[
		z+\tau\nabla B(s)=0.
		\]
		Thus the Jordan system has the same zero set as the barrier optimality
		system and is admissible.
	\end{proof}
	
	\subsection{Nesterov-Todd scaling}
	
	\begin{definition}[Nesterov-Todd scaling]
		Let \(K\) be symmetric and let \(s,z\in\intr K\). The
		\emph{Nesterov-Todd scaling} between \(s\) and \(z\) is the unique
		self-adjoint positive-definite automorphism \(W\in\Aut(K)\) satisfying
		\[
		Ws=z.
		\]
		It is given by
		\begin{equation}
			W=Q(w),
			\qquad
			w
			=
			Q(s^{-1/2})
			\left(
			Q(s^{1/2})z
			\right)^{1/2}.
		\end{equation}
	\end{definition}
	
	\begin{definition}[Nesterov-Todd Newton system]
		Define the scaled variables
		\[
		\widetilde s=W^{1/2}s,
		\qquad
		\widetilde z=W^{-1/2}z.
		\]
		Since \(Ws=z\),
		\[
		\widetilde s=\widetilde z.
		\]
		Linearizing the Jordan complementarity equation in these scaled variables,
		while holding \(W\) fixed, produces the symmetric block
		\[
		H=W.
		\]
		The resulting Newton system is
		\begin{equation}
			\begin{pmatrix}
				P & 0 & C^\top & 0\\
				0 & 0 & I      & -I\\
				C & I & 0      & 0\\
				0 & W & 0      & I
			\end{pmatrix}
			\begin{pmatrix}
				\Delta x\\
				\Delta s\\
				\Delta y\\
				\Delta z
			\end{pmatrix}
			=
			-
			\begin{pmatrix}
				Px+q+C^\top y\\
				y-z\\
				s+c(x)\\
				z+\tau\nabla B(s)
			\end{pmatrix}.
			\label{eq:nt-newton-system}
		\end{equation}
		After eliminating \(\Delta z\),
		\begin{equation}
			\begin{pmatrix}
				P & 0 & C^\top\\
				0 & W & I\\
				C & I & 0
			\end{pmatrix}
			\begin{pmatrix}
				\Delta x\\
				\Delta s\\
				\Delta y
			\end{pmatrix}
			=
			-
			\begin{pmatrix}
				Px+q+C^\top y\\
				y+\tau\nabla B(s)\\
				s+c(x)
			\end{pmatrix}.
		\end{equation}
	\end{definition}
	
	\begin{proposition}[Properties of the Nesterov-Todd system]
		The Nesterov-Todd system is admissible, its block \(W\) is symmetric positive
		definite, and it is tangent on the central path.
	\end{proposition}
	
	\begin{proof}
		The nonlinear residuals are the barrier optimality residuals, so the system
		is admissible. By construction, \(W\) is self-adjoint and positive definite.
		
		On the central path,
		\[
		z=C\tau s^{-1}.
		\]
		Logarithmic homogeneity differentiated twice gives at $\lambda =1$:
		\[
		\nabla^2B(s)s= -\nabla B(s)= Cs^{-1},
		\]
		and hence
		\[
		\tau\nabla^2B(s)s=z.
		\]
		The operator \(\tau\nabla^2B(s)\) is a self-adjoint positive-definite
		automorphism of \(K\) mapping \(s\) to \(z\). By uniqueness of the
		Nesterov-Todd scaling,
		\[
		W=\tau\nabla^2B(s).
		\]
		Therefore, the system is tangent.
	\end{proof}
	
	\begin{proposition}[Short-step complexity with Nesterov-Todd scaling,
		\cite{NesterovTodd1997}]
		For a symmetric cone, the short-step method using the Nesterov-Todd system
		reaches a duality gap of \(\varepsilon\) in
		\[
		O\left(
		\sqrt{\nu}\log\frac1{\varepsilon}
		\right)
		\]
		iterations, with one Newton step per parameter update.
	\end{proposition}
	
	\section{Predictor-corrector methods}
	
	The short-step schedule has a polynomial complexity bound, but its reduction
	factor approaches \(1\) as \(\nu\) grows. Practical solvers therefore use
	more aggressive parameter updates.
	
	For example, if \(K=\R_+^{1000}\), then \(\nu=1000\). With
	\(\gamma=1/16\), the short-step reduction factor is
	\[
	1-\frac{\gamma}{\sqrt{\nu}}
	\approx
	0.998.
	\]
	Reducing the parameter by eight orders of magnitude then requires about
	\(9300\) iterations. Practical primal-dual solvers usually require only a
	few tens of iterations, although this faster behavior is not described by
	the short-step worst-case bound.
	
	\begin{definition}[Fast Newton system]
		A tangent Newton system is called \emph{fast} if its matrix block \(H\)
		does not depend explicitly on the target parameter \(\tau\). Changing
		\(\tau\) then changes only the right-hand side, so several directions can be
		computed using the same matrix factorization.
		
		The averaged, Jordan, and Nesterov-Todd systems are fast. The exact barrier
		system, whose block is
		\[
		H=\tau\nabla^2B(s),
		\]
		is not fast.
	\end{definition}
	
	\begin{definition}[Affine predictor step]
		The \emph{affine predictor step} is obtained by setting the target
		complementarity parameter to
		\[
		\tau=0
		\]
		in a fast Newton system.
	\end{definition}
	
	\begin{definition}[Maximum affine step lengths]
		Let
		\[
		(\Delta x^{\mathrm a},
		\Delta s^{\mathrm a},
		\Delta y^{\mathrm a},
		\Delta z^{\mathrm a})
		\]
		be the affine predictor direction. Define
		\begin{align}
			\alpha_{\mathrm p}^{\mathrm a}
			&=
			\sup
			\left\{
			\alpha\in[0,1]:
			s+\alpha\Delta s^{\mathrm a}\in K
			\right\},\\
			\alpha_{\mathrm d}^{\mathrm a}
			&=
			\sup
			\left\{
			\alpha\in[0,1]:
			z+\alpha\Delta z^{\mathrm a}\in K^*
			\right\}.
		\end{align}
	\end{definition}
	
	\begin{definition}[Affine complementarity ratio]
		Let
		\[
		\mu=\frac{\inner{s}{z}}{\nu}
		\]
		and define
		\[
		\mu^{\mathrm a}
		=
		\frac1\nu
		\inner{
			s+\alpha_{\mathrm p}^{\mathrm a}\Delta s^{\mathrm a}
		}{
			z+\alpha_{\mathrm d}^{\mathrm a}\Delta z^{\mathrm a}
		}.
		\]
		The \emph{affine complementarity ratio} is
		\begin{equation}
			r=\frac{\mu^{\mathrm a}}{\mu}.
		\end{equation}
	\end{definition}
	
	If $r$ is small intuitively the affine step worked and we should trust it otherwise when it is near $1$ we shouldn't trust it  as it didn't made the situation better. A way to transform this intuition into calculation is Mehrotra's heuristic:
	
	\begin{definition}[Mehrotra centering rule, \cite{Mehrotra1992}]
		Let
		\[
		0<\sigma_{\min}<\sigma_{\max}<1.
		\]
		The Mehrotra centering parameter is
		\begin{equation}
			\sigma
			=
			\min
			\left\{
			\sigma_{\max},
			\max\{\sigma_{\min},r^3\}
			\right\},
		\end{equation}
		and the target complementarity is
		\[
		\tau=\sigma\mu.
		\]
	\end{definition}
	
	
	\begin{definition}[Corrector term]
		For a symmetric cone the exact expansion is
		\[
		(s+\Delta s)\circ(z+\Delta z)
		= s\circ z + L(z)\Delta s + L(s)\Delta z + \Delta s\circ\Delta z ,
		\]
		in which no inner product convention appears. The Newton linearisation omits
		the final quadratic term, and the \emph{corrector} estimates it by the affine
		predictor $\Delta s^{\mathrm a}\circ\Delta z^{\mathrm a}$. Since the centrality
		target is $s\circ z=C\tau e$ , the
		corrected residual is
		\begin{equation}
			r_4 = s\circ z + \Delta s^{\mathrm a}\circ\Delta z^{\mathrm a}
			- C\sigma\tau\,e ,
		\end{equation}
		the constant $C$ entering through the target alone.
		
		For a nonsymmetric cone there is no Jordan product, so neither the corrector nor
		$C$ is defined; those blocks carry the centring target in the barrier form,
		\begin{equation}
			r_4 = z + \sigma\tau\,\nabla B(s) ,
		\end{equation}
		and receive pure centring.
	\end{definition}
	
	The barrier is finite only on the interior of the cone, so a step must stop
	short of the boundary. The rule that does so is the following, and it is used
	by every interior point method below.

	\begin{definition}[Largest feasible step and fraction to the
		boundary]\label{def:fraction}
		Let $K$ be a closed convex cone, $v\in\intr K$ and $\Delta v$ a direction.
		The \emph{largest feasible step} along $\Delta v$ is
		\begin{equation}
			\bar\alpha(v,\Delta v)
			=\sup\bigl\{\alpha\geq0\ :\ v+\alpha\Delta v\in K\bigr\}
			\ \in\ (0,+\infty] ,
		\end{equation}
		infinite exactly when the ray never leaves the cone. For a parameter
		$\theta\in(0,1)$, the \emph{fraction to the boundary} rule takes the step
		\begin{equation}
			\alpha=\min\bigl(1,\ \theta\,\bar\alpha(v,\Delta v)\bigr) .
		\end{equation}
		On a product cone the step is the smallest over the blocks, and the primal
		and the dual variables each take their own, giving a pair
		$(\alpha_{\mathrm p},\alpha_{\mathrm d})$.
	\end{definition}

	\begin{proposition}[The rule keeps the iterate interior]\label{prop:fraction}
		With $v\in\intr K$, one has $\bar\alpha>0$, and
		\begin{equation}
			v+\alpha\Delta v\in\intr K
			\qquad\text{for every }0\leq\alpha<\bar\alpha .
		\end{equation}
		In particular the step of Definition~\ref{def:fraction} is strictly
		positive and lands strictly inside the cone, since $\theta<1$.
	\end{proposition}


	Taking $\theta$ strictly below one is therefore not a safeguard against a
	numerical accident but the condition under which the next iteration exists at
	all: at $\alpha=\bar\alpha$ the barrier is infinite and its Hessian singular,
	and there would be no scaling to build.

	\begin{algorithm}[H]
		\caption{Predictor-corrector method}
		\begin{algorithmic}[1]
			\Require A fast admissible Newton system, an interior initial point
			\(s^0\in\relint K\), \(z^0\in\relint K^*\), initial points \(x^0,y^0\),
			constants \(0<\sigma_{\min}<\sigma_{\max}<1\), a fraction-to-the-boundary
			parameter \(\theta\in(0,1)\), and a tolerance \(\varepsilon>0\)
			\Ensure An approximate primal-dual solution
			\For{\(k=0,1,2,\ldots\)}
			\State Compute
			\[
			\mu_k=\frac{\inner{s^k}{z^k}}{\nu}
			\]
			\If{\(\nu\mu_k\leq\varepsilon\)}
			\State \Return \((x^k,s^k,y^k,z^k)\)
			\EndIf
			\State Assemble and factorize the Newton matrix
			\State Solve the system with target \(\tau=0\) to obtain the affine
			predictor direction
			\State Compute
			\(\alpha_{\mathrm p}^{\mathrm a}\),
			\(\alpha_{\mathrm d}^{\mathrm a}\),
			\(\mu^{\mathrm a}\), and
			\(r=\mu^{\mathrm a}/\mu_k\)
			\State Set
			\[
			\sigma
			=
			\min
			\left\{
			\sigma_{\max},
			\max\{\sigma_{\min},r^3\}
			\right\}
			\]
			\State Reuse the factorization to solve the corrected system with target
			\(\tau=\sigma\mu_k\)
			\State Compute the maximum primal and dual step lengths
			\(\alpha_{\mathrm p}\) and \(\alpha_{\mathrm d}\)
			\State Update
			\[
			x^{k+1}
			=
			x^k+\theta\alpha_{\mathrm p}\Delta x,
			\qquad
			s^{k+1}
			=
			s^k+\theta\alpha_{\mathrm p}\Delta s
			\]
			\State Update
			\[
			y^{k+1}
			=
			y^k+\theta\alpha_{\mathrm d}\Delta y,
			\qquad
			z^{k+1}
			=
			z^k+\theta\alpha_{\mathrm d}\Delta z
			\]
			\EndFor
		\end{algorithmic}
	\end{algorithm}
	
	\begin{remark}[Theory and practice]
		The short-step method has a polynomial complexity bound but uses conservative
		parameter reductions. The Mehrotra predictor-corrector method uses much more
		aggressive reductions and is usually far faster in practice.
		
		The centering rule and corrector are heuristic. Consequently, the short-step
		complexity proof does not apply directly to the practical predictor-corrector
		algorithm. Its effectiveness is empirical and is supported by extensive
		solver experience rather than by the short-step iteration bound.
	\end{remark}
	
	\chapter{Augmented Lagrangian method basics}
	
	\section{Augmented Lagrangian smoothing family}
	
	Throughout this chapter the problem is
	\begin{equation}
		\min_{x\in\R^n}\ f(x)
		\qquad\text{subject to}\qquad
		c(x)\preceq_K0 ,
	\end{equation}
	with $f(x)=\tfrac12\inner{x}{Px}+\inner{q}{x}$ and $c(x)=Cx+d$, and $s=-c(x)$
	denotes the slack, so that the constraint is $s\in K$.
	
	\begin{definition}[Augmented Lagrangian smoothing]
		Let $\theta=(M,z_c)$ with $M\in\mathcal{A}(K)$ and $z_c\in\R^m$. The
		\emph{augmented Lagrangian smoothing} of $\delta_K$ is
		\begin{equation}
			\phi^{\mathrm{AL}}_\theta(s)
			=\sup_{z\in K^*}\Bigl\{-\inner{z}{s}-\tfrac12\norm{z-z_c}^2_M\Bigr\}
			=\tfrac12\norm{\proj{Mz_c-s}{K^*}}^2_{M^{-1}}
			-\tfrac12\norm{z_c}^2_M .
		\end{equation}
	\end{definition}
	
	\begin{definition}[Moreau smoothing]
		With $\theta=(M,z_c)$ as above, the \emph{Moreau smoothing} of $\delta_K$ is
		the Moreau envelope of $\delta_K$ in the metric $M^{-1}$ evaluated at the
		shifted point,
		\begin{equation}
			\phi^{\mathrm{Mor}}_\theta(s)
			=e^{M^{-1}}_{\delta_K}\bigl(s-Mz_c\bigr)
			=\tfrac12\dist_{M^{-1}}\bigl(s-Mz_c,\ K\bigr)^2 .
		\end{equation}
	\end{definition}
	
	
	\begin{proposition}[The two smoothings differ by a constant]
		Let \(\theta=(M,z_c)\). Then, for every \(s\in\R^m\),
		\begin{equation}
			\phi^{\mathrm{Mor}}_\theta(s)
			=
			\phi^{\mathrm{AL}}_\theta(s)
			+
			\frac12\norm{z_c}_M^2.
		\end{equation}
	\end{proposition}
	
	\begin{proof}
		Since M is admissible $M^{-1}$ too hence in this metric: 
		$$ \dist(s-Mz_c,K) = \|\Pi_{K^*}(Mz_c-s)\|_{M^{-1}}$$
		
	\end{proof}
	
	\begin{proposition}[Optimality conditions for the smoothed subproblem]
		Let
		\[
		\theta=(M,z_c),
		\]
		and consider
		\begin{equation}
			\min_x
			\left\{
			f(x)+\phi_\theta\bigl(-c(x)\bigr)
			\right\},
			\label{eq:smoothed-primal-subproblem}
		\end{equation}
		where \(\phi_\theta\) is either
		\(\phi_\theta^{\mathrm{Mor}}\) or
		\(\phi_\theta^{\mathrm{AL}}\).
		
		Then \(x\) minimizes \eqref{eq:smoothed-primal-subproblem} if and only if
		there exists \(z\) such that
		\begin{equation}
			\nabla f(x)+C^\top z=0,
			\qquad
			z
			=
			M^{-1}
			\proj{Mz_c+c(x)}{K^*}.
		\end{equation}
	\end{proposition}
	
	\begin{proof}
		Consequence of the general smoothing optimality condition, combined with the fact that :
		$$ \nabla \phi_{\theta}(u) = \nabla (\frac{1}{2} \|\Pi_{K^*}\|^2_{M^{-1}}(Mz_c-.)+ cst)(u) = -M^{-1}\Pi_{K^*}(Mz_c-u)$$
	\end{proof}
	\begin{algorithm}[H]
		\caption{Augmented Lagrangian method}
		\begin{algorithmic}[1]
			\Require a sequence $(M_k)$ in $\mathcal{A}(K)$, an initial $z_c^0\in K^*$,
			inner tolerances $(\varepsilon_k)$
			\Ensure an approximate solution of the problem
			\For{$k=0,1,2,\dots$}
			\State solve, to accuracy $\varepsilon_k$ and by the semismooth Newton method,
			\[
			x^{k+1}\in\argmin_x\ f(x)+\phi_{(M_k,z_c^k)}\bigl(-c(x)\bigr)
			\]
			\State $z^{k+1}\gets M_k^{-1}\proj{M_kz_c^k+c(x^{k+1})}{K^*}$
			\State $z_c^{k+1}\gets z^{k+1}$
			\EndFor
		\end{algorithmic}
	\end{algorithm}
	
	\section{The augmented Lagrangian method as a proximal point algorithm}
	
	\begin{definition}[Dual optimality operator]
		Let \(g\) be the dual function, and assume that the dual problem admits at
		least one point with finite dual value. Define
		\begin{equation}
			h(z)
			=
			-g(z)+\delta_{K^*}(z).
		\end{equation}
		Then \(h\) is closed, proper, and convex, and its minimizers are precisely
		the solutions of the dual problem.
		
		The \emph{dual optimality operator} is
		\begin{equation}
			T_{\mathrm D}
			=
			\partial h
			=
			\partial\bigl(-g+\delta_{K^*}\bigr).
		\end{equation}
	\end{definition}
	
	\begin{proposition}[Maximal monotonicity of the dual optimality operator]
		The operator \(T_{\mathrm D}\) is maximal monotone, and
		\begin{equation}
			T_{\mathrm D}^{-1}(0)
			=
			\argmax_{z\in K^*}g(z).
		\end{equation}
	\end{proposition}
	
	\begin{proof}
		The dual function \(g\) is concave and upper semicontinuous. Therefore,
		\(-g\) is closed and convex. Since \(K^*\) is closed and convex,
		\(\delta_{K^*}\) is also closed and convex. By the assumed existence of a
		dual feasible point with finite dual value,
		\[
		h=-g+\delta_{K^*}
		\]
		is proper.
		
		The subdifferential of a closed proper convex function is maximal monotone.
		Hence
		\[
		T_{\mathrm D}=\partial h
		\]
		is maximal monotone.
		
		Finally, Fermat's rule gives
		\[
		0\in T_{\mathrm D}(z)
		\quad\Longleftrightarrow\quad
		z\in\argmin h.
		\]
		Since
		\[
		h(z)
		=
		\begin{cases}
			-g(z),&z\in K^*,\\
			+\infty,&z\notin K^*,
		\end{cases}
		\]
		minimizing \(h\) is equivalent to maximizing \(g\) over \(K^*\).
	\end{proof}
	
	\begin{proposition}[The method of multipliers as a proximal point method,
		\cite{Rockafellar1976b}]
		Let \(M_k\in\mathcal A(K)\). Suppose that the augmented Lagrangian
		subproblem is solved exactly at iteration \(k\), producing
		\((x^{k+1},z^{k+1})\). Then the multiplier update satisfies
		\begin{equation}
			z^{k+1}
			=
			J_{T_{\mathrm D}}^{M_k}(z^k).
		\end{equation}
		Thus the multiplier sequence is the variable-metric proximal point
		iteration applied to the dual optimality operator \(T_{\mathrm D}\).
	\end{proposition}
	
	\begin{proof}
		Fix an iteration and write
		\[
		M=M_k,
		\qquad
		z_c=z^k,
		\qquad
		x^+=x^{k+1},
		\qquad
		z^+=z^{k+1}.
		\]
		The optimality conditions of the augmented Lagrangian subproblem give
		\begin{equation}
			\nabla f(x^+)+C^\top z^+=0
			\label{eq:alm-stationarity}
		\end{equation}
		and
		\begin{equation}
			z^+
			=
			M^{-1}
			\proj{Mz_c+c(x^+)}{K^*}.
			\label{eq:alm-multiplier-update}
		\end{equation}
		
		Equation \eqref{eq:alm-stationarity} says that \(x^+\) minimizes
		\(L(\cdot,z^+)\). Consequently,
		\[
		g(z^+)=L(x^+,z^+).
		\]
		For every \(z\), the definition of the dual function gives
		\begin{align}
			g(z)
			&\leq
			L(x^+,z)\\
			&=
			L(x^+,z^+)
			+
			\inner{z-z^+}{c(x^+)}\\
			&=
			g(z^+)
			+
			\inner{z-z^+}{c(x^+)}.
		\end{align}
		Thus \(c(x^+)\) is a supergradient of \(g\) at \(z^+\), or equivalently,
		\begin{equation}
			-c(x^+)\in\partial(-g)(z^+).
			\label{eq:dual-function-subgradient}
		\end{equation}
		
		By admissibility,
		\eqref{eq:alm-multiplier-update} is equivalent to
		\[
		z^+
		=
		\proj{z_c+M^{-1}c(x^+)}{K^*}^{M}.
		\]
		The metric projection characterization therefore gives
		\[
		M\bigl(
		z_c+M^{-1}c(x^+)-z^+
		\bigr)
		\in
		N_{K^*}(z^+).
		\]
		Equivalently,
		\begin{equation}
			M(z_c-z^+)+c(x^+)
			\in
			\partial\delta_{K^*}(z^+).
			\label{eq:normal-component}
		\end{equation}
		
		Adding \eqref{eq:dual-function-subgradient} and
		\eqref{eq:normal-component}, we obtain
		\[
		M(z_c-z^+)
		\in
		\partial(-g)(z^+)
		+
		\partial\delta_{K^*}(z^+)
		\subseteq
		\partial h(z^+).
		\]
		Therefore,
		\[
		M(z_c-z^+)\in T_{\mathrm D}(z^+),
		\]
		which is equivalent to
		\[
		z^+
		=
		J_{T_{\mathrm D}}^M(z_c).
		\]
	\end{proof}
	
	\section{The proximal point method on the KKT operator}
	
	\begin{definition}[KKT operator]
		For \(v=(x,z)\in\R^{n+m}\), define
		\begin{equation}
			T(x,z)
			=
			\begin{pmatrix}
				Px+q+C^\top z\\[1mm]
				-c(x)+N_{K^*}(z)
			\end{pmatrix}.
			\label{eq:kkt-operator}
		\end{equation}
		Equivalently,
		\[
		T(x,z)
		=
		\begin{pmatrix}
			Px+q+C^\top z\\[1mm]
			-c(x)+\partial\delta_{K^*}(z)
		\end{pmatrix}.
		\]
	\end{definition}
	
	\begin{proposition}[Maximal monotonicity and zeros]
		The KKT operator \(T\) is maximal monotone. Moreover,
		\[
		0\in T(x,z)
		\]
		if and only if \((x,z)\) satisfies the KKT conditions
		\begin{align}
			Px+q+C^\top z&=0,\\
			-c(x)&\in K,\\
			z&\in K^*,\\
			\inner{z}{c(x)}&=0.
		\end{align}
	\end{proposition}
	
	\begin{proof}
		Write
		\[
		T=T_1+T_2,
		\]
		where
		\[
		T_1(x,z)
		=
		\begin{pmatrix}
			P&C^\top\\
			-C&0
		\end{pmatrix}
		\begin{pmatrix}
			x\\z
		\end{pmatrix}
		+
		\begin{pmatrix}
			q\\-d
		\end{pmatrix}
		\]
		and
		\[
		T_2(x,z)
		=
		\begin{pmatrix}
			0\\
			N_{K^*}(z)
		\end{pmatrix}.
		\]
		
		The symmetric part of the linear map defining \(T_1\) is
		\[
		\begin{pmatrix}
			P&0\\
			0&0
		\end{pmatrix}
		\succeq0.
		\]
		Indeed,
		\[
		\inner{T_1(v)-T_1(v')}{v-v'}
		=
		\inner{x-x'}{P(x-x')}
		\geq0.
		\]
		Thus \(T_1\) is monotone. Since it is affine and defined on the whole space,
		it is maximal monotone.
		
		The function
		\[
		(x,z)\longmapsto\delta_{K^*}(z)
		\]
		is closed, proper, and convex, and its subdifferential is \(T_2\). Hence
		\(T_2\) is maximal monotone. Since \(T_1\) has full domain, the maximal
		monotone sum theorem implies that \(T=T_1+T_2\) is maximal monotone.
		
		Now \(0\in T(x,z)\) means
		\[
		Px+q+C^\top z=0
		\]
		and
		\[
		c(x)\in N_{K^*}(z).
		\]
		For \(z\in K^*\),
		\[
		N_{K^*}(z)
		=
		(K^*)^\circ\cap z^\perp
		=
		-K\cap z^\perp.
		\]
		Therefore,
		\[
		c(x)\in N_{K^*}(z)
		\]
		is equivalent to
		\[
		-c(x)\in K,
		\qquad
		z\in K^*,
		\qquad
		\inner{z}{c(x)}=0.
		\]
		These are precisely the KKT conditions.
	\end{proof}
	
	\begin{algorithm}[H]
		\caption{Proximal method of multipliers}
		\begin{algorithmic}[1]
			\Require Sequences \((\rho_k)\) with \(\rho_k>0\), admissible metrics
			\(M_k\in\mathcal A(K)\), an initial center
			\((x_c^0,z_c^0)\), and summable errors \((\varepsilon_k)\)
			\Ensure An approximate point of \(T^{-1}(0)\)
			\For{\(k=0,1,2,\ldots\)}
			\State Set
			\[
			\mathcal M_k
			=
			\diag(\rho_kI_n,M_k)
			\]
			\State Compute \((x^{k+1},z^{k+1})\) satisfying
			\[
			\norm{
				(x^{k+1},z^{k+1})
				-
				J_T^{\mathcal M_k}(x_c^k,z_c^k)
			}_{\mathcal M_k}
			\leq
			\varepsilon_k
			\]
			\State Update
			\[
			(x_c^{k+1},z_c^{k+1})
			\gets
			(x^{k+1},z^{k+1})
			\]
			\EndFor
		\end{algorithmic}
	\end{algorithm}
	
	\section{Primal formulation of a proximal step}
	
	\begin{definition}[Augmented Lagrangian penalty]
		Let \(M\in\mathcal A(K)\) and \(z_c\in\R^m\). Define
		\begin{equation}
			\phi^{\mathrm{AL}}_{(M,z_c)}(s)
			=
			\frac12
			\norm{
				\proj{Mz_c-s}{K^*}
			}_{M^{-1}}^2
			-
			\frac12\norm{z_c}_M^2.
			\label{eq:augmented-lagrangian-penalty}
		\end{equation}
		Then
		\begin{equation}
			\nabla
			\phi^{\mathrm{AL}}_{(M,z_c)}(s)
			=
			-
			M^{-1}
			\proj{Mz_c-s}{K^*}.
		\end{equation}
	\end{definition}
	
	\begin{proposition}[Primal formulation of one proximal step]
		Let
		\[
		\rho>0,
		\qquad
		M\in\mathcal A(K),
		\qquad
		\mathcal M=\diag(\rho I_n,M),
		\]
		and let
		\[
		v_c=(x_c,z_c).
		\]
		Define
		\begin{equation}
			\varphi_c(x)
			=
			f(x)
			+
			\frac{\rho}{2}\norm{x-x_c}^2
			+
			\phi^{\mathrm{AL}}_{(M,z_c)}\bigl(-c(x)\bigr).
			\label{eq:primal-proximal-objective}
		\end{equation}
		Then
		\[
		v^+=(x^+,z^+)=J_T^{\mathcal M}(v_c)
		\]
		if and only if \(x^+\) minimizes \(\varphi_c\) and
		\begin{equation}
			z^+
			=
			M^{-1}
			\proj{Mz_c+c(x^+)}{K^*}.
			\label{eq:recovered-dual-variable}
		\end{equation}
	\end{proposition}
	
	\begin{proof}
		The resolvent inclusion
		\[
		\mathcal M(v_c-v^+)\in T(v^+)
		\]
		is equivalent to
		\begin{align}
			\rho(x_c-x^+)
			&=
			Px^++q+C^\top z^+,
			\label{eq:proximal-first-row}\\
			M(z_c-z^+)
			&\in
			-c(x^+)+N_{K^*}(z^+).
			\label{eq:proximal-second-row}
		\end{align}
		
		Set
		\[
		w=z_c+M^{-1}c(x^+).
		\]
		The second inclusion becomes
		\[
		M(w-z^+)\in N_{K^*}(z^+).
		\]
		By the characterization of the metric projection,
		\[
		z^+
		=
		\proj{w}{K^*}^{M}.
		\]
		Admissibility gives
		\[
		\proj{w}{K^*}^{M}
		=
		\proj{w}{K^*}
		=
		M^{-1}\proj{Mw}{K^*}.
		\]
		Since
		\[
		Mw=Mz_c+c(x^+),
		\]
		we obtain \eqref{eq:recovered-dual-variable}.
		
		Now set
		\[
		\sigma(x)=Mz_c+c(x).
		\]
		Differentiating \eqref{eq:primal-proximal-objective} gives
		\[
		\nabla\varphi_c(x)
		=
		Px+q
		+
		\rho(x-x_c)
		+
		C^\top M^{-1}\proj{\sigma(x)}{K^*}.
		\]
		Using \eqref{eq:recovered-dual-variable}, the condition
		\[
		\nabla\varphi_c(x^+)=0
		\]
		is precisely \eqref{eq:proximal-first-row}. Since \(\varphi_c\) is strongly
		convex, this stationarity condition is equivalent to \(x^+\) being its
		unique minimizer.
	\end{proof}
	
	\begin{proposition}[Properties of the primal objective]
		The function \(\varphi_c\) is finite and continuously differentiable on
		\(\R^n\). It is strongly convex with modulus \(\rho\), and therefore has a
		unique minimizer.
		
		Its gradient is
		\begin{equation}
			\nabla\varphi_c(x)
			=
			Px+q
			+
			\rho(x-x_c)
			+
			C^\top M^{-1}\proj{\sigma(x)}{K^*},
			\qquad
			\sigma(x)=Mz_c+c(x).
			\label{eq:primal-objective-gradient}
		\end{equation}
		For every
		\[
		F\in\partial\proj{\cdot}{K^*}(\sigma(x)),
		\]
		the matrix
		\begin{equation}
			G_F(x)
			=
			P+\rho I+C^\top M^{-1}FC
			\label{eq:primal-generalized-hessian}
		\end{equation}
		belongs to
		\[
		\partial\bigl(\nabla\varphi_c\bigr)(x)
		\]
		and satisfies
		\begin{equation}
			G_F(x)\succeq\rho I.
		\end{equation}
		At a point where the projection is differentiable,
		\[
		\nabla^2\varphi_c(x)
		=
		P+\rho I+C^\top M^{-1}FC,
		\]
		with \(F=D\proj{\cdot}{K^*}(\sigma(x))\).
	\end{proposition}
	
	\begin{proof}
		The squared projection potential is finite and continuously differentiable.
		Consequently, \(\varphi_c\) is finite and continuously differentiable, with
		gradient \eqref{eq:primal-objective-gradient}.
		
		The function \(f\) is convex, the augmented Lagrangian penalty is convex,
		and
		\[
		x\longmapsto\frac{\rho}{2}\norm{x-x_c}^2
		\]
		is strongly convex with modulus \(\rho\). Thus \(\varphi_c\) is strongly
		convex with modulus \(\rho\).
		
		The Clarke chain rule gives
		\[
		G_F(x)
		\in
		\partial\bigl(\nabla\varphi_c\bigr)(x).
		\]
		Admissibility implies that \(M^{-1}\) commutes with \(F\), and
		\[
		M^{-1}F
		\]
		is symmetric positive semidefinite. Hence
		\[
		C^\top M^{-1}FC\succeq0.
		\]
		Since \(P\succeq0\),
		\[
		G_F(x)
		=
		P+\rho I+C^\top M^{-1}FC
		\succeq
		\rho I.
		\]
	\end{proof}
	
	The semismooth Newton direction is therefore defined by
	\begin{equation}
		\bigl(
		P+\rho I+C^\top M^{-1}FC
		\bigr)\Delta x
		=
		-\nabla\varphi_c(x),
		\qquad
		F\in\partial\proj{\cdot}{K^*}(\sigma(x)).
		\label{eq:primal-semismooth-newton-system}
	\end{equation}
	
	\begin{algorithm}[H]
		\caption{Semismooth Newton method for the primal proximal problem}
		\begin{algorithmic}[1]
			\Require Parameters \(\rho,M\), centers \((x_c,z_c)\), an initial point
			\(x^0\), a merit function, a line-search rule, and a tolerance
			\(\varepsilon>0\)
			\Ensure A point \(x\) satisfying
			\(\norm{\nabla\varphi_c(x)}\leq\varepsilon\)
			\For{\(j=0,1,2,\ldots\)}
			\If{\(\norm{\nabla\varphi_c(x^j)}\leq\varepsilon\)}
			\State \Return \(x^j\)
			\EndIf
			\State Set
			\[
			\sigma^j=Mz_c+c(x^j)
			\]
			\State Choose
			\[
			F_j\in\partial\proj{\cdot}{K^*}(\sigma^j)
			\]
			\State Solve
			\[
			\bigl(
			P+\rho I+C^\top M^{-1}F_jC
			\bigr)\Delta x^j
			=
			-\nabla\varphi_c(x^j)
			\]
			\State Compute a step length \(\alpha_j\) using the line-search rule
			\State Set
			\[
			x^{j+1}=x^j+\alpha_j\Delta x^j
			\]
			\EndFor
		\end{algorithmic}
	\end{algorithm}
	
	\section{Infeasibility certificates}
	
	\begin{proposition}[Primal alternative]
		Exactly one of the following statements holds:
		
		\begin{enumerate}
			\item
			\[
			-d\in\operatorname{cl}\bigl(K+\range C\bigr).
			\]
			
			\item There exists \(z\in\R^m\) such that
			\[
			z\in K^*,
			\qquad
			C^\top z=0,
			\qquad
			\inner{z}{d}>0.
			\]
		\end{enumerate}
		
		If the primal problem is feasible, then the first statement holds.
		Consequently, the second statement certifies primal infeasibility.
	\end{proposition}
	
	\begin{proof}
		Set
		\[
		S=\operatorname{cl}\bigl(K+\range C\bigr).
		\]
		Then \(S\) is a closed convex cone. We first observe that
		\begin{equation}
			S^*
			=
			\left\{
			z\in K^*:C^\top z=0
			\right\}.
			\label{eq:dual-sum-cone}
		\end{equation}
		Indeed, if \(z\in S^*\), then \(K\subseteq S\) implies \(z\in K^*\).
		Moreover, \(\range C\subseteq S\) is a subspace. Hence both \(Cx\) and
		\(-Cx\) belong to \(S\), so
		\[
		\inner{z}{Cx}=0
		\qquad
		\text{for every }x.
		\]
		Thus \(C^\top z=0\). The converse follows immediately from
		\[
		\inner{z}{k+Cx}
		=
		\inner{z}{k}
		\geq0
		\]
		for \(k\in K\).
		
		The two alternatives cannot hold simultaneously. Indeed, if
		\(-d\in S\) and \(z\in S^*\), then
		\[
		\inner{z}{-d}\geq0,
		\]
		so
		\[
		\inner{z}{d}\leq0.
		\]
		
		Suppose now that \(-d\notin S\). Suppose that for any $z\in S^*$ we have that $\langle z,-d\rangle \geq 0$ hence that $S^* \subset  \{-d\}^*$ thus taking the dual given that $S$ is closed and convex and dual inverse monotonicity:
		$$ \{-d\} \in S^{**} =  S  $$
		
		Which is contradictory.Thus there exist $z$ such that:
		\[
		z\in S^*
		\]
		and
		\[
		\inner{z}{d}>0.
		\]
		Using \eqref{eq:dual-sum-cone}, we obtain
		\[
		z\in K^*,
		\qquad
		C^\top z=0.
		\]
		
		Finally, if \(x\) is primal feasible, then
		\[
		-(Cx+d)\in K,
		\]
		and therefore
		\[
		-d
		=
		-(Cx+d)+Cx
		\in
		K+\range C
		\subseteq S.
		\]
	\end{proof}
	
	\begin{proposition}[Improving ray]
		Suppose that there exists \(\delta\in\R^n\) satisfying
		\begin{equation}
			P\delta=0,
			\qquad
			-C\delta\in K,
			\qquad
			\inner{q}{\delta}<0.
			\label{eq:dual-infeasibility-ray}
		\end{equation}
		Then the dual problem is infeasible. If the primal problem is feasible, then
		\[
		p^\star=-\infty.
		\]
	\end{proposition}
	
	\begin{proof}
		Suppose, for contradiction, that \(z\in K^*\) is dual feasible. Then there
		exists \(x_z\) satisfying
		\[
		Px_z+q+C^\top z=0.
		\]
		Taking the inner product with \(\delta\), and using \(P\delta=0\), gives
		\begin{align}
			0
			&=
			\inner{Px_z+q+C^\top z}{\delta}\\
			&=
			\inner{q}{\delta}
			+
			\inner{z}{C\delta}\\
			&=
			\inner{q}{\delta}
			-
			\inner{z}{-C\delta}.
		\end{align}
		Since
		\[
		z\in K^*
		\qquad\text{and}\qquad
		-C\delta\in K,
		\]
		we have
		\[
		\inner{z}{-C\delta}\geq0.
		\]
		It follows that
		\[
		0
		\leq
		\inner{q}{\delta},
		\]
		contradicting \(\inner{q}{\delta}<0\). Thus the dual problem is infeasible.
		
		Now suppose that \(x\) is primal feasible. For every \(t\geq0\),
		\begin{align}
			-\bigl(C(x+t\delta)+d\bigr)
			&=
			-(Cx+d)+t(-C\delta)\\
			&\in K,
		\end{align}
		because \(K\) is a convex cone. Hence \(x+t\delta\) remains feasible.
		
		Since \(P\delta=0\) and \(P\) is symmetric,
		\[
		f(x+t\delta)
		=
		f(x)+t\inner{q}{\delta}.
		\]
		Because \(\inner{q}{\delta}<0\),
		\[
		f(x+t\delta)\longrightarrow-\infty
		\qquad
		\text{as }t\to+\infty.
		\]
		Therefore,
		\[
		p^\star=-\infty.
		\]
	\end{proof}
	
	\begin{definition}[Infeasibility certificates]\label{def:certificates}
		A \emph{primal infeasibility certificate} is a vector \(z\in\R^m\)
		satisfying
		\begin{equation}
			z\in K^*,
			\qquad
			C^\top z=0,
			\qquad
			\inner{z}{d}>0.
		\end{equation}
		
		A \emph{dual infeasibility certificate} is a vector
		\(\delta\in\R^n\) satisfying
		\begin{equation}
			P\delta=0,
			\qquad
			-C\delta\in K,
			\qquad
			\inner{q}{\delta}<0.
		\end{equation}
		
		They are the certificates used in the algorithm for testing infeasibility.
	\end{definition}
	
	\section{Asymptotic displacement and certificates}
	
	\begin{proposition}[Asymptotic displacement, \cite{Pazy1971}]
		Let \(T\) be maximal monotone, let \(\mathcal M\succ0\), and define
		\[
		v^{k+1}=J_T^{\mathcal M}(v^k).
		\]
		Then the displacement sequence
		\[
		v^{k+1}-v^k
		\]
		converges to a limit \(\delta\) that is independent of the initial point
		\(v^0\). If \(T^{-1}(0)\neq\emptyset\), then
		\[
		\delta=0.
		\]
	\end{proposition}
	
	\begin{proposition}[Decomposition of the asymptotic displacement]
		Let
		\[
		\mathcal M=\diag(\rho I_n,M),
		\qquad
		\rho>0,
		\qquad
		M\succ0,
		\]
		and let \(T\) be the KKT operator. Write
		\[
		\delta=(\delta_x,\delta_z)
		\]
		for the asymptotic displacement of the proximal iteration. Then
		\begin{align}
			P\delta_x&=0,
			&
			-C\delta_x&\in K,
			\label{eq:displacement-primal}\\
			C^\top\delta_z&=0,
			&
			\delta_z&\in K^*,
			\label{eq:displacement-dual}\\
			\inner{C\delta_x}{\delta_z}&=0,
			\label{eq:displacement-orthogonality}\\
			\inner{q}{\delta_x}
			&=
			-\rho\norm{\delta_x}^2,
			\label{eq:displacement-q}\\
			\inner{d}{\delta_z}
			&=
			\norm{\delta_z}_M^2.
			\label{eq:displacement-d}
		\end{align}
		Consequently:
		
		\begin{enumerate}
			\item If \(\delta_x\neq0\), then \(\delta_x\) is a dual infeasibility
			certificate.
			
			\item If \(\delta_z\neq0\), then \(\delta_z\) is a primal infeasibility
			certificate.
			
			\item If \(\delta\neq0\), then at least one infeasibility certificate is
			obtained.
		\end{enumerate}
	\end{proposition}
	
	\begin{remark}
		The theorem is only expressed here for fixed metric but we use it for variable metric.
	\end{remark}
	\begin{remark}[Vanishing displacement does not imply solvability]
		The preceding proposition gives certificates only when the asymptotic
		displacement is nonzero. Although
		\[
		T^{-1}(0)\neq\emptyset
		\quad\Longrightarrow\quad
		\delta=0,
		\]
		the converse is false in general. It is possible to have
		\[
		\delta=0
		\qquad\text{and}\qquad
		T^{-1}(0)=\emptyset.
		\]
		This situation may occur for weakly infeasible problems or for problems
		that are unbounded along a curved feasible path but admit no improving
		recession direction. In the algorithm we use the implication:
		$$\delta \neq 0\implies T^{-1}(0) = \emptyset$$
		To detect infeasibility. 
	\end{remark}
	
	\part{Method construction}
	
	\chapter{Getting a Newton}
	
	\section{First systems}
	This section derives the Newton system of the solver's semismooth method. We
	first recall the optimality conditions of the problem.
	
	\begin{proposition}[Optimality conditions of one proximal step]
		Let $\rho>0$, $M\in\mathcal{A}(K)$, $\mathcal{M}=\diag(\rho I_n,M)$ and
		$v_c=(x_c,z_c)$. Then $v=(x,z)=J^{\mathcal{M}}_T(v_c)$ if and only if
		\begin{align}
			r_{\mathrm{stat}}&=Px+q+\rho(x-x_c)+C^\top z=0, \\
			r_{\mathrm{cone}}&=Mz-\proj{\sigma}{K^*}=0,
			\qquad \sigma=Mz_c+c(x) .
		\end{align}
	\end{proposition}
	
	\begin{equation}
		Px+q+\rho(x-x_c)+C^\top M^{-1}\proj{Mz_c+c(x)}{K^*}=0 .
	\end{equation}
	
	\begin{proposition}[The system obtained directly from the two residuals]
		Let $F\in\partial\proj{\cdot}{K^*}(\sigma)$. The semismooth Newton step for
		$\bigl(r_{\mathrm{stat}},r_{\mathrm{cone}}\bigr)$ solves
		\begin{equation}
			\begin{pmatrix}
				P+\rho I & C^\top\\
				-FC & M
			\end{pmatrix}
			\begin{pmatrix}\Delta x\\ \Delta z\end{pmatrix}
			=-\begin{pmatrix}r_{\mathrm{stat}}\\ r_{\mathrm{cone}}\end{pmatrix} ,
		\end{equation}
	\end{proposition}
	
	\begin{proof}
		Differentiating gives $$\partial r_{\mathrm{stat}}/\partial x=P+\rho I$$ and
		$$\partial r_{\mathrm{stat}}/\partial z=C^\top$$; for the second residual
		$$\partial\sigma/\partial x=C$$, so
		$$\partial r_{\mathrm{cone}}/\partial x=-FC$$ and
		$$\partial r_{\mathrm{cone}}/\partial z=M$$. The diagonal blocks are symmetric.
		
	\end{proof}
	
	The trouble with the preceding matrix is that it is not symmetric, and that
	its sparsity pattern changes at every step because of the product $FC$. The
	first of these is repaired by a residual carrying $F$ in its own definition.
	
	
	\begin{definition}[Symmetrised stationarity residual]
		With $\sigma=Mz_c+c(x)$ and $F\in\partial\proj{\cdot}{K^*}(\sigma)$,
		\begin{equation}
			r^{*}_{\mathrm{stat}}=r_{\mathrm{stat}}+C^\top(F-I)z
			=Px+q+\rho(x-x_c)+C^\top Fz .
		\end{equation}
	\end{definition}
	Let recall the following lemma that we proved before:
	\begin{lemma}[$F$ fixes the projection]
		Let $\sigma\in\R^m$, $p=\proj{\sigma}{K^*}$ and
		$F\in\partial\proj{\cdot}{K^*}(\sigma)$. Then $Fp=p$.
	\end{lemma}
	
	\begin{proposition}[The two stationarity residuals agree at a cone zero]
		If $r_{\mathrm{cone}}=0$ then $Fz=z$, and consequently
		$r^{*}_{\mathrm{stat}}=r_{\mathrm{stat}}$.
	\end{proposition}
	
	\begin{proof}
		$r_{\mathrm{cone}}=0$ reads $Mz=p$. Since $M\in\mathcal{A}(K)$ commutes with
		$F$, so does $M^{-1}$, and the lemma gives
		$Fz=FM^{-1}p=M^{-1}Fp=M^{-1}p=z$. Then
		$r^{*}_{\mathrm{stat}}-r_{\mathrm{stat}}=C^\top(F-I)z=0$.
	\end{proof}
	
	\begin{proposition}[The symmetrised system]
		Let $V=C^\top D^2\proj{\cdot}{K^*}(\sigma)[\,\cdot\,,\,\cdot\,,z]\,C$. The
		semismooth Newton step for
		$\bigl(r^{*}_{\mathrm{stat}},r_{\mathrm{cone}}\bigr)$, with the second block
		row negated, solves
		\begin{equation}
			\begin{pmatrix}
				P+\rho I+V & C^\top F\\
				FC & -M
			\end{pmatrix}
			\begin{pmatrix}\Delta x\\ \Delta z\end{pmatrix}
			=\begin{pmatrix}-r^{*}_{\mathrm{stat}}\\ r_{\mathrm{cone}}\end{pmatrix} ,
		\end{equation}
		and this matrix is symmetric.
	\end{proposition}
	
	\begin{proof}
		The $z$ derivative of $r^{*}_{\mathrm{stat}}$ is $C^\top F$, since $\sigma$
		does not depend on $z$. Its $x$ derivative is
		$$P+\rho I+C^\top\bigl(DF(\sigma)[C\,\cdot\,]\bigr)z$$, and
		$DF(\sigma)=D^2\proj{\cdot}{K^*}(\sigma)$ contracted against $z$ in its last
		argument gives $V$. The second row is unchanged from the previous proposition
		up to the negation. Symmetry: $C^\top F$ and $FC$ are transposes since
		$F=F^\top$, $M$ is symmetric, and $V$ is symmetric because a second derivative
		is symmetric in its first two arguments.
	\end{proof}
	
	\begin{remark}
		Because \(r_{\mathrm{stat}}^*\) contains a term involving a generalized Jacobian
		\(F\) of the projection, it is not necessarily semismooth for an arbitrary
		closed convex cone \(K\). To justify a classical Newton linearisation at this
		stage, we therefore assume that \(\Pi_K\) is twice differentiable at the chosen
		point \(\sigma\). This assumption is restrictive, but it will be removed later.
		Indeed, the resulting linear system will be shown to be equivalent to a
		semismooth Newton step applied to \(\nabla\phi_c\), which provides a rigorous
		interpretation even when \(\Pi_K\) is not twice differentiable.
	\end{remark}
	
	This system still has a sparsity pattern that changes; it carries a product
	$C^\top C$, which squares the conditioning; $V$ is hard to compute for most
	curved cones, though it vanishes for the orthant; and no merit function is
	attached to it.
	
	\section{A family of merit functions}
	
	\begin{definition}[Merit family]
		Write $w=(x,z)$, $R(w)=\bigl(r^{*}_{\mathrm{stat}},r_{\mathrm{cone}}\bigr)$ and
		$J=\partial R$. A function $\mathcal{N}$ belongs to the \emph{family} if
		\begin{equation}
			\nabla\mathcal{N}(w)=U(w)\,R(w)
		\end{equation}
		for some $U(w)$ invertible at every $w$.
	\end{definition}
	
	\begin{proposition}[Properties of the family]
		Let $\mathcal{N}$ belong to the family. Then $\nabla\mathcal{N}(w)=0$ if and only if $R(w)=0$ and
		$\nabla^2\mathcal{N}=UJ+\mathcal{E}$ with
		$\mathcal{E}=(\partial U)[\,\cdot\,,R]$, so that Newton's method on
		$\mathcal{N}$ solves
		\begin{equation}
			\bigl(UJ+\mathcal{E}\bigr)\Delta w=-UR,
			\qquad\text{equivalently}\qquad
			\bigl(J+U^{-1}\mathcal{E}\bigr)\Delta w=-R \label{eq:displaced} .
		\end{equation}
	\end{proposition}
	
	\begin{proof}
		The first is immediate from invertibility of $U$. For the second, differentiating
		$\nabla\mathcal{N}=U R$ by the product rule gives
		$$\nabla^2\mathcal{N}=U\,\partial R+(\partial U)[\,\cdot\,,R]$$, and Newton's
		method on $\mathcal{N}$ is $\nabla^2\mathcal{N}\Delta w=-\nabla\mathcal{N}$,
		which is the displayed system; multiplying by $U^{-1}$ gives the second form.
	\end{proof}
	
	We use the term merit family here not in the strict sense given in the Newton method chapter, even through we just saw that element of the family have gradient that are zero at optimility point only. The one giving the symetrised  is such an exemple:
	\begin{definition}[Proximal merit]
		With $\sigma=Mz_c+c(x)$,
		\begin{equation}
			\mathcal{N}_L(x,z)=f(x)+\tfrac{\rho}{2}\norm{x-x_c}^2
			+\inner{z}{\proj{\sigma}{K^*}}-\tfrac12\norm{z}^2_M .
		\end{equation}
	\end{definition}
	
	\begin{proposition}[The proximal merit is a saddle function of the family]
		Let $E=\diag(I_n,-I_m)$. Then $$\nabla\mathcal{N}_L=E\,R$$ and
		$\nabla^2\mathcal{N}_L$ is the symmetrised matrix of the preceding section. In
		particular $\mathcal{N}_L$ belongs to the family with $U=E$, and its second
		block diagonal is $-M\prec0$, so it is not bounded below.
	\end{proposition}
	
	\begin{proof}
		Differentiating in $z$ gives $\proj{\sigma}{K^*}-Mz=-r_{\mathrm{cone}}$.
		Differentiating in $x$ gives
		$Px+q+\rho(x-x_c)+C^\top F^\top z$, and $F=F^\top$, so this is
		$r^{*}_{\mathrm{stat}}$. Hence $\nabla\mathcal{N}_L=(r^{*}_{\mathrm{stat}},
		-r_{\mathrm{cone}})=E R$. Differentiating once more reproduces the blocks of
		the symmetrised system, the $(1,1)$ entry acquiring $V$ from the derivative of
		$F$ contracted against $z$.
	\end{proof}
	
	\begin{definition}[The proximal second order block]
		\begin{equation}
			V_{\mathrm{prox}}
			=C^\top D^2\proj{\cdot}{K^*}(\sigma)[\,\cdot\,,\,\cdot\,,z]\,C .
		\end{equation}
	\end{definition}
	
	An actual merit can be found in the family i.e for which the optimum is a minimum.
	\begin{definition}[Gill and Robinson merit]
		For $Q\in\mathcal{A}(K)$,
		\begin{equation}
			\mathcal{N}_c(x,z)=\varphi_c(x)
			+\tfrac12\norm{r_{\mathrm{cone}}}^2_{Q^{-1}} .
		\end{equation}
	\end{definition}
	
	\begin{proposition}[The weight is forced]
		$\nabla_z\mathcal{N}_c=MQ^{-1}r_{\mathrm{cone}}$, so
		$\nabla_z\mathcal{N}_c=r_{\mathrm{cone}}$ for every $z$ if and only if $Q=M$.
	\end{proposition}
	
	\begin{proof}
		$\varphi_c$ does not depend on $z$, and
		$$r_{\mathrm{cone}}=Mz-\proj{\sigma}{K^*}$$ has $z$ derivative $M$, so
		$$\nabla_z\tfrac12\norm{r_{\mathrm{cone}}}^2_{Q^{-1}}
		=M^\top Q^{-1}r_{\mathrm{cone}}=MQ^{-1}r_{\mathrm{cone}}$$. This equals
		$r_{\mathrm{cone}}$ for all $z$ exactly when $MQ^{-1}=I$.
	\end{proof}
	
	\begin{proposition}[The Gill and Robinson merit belongs to the family]
		With $Q=M$, :
		\begin{equation}
			\nabla\mathcal{N}_c
			=\bigl(r^{*}_{\mathrm{stat}}-2C^\top M^{-1}Fr_{\mathrm{cone}},\
			r_{\mathrm{cone}}\bigr)=U\,R,
			\qquad
			U=\begin{pmatrix}I & -2C^\top M^{-1}F\\ 0 & I\end{pmatrix} ,
		\end{equation}
		which is block upper triangular and invertible.
	\end{proposition}
	
	\begin{proposition}[The Gill and Robinson system]
		Using \eqref{eq:displaced} and scaling the second block row gives the
		symmetrised matrix with $V$ replaced by $V_{\mathrm{GR}}$ and the same right
		hand side structure,
		\begin{equation}
			\begin{pmatrix}
				P+\rho I+V_{\mathrm{GR}} & C^\top F\\
				FC & -M
			\end{pmatrix}
			\begin{pmatrix}\Delta x\\ \Delta z\end{pmatrix}
			=\begin{pmatrix}-r^{*}_{\mathrm{stat}}\\ r_{\mathrm{cone}}\end{pmatrix} .
		\end{equation}
		The right hand side is the one of the symmetrised system, unchanged.
	\end{proposition}
	
	\begin{definition}[The Gill and Robinson second order block]
		\begin{equation}
			V_{\mathrm{GR}}=C^\top D^2\proj{\cdot}{K^*}(\sigma)
			\bigl[\,\cdot\,,\,\cdot\,,M^{-1}(p-r_{\mathrm{cone}})\bigr]\,C ,
			\qquad p=\proj{\sigma}{K^*} .
		\end{equation}
	\end{definition}
	
	The Gill and Robinson system still suffers from the changing sparsity pattern
	and from the cost of computing $V$. Let recall the following proposition points , that will point us in the right
	direction. 
	
	\begin{proposition}[The second derivative contracted at $M^{-1}p$]
		At a $\sigma$ where $\proj{\cdot}{K^*}$ is twice differentiable, and with
		$M\in\mathcal{A}(K)$,
		\begin{equation}
			D^2\proj{\cdot}{K^*}(\sigma)\bigl[\,\cdot\,,\,\cdot\,,M^{-1}p\bigr]
			=M^{-1}F(I-F)\succeq0 .
		\end{equation}
	\end{proposition}
	
	
	
	This naturally suggests the following definition.
	\begin{definition}[The block at the projection]
		\begin{equation}
			V_p=C^\top M^{-1}F(I-F)\,C\ \succeq0 .
		\end{equation}
	\end{definition}
	
	\begin{equation}
		\begin{pmatrix}
			P+\rho I+C^\top M^{-1}F(I-F)C & C^\top F\\[1mm]
			FC & -M
		\end{pmatrix}
		\begin{pmatrix}\Delta x\\ \Delta z\end{pmatrix}
		=\begin{pmatrix}-r^{*}_{\mathrm{stat}}\\ r_{\mathrm{cone}}\end{pmatrix} .
	\end{equation}
	
	
	This system is not in the family, so the demand on $U$ has to be relaxed.
	Nothing is lost by doing so, because the system is in fact the Newton system of
	$\varphi_c$ taken as merit.
	\begin{proposition}[The system with $V_p$ is the Newton system of the primal function]
		Eliminating $\Delta z$ from the system above leaves
		\begin{equation}
			\nabla^2\varphi_c(x)\,\Delta x=-\nabla\varphi_c(x) ,
		\end{equation}
		so that system is a representation of the primal Newton step in which $\Delta z$
		is obtained from the same factorisation, and its $\Delta x$ inherits
		$\inner{\nabla\varphi_c}{\Delta x}\leq-\rho\norm{\Delta x}^2$.
	\end{proposition}
	
	\begin{proof}
		The second row gives $\Delta z=M^{-1}\bigl(FC\Delta x-r_{\mathrm{cone}}\bigr)$.
		Substituting into the first leaves
		\[
		\bigl(P+\rho I+C^\top M^{-1}F(I-F)C+C^\top FM^{-1}FC\bigr)\Delta x
		=-r^{*}_{\mathrm{stat}}+C^\top FM^{-1}r_{\mathrm{cone}} .
		\]
		On the left by admissibility of $M$: $$C^\top FM^{-1}FC=C^\top M^{-1}F^2C$$, which added to
		$$C^\top M^{-1}F(I-F)C$$ gives $$C^\top M^{-1}FC$$, so the matrix is
		$ \nabla^2\varphi_c(x)$. On the right $FM^{-1}=M^{-1}F$ and
		$r^{*}_{\mathrm{stat}}=r_{\mathrm{stat}}+C^\top(F-I)z$ combine with
		$r_{\mathrm{cone}}=Mz-p$ and $Fp=p$ to give
		$$-r_{\mathrm{stat}}-C^\top(F-I)z+C^\top M^{-1}F(Mz-p)
		=-r_{\mathrm{stat}}+C^\top(z-M^{-1}p)-C^\top z
		=-\nabla\varphi_c(x)$$. The descent bound follows from
		$\inner{\nabla\varphi_c}{\Delta x}
		=-\inner{\nabla^2\varphi_c\Delta x}{\Delta x}$ and
		$\nabla^2\varphi_c\succeq\rho I$.
	\end{proof}
	
	Two problems remain: the sparsity pattern and the product.
	
	\section{Changing the representation}
	
	Most general conic methods use the following form.
	\begin{definition}[Split form]
		The \emph{split form} of the problem is
		\begin{equation}
			\min_{x\in\R^n,\,s\in\R^m}\ f(x)
			\qquad\text{subject to}\qquad
			Ex+s=b,\quad s\in K .
		\end{equation}
	\end{definition}
	
	It the immediate that :
	\begin{proposition}[The two forms are the same problem]
		With $E=C$ and $b=-d$, the map $x\mapsto(x,-c(x))$ is a bijection from the
		feasible set of the original problem onto the feasible set of the split form
		which preserves the objective. Consequently the two problems have the same
		optimal value, and $x$ is optimal for one if and only if $(x,-c(x))$ is optimal
		for the other.
	\end{proposition}
	
	
	Let introduce some notation:
	\begin{definition}[Residual splitting]
		Let $\rho>0$, $\rho_p>0$, $M_s\succ0$ and $M\in\mathcal{A}(K)$, with centres
		$(x_c,s_c,y_c,z_c)$. With $\sigma=Mz_c-s$, $p=\proj{\sigma}{K^*}$ and
		$F\in\partial\proj{\cdot}{K^*}(\sigma)$, set
		\begin{align}
			r_{\mathrm{stat}}^x
			&:= Px+q+\rho(x-x_c)+E^\top y, \\
			r_{\mathrm{stat}}^{*s}
			&:= \rho_p(s-s_c)+y-Fz, \\
			r_{\mathrm{eq}}
			&:= M_s(y-y_c)-(Ex+s-b), \\
			r_{\mathrm{cone}}
			&:= Mz-p.\\	
			\varrho_s
			&:= r_{\mathrm{stat}}^{*s}+M^{-1}Fr_{\mathrm{cone}}.
		\end{align}
		
	\end{definition}
	
	What this form buys is the system below, which settles both the sparsity
	pattern and the product.
	\begin{equation}\label{eq:k4}
		\begin{pmatrix}
			P+\rho I & 0 & E^\top & 0\\[1mm]
			0 & \rho_pI+M^{-1}F(I-F) & I & -F\\[1mm]
			E & I & -M_s & 0\\[1mm]
			0 & -F & 0 & -M
		\end{pmatrix}
		\begin{pmatrix}\Delta x\\ \Delta s\\ \Delta y\\ \Delta z\end{pmatrix}
		=\begin{pmatrix}
			-r^{x}_{\mathrm{stat}}\\ -r^{*s}_{\mathrm{stat}}\\
			r_{\mathrm{eq}}\\ r_{\mathrm{cone}}
		\end{pmatrix} .
	\end{equation}
	
	\begin{definition}[Cone operator]
		\begin{equation}
			G=\rho_pI+M^{-1}F\ \succ0 .
		\end{equation}
	\end{definition}
	
	In general the three by three version, of size $n+2m$, is the one solved.
	\begin{proposition}[Reduction to the three by three system]
		Eliminating $\Delta z$ from \eqref{eq:k4} through
		$\Delta z=-M^{-1}\bigl(F\Delta s+r_{\mathrm{cone}}\bigr)$ leaves
		\begin{equation}\label{eq:k3}
			\begin{pmatrix}
				P+\rho I & 0 & E^\top\\
				0 & G & I\\
				E & I & -M_s
			\end{pmatrix}
			\begin{pmatrix}\Delta x\\ \Delta s\\ \Delta y\end{pmatrix}
			=\begin{pmatrix}
				-r^{x}_{\mathrm{stat}}\\
				-\varrho_s\\
				r_{\mathrm{eq}}
			\end{pmatrix} ,
		\end{equation}
		which is symmetric quasi-definite.
	\end{proposition}
	
	
	\begin{proposition}[Reduction to the two by two system]
		Eliminating $\Delta s$ from \eqref{eq:k3} through
		$\Delta s=-G^{-1}\bigl(\varrho_s+\Delta\nu\bigr)$, leaves
		\begin{equation}\label{eq:k2}
			\begin{pmatrix}
				P+\rho I & E^\top\\
				E & -\bigl(M_s+G^{-1}\bigr)
			\end{pmatrix}
			\begin{pmatrix}\Delta x\\ \Delta y\end{pmatrix}
			=\begin{pmatrix}
				-r^{x}_{\mathrm{stat}}\\
				r_{\mathrm{eq}}+G^{-1}\varrho_s
			\end{pmatrix} ,
		\end{equation}
		again symmetric quasi-definite.
	\end{proposition}
	
	
	The two by two version is not implemented, since it requires inverting $G$ and
	which is ill conditioned.
	
	The Schur complement, of size $n$, is used as well. Written as
	$\Sigma=G(I+M_sG)^{-1}$ it needs no inverse of $G$ anywhere, which is what makes
	it usable: $G=\rho_pI+M^{-1}F$ has an eigenvalue exactly $\rho_p$ wherever the
	Clarke Jacobian $F$ has a zero one, so $\kappa(G)$ is of order
	$1/(\rho_p\mu)$, and both parameters are driven to their floors. Any route
	through $G^{-1}$ therefore carries a condition number the working precision
	cannot hold, while $I+M_sG$ has condition number of order
	$1+\norm{G}\norm{M_s}$. The reduction involves a product, which increases the
	fill in the sparse case; with $G^{-1}$ we are able to make it reach the same accuracy as the 3x3 using Ruiz(we will talk about it later) to minimize the conditioning of product appearing in it.
	
	\begin{proposition}[Schur complement]
		Let $\Sigma=\bigl(M_s+G^{-1}\bigr)^{-1}=\bigl(I+GM_s\bigr)^{-1}G
		=G\bigl(I+M_sG\bigr)^{-1}$ and
		\begin{equation}\label{eq:schur}
			S=P+\rho I+E^{\top}\Sigma E\ \succ0 .
		\end{equation}
		Then $\Delta x$ solves
		$S\Delta x=-r^{x}_{\mathrm{stat}}+E^{\top}\Sigma\bigl(r_{\mathrm{eq}}
		+G^{-1}\varrho_s\bigr)$, after which
		$\Delta y=\Sigma\bigl(E\Delta x-r_{\mathrm{eq}}-G^{-1}\varrho_s\bigr)$,
		$\Delta s=-G^{-1}(\varrho_s+\Delta y)$ and
		$\Delta z=-M^{-1}(F\Delta s+r_{\mathrm{cone}})$.
	\end{proposition}
	
	We never invert $G$. Writing $t=r_{\mathrm{eq}}-M_s\varrho_s$ and using
	\begin{equation}
		(I+M_sG)^{-1}=I-M_s\Sigma,
		\qquad
		(I+GM_s)^{-1}=I-\Sigma M_s=\Sigma G^{-1},
	\end{equation}
	the three steps become
	\begin{equation}
		S\Delta x=-r^{x}_{\mathrm{stat}}+E^{\top}\bigl(\Sigma t+\varrho_s\bigr),
		\qquad
		\Delta y=\Sigma e-\varrho_s,
		\qquad
		\Delta s=M_s\Sigma e-e ,
	\end{equation}
	with $e=E\Delta x-t$, in which $G$ appears only through $\Sigma$.
	Applying KKT conditions we get:
	\begin{proposition}[Optimality conditions in the split variables]
		$(x,s,z)$ is optimal for the split form together with a multiplier if and only
		if
		\begin{equation}
			Px+q+E^\top z=0,\qquad Ex+s=b,\qquad s\in K,\qquad z\in K^*,\qquad
			\inner{s}{z}=0 .
		\end{equation}
		Moreover the multiplier $y$ of the equality constraint equals $z$ at such a
		point.
	\end{proposition}
	
	\begin{proof}[Proof that $y=z$]
		The Lagrangian of the split form is
		\begin{equation}
			L(x,s,y,z)=f(x)+\inner{y}{Ex+s-b}-\inner{z}{s},
			\qquad z\in K^* ,
		\end{equation}
		the term $-\inner{z}{s}$ carrying the constraint $s\in K$. Stationarity in $s$
		reads $\nabla_sL=y-z=0$, since $s$ enters only through $\inner{y}{s}$ and
		$-\inner{z}{s}$. Hence $y=z$, and substituting into $\nabla_xL=Px+q+E^\top y=0$
		gives the first displayed condition.
	\end{proof}
	
	The following proposition is used as the duality gap  convergence criteria:
	
	\begin{proposition}[The gap is what the constraints leave]\label{prop:gap}
		Let $(x,s)$ satisfy $Ex+s=b$ with $s\in K$, and let $z\in K^*$ satisfy
		$Px+q+E^\top z=0$. Then
		\begin{equation}
			\inner{s}{z}=\inner{x}{Px}+\inner{q}{x}+\inner{b}{z}\ \geq0 ,
		\end{equation}
		and it vanishes if and only if $(x,s,z)$ is optimal.
	\end{proposition}
	
	\begin{proof}
		$\inner{s}{z}=\inner{b-Ex}{z}=\inner{b}{z}-\inner{x}{E^\top z}
		=\inner{b}{z}+\inner{x}{Px+q}$, which is the displayed identity. It is
		nonnegative because $s\in K$ and $z\in K^*$. If it vanishes, the five
		conditions of the preceding proposition all hold, so the point is optimal; if
		it is optimal, complementarity gives $\inner{s}{z}=0$.
	\end{proof}
	
	
	\begin{definition}[Termination test]\label{def:term}
		Given $(x,y,z)$, let $\widehat s=b-Ex$ be the slack recomputed from the data,
		and let
		\begin{equation}
			\begin{aligned}
				\pi &= \norm{\proj{-\widehat s}{K}}_\infty,
				& \pi_{\mathrm{scale}} &=
				\max\{\norm{Ex}_\infty,\norm{\widehat s}_\infty,\norm{b}_\infty\},\\
				\delta &= \norm{Px+q+E^\top z}_\infty,
				& \delta_{\mathrm{scale}} &=
				\max\{\norm{Px}_\infty,\norm{q}_\infty,\norm{E^\top z}_\infty\},\\
				\gamma &= \norm{z-\proj{z}{K^*}}_\infty,
				& \gamma_{\mathrm{scale}} &= \norm{z}_\infty,\\
				\eta &= \bigl|\inner{\widehat s}{z}\bigr|,
				& \eta_{\mathrm{scale}} &= \max\{|p^\star|,|d^\star|\},
			\end{aligned}
		\end{equation}
		with $p^\star=\tfrac12\inner{x}{Px}+\inner{q}{x}$ and
		$d^\star=-\tfrac12\inner{x}{Px}-\inner{b}{z}$. The test passes when
		\begin{equation}
			\pi\leq\varepsilon_{\mathrm{abs}}
			+\varepsilon_{\mathrm{rel}}\pi_{\mathrm{scale}},
			\quad
			\delta\leq\varepsilon_{\mathrm{abs}}
			+\varepsilon_{\mathrm{rel}}\delta_{\mathrm{scale}},
			\quad
			\gamma\leq\varepsilon_{\mathrm{abs}}
			+\varepsilon_{\mathrm{rel}}\gamma_{\mathrm{scale}},
			\quad
			\eta\leq\varepsilon_{\mathrm{abs}}
			+\varepsilon_{\mathrm{rel}}\eta_{\mathrm{scale}}
		\end{equation}
		all hold. The slack is recomputed so that $Ex+\widehat s=b$ holds exactly and
		only the four displayed quantities can fail. The fourth is the duality gap,
		since $p^\star-d^\star=\inner{\widehat s}{z}$ by
		Proposition~\ref{prop:gap}.
	\end{definition}
	
	\begin{algorithm}[H]
		\caption{Proximal method of multipliers in the split variables}
		\begin{algorithmic}[1]
			\Require data $(P,q,E,b,K)$, weights $\rho,\rho_p>0$, $M_s\succ0$,
			$M\in\mathcal{A}(K)$, a penalty update rule, tolerances
			$\varepsilon_{\mathrm{abs}},\varepsilon_{\mathrm{rel}},
			\varepsilon_{\mathrm{inf}}$ and inner tolerances $(\varepsilon_k)$
			\Ensure $(x,s,y,z)$ passing the test of Definition~\ref{def:term}, or an
			infeasibility certificate
			\State $(x_c,s_c,y_c,z_c)\gets$ a starting point, and
			$(x,s,y,z)\gets(x_c,s_c,y_c,z_c)$
			\For{$k=0,1,2,\dots$}
			\State $(x^-,z^-)\gets(x,z)$
			\Comment{kept for the displacement}
			\While{$\norm{(r^{x}_{\mathrm{stat}},r^{s}_{\mathrm{stat}},
					r_{\mathrm{eq}},r_{\mathrm{cone}})}_\infty>\min_i(\varepsilon_k)_i$}
			\Comment{$\varepsilon_k$ is per row; the loop stops against its
				smallest entry}
			\State $\sigma\gets Mz_c-s$, $p\gets\proj{\sigma}{K^*}$, choose
			$F\in\partial\proj{\cdot}{K^*}(\sigma)$ and form $G=\rho_pI+M^{-1}F$
			\State solve \eqref{eq:k3}, or its reduction \eqref{eq:k2} or
			\eqref{eq:schur}, for $(\Delta x,\Delta s,\Delta y)$
			\State $\Delta z\gets-M^{-1}\bigl(F\Delta s+r_{\mathrm{cone}}\bigr)$
			\State obtain $\alpha$ from the line search rule applied to
			$\varphi_c$
			\Comment{the subproblem's objective, not $\mathcal{N}_c$}
			\State $(x,s,y,z)\gets(x,s,y,z)
			+\alpha(\Delta x,\Delta s,\Delta y,\Delta z)$
			\EndWhile
			\If{the test of Definition~\ref{def:term} passes}
			\State \Return $(x,b-Ex,y,z)$
			\EndIf
			\State $\delta x\gets x-x^-$, \quad $\delta z\gets z-z^-$
			\Comment{the displacement over the outer iteration}
			\If{$\delta z$ certifies primal infeasibility to
				$\varepsilon_{\mathrm{inf}}$}
			\State \Return \textsc{PrimalInfeasible}
			\EndIf
			\If{$\delta x$ certifies dual infeasibility to
				$\varepsilon_{\mathrm{inf}}$}
			\State \Return \textsc{DualInfeasible}
			\EndIf
			\State apply the penalty update rule, which moves the centres
			$(x_c,s_c,y_c,z_c)$, the weights $\rho_p,M_s,M$ and
			$\varepsilon_{k+1}$
			\EndFor
		\end{algorithmic}
	\end{algorithm}
	
	
	\chapter{Smoothing point of view and interior point version}
	
	\section{The proximal step as a smoothing of the slack}
	
	\begin{equation}
		\min_{x\in\R^n,\ s\in\R^m}\ \varphi_c(x,s) .
	\end{equation}
	
	\begin{proposition}[The hard constraint on $s$ problem]
		Let
		\begin{equation}
			h(x,s)=f(x)+\tfrac{\rho}{2}\norm{x-x_c}^2
			+\tfrac{\rho_p}{2}\norm{s-s_c}^2
			+\inner{y_c}{\ell}+\tfrac12\norm{\ell}^2_{M_s^{-1}},
			\qquad \ell=Ex+s-b .
		\end{equation}
		
		The problem with hard constraint on $s$ is: 
		
		$$ \min_{x, s\in K} h(x,s)$$
		
		
		
	\end{proposition}
	
	
	
	\section{General theory for a general smoothing}
	
	\begin{definition}[Smoothed proximal problem]
		Let
		$(\phi_{\theta})$ be a smoothing of $\delta_K$ , for a fixed $\theta$ use  $\phi = \phi_{\theta}$ then the \emph{smoothed proximal problem}  is
		\begin{equation}
			\min_{x\in\R^n,\ s\in D^\theta}\ h(x,s)+\phi(s) .
		\end{equation}
	\end{definition}
	
	\begin{definition}[Smoothing stationarity residual]
		Introducing the multipliers $y$ and $z$, define
		\begin{align*}
			r_{\mathrm{stat}}^x
			&:= Px+q+\rho(x-x_c)+E^{\top}y, \\
			r_\phi^s
			&:= \rho_p (s-s_c)+y-z, \\
			r_\phi
			&:= z+\nabla\phi(s), \\
			r_{\mathrm{eq}}
			&:= M_s(y-y_c)-(Ex+s-b).
		\end{align*}
	\end{definition}
	\begin{proposition}[Stationary condition of the smoothed problem]
		A pair $(x,s)$ solves the smoothed problem if and only if
		there exist multipliers $y$ and $z$ such that
		\[
		r_{\mathrm{stat}}^x=0,
		\qquad
		r_\phi^s=0,
		\qquad
		r_\phi=0,
		\qquad
		r_{\mathrm{eq}}=0.
		\]
	\end{proposition}
	\begin{proof}
		Consequence of the general smoothing stationary conditions.
	\end{proof}
	
	
	\begin{proposition}[Newton on the proximal smoothed problem]
		Let $H\in\partial(\nabla\phi)(s)$. Then the Newton system is
		\begin{equation}
			\begin{pmatrix}
				P+\rho I & 0 & E^\top & 0\\
				0 & \rho_p I & I & -I\\
				-E & -I & M_s & 0\\
				0 & H & 0 & I
			\end{pmatrix}
			\begin{pmatrix}
				\Delta x\\ \Delta s\\ \Delta y\\ \Delta z
			\end{pmatrix}
			=-
			\begin{pmatrix}
				r_{\mathrm{stat}}^x\\ r_\phi^s\\ r_{\mathrm{eq}}\\ r_\phi
			\end{pmatrix},
		\end{equation}
		equivalently, negating the third block row,
		\begin{equation}
			\begin{pmatrix}
				P+\rho I & 0 & E^\top & 0\\
				0 & \rho_p I & I & -I\\
				E & I & -M_s & 0\\
				0 & H & 0 & I
			\end{pmatrix}
			\begin{pmatrix}
				\Delta x\\ \Delta s\\ \Delta y\\ \Delta z
			\end{pmatrix}
			=-
			\begin{pmatrix}
				r_{\mathrm{stat}}^x\\ r_\phi^s\\ -r_{\mathrm{eq}}\\ r_\phi
			\end{pmatrix}.
		\end{equation}
	\end{proposition}
	
	\begin{definition}[Perturbed residual]
		Let $\Upsilon$ satisfy $\Upsilon(s,z)=0$ at the optimum of the smoothed
		problem. The perturbed residual is
		\[
		r_{\phi}^{\Upsilon,s}=r^s_{\phi}+\Upsilon(s,z).
		\]
	\end{definition}
	
	\begin{definition}[Smoothed system with perturbed right hand side]
		\begin{equation}\label{eq:k4}
			\begin{pmatrix}
				P+\rho I & 0 & E^\top & 0\\
				0 & \rho_p I & I & -I\\
				E & I & -M_s & 0\\
				0 & H & 0 & I
			\end{pmatrix}
			\begin{pmatrix}
				\Delta x\\ \Delta s\\ \Delta y\\ \Delta z
			\end{pmatrix}
			=-
			\begin{pmatrix}
				r_{\mathrm{stat}}^x\\ r_{\phi}^s\\ -r_{\mathrm{eq}}\\
				r_{\phi}^{\Upsilon}
			\end{pmatrix}.
		\end{equation}
	\end{definition}
	
	\begin{definition}[Cone operator and reduced residual]
		
		For
		\[
		H \in \partial\!\left(\nabla \phi\right)(s),
		\]
		define
		\begin{equation}
			G_{\phi} := \rho_p I + H,
			\label{eq:G_phi}
		\end{equation}
		and
		\begin{equation}
			\varrho^\Upsilon_s := r_{\phi}^{\Upsilon,s} + r_{\phi}.
			\label{eq:varrho_s}
		\end{equation}
		
	\end{definition}
	
	\begin{proposition}[3x3 system of proximal smoothing subproblem]
		Eliminate $\Delta z$ in the preceding system you get:
		\begin{equation}
			\begin{pmatrix}
				P+\rho I & 0 & E^{\top}\\
				0 & G_\phi & I\\
				E & I & -M_s
			\end{pmatrix}
			\begin{pmatrix}\Delta x\\ \Delta s\\ \Delta y\end{pmatrix}
			=-\begin{pmatrix}r^{x}_{\mathrm{stat}}\\ \varrho^\Upsilon_s\\ -r_{\mathrm{eq}}\end{pmatrix} .
		\end{equation}
	\end{proposition}
	
	For the augmented Lagrangian smoothing we have:
	\begin{equation}
		G_{\phi^{\mathrm{AL}}}
		=
		\rho_p I + M^{-1}F.
		\label{eq:G_phi_AL}
	\end{equation}
	
	\begin{proposition}[Proposition (Our problem is this problem)]
		$\varphi_c=h+\phi^{\mathrm{AL}}_{(M,z_c)}$, that is the proximal
		subproblem is exactly the augmented Lagrangian smoothing, with parameter
		$\theta=(M,z_c)$, of the constrained problem
		\begin{equation}
			\min_{x,s}\ h(x,s)\qquad\text{subject to}\qquad s\in K .
		\end{equation}
		And using $\Upsilon(s,z) = z-Fz $.
	\end{proposition}
	
	A natural question is what if we use an interior smoothing? This lead to the following:
	\begin{definition}[Proximal interior smoothing problem]
		\begin{align}
			\nabla\phi_\tau(s) &= \tau\nabla B(s),
			& \partial(\nabla\phi_\tau)(s) &= \{\tau\nabla^{2}B(s)\},\\
			G_\phi &= \rho_pI+\tau\nabla^{2}B(s),
			& \Upsilon &= 0,\\
			r_\phi &= z+\tau\nabla B(s),
			& r^{\Upsilon,s}_\phi &= r^s_\phi,\\
			\varrho^{\Upsilon}_s &= \rho_p(s-s_c)+y+\tau\nabla B(s).
		\end{align}
	\end{definition}
	
	Which give the following algorithm:
	
	\begin{algorithm}[H]
		\caption{Proximal interior point method in the split variables}
		\begin{algorithmic}[1]
			\Require data $(P,q,E,b,K)$, weights $\rho,\rho_p>0$, $M_s\succ0$, a
			barrier $B$ of degree $\nu$, centring bounds
			$0<\sigma_{\min}\leq\sigma_{\max}<1$, a penalty update rule, tolerances
			$\varepsilon_{\mathrm{abs}},\varepsilon_{\mathrm{rel}},
			\varepsilon_{\mathrm{inf}}$
			\Ensure $(x,s,y,z)$ passing the test of Definition~\ref{def:term}, or an
			infeasibility certificate
			\State $(x_c,s_c,y_c,z_c)\gets$ a starting point with
			$s\in\intr K$ and $z\in\intr K^*$, and
			$(x,s,y,z)\gets(x_c,s_c,y_c,z_c)$
			\Comment{the smoothing is interior, so the start must be}
			\For{$k=0,1,2,\dots$}
			\State $(x^-,z^-)\gets(x,z)$
			\Comment{kept for the displacement}
			\State $\tau\gets\inner{s}{z}/\nu$
			\Comment{Euler's identity read backwards; one $\tau$ for every block}
			\State form $H\gets\tau\nabla^{2}B(s)$ and
			$G_\phi\gets\rho_pI+H$
			\State solve \eqref{eq:k4} with target $0$ for
			$(\Delta x^{\mathrm a},\Delta s^{\mathrm a},\Delta y^{\mathrm a})$, then
			$\Delta z^{\mathrm a}\gets-(H\Delta s^{\mathrm a}+r_\phi)$
			\Comment{the affine predictor: no centring}
			\State $(\alpha_p,\alpha_d)\gets$ fraction to the boundary along
			$(\Delta s^{\mathrm a},\Delta z^{\mathrm a})$
			\State $\tau_{\mathrm a}\gets
			\inner{s+\alpha_p\Delta s^{\mathrm a}}
			{z+\alpha_d\Delta z^{\mathrm a}}/\nu$
			\State $\sigma\gets\min\bigl\{\sigma_{\max},
			\max\{\sigma_{\min},(\tau_{\mathrm a}/\tau)^3\}\bigr\}$
			\Comment{the affine step reports how much centring is needed}
			\State solve \eqref{eq:k4} again with target $\sigma\tau$ and the
			residual corrected by
			$\Delta s^{\mathrm a}\!\circ\!\Delta z^{\mathrm a}$, for
			$(\Delta x,\Delta s,\Delta y)$, then
			$\Delta z\gets-(H\Delta s+r_\phi^{\Upsilon})$
			\Comment{same matrix, so one factorisation serves both}
			\State $(\alpha_p,\alpha_d)\gets$ fraction to the boundary along
			$(\Delta s,\Delta z)$
			\State $(x,s)\gets(x,s)+\alpha_p(\Delta x,\Delta s)$, \quad
			$(y,z)\gets(y,z)+\alpha_d(\Delta y,\Delta z)$
			\Comment{one Newton step per outer iteration}
			\If{the test of Definition~\ref{def:term} passes}
			\State \Return $(x,b-Ex,y,z)$
			\EndIf
			\State $\delta x\gets x-x^-$, \quad $\delta z\gets z-z^-$
			\Comment{the displacement over the outer iteration}
			\If{$\delta z$ certifies primal infeasibility to
				$\varepsilon_{\mathrm{inf}}$}
			\State \Return \textsc{PrimalInfeasible}
			\EndIf
			\If{$\delta x$ certifies dual infeasibility to
					$\varepsilon_{\mathrm{inf}}$}
			\State \Return \textsc{DualInfeasible}
			\EndIf
			\State apply the penalty update rule, which moves the centres
			$(x_c,s_c,y_c,z_c)$ and the weights $\rho,\rho_p,M_s$
			\EndFor
		\end{algorithmic}
	\end{algorithm}
	
	
	
	
	\part{Implementation detail}

	\chapter{The map}

	The theory part ends on a claim: the smoothing enters the Newton system only
	through the cone operator $G_\phi$ and the residuals, and everything
	downstream is common to both methods. The code is arranged so that this claim
	is not a remark but a compilable statement.

	The solver is small enough to read in full, but raw file order hides its
	architecture. It is laid out in five layers. The cone kernel is a shared
	service used by smoothing, linear algebra, initialisation, globalisation and
	termination; it is not itself another stage of the iteration.

	\begin{figure}[H]
		\centering
		\begin{tikzpicture}[
			node distance=4mm,
			box/.style={draw, rounded corners=1pt, align=center,
				minimum height=8mm, font=\small, inner xsep=3mm},
			lab/.style={font=\footnotesize\itshape, align=left}
			]
			\node[box, minimum width=88mm] (solve) {\texttt{solve.cpp}\ \ the outer and inner loops};
			\node[box, minimum width=88mm, below=of solve] (smooth) {\texttt{smoothing/}\ \ $\phi$, $G_\phi$, the residuals};
			\node[box, minimum width=88mm, below=of smooth] (road) {\texttt{road/}\ \ the Newton system \eqref{eq:k3} and its reduction};
			\node[box, minimum width=88mm, below=of road] (back) {\texttt{backend/}\ \ the factorisation};
			\node[box, minimum width=20mm, right=8mm of smooth, minimum height=30mm] (cone)
			{\texttt{cone/}\\[2pt] $\proj{\cdot}{K^*}$\\ $F$, $\nabla^2B$\\ scalings};
			\draw[-{Latex[length=2mm]}] (solve) -- (smooth);
			\draw[-{Latex[length=2mm]}] (smooth) -- (road);
			\draw[-{Latex[length=2mm]}] (road) -- (back);
			\draw[-{Latex[length=2mm]}, dashed] (cone.west |- smooth.east) -- (smooth.east);
			\draw[-{Latex[length=2mm]}, dashed] (cone) -- (road.east);
		\end{tikzpicture}
		\caption{The five layers. Each solid arrow is a narrow interface: the layer
			above holds a pointer to an abstract type and never learns which
			implementation it has. The dashed arrows are the cone kernel, which every
			layer calls and which owns all six cones.}
	\end{figure}

	\section{What is where}

	Every directory implements one object of the theory part, and the table names
	which.

	\begin{table}[H]
		\centering\small
		\begin{tabular}{ll}
			\toprule
			directory & what it implements\\
			\midrule
			\texttt{solve.cpp} and the root headers
			& the two algorithms and the iterate\\
			\texttt{ruiz/} & equilibration and scaling maps\\
			\texttt{cone/} & projections, linearisations, barriers and cone geometry\\
			\texttt{smoothing/} & $G_\phi$, residuals and recovery of $\Delta z$\\
			\texttt{road/} & assembly and solution of \eqref{eq:k3} or \eqref{eq:schur}\\
			\texttt{backend/} & symbolic and numeric factorisation\\
			\texttt{globalisation/} & step selection and the merit $\varphi_c$\\
			\texttt{schedule/} & centres, penalties, tolerances and regularisation\\
			\texttt{initialisation/} & cold and warm starting points\\
			\texttt{termination/} & Definition~\ref{def:term} and certificates\\
			\texttt{bindings/} & the Python surface\\
			\bottomrule
		\end{tabular}
		\caption{The implementation components. Counts are omitted deliberately:
			the contracts are stable while file lengths are not.}
	\end{table}

	Only two of the three systems of the theory part are built. The two by two
	form \eqref{eq:k2} is never assembled, because eliminating $\Delta s$ requires
	$G^{-1}$, and $G=\rho_pI+M^{-1}F$ has an eigenvalue exactly $\rho_p$ wherever
	$F$ has a zero one, so its conditioning grows as $1/(\rho_p\mu)$. The Schur
	road reaches \eqref{eq:schur} from \eqref{eq:k3} directly, through $\Sigma$,
	without ever passing through \eqref{eq:k2}.

	\section{A reading order}

	The shortest path through the method is the following sequence.

	\begin{enumerate}[leftmargin=*]
		\item \texttt{solve.cpp}, the outer and inner loops. Read it first
		and once, without following any call; it names every phase in order and
		nothing else in the repository has that view.
		\item \texttt{smoothing/smoothing.hpp}, the interface. This is the
		claim of the theory part written as a type: eleven pure virtual functions,
		of which \texttt{blocks()} and \texttt{road\_residuals()} are the only two
		the linear algebra ever calls.
		\item \texttt{smoothing/projection.cpp}, the augmented Lagrangian
		smoothing.
		\item \texttt{road/three\_by\_three.cpp}, the assembly and solve of
		\eqref{eq:k3}, which is the default road. The pattern is built once and the values refilled every
		iteration; that split is the reason the road is fast and the reason a cone
		may not change the sparsity of its block between iterations.
		\item \texttt{cone/nonneg.cpp}, the smallest complete cone. Every
		function a cone must provide appears here in its simplest form, so it is
		the template against which the other five are read.
	\end{enumerate}

	After those, \texttt{smoothing/barrier.cpp} gives the interior method in the
	same shape as the third file.

	\chapter{One solve, end to end}\label{ch:onesolve}

	This chapter follows \texttt{solve.cpp} from the first line to the last. It is
	the only file with a view of the whole method, and the order of its phases is
	the order of this chapter. At each phase the object of the theory part being
	computed is named, so that the reader who has the theory can check the
	correspondence rather than take it on trust.

	The point to carry through the chapter is that both algorithms of the theory
	part are this one loop. 

	\section{Setup}

	The settings are split by \emph{schedule}, not merely by road: \texttt{bcl},
	\texttt{gbcl}, \texttt{semismooth} and \texttt{interior} each own a block, so
	that tuning one cannot move another. The four blocks that carry a road also
	inherit a shared part holding $\rho$, $\rho_p$, the equality penalty, the
	refinement budget and the refactorisation limits, which every road needs and
	none of them needs differently. The split is what makes \texttt{bcl} usable as
	a fixed reference while \texttt{gbcl} is tuned against it.
	The problem is then equilibrated.



	\section{The starting point}

	The iterate carries $x,s,y,z$, their four centres, the two penalty vectors and
	the scalars $\rho$ and $\rho_p$. The penalties are set per row rather than
	globally, and a block with empty interior, that is a zero cone block, receives
	the equality penalty instead of the cone penalty, since there is no
	complementarity to relax there.

	A warm start scales the given point into the working problem and sets
	$y\gets z$. The two multipliers are distinct variables and are stepped
	separately, but they agree at the solution, by the stationarity in $s$, so a
	point inherited from a neighbouring problem has nothing better to offer for
	$y$. A cold start is delegated to the initialisation, which is where the two
	methods first diverge: the barrier requires
	$s\in\intr K$ and $z\in\intr K^*$ and reports this through
	\texttt{requires\_interior()}, so it is given an interior point, while the
	projection smoothing is exterior and accepts any point at all.

	\section{The outer loop}

	The outer loop is the proximal iteration. Each pass performs the following, in
	order.

	The numeric values contributed by $P$, $E$ and $\rho$ are refilled, but only
	if $\rho$ has moved since they were last written. The pattern was built once
	at setup and never changes, and since $\rho$ is usually held fixed for many
	iterations the guard saves a full refill. The guard covers only these three:
	the cone operator and the equality penalty are written by
	\texttt{assemble}, which runs on every inner iteration without a guard, so a
	change in $M$, $M_s$ or $\rho_p$ is picked up immediately. The split is by
	rate of change, not by importance.

	The inner loop then runs, and is described below. Afterwards the iterate is
	unscaled and graded by \texttt{termination.evaluate} against the original
	data, giving the four quantities of Definition~\ref{def:term}.

	The best iterate seen is tracked, ranked by the largest of those four. This is
	not cosmetic. On exit, if the
	current iterate is worse than the best recorded one, the best is restored and
	graded again, and a point that fails the test as the last iterate may pass it
	as the best.

	The displacement over the outer iteration is then formed, $\delta x=x-x^-$ and
	$\delta z=z-z^-$, and each is offered to the certificate it can prove. The
	pairing is crossed: the dual displacement certifies primal infeasibility, and
	the primal displacement certifies dual infeasibility. Finally the schedule is
	given the violation per row, the reference value the smoothing supplies, the
	two residuals and whether the inner loop converged, and it moves the centres,
	the penalties and the tolerances.

	\section{How the metrics are represented}

	A natural question at this point is what $M$ and $M_s$ are in the code, since
	the theory allows any $M\in\mathcal{A}(K)$ and any $M_s\succ0$. Both are
	penalties, and both are stored as vectors:
	\texttt{iterate.cone\_penalty} is $M$ and
	\texttt{iterate.equality\_penalty} is $M_s$, one entry per row. They are
	therefore diagonal, and the products $M^{-1}v$ and $M_sv$ that appear
	throughout are elementwise.

	Diagonal is not a restriction on $M_s$, which only has to be positive
	definite, and a diagonal choice is what makes \eqref{eq:k3} cheap. On $M$ it
	is a restriction, and admissibility says exactly which one. A diagonal metric
	commutes with the projection onto the orthant for any positive entries, so
	there the per row freedom is real and the schedule uses it. On a curved block
	it does not: admissibility forces $M$ to be a positive multiple of the
	identity on that block. The code enforces this rather than assuming it, in
	\texttt{flatten\_on\_curved}, which sets every entry of a curved block to the
	smallest entry in that block whenever the schedule has moved the penalties.
	The equilibration does the same to its row scaling.

	So $M$ is diagonal as stored, uniform per block where the cone is curved, and
	free per row where it is not. That is $\mathcal{A}(K)$, written as a vector
	and one flattening pass.

	\section{The inner loop}

	The inner budget :
	\begin{equation}
		\text{budget}=
		\begin{cases}
			1 & \text{if \texttt{one\_step\_per\_outer()}},\\
			\texttt{max\_iter\_inner} & \text{otherwise},
		\end{cases}
	\end{equation}
	and the inner loop stops in any case once the running total of Newton steps
	reaches \texttt{max\_newton}. The three limits answer different questions:
	\texttt{max\_iter\_outer} bounds the proximal iteration,
	\texttt{max\_iter\_inner} bounds one subproblem, and \texttt{max\_newton}
	bounds the whole solve in the only unit that costs anything, factorisations.
	The barrier returns true and therefore takes exactly one Newton step per
	outer iteration, which is the interior point method; the projection smoothing
	returns false and iterates to a tolerance, which is the semismooth method.

	The residuals are computed by the smoothing, and the inner loop stops when
	their largest entry falls below the smallest entry of the per row tolerance.
	The tolerance is a vector because the schedule sets it per row; the loop needs
	a scalar, and the smallest entry is the conservative choice.

	The matrix is then assembled from \texttt{blocks()} and \texttt{scalings()}
	and factored. A failed factorisation is not fatal: $\rho_p$ and $\rho$ are
	multiplied by a fixed factor, the values are rewritten, the residuals are
	recomputed at the new weights and the factorisation is retried, up to a
	budget. Only if the budget is exhausted does the solve stop, with the best
	point recovered as above.

	A factorisation that \emph{succeeds} can be just as bad, and is treated the
	same way. A backend may report success and still return a direction that is
	not finite, or one so large that the fraction to the boundary admits almost no
	step; the outer iteration is then spent taking no step at all, and the next
	repeats it at the same regularisation. A non finite direction triggers the
	same response as a failed factorisation, with one addition: it also ratchets a
	floor under the regularisation that the schedule is not allowed to undo. A
	short step is recognised the same way, against \texttt{degenerate\_step},
	but that test is off by default: it is left at zero, since recovering from a
	short step cost more solves than it saved on every set. Without the ratchet a
	bump lasts exactly one attempt, because the schedule drives the regularisation
	straight back down on the next outer iteration. 

	The system is solved for $(\Delta x,\Delta s,\Delta y)$, and
	\texttt{recover\_dual\_step} reconstructs $\Delta z$ from $\Delta s$. The cone
	multiplier is eliminated before the factorisation and rebuilt afterwards, so
	it costs nothing in the linear algebra while remaining a genuine increment
	rather than a substitution.

	The step lengths come from the globalisation, one for the primal pair and one
	for the dual pair. The smoothing is then offered the chance to correct: if
	\texttt{corrector} returns true it has rewritten its own residual, using the
	step just computed as an affine predictor, and the system is solved a second
	time and the step recomputed. The matrix is untouched, so the second solve
	reuses the factorisation. The projection smoothing declines and this costs
	nothing; the barrier accepts, which is the Mehrotra predictor and corrector.

	Finally the four variables are advanced, the primal pair by $\alpha_p$ and the
	dual pair by $\alpha_d$. Every step is guarded: a non finite direction, a non
	finite step length or a non finite iterate ends the inner loop, and in the
	last case the best recorded point is restored first.

	\section{What leaves the solver}

	The reported primal and dual variables are the unscaled $x$ and $z$. The slack
	is not the stored $s$ but is recomputed as $b-Ex$ from the original data, so
	that $Ex+s=b$ holds exactly in the answer rather than to the accuracy the
	iteration reached. The multiplier $y$ never leaves; it exists for the linear
	algebra and has no meaning to the caller.

	\chapter{Equilibration}

	\section{Where it happens}

	Equilibration is \texttt{ruiz/ruiz.cpp}, $172$ lines, and it is called
	exactly once, from \texttt{solve.cpp} before anything else touches the data:

	\begin{quote}
		\texttt{equilibrate(P, q, E, b, cones, ruiz\_settings, scaled)}
	\end{quote}

	It is guarded by \texttt{settings.equilibrate}; when that is false the
	scaled problem is a copy of the original with both scalings set to one and
	the cost scale set to one, so every downstream formula is unchanged and the
	branch costs nothing but a copy. The number of passes comes from the shared
	tuning as \texttt{ruiz\_iter}, and whether the objective is rescaled from
	\texttt{settings.scale\_cost}.

	Nothing else in the solver calls it, and nothing else writes to the scaled
	problem. From that call until the final unscaling every quantity in the
	solver lives in the scaled problem, which is the contract stated in
	Chapter~\ref{ch:onesolve}.

	\section{What is produced}

	The routine returns a \texttt{ScaledProblem}: the four scaled data items and
	the three scalings that generated them. Writing
	\[
	D=\diag(\text{\texttt{column\_scale}})\in\R^{n\times n},
	\qquad
	R=\diag(\text{\texttt{row\_scale}})\in\R^{m\times m},
	\qquad
	c=\text{\texttt{cost\_scale}}>0 ,
	\]
	the scaled data are
	\[
	\widetilde P=c\,DPD,
	\qquad
	\widetilde q=c\,Dq,
	\qquad
	\widetilde E=RED,
	\qquad
	\widetilde b=Rb .
	\]

	The variables travel with them. The solver works with $\widetilde x$,
	$\widetilde s$, $\widetilde z$ related to the original by
	\[
	x=D\widetilde x,
	\qquad
	s=R^{-1}\widetilde s,
	\qquad
	z=\frac1c\,R\,\widetilde z ,
	\]
	Two of the three are \texttt{unscale\_primal} and \texttt{unscale\_dual}; the
	third has no implementation, because the slack is never unscaled. What leaves
	the solver is $b-Ex$ recomputed from the original data, as the termination
	chapter describes, so the map for $s$ is stated here for completeness and used
	nowhere. The three inverse maps \texttt{scale\_primal}, \texttt{scale\_slack}
	and \texttt{scale\_dual} do exist, and are used only by the warm start, which
	receives a point in the original problem and must push it in.

	\begin{proposition}[The scaled problem is the same problem]
		\label{prop:ruizsame}
		Suppose $RK=K$. Then $\widetilde x$ solves
		\begin{equation}
			\min_{\widetilde x,\widetilde s}\
			\tfrac12\inner{\widetilde x}{\widetilde P\widetilde x}
			+\inner{\widetilde q}{\widetilde x}
			\quad\text{subject to}\quad
			\widetilde E\widetilde x+\widetilde s=\widetilde b,
			\quad \widetilde s\in K
		\end{equation}
		if and only if $x=D\widetilde x$ solves the original problem, and the two
		objectives differ by the factor $c$.
	\end{proposition}

	\begin{proof}
		For the objective,
		\[
		\tfrac12\inner{\widetilde x}{\widetilde P\widetilde x}
		=\tfrac{c}{2}\inner{\widetilde x}{DPD\widetilde x}
		=\tfrac{c}{2}\inner{D\widetilde x}{P\,D\widetilde x}
		=\tfrac{c}{2}\inner{x}{Px},
		\qquad
		\inner{\widetilde q}{\widetilde x}
		=c\inner{Dq}{\widetilde x}
		=c\inner{q}{x} ,
		\]
		using that $D$ is symmetric. So the scaled objective is $c$ times the
		original, and $c>0$ does not move the minimiser.

		For the constraint, $\widetilde E\widetilde x=RED\widetilde x=REx$, so
		\[
		\widetilde E\widetilde x+\widetilde s=\widetilde b
		\iff
		REx+\widetilde s=Rb
		\iff
		Ex+R^{-1}\widetilde s=b ,
		\]
		which is $Ex+s=b$ for $s=R^{-1}\widetilde s$. Finally $\widetilde s\in K$
		is equivalent to $s=R^{-1}\widetilde s\in K$ precisely because $RK=K$.
	\end{proof}

	

	\section{The hypothesis on the row scaling}

	The solver hands the \emph{same} cone objects to the scaled problem that it
	received for the original one. It never builds a scaled cone. That is only
	sound if the row scaling leaves each cone where it was.

	\begin{quote}
		The row scaling must satisfy $RK=K$, that is $R\in\Aut(K)$.
	\end{quote}

	This is a real restriction, and it is exactly the restriction admissibility
	imposes on the metric, for a different reason. The two are connected.

    It is clear that $M$ admissible implies $MK = K$ but the reverse isn't true. In the solver we take as hypothesis that $R$ is admissible.


	
	This is what \texttt{flatten\_curved} enforces, and it imposes the same value on the same block in $R$. After each pass computes a candidate row scaling
	entry by entry, every block whose cone reports \texttt{is\_curved} has its
	entries replaced by a single shared value, the geometric mean of what the
	pass proposed:
	\[
	r_i\ \longleftarrow\ \exp\Bigl(\frac1{|J|}\sum_{j\in J}\log r_j\Bigr)
	\qquad\text{for every }i\in J .
	\]
	The geometric mean is the natural choice here: it is the arithmetic mean in
	the logarithm, which is the space the iteration actually works in, and it
	leaves a block that is already uniform untouched.

	\section{Ruiz's iteration}

	The scaling itself is the algorithm of \cite{Ruiz2001}, applied to the pair
	$(P,E)$ rather than to a single matrix. The two share their column index,
	since both act on $x$, and only $E$ has rows of its own; $P$ is symmetric and
	so contributes to a column and to the matching row of $D$ at once.

	\begin{algorithm}[H]
		\caption{Ruiz equilibration of the pair $(P,E)$}
		\begin{algorithmic}[1]
			\Require $P$, $q$, $E$, $b$, the cones, a pass budget and a tolerance
			\Ensure $\widetilde P,\widetilde q,\widetilde E,\widetilde b$ and the
			scalings $D$, $R$
			\State $D\gets I$,\quad $R\gets I$
			\For{$\text{pass}=1,\dots,\text{budget}$}
			\State $\gamma_j\gets\max_i|\widetilde P_{ij}|$ over the column, then
			$\gamma_j\gets\max\bigl(\gamma_j,\ \max_i|\widetilde E_{ij}|\bigr)$
			\State $\eta_i\gets\max_j|\widetilde E_{ij}|$
			\Comment{$\infty$ norms of the columns and of the rows}
			\State $\omega\gets\max\bigl\{|\gamma_j-1|,\ |\eta_i-1|\bigr\}$ over the
			nonempty rows and columns
			\State $\gamma_j\gets1/\sqrt{\gamma_j}$,\quad
			$\eta_i\gets1/\sqrt{\eta_i}$
			\Comment{clamped, see below}
			\State flatten $\eta$ on every curved block
			\Comment{so that $R$ stays admissible}
			\State $\widetilde P_{ij}\gets\gamma_i\gamma_j\widetilde P_{ij}$,\quad
			$\widetilde E_{ij}\gets\eta_i\gamma_j\widetilde E_{ij}$
			\State $\widetilde q\gets\gamma\odot\widetilde q$,\quad
			$\widetilde b\gets\eta\odot\widetilde b$
			\State $D\gets\gamma D$,\quad $R\gets\eta R$
			\If{$\omega\leq\text{tolerance}$}
			\State \textbf{break}
			\Comment{every nonempty row and column is within tolerance of unit norm}
			\EndIf
			\EndFor
			\State rescale the objective, if enabled
		\end{algorithmic}
	\end{algorithm}


	The convergence is fast enough that the pass budget is small: the default is
	ten, and the loop usually leaves early on the tolerance $10^{-3}$. 
	
	\paragraph{Three guards.}
	The square root is taken through
	\[
	\text{guarded\_root}(v)=
	\begin{cases}
		1 & v\leq10^{-12},\\[1mm]
		1/\sqrt{\min\bigl(10^{8},\max(10^{-8},v)\bigr)} & \text{otherwise},
	\end{cases}
	\]
	so an empty row or column is left at scale one rather than producing an
	infinite scaling, and a row whose entries span more than sixteen orders of
	magnitude is clamped instead of being trusted. The convergence test ignores
	rows and columns whose norm is below $10^{-12}$, so an empty one cannot keep
	the loop running to its budget.

	\section{The cost scaling}

	After the passes, and only if \texttt{scale\_cost} is set, the objective is
	brought to unit size:
	\[
	c=\frac{1}{\max\bigl(\norm{\widetilde P}_{\max},\
		\norm{\widetilde q}_\infty\bigr)},
	\qquad
	\widetilde P\gets c\widetilde P,
	\qquad
	\widetilde q\gets c\widetilde q ,
	\]
	skipped entirely when that maximum falls outside $[10^{-8},10^{8}]$. This
	touches the objective alone, so the primal solution is unchanged and only the
	dual is scaled, which is why $c$ appears in \texttt{unscale\_dual} and in no
	other unscaling map.

\begin{remark}
	The solver is tuned on the scaled problem, and the tolerances, the initial
	penalties and the restart thresholds are all calibrated against data whose
	rows and columns have unit norm. Turning \texttt{equilibrate} off is
	supported, and useful for isolating a numerical question, but it is not a
	configuration the schedules were designed for.
\end{remark}

	\chapter{Initialisation}

	\section{Where it happens}

	Initialisation is \texttt{initialisation/}, $295$ lines over four files:
	an interface and a dispatch in \texttt{initialisation.cpp}, and one
	implementation for each road.

	\begin{table}[H]
		\centering\small
		\begin{tabular}{lll}
			\toprule
			file & class & used by\\
			\midrule
			\texttt{feasible\_start.cpp} & \texttt{FeasibleStart}
			& the semismooth road\\
			\texttt{interior\_start.cpp} & \texttt{InteriorStart}
			& the barrier road\\
			\texttt{initialisation.cpp} & the shared helpers & both\\
			\bottomrule
		\end{tabular}
	\end{table}

	The choice is not a setting. It is made in \texttt{solve.cpp} by asking the
	smoothing which it needs,
	\begin{quote}
		\texttt{make\_initialisation(smoothing->requires\_interior(), settings)}
	\end{quote}
	and the factory returns the interior start when that answer is true and the
	feasible start when it is false. A road therefore cannot be given the wrong
	initialisation, because the question is answered by the object that will
	consume the point.

	Both classes provide \texttt{cold} and \texttt{warm}, and
	\texttt{solve.cpp} calls exactly one of the two, immediately after the
	equilibration and before the outer loop begins. Everything below happens in
	the scaled problem.

	\section{The semismooth start}

	The projection smoothing is exterior: it accepts any point at all, and
	nothing about $s$ has to lie in the cone for the residuals to be defined.
	That freedom makes the question interesting rather than trivial, since some
	choice still has to be made and there is no feasibility requirement to force
	it.

	The choice made is the following. Set
	\[
	x=0,
	\qquad
	y=0,
	\qquad
	z=0,
	\qquad
	s=b+\proj{-b}{K^*} ,
	\]
	then anchor the four centres on that point. In the code the slack is built in
	two lines, the projection being taken onto the dual cone because that is the
	only projection the kernel exposes.

	\begin{proposition}[The semismooth start is the nearest feasible slack]
		\label{prop:feasiblestart}
		The slack above satisfies
		\begin{equation}
			s=\proj{b}{K} ,
		\end{equation}
		and consequently, among all slacks lying in $K$, it is the one that
		minimises the initial equality residual: $s\in K$ and
		\begin{equation}
			\norm{r_{\mathrm{eq}}}=\norm{b-s}
			=\min_{\sigma\in K}\norm{b-\sigma} .
		\end{equation}
	\end{proposition}

	\begin{proof}
		Moreau's decomposition applied at $b$ gives
		\[
		b=\proj{b}{K}+\proj{b}{K^{\circ}},
		\qquad
		K^{\circ}=-K^* .
		\]
		Projection commutes with negating both the point and the cone, so
		\[
		\proj{b}{K^{\circ}}=\proj{b}{-K^*}=-\proj{-b}{K^*} ,
		\]
		and substituting gives $b=\proj{b}{K}-\proj{-b}{K^*}$, that is
		$\proj{b}{K}=b+\proj{-b}{K^*}$, which is the displayed $s$.

		For the second claim, the equality residual is
		$r_{\mathrm{eq}}=M_s(y-y_c)-(Ex+s-b)$, and at this point $x=0$ and
		$y=y_c=0$, so $r_{\mathrm{eq}}=b-s$. The projection onto a closed convex
		set is by definition the unique minimiser of the distance to it, so
		$\proj{b}{K}$ minimises $\norm{b-\sigma}$ over $\sigma\in K$.
	\end{proof}

	So the start is feasible for the cone and as close as it can be to feasible
	for the equality, which is what the class name records. Note what it does
	\emph{not} do: it makes no attempt to satisfy the stationarity, and with
	$z=0$ the initial dual residual is simply $q$.

	\begin{algorithm}[H]
		\caption{The semismooth cold start}
		\begin{algorithmic}[1]
			\Require the scaled data, the cones
			\Ensure an iterate and its centres
			\State $x\gets0$,\quad $y\gets0$,\quad $z\gets0$
			\State $w\gets-b$
			\State $w\gets\proj{w}{K^*}$
			\Comment{blockwise, the only projection the kernel exposes}
			\State $s\gets b+w$
			\Comment{$=\proj{b}{K}$ by Proposition~\ref{prop:feasiblestart}}
			\State $(x_c,s_c,y_c,z_c)\gets(x,s,y,z)$
		\end{algorithmic}
	\end{algorithm}

	The warm path does none of this. It takes the point it was handed, already
	scaled into the working problem by the caller, and only anchors the centres
	on it. There is nothing to repair: an exterior smoothing has no feasibility
	for a warm point to violate.

	\section{The interior start}

	The barrier road cannot begin anywhere. It requires
	\[
	s\in\intr K,
	\qquad
	z\in\intr K^* ,
	\]
	since the barrier and its derivatives are undefined outside, and it requires
	the two to be far enough inside that the first step of
	Definition~\ref{def:fraction} is
	not immediately tiny. The construction has three stages, and the second may
	be skipped.

	\subsection{A reference point}

	Every cone supplies a fixed strictly feasible direction through
	\texttt{interior\_direction}, and the first stage simply takes it:
	\[
	x\gets0,
	\qquad
	y\gets0,
	\qquad
	s_j\gets d_j^{\mathrm{primal}},
	\qquad
	z_j\gets d_j^{\mathrm{dual}} ,
	\]
	on every block with an interior, and $s_j=z_j=0$ on the blocks without one.

	\begin{table}[H]
		\centering\small
		\begin{tabular}{lll}
			\toprule
			cone & primal direction & dual direction\\
			\midrule
			\texttt{Nonneg} & $\mathbf{1}$ & $\mathbf{1}$\\
			\texttt{SecondOrder} & $e_0$ & $e_0$\\
			\texttt{PSDTriangle} & $\svec(I)$ & $\svec(I)$\\
			\texttt{Exponential} & $(-1,1,2)$ & $(-1,0,1)$\\
			\texttt{Power} & $\alpha$ dependent & $\alpha$ dependent\\
			\bottomrule
		\end{tabular}
		\caption{The interior directions. Each has margin exactly one, by
			Proposition~\ref{prop:marginprop}, which is what the shift relies on.}
	\end{table}

	This point is strictly feasible and nothing else. It ignores the problem
	entirely, and if the seeded stage is disabled it is what the solver starts
	from.

	\subsection{A seeded point}

	When \texttt{seeded\_start} is set, which it is by default, the reference
	point is used only as a place to build a scaling, and the actual start comes
	from one linear solve.

	The barrier scaling is built at the reference point, the cone operator blocks
	are materialised from it, the road is assembled and factored, and the system
	is solved once with the problem data itself on the right hand side:
	\[
	r^{x}_{\mathrm{stat}}=q,
	\qquad
	\varrho_s=0,
	\qquad
	r_{\mathrm{eq}}=b .
	\]
	The three components of the solution become $x$, $s$ and $y$ directly, and
	the dual is recovered from the slack by
	\[
	z=-\bigl(G-\rho_pI\bigr)s=-Hs .
	\]

	There is no theorem here; the construction is the heuristic of
	\cite{Mehrotra1992} carried over to the conic case, and what it computes is
	worth stating plainly.

	\begin{proposition}[What the seed solves]\label{prop:seed}
		Let $G=\rho_pI+H$ be the cone operator built at the reference point. The
		seeded start is the exact solution of the problem obtained from the
		original one by replacing the cone constraint with the quadratic form of
		$G$,
		\begin{equation}
			\min_{x,s}\
			\tfrac12\inner{x}{(P+\rho I)x}+\inner{q}{x}+\tfrac12\inner{s}{Gs}
			\quad\text{subject to}\quad
			Ex+s-M_sy=b ,
		\end{equation}
		and $z=-Hs$ is the dual it induces.
	\end{proposition}

	\begin{proof}
		The Newton system of the road, with the right hand side above, reads
		\[
		(P+\rho I)\Delta x+E^{\top}\Delta y=-q,
		\qquad
		G\Delta s+\Delta y=0,
		\qquad
		E\Delta x+\Delta s-M_s\Delta y=b .
		\]
		These are exactly the stationarity conditions of the displayed problem in
		$(x,s)$ with multiplier $y$: the first is stationarity in $x$, the second
		is stationarity in $s$, and the third is the constraint. Since $P+\rho
		I\succ0$ and $G\succ0$, the problem is strictly convex and its solution is
		unique, so the Newton step from the origin is that solution rather than an
		approximation to it. The second equation gives $\Delta y=-G\Delta s$, and
		the dual assigned to the cone is the part of that coming from $H$ alone,
		$z=-Hs$, which is the gradient at $s$ of the quadratic model $\tfrac12
		\inner{s}{Hs}$ that replaced the barrier.
	\end{proof}

	In words: the seed is the point the solver would reach if the barrier were
	exactly its own second order model at the reference point. That model is a
	poor approximation far from the reference, which is why the result is not
	trusted and is repaired in the next stage. If the factorisation fails, the
	whole stage is abandoned and the reference point stands.

	\subsection{Pushing the point inside}

	Nothing so far guarantees that the seeded point is interior at all: it solves
	a quadratic problem that has forgotten the cone. The last stage restores the
	requirement, on the primal and on the dual separately.

	The quantity both stages measure is the margin. It is a property of a cone together
	with a chosen direction inside it, and every cone in the kernel supplies both.

	\begin{definition}[Margin]\label{def:margin}
		Let $K\subseteq\R^d$ be a pointed closed convex cone and let
		$d\in\intr K$, $d\neq0$, be a distinguished interior direction. The
		\emph{margin} of a point $v\in\R^d$ is
		\begin{equation}
			m_K(v)=\sup\bigl\{t\in\R\ :\ v-t\,d\in K\bigr\} .
		\end{equation}
	\end{definition}

	The supremum is over all of $\R$, not only over $t\geq0$, so the margin is
	signed: it is the amount of $d$ that can be removed from $v$ before leaving
	the cone when $v$ is inside, and minus the amount that must be added to enter
	it when $v$ is outside. That the definition is signed is what allows one
	formula to repair a point that is merely too close to the boundary and a
	point that is not in the cone at all, which is the situation the seeded start
	leaves behind.

	\begin{proposition}[Properties of the margin]\label{prop:marginprop}
		With $K$ and $d$ as above, the supremum is finite and attained, and
		\begin{enumerate}[label=(\roman*)]
			\item $m_K(v)\geq0$ if and only if $v\in K$, and $m_K(v)>0$ if and only
			if $v\in\intr K$;
			\item $m_K$ is positively homogeneous of degree one;
			\item $m_K$ is concave, and therefore superadditive;
			\item $m_K(d)=1$.
		\end{enumerate}
	\end{proposition}

	\begin{proof}
		Write $T(v)=\{t:v-td\in K\}$. It is closed, because $K$ is closed and the
		map $t\mapsto v-td$ is continuous. It is bounded above: were it not, then
		$v-td\in K$ for arbitrarily large $t$, hence $v/t-d\in K$ by the cone
		property, and letting $t\to\infty$ and using closedness gives $-d\in K$, so
		$d\in K\cap(-K)=\{0\}$ by pointedness, contradicting $d\neq0$. A closed set  in $\mathbb{R}$
		bounded above contains its supremum.

		\emph{(i)} If $v\in K$ then $0\in T(v)$, so $m_K(v)\geq0$; conversely
		$m_K(v)\geq0$ with the supremum attained gives $v-m_K(v)d\in K$ and hence
		$v\in K$, being the sum of that element with $m_K(v)d\in K$. If
		$v\in\intr K$ then $v-td\in K$ for all small $t>0$, so $m_K(v)>0$;
		conversely if $m_K(v)=t>0$ then
		\[
		v=\bigl(v-td\bigr)+td\in K+\intr K\subseteq\intr K .
		\]
		
		Where the last inequality is due to the fact that if $O$ is convex then   $\alpha x +(1-\alpha) y \in \intr O $ for any $x\in O,y \in \intr O,\alpha \in [0,1[$. Used at $\alpha = \frac{1}{2}$ combined with the fact that $K$ is a cone.
		
		\emph{(ii)} For $\lambda>0$,
		\[
		T(\lambda v)=\{t:\lambda v-td\in K\}
		=\{t:v-(t/\lambda)d\in K\}=\lambda\,T(v) ,
		\]
		dividing by $\lambda$ inside $K$, so the suprema scale by $\lambda$.

		\emph{(iii)} Let $t_1=m_K(v_1)$, $t_2=m_K(v_2)$ and $\theta\in[0,1]$. Both
		$v_1-t_1d$ and $v_2-t_2d$ lie in $K$, so their convex combination does too:
		\[
		\theta v_1+(1-\theta)v_2-\bigl(\theta t_1+(1-\theta)t_2\bigr)d\in K ,
		\]
		whence $m_K(\theta v_1+(1-\theta)v_2)\geq\theta t_1+(1-\theta)t_2$.
		Superadditivity follows from concavity and homogeneity:
		\[
		m_K(a+b)=2\,m_K\Bigl(\frac{a+b}{2}\Bigr)
		\geq2\Bigl(\tfrac12m_K(a)+\tfrac12m_K(b)\Bigr)=m_K(a)+m_K(b) .
		\]

		\emph{(iv)} $T(d)=\{t:(1-t)d\in K\}$. For $t\leq1$ the factor $1-t$ is
		nonnegative and $(1-t)d\in K$ since $K$ is a cone containing $d$. For $t>1$
		membership would give $-d\in K$, excluded as above. Hence
		$m_K(d)=1$.
	\end{proof}

	Part \emph{(iv)} holds for every pointed cone and every nonzero interior
	direction, so it is not a property the five directions of the table were
	chosen to have. On the three symmetric cones the margin has a closed form,
	and each is the definition evaluated:
	\begin{equation}
		m_{\R^d_+}(v)=\min_iv_i,
		\qquad
		m_{\mathcal{Q}^d}(v)=v_0-\norm{v_{1:}},
		\qquad
		m_{\Sp^k_+}(X)=\lambda_{\min}(X) .
	\end{equation}
	 The
	exponential and power cones have no closed form and bisect on the membership
	predicate instead, bracketing by doubling and then halving a fixed number of
	times, so what they return is Definition~\ref{def:margin} to the bisection
	tolerance rather than exactly.

	\begin{algorithm}[H]
		\caption{\texttt{shift\_inside}, on one side}
		\begin{algorithmic}[1]
			\Require a point $v$, the cones, the constants
			$\pi_{\mathrm{abs}}$ and $\pi_{\mathrm{rel}}$
			\Ensure a point whose every block has margin at least the target
			\State $m\gets\min_j\ \text{margin}\bigl(K_j,v_j\bigr)$
			over the blocks with an interior
			\State $\varsigma\gets\max_j\ \norm{v_j}_\infty$
			over the same blocks
			\State $t^{\star}\gets\pi_{\mathrm{abs}}+\pi_{\mathrm{rel}}\,\varsigma$
			\Comment{the target margin, scaled to the point}
			\State $t\gets\max\bigl(0,\ t^{\star}-m\bigr)$
			\If{$t=0$}
			\State \Return $v$
			\Comment{already inside by enough}
			\EndIf
			\For{every block $j$ with an interior}
			\State $v_j\gets v_j+t\,d_j$
			\Comment{$d_j$ the interior direction of that cone}
			\EndFor
			\State \Return $v$
		\end{algorithmic}
	\end{algorithm}

	The defaults are $\pi_{\mathrm{abs}}=1$ and $\pi_{\mathrm{rel}}=0.1$, so the
	target is one plus a tenth of the largest entry of the point. Making the
	target grow with the point matters: a fixed absolute margin would be
	generous on a problem whose data is of order one and negligible on one whose
	data is of order $10^6$.

	That the shift achieves what it claims is elementary, and rests on one
	property of the margin.

	\begin{proposition}[The shift restores the margin]\label{prop:shift}
		Let $t^{\star}$ be the target, let $m^{\min}$ be the smallest block margin
		before the shift, and let $t=\max(0,t^{\star}-m^{\min})$. Then after the
		shift every block has margin at least $t^{\star}$.
	\end{proposition}

	\begin{proof}
		Let $j$ be any block. By superadditivity and $m(d_j)=1$, both from
		Proposition~\ref{prop:marginprop},
		\[
		m\bigl(v_j+t\,d_j\bigr)\ \geq\ m(v_j)+t\,m(d_j)=m(v_j)+t .
		\]
		If $t>0$ then $m(v_j)\geq m^{\min}$ gives
		$m(v_j)+t\geq m^{\min}+t^{\star}-m^{\min}=t^{\star}$. If $t=0$ then
		$t^{\star}\leq m^{\min}\leq m(v_j)$ already.
	\end{proof}

	Note that the proposition needs no hypothesis beyond
	Proposition~\ref{prop:marginprop}, so it covers the exponential and power
	blocks as well, up to the accuracy with which their margins are bisected.

	\begin{algorithm}[H]
		\caption{The interior cold start}
		\begin{algorithmic}[1]
			\Require the scaled data, the cones, a factored road
			\Ensure a strictly feasible iterate and its centres
			\State $x\gets0$,\quad $y\gets0$
			\For{every block $j$}
			\State $s_j\gets d_j^{\mathrm{primal}}$,\quad
			$z_j\gets d_j^{\mathrm{dual}}$ if the cone has an interior, else $0$
			\EndFor
			\If{\texttt{seeded\_start}}
			\State build the scalings and the blocks of $G$ at this point
			\State assemble and factor the road
			\If{the factorisation succeeded}
			\State solve once with right hand side $(q,\,0,\,b)$
			\State $(x,s,y)\gets$ the solution,\quad $z\gets-(G-\rho_pI)s$
			\State \texttt{shift\_inside} the primal, then the dual
			\EndIf
			\EndIf
			\State $(x_c,s_c,y_c,z_c)\gets(x,s,y,z)$
		\end{algorithmic}
	\end{algorithm}

	Note where the shift sits: inside the seeded branch. If the seed is skipped,
	or if the factorisation fails, no shift is applied, because the reference
	point is already interior by construction and there is nothing to repair.

	The warm path applies the two shifts and nothing else. A point inherited from
	a neighbouring problem is usually interior already, but it may sit very close
	to the boundary of the new one, and the shift is what prevents the first
	step of Definition~\ref{def:fraction} from being vanishingly small.

	\section{The centres}

	Both roads finish by calling \texttt{anchor\_centres}, which is four
	assignments:
	\[
	x_c\gets x,
	\qquad
	s_c\gets s,
	\qquad
	y_c\gets y,
	\qquad
	z_c\gets z .
	\]
	This is the only place the centres are set outside the schedule.


	\begin{remark}[The warm path trusts its input.]
	Neither class validates the point it is handed. The semismooth road has
	nothing to validate, but the interior road assumes the caller has already
	scaled the point into the working problem, and the shift will happily push an
	unscaled point further away from where it should be rather than towards it.
\end{remark}
	\chapter{Termination}

	\section{Where it happens}

	Termination is \texttt{termination/}, $183$ lines over two files, and it is
	the smallest directory in the solver. It is reached from three places in
	\texttt{solve.cpp}, all of them in the outer loop.

	\begin{table}[H]
		\centering\small
		\begin{tabular}{lll}
			\toprule
			call & when & decides\\
			\midrule
			\texttt{evaluate} & after the inner loop & \textsc{Solved}\\
			\texttt{primal\_infeasible} & after the displacement
			& \textsc{PrimalInfeasible}\\
			\texttt{dual\_infeasible} & after the displacement
			& \textsc{DualInfeasible}\\
			\texttt{evaluate} & once more, on exit & the rescue below\\
			\bottomrule
		\end{tabular}
	\end{table}

	The remaining outcomes need no call. \textsc{NumericalFailure} is set where a
	factorisation exhausts its retries, \textsc{MaxNewton} where the global count
	of Newton steps is reached, and \textsc{MaxIterOuter} is the initial value of
	\texttt{results.status}, so a loop that simply runs out of outer iterations
	reports it without anyone assigning it. The three limits are distinguished on
	purpose: a solve stopped by the Newton budget and one stopped by the outer
	budget have failed for different reasons and are fixed by different settings.

	Two things about the call are worth stating before the contents. It is given
	\texttt{original} rather than \texttt{working}, and explicitly unscaled copies
	of $x$ and $z$: the verdict is pronounced on the problem the caller posed,
	never on the equilibrated one. And it is a free function of the data and the
	point, holding no state between calls, so the same point always receives the
	same verdict.

	\section{What is measured}

	The four quantities of Definition~\ref{def:term} are formed as follows.

	The slack is not read from the iterate but rebuilt from the data,
	\[
	\widehat s=b-Ex ,
	\]
	so that $Ex+\widehat s=b$ holds exactly and the equality constraint cannot
	fail. Only four things remain that can, and they are the four tested.

	The primal residual is the distance from that slack to the cone. It is
	computed through the dual projection, which is the only projection the kernel
	exposes, and they are the same number by the theory part.


	With that, the four quantities are
	\[
	\begin{aligned}
		\text{primal residual} &= \norm{\proj{-\widehat s}{K^*}}_\infty
		=\dist(\widehat s,K),
		&\quad
		\text{dual residual} &= \norm{Px+q+E^{\top}z}_\infty,\\[1mm]
		\text{cone violation} &= \norm{z-\proj{z}{K^*}}_\infty
		=\dist(z,K^*),
		&\quad
		\text{complementarity} &= \bigl|\inner{\widehat s}{z}\bigr| ,
	\end{aligned}
	\]
	The first three are compared against $\varepsilon_{\mathrm{abs}}
	+\varepsilon_{\mathrm{rel}}$ times their own scale, and the gap against its
	own pair $\varepsilon^{\mathrm{gap}}_{\mathrm{abs}}
	+\varepsilon^{\mathrm{gap}}_{\mathrm{rel}}$ times its own scale,
	\[
	\begin{aligned}
		\pi_{\mathrm{scale}}
		&=\max\bigl\{\norm{Ex}_\infty,\norm{\widehat s}_\infty,
		\norm{b}_\infty\bigr\},
		&\qquad
		\delta_{\mathrm{scale}}
		&=\max\bigl\{\norm{Px}_\infty,\norm{q}_\infty,
		\norm{E^{\top}z}_\infty\bigr\},\\[1mm]
		\gamma_{\mathrm{scale}}&=\norm{z}_\infty,
		&\qquad
		\eta_{\mathrm{scale}}&=\max\bigl\{|p^{\star}|,|d^{\star}|\bigr\} .
	\end{aligned}
	\]

	The gap carries its own pair, as it does in \textsc{piqp} and in Clarabel.

	The scale of the fourth is the pair of objectives, and that is only meaningful
	because the complementarity \emph{is} the duality gap: by the proposition on
	the gap, $p^{\star}-d^{\star}=\inner{\widehat s}{z}$, so measuring
	$|\inner{\widehat s}{z}|$ against $\max(|p^{\star}|,|d^{\star}|)$ is a
	relative gap test and not an arbitrary rescaling.

	\section{Termination of the inner loop}

	The test above ends the solve. A second, quite different test ends each inner
	loop, and the two should not be confused: the first is about the original
	problem, the second about the current subproblem.

	The inner loop stops when the residuals of the \emph{smoothing} fall below the
	tolerance the schedule currently holds,
	\[
	\max\bigl\{\norm{r^{x}_{\mathrm{stat}}}_\infty,
	\norm{r^{s}_{\mathrm{stat}}}_\infty,
	\norm{r_{\mathrm{eq}}}_\infty,
	\norm{r_{\mathrm{cone}}}_\infty\bigr\}
	\ \leq\ \min_i\,(\varepsilon_k)_i .
	\]
	The tolerance is a vector because the schedule sets it per row; the loop needs
	a scalar, and the smallest entry is the conservative reading. Nothing here is
	compared against $\varepsilon_{\mathrm{abs}}$: the subproblem is allowed to be
	solved loosely, and how loosely is the schedule's decision.

	\paragraph{The two roads differ only in the budget.}
	The test is the same for both, but
	\[
	\text{budget}=
	\begin{cases}
		1 & \text{barrier, since \texttt{one\_step\_per\_outer} is true},\\
		\texttt{max\_inner} & \text{projection}.
	\end{cases}
	\]
	The barrier road therefore evaluates the test once, and if it fails takes
	exactly one step and leaves; it never iterates to a tolerance, because an
	interior point method has no subproblem to solve to one. The projection road
	iterates until the test passes or the budget is exhausted, and reports which
	through \texttt{inner\_converged}, a flag the schedule reads when it decides
	how to move the penalties.

	\section{The certificates}

	Neither certificate is derived here; both are
	Definition~\ref{def:certificates} of the theory part, and what follows is how
	the code tests for them.

	Both are tested on the displacement over one outer iteration, $\delta x$ and
	$\delta z$, which is the object the asymptotic result of the theory part
	identifies as converging to a certificate when one exists. The pairing is
	crossed, the dual displacement certifying primal infeasibility and the primal
	displacement certifying dual infeasibility.

	\paragraph{Normalisation comes first.}
	The three conditions of each certificate are positively homogeneous: if $w$
	satisfies them so does $\lambda w$ for every $\lambda>0$. A certificate is
	therefore only ever determined up to positive scale, and a tolerance applied
	to an unnormalised vector would be meaningless, passing for a large
	displacement and failing for a small one that points the same way. Both tests
	begin by dividing through by $\norm{\cdot}_\infty$, and return false at once
	if that norm vanishes.

	\begin{algorithm}[H]
		\caption{The primal infeasibility test}
		\begin{algorithmic}[1]
			\Require the dual displacement $\delta z$, a tolerance
			$\varepsilon_{\mathrm{inf}}$
			\Ensure whether $\delta z$ certifies that the primal is infeasible
			\If{$\norm{\delta z}_\infty=0$}
			\State \Return false
			\EndIf
			\State $w\gets\delta z/\norm{\delta z}_\infty$
			\If{$\norm{w-\proj{w}{K^*}}_\infty>\varepsilon_{\mathrm{inf}}$}
			\State \Return false
			\Comment{$w\notin K^*$}
			\EndIf
			\If{$\norm{E^{\top}w}_\infty>\varepsilon_{\mathrm{inf}}$}
			\State \Return false
			\Comment{$E^{\top}w\neq0$}
			\EndIf
			\State \Return $\inner{b}{w}<-\varepsilon_{\mathrm{inf}}$
		\end{algorithmic}
	\end{algorithm}

	The three checks are the three conditions of the definition, in the split
	notation where $b=-d$, so $\inner{z}{d}>0$ reads $\inner{b}{w}<0$. That they
	suffice is one line, worth repeating because it is the whole content of the
	test.

	\begin{proposition}[The primal certificate is sound]
		If $w\in K^*$, $E^{\top}w=0$ and $\inner{b}{w}<0$, then no $(x,s)$ satisfies
		$Ex+s=b$ with $s\in K$.
	\end{proposition}


	\begin{algorithm}[H]
		\caption{The dual infeasibility test}
		\begin{algorithmic}[1]
			\Require the primal displacement $\delta x$, a tolerance
			$\varepsilon_{\mathrm{inf}}$
			\Ensure whether $\delta x$ certifies that the dual is infeasible
			\If{$\norm{\delta x}_\infty=0$}
			\State \Return false
			\EndIf
			\State $u\gets\delta x/\norm{\delta x}_\infty$
			\If{$\norm{Pu}_\infty>\varepsilon_{\mathrm{inf}}$}
			\State \Return false
			\Comment{$Pu\neq0$}
			\EndIf
			\If{$\inner{q}{u}>-\varepsilon_{\mathrm{inf}}$}
			\State \Return false
			\Comment{the objective does not decrease along $u$}
			\EndIf
			\State \Return
			$\norm{\proj{Eu}{K^*}}_\infty\leq\varepsilon_{\mathrm{inf}}$
			\Comment{$\proj{v}{K^*}=0$ exactly when $-v\in K$}
		\end{algorithmic}
	\end{algorithm}


	\begin{proposition}[The dual certificate is sound]
		If $Pu=0$, $-Eu\in K$ and $\inner{q}{u}<0$, then the primal objective is
		unbounded below on the feasible set whenever that set is nonempty, so the
		dual is infeasible.
	\end{proposition}



	\section{The verdict, and the rescue}

	After the loop the status is whatever the loop left, and one further step is
	taken before returning. If the status is not \textsc{Solved}, the largest of
	the four quantities at the final iterate is compared with the best value seen
	during the solve; if the final iterate is the worse of the two, the best point
	is restored, graded again, and may be declared \textsc{Solved} even though the
	last iterate was not.


	\begin{remark}[The slack is rebuilt, so one residual can never be seen.]
	Because $\widehat s=b-Ex$ is recomputed, the equality residual is identically
	zero at the point being graded, whatever the iterate's own $s$ was doing. That
	is the correct choice, since the answer returned to the caller carries that
	same recomputed slack. But it means a bug that lets $s$ drift away from
	$b-Ex$ inside the solver will never show up in the verdict, and must be caught
	by the inner residuals instead.
\end{remark}
	\begin{remark}[A certificate is only ever approximate.]
	Every test is up to $\varepsilon_{\mathrm{inf}}$, so what is reported is a
	vector that nearly satisfies the three conditions, not one that satisfies
	them. On a problem that is feasible but very badly conditioned, a displacement
	can pass the three checks without a certificate existing, and the solver will
	report infeasibility for a problem that has an answer. The tolerance is set to
	$10^{-9}$ for that reason, and lowering it trades one kind of wrong answer for
	another.
\end{remark}

	\chapter{The Newton systems}

	\section{Where it happens}

	The two roads are \texttt{road/}, $847$ lines over four files: an interface
	and the selection rule in \texttt{road.cpp}, and one implementation each in
	\texttt{three\_by\_three.cpp} and \texttt{schur.cpp}. They are the only place
	in the solver that knows what a sparse matrix is.

	A road is used in four steps, and the separation between them is the reason
	the solver is fast.

	\begin{table}[H]
		\centering\small
		\begin{tabular}{lll}
			\toprule
			call & how often & what it does\\
			\midrule
			\texttt{setup} & once per solve & plans the sparsity pattern and
			factorises symbolically\\
			\texttt{values} & when $\rho$ moves & writes the entries coming from
			$P$, $E$ and $\rho$\\
			\texttt{assemble} & every inner iteration & writes the entries coming
			from $G$ and $M_s$\\
			\texttt{factor} & every inner iteration & the numeric factorisation\\
			\texttt{solve} & once or twice per iteration & the triangular solves
			and refinement\\
			\bottomrule
		\end{tabular}
		\caption{The road interface. \texttt{setup} is the only call that decides
			where a nonzero may appear; everything after it writes into positions
			that already exist.}
	\end{table}

	That split is what \texttt{PatternBuilder} and \texttt{position\_of} exist
	for. The pattern is built once from triplets, compressed, and then every
	later write goes through an index recorded at setup. If a write ever asks for
	a position the pattern does not contain, \texttt{position\_of} raises rather
	than corrupting a neighbouring entry, with a message naming the disagreement
	between setup and assembly.

	\section{Choosing a road}

	When the caller asks for \texttt{Auto}, the choice is made once, before
	anything is built, from three numbers: the column count $n$, the row count
	$m$, and the number of nonzeros in $E$.

	\begin{algorithm}[H]
		\caption{\texttt{resolve\_road}}
		\begin{algorithmic}[1]
			\Require the requested kind, $n$, $m$, and $\operatorname{nnz}(E)$
			\Ensure a concrete road
			\If{the request is not \texttt{Auto}}
			\State \Return it
			\Comment{the caller's choice is never overridden}
			\EndIf
			\If{$m\leq n$}
			\State \Return \textsc{ThreeByThree}
			\Comment{too few rows for the reduction to pay}
			\EndIf
			\State $d\gets\operatorname{nnz}(E)/(nm)$
			\Comment{the density of the constraint}
			\State $\hat d\gets\min\bigl(1,\ d^2m\bigr)$
			\Comment{the expected density of $E^{\top}\Sigma E$}
			\If{$\hat d<\hat d_{\min}$}
			\State \Return \textsc{ThreeByThree}
			\EndIf
			\State \Return \textsc{Schur}
		\end{algorithmic}
	\end{algorithm}

	\begin{remark}[The first guard is not about size.]
	The reduced system is $n\times n$ and the three by three system is
	$(n+2m)\times(n+2m)$, so the reduction is always the smaller of the two and
	the guard $m\leq n$ is not a statement about dimension. It is a statement
	about flops, and the two counts can be written down.
    Assuming that $G$ is diagonal , $P$ is sparse and we are using a minimum degree ordering , one can give the following estimate of the number flops of the 3x3 road:
	\[
	\underbrace{F_P}_{\text{the }n\text{ sparse columns}}
	\;+\;
	\underbrace{\frac{nm^2}{2}}_{d_k=m\text{ for each of them}}
	\;+\;
	\underbrace{\frac{m^3}{6}}_{\text{the dense }m\times m\text{ corner}} ,
	\]
	where $F_P$ is the cost of factoring $P$ alone. 

	The Schur road pays instead assuming that $E$ is dense:
	\[
	\underbrace{\frac{mn^2}{2}}_{\text{forming }E^{\top}\Sigma E}
	\;+\;
	\underbrace{\frac{n^3}{6}}_{\text{factoring the dense }S} ,
	\]
	the second term because adding $E^{\top}\Sigma E$ to $P$ makes every entry of
	$S$ nonzero when $E$ is dense, so $d_k=n-k$ and no ordering can help.

	The two expressions are exchanged by swapping $m$ and $n$. With $F_P$
	negligible they therefore cross exactly at $m=n$, which is where the guard
	sits.

	\begin{table}[H]
		\centering\small
		\begin{tabular}{lrrrr}
			\toprule
			$m$ & three by three & Schur & measured ratio & from the counts\\
			\midrule
			$n$ & $6.64\cdot10^{5}$ & $6.64\cdot10^{5}$ & $1.00$ & $1.00$\\
			$n/2$ & $1.55\cdot10^{5}$ & $4.14\cdot10^{5}$ & $2.67$ & $2.86$\\
			$n/4$ & $3.88\cdot10^{4}$ & $2.89\cdot10^{5}$ & $7.45$ & $8.62$\\
			$n/10$ & $7.31\cdot10^{3}$ & $2.14\cdot10^{5}$ & $29.3$ & $41.9$\\
			\bottomrule
		\end{tabular}
		\caption{Flops by symbolic elimination on the two patterns, $n=100$, $P$
			banded and $E$ dense, under a minimum degree ordering computed for each.
			The crossover at $m=n$ is exact. The counts overstate the advantage at
			small $m$, where $F_P$ is no longer negligible against a border term that
			has become small.}
	\end{table}

The argument does not need $P$
	sparse: a dense $P$ adds $n^3/6$ to the three by three road, but the Schur
	road pays that same $n^3/6$ \emph{and} the $mn^2/2$ to form the product, so it
	still does not gain. Sparsity in $P$ widens the gap rather than creating it.
	What creates it is that the reduction always pays $O(mn^2)$ to build $S$ and
	recoups it only by shrinking the factorisation from $n+2m$ to $n$, which is no
	great shrink when $m\leq n$.
\end{remark}


	\begin{remark}[The guard is coarser than the count.]
	The counts assume $E$ dense, and the guard does not: it returns on $m\leq n$
	alone, before the density of $E$ is ever looked at. The argument above
	therefore justifies it on only half of its domain, and that is deliberate,
	because the two halves are not symmetric.

	When $E$ is dense and $m\leq n$, taking the Schur road is expensive, and the
	table gives the price: a factor of thirty at $m=n/10$. When $E$ is sparse and
	$m\leq n$, the product $E^{\top}\Sigma E$ may itself stay sparse, the $n^3/6$
	never appears, and the Schur road might well have been the cheaper of the two;
	the guard is then conservative and gives something up. But both systems are
	sparse in that case and the larger has size at most $3n$, so what is given up
	is bounded and small in absolute terms. The guard accepts a bounded loss on
	sparse problems in order to avoid an unbounded one on dense problems.
\end{remark}


	\begin{remark}[On a curved block none of this holds.]
	A semidefinite block of side $k$ contributes a dense $G$ of dimension
	$k(k+1)/2$, and is treated differently then a pure diagonal block. The guard $m\leq n$ is applied to curved
	problems all the same, on the grounds that argument is the right
	intuition even where the count is not available, and it should be read there
	as a heuristic rather than as the conclusion of the paragraphs above.
	\end{remark}

	\begin{remark}[The second guard is about the density of the reduction.]
	By the argument above $m\geq n$ and $E$ is dense we should take the Schur road. The quantity $\hat d$ is the estimate that decides it. If $E$ has
	density $d$, then a given entry of $E^{\top}\Sigma E$ is a sum over $m$ terms,
	each present when two independent entries of $E$ are, so under an
	independence assumption it is nonzero with probability about $1-(1-d^2)^m \approx d^2m$. The rule
	sends the problem to the Schur road when that estimate exceeds
	$\hat d_{\min}$, a fixed threshold of the selection rule rather than a
	setting. Another reason why the Schur road is preferred when the reduced system is that its matrix is
	positive definite, so the backend may use a supernodal factorisation, which
	spends its time in dense level three kernels; on a dense reduced system that
	is far faster than the simplicial factorisation the indefinite three by three
	system is restricted to. A sparse reduced system gives the supernodal
	algorithm nothing to work with, and then the larger but sparser three by
	three matrix wins.
\end{remark}
  
	\section{The three by three road}

	The system is \eqref{eq:k3}, of size $n+2m$, and it is assembled directly.
	Its virtue is stated in the theory part: it is symmetric quasi-definite, so
	an $LDL^{\top}$ factorisation exists for any symmetric permutation, and the
	ordering may therefore be chosen for sparsity alone without any pivoting for
	stability.

	\subsection{Three ways to hold a block}

	The cone operator occupies the middle diagonal block, and how it is stored
	depends on what shape the cone will produce. The decision is taken at setup,
	from \texttt{widest\_operator\_kind}, and never revisited.

	\begin{table}[H]
		\centering\small
		\begin{tabular}{lll}
			\toprule
			allocation & when & entries for a block of size $\ell$\\
			\midrule
			\textsc{Diagonal} & the widest kind is diagonal & $\ell$\\
			\textsc{Expanded} & diagonal plus low rank, and
			$\ell\geq8$ & $3\ell+2$\\
			\textsc{Dense} & otherwise & $\ell(\ell+1)/2$\\
			\bottomrule
		\end{tabular}
	\end{table}

	The middle row is the interesting one. A diagonal plus rank two block could
	always be stored densely, but it need not be: introducing two auxiliary
	variables carries the low rank part exactly while keeping the matrix sparse.

	\begin{proposition}[The expansion is exact]\label{prop:expansion}
		Let $G=\delta I+UTU^{\top}$ with $T\in\Sp^2$ invertible, and let
		$T=Q\Lambda Q^{\top}$ with $V=UQ$, so that
		$G=\delta I+V\Lambda V^{\top}$. Then the augmented matrix
		\begin{equation}
			\begin{pmatrix}
				\delta I & V\\
				V^{\top} & -\Lambda^{-1}
			\end{pmatrix}
		\end{equation}
		has Schur complement $G$ onto its leading block, and its inertia is that of
		$\delta I$ together with one negative eigenvalue for each positive
		$\lambda_k$ and one positive for each negative $\lambda_k$.
	\end{proposition}

	\begin{proof}
		Eliminating the trailing block gives
		\[
		\delta I-V\bigl(-\Lambda^{-1}\bigr)^{-1}V^{\top}
		=\delta I+V\Lambda V^{\top}=G .
		\]
		The inertia statement is the inertia additivity of a symmetric block
		elimination: the trailing block $-\Lambda^{-1}$ contributes the signs of
		$-1/\lambda_k$, which are opposite to those of $\lambda_k$.
	\end{proof}

	The sign bookkeeping is not decoration. The three by three matrix is
	quasi-definite because its leading block is positive definite and its trailing
	block negative definite, and an auxiliary variable must be placed on the side
	matching the sign of $-1/\lambda_k$ or the property is lost. The code counts
	the two cases separately as it writes them, which is what
	\texttt{auxiliary\_positive} and \texttt{auxiliary\_negative} record.

	\begin{algorithm}[H]
		\caption{\texttt{write\_expanded}, one block}
		\begin{algorithmic}[1]
			\Require a block $G=\delta I+UTU^{\top}$ of size $\ell$
			\Ensure its entries, written into the pattern reserved at setup
			\State $(Q,\Lambda)\gets$ eigendecomposition of the $2\times2$ middle
			$T$
			\State $V\gets UQ$
			\State write $\delta$ on the $\ell$ diagonal positions of the block
			\For{$k=1,2$}
			\If{$|\lambda_k|\leq\varepsilon_{\mathrm{pin}}$}
			\State write a zero column and $-1$ on the auxiliary diagonal
			\Comment{the correction has rank below two here}
			\Else
			\State write the column $V_{\cdot k}$ and $-1/\lambda_k$ on the
			auxiliary diagonal
			\EndIf
			\State record the sign of that diagonal
			\EndFor
		\end{algorithmic}
	\end{algorithm}

	The pinning branch is what makes the expansion safe when the rank two
	correction degenerates to rank one or zero, which happens on a second order
	block whose linearisation point sits on the axis. Writing $-1/\lambda_k$ there
	would overflow; instead the column is zeroed and the diagonal set to $-1$, so
	that the auxiliary variable decouples and contributes nothing.

	\begin{remark}[Why eight.]
	Dense storage of a block of size $\ell$ costs $\ell(\ell+1)/2$ entries in the
	lower triangle, and the expansion costs $\ell$ on the diagonal, $2\ell$ in the
	two columns and $2$ on the auxiliary diagonal. The expansion is cheaper when
	\[
	\frac{\ell(\ell+1)}{2}>3\ell+2
	\qquad\Longleftrightarrow\qquad
	\ell^2-5\ell-4>0
	\qquad\Longleftrightarrow\qquad
	\ell\geq6 ,
	\]
	so the default threshold of eight is that crossover with a small margin, and a
	second order block of dimension three or four is stored densely because it is
	genuinely cheaper that way.
\end{remark}

	\subsection{Solving}

	The right hand side is assembled from the three residuals, with zeros on the
	auxiliary rows, and the sign of the third block matches the theory part:
	\[
	\text{rhs}=\bigl(-r^{x}_{\mathrm{stat}},\ -\varrho_s,\ +r_{\mathrm{eq}},\
	0\bigr) .
	\]
	After the backend's solve, the direction is read off the first three segments
	and the auxiliary entries are discarded; they exist only to carry the low rank
	part through the factorisation.

	\paragraph{Iterative refinement.}
	Every solve is refined, at most \texttt{refinement\_budget} times, which is
	the shared setting \texttt{refine} and defaults to eight. Each pass forms the true residual against the assembled
	matrix, solves with it, and adds the correction. The loop stops early on
	either of two conditions, and the second matters as much as the first:
	\[
	\norm{r}_\infty\ \geq\ \text{the previous pass's value},
	\qquad\text{or}\qquad
	\norm{r}_\infty\ \leq\ \varepsilon_{\mathrm{mach}}\max\bigl(
	\norm{\text{rhs}}_\infty,1\bigr) .
	\]
	Refinement that has stopped improving is refinement that has started chasing
	rounding, and continuing costs a solve per pass for nothing.

	\section{The Schur road}

	The reduced system is \eqref{eq:schur}, of size $n$ and positive definite. Two
	things have to be arranged: the weight $\Sigma$, and the reduction itself
	without ever forming $G^{-1}$.

	\subsection{The weight}

	The weight is $\Sigma=G(I+M_sG)^{-1}$, computed once per block per iteration
	by \texttt{schur\_weight}, and its form follows the block's.

	\begin{table}[H]
		\centering\small
		\begin{tabular}{ll}
			\toprule
			block & $\Sigma$\\
			\midrule
			diagonal & $g/(1+g\,\mu_s)$, elementwise\\
			diagonal plus low rank & diagonal plus low rank, one $r\times r$ solve\\
			dense & one dense solve of $(I+M_sG)X=G$\\
			\bottomrule
		\end{tabular}
		\caption{The Schur weight preserves the shape of the block. That it does so
			for the middle row is Proposition~\ref{prop:closedweight}.}
	\end{table}

	\begin{proposition}[The low rank form is closed under the weight]
		\label{prop:closedweight}
		Let $G=\delta I+UTU^{\top}$ and let $M_s=\mu_sI$ be a multiple of the
		identity on the block. Then $\Sigma=G(I+M_sG)^{-1}$ is again of the form
		$\delta'I+UT'U^{\top}$ with the same columns $U$, and
		\begin{equation}
			\delta'=\frac{\delta}{1+\mu_s\delta} .
		\end{equation}
	\end{proposition}

	\begin{proof}
		Write $\sigma=1+\mu_s\delta$. Then
		$I+M_sG=\sigma I+\mu_sUTU^{\top}$, and by the
		Sherman and Morrison and Woodbury identity its inverse is
		\[
		\frac1\sigma I-\frac{1}{\sigma^2}U
		\Bigl(I+\frac{\mu_s}{\sigma}TU^{\top}U\Bigr)^{-1}\mu_sT\,U^{\top} ,
		\]
		which is a multiple of the identity plus a correction whose columns are
		again $U$. Multiplying on the left by $G=\delta I+UTU^{\top}$, every term is
		either a multiple of the identity or has $U$ on the left and $U^{\top}$ on
		the right, so the product has the stated form, and the coefficient of the
		identity is $\delta/\sigma$.
	\end{proof}

	This is why the metric must be uniform on a curved block, and the code says so
	where the requirement bites: \texttt{add\_penalty} raises if a diagonal plus
	low rank block is handed a penalty that varies across its rows, because the
	scalar $\delta$ cannot absorb it. Admissibility guarantees the uniformity, and
	the proposition is what turns that guarantee into a closed representation.

	\subsection{The reduction}

	The elimination itself is where the road was once wrong, and where the
	identity that fixed it lives.

	\begin{algorithm}[H]
		\caption{\texttt{reduce}, one solve of the Schur road}
		\begin{algorithmic}[1]
			\Require the three residuals, the stored weight $\Sigma$ and penalty
			$M_s$
			\Ensure $(\Delta x,\Delta s,\Delta y)$
			\State $t\gets r_{\mathrm{eq}}-M_s\varrho_s$
			\State $\text{rhs}\gets-r^{x}_{\mathrm{stat}}
			+E^{\top}\bigl(\Sigma t+\varrho_s\bigr)$
			\State solve $S\,\Delta x=\text{rhs}$
			\Comment{the only factorised solve}
			\State $e\gets E\Delta x-t$
			\State $\Delta y\gets\Sigma e-\varrho_s$
			\State $\Delta s\gets M_s\Sigma e-e$
			\Comment{$=-(I-M_s\Sigma)e$}
		\end{algorithmic}
	\end{algorithm}

	Every appearance of $G$ in that algorithm is inside $\Sigma$. The last line is
	the one that matters, and it is the push through identity of the theory part
	written out: $(I+M_sG)^{-1}=I-M_s\Sigma$. Eliminating $\Delta s$ the obvious
	way would ask for $G^{-1}$ instead, and
	\[
	G=\rho_pI+M^{-1}F
	\]
	has an eigenvalue exactly $\rho_p$ wherever $F$ has a zero one, so its
	condition number grows like $1/(\rho_p\mu)$ and reaches $10^{25}$ late in a
	solve. The identity is not an optimisation; it is the difference between a
	road that converges on curved cones and one that does not.

	\paragraph{Refinement, on the full system.}
	The Schur road refines like the other, but against the \emph{unreduced} $3
	\times3$ residual rather than against $S\Delta x-\text{rhs}$. The correction
	is obtained by calling \texttt{reduce} again on that residual, so each pass
	costs one triangular solve and no refactorisation. Refining against the
	reduced residual instead would measure only the error the reduction did not
	introduce.

	\begin{remark}[The pattern is a contract, and it is one sided.]
		\texttt{setup} promises that every position a later write can ask for
		exists, and it keeps that promise by asking each cone for the widest block
		it will ever emit. A cone that under-reported would produce a
		\texttt{position\_of} failure, which is loud. A cone that over-reported
		would produce a matrix with structural zeros, which is silent and merely
		slower. Only one of the two mistakes is caught.
	\end{remark}

	\begin{remark}[Quasi-definiteness is assumed, not verified.]
		The three by three road never checks the inertia of what it assembled. The
		property is guaranteed by construction, from $P+\rho I\succ0$, $G\succ0$
		and $M_s\succ0$, and by the sign discipline of
		Proposition~\ref{prop:expansion} on the auxiliary variables. If a change
		broke that discipline the factorisation would still run, because a backend
		asked for an $LDL^{\top}$ without pivoting will produce one; it would simply
		be the wrong one.
	\end{remark}

	\begin{remark}[The two roads can disagree.]
		They solve the same system and return directions that agree to the accuracy
		of the refinement, which is not the same as agreeing exactly. On a problem
		near the edge of convergence the two can take different numbers of
		iterations, and one can converge where the other does not. 
	\end{remark}

	\chapter{Backends}

	\section{Where it happens}

	A backend is a factorisation, and nothing else. The directory is
	\texttt{backend/}, $769$ lines over nine files: an interface and a dispatch,
	five implementations, and the thread control they share. A road holds one
	backend and calls it; no other part of the solver knows a backend exists.

	The interface is deliberately narrow.

	\begin{table}[H]
		\centering\small
		\begin{tabular}{lll}
			\toprule
			call & how often & what it does\\
			\midrule
			\texttt{capabilities} & never, at present & declares what the backend
			can do\\
			\texttt{symbolic} & once per solve & takes the pattern, orders it,
			plans the factor\\
			\texttt{value\_count} & once & how many values the pattern holds\\
			\texttt{numeric} & every inner iteration & factors the values into
			that plan\\
			\texttt{update} & never, at present & applies a low rank change to the
			factor\\
			\texttt{solve} & once or twice per iteration & the triangular solves\\
			\texttt{set\_threads} & once & how many threads to ask for\\
			\bottomrule
		\end{tabular}
		\caption{The backend interface. The split between \texttt{symbolic} and
			\texttt{numeric} is Definitions~\ref{def:symbolic} and
			\ref{def:numeric}, and it is the reason a road may factor a hundred times
			after ordering once.}
	\end{table}

	\texttt{numeric} returns a boolean rather than raising. A factorisation that
	fails is an ordinary event, handled by the retry loop of
	Chapter~\ref{ch:onesolve}, not an error.

	Two members of the contract are declared and not yet called.
	\texttt{update} applies a low rank change to an existing factorisation rather
	than recomputing it, and takes a \texttt{Delta} naming the columns and the
	sign. Only \texttt{eigen\_dense} implements it, through a rank update of its
	factorisation; the other four accept the call and return false. It is kept
	because the place it would pay is identifiable rather than hypothetical: on
	the Schur road applied to a quadratic program, a row changing activity moves
	one entry of $\Sigma$ and so changes $S$ by a rank one term, which an update
	would absorb in $O(\operatorname{nnz}L)$ against a refactorisation's $n^3/6$.

	\texttt{capabilities} is how a caller would find out whether that is
	available, reporting \texttt{supports\_update} and the thread count the
	backend can use. Nothing reads it yet either, for the same reason: with no
	caller of \texttt{update} there is nothing to condition on. The pair is the
	extension point for a future active set update, left in place and honestly
	inert.

	\section{The pattern contract}

	Everything a backend receives is a pattern in compressed sparse column form,
	and \texttt{validate\_pattern} checks five conditions before any of them sees
	it:
	\begin{itemize}
		\item the column starts begin at zero and do not decrease;
		\item no column is empty, so every diagonal entry is present;
		\item each column begins with its own diagonal;
		\item no entry lies above the diagonal, the matrix being stored lower
		triangular;
		\item within a column the row indices increase.
	\end{itemize}

	The conditions are not arbitrary. Storing only the lower triangle halves the
	memory and is what every sparse factorisation expects; requiring the diagonal
	present means no backend has to consider a structurally singular pivot;
	requiring it first, and the rest sorted, means a backend can locate an entry
	by a binary search rather than a scan, which is what \texttt{position\_of} in
	the road does on every write.

	Validation happens once, at \texttt{symbolic}. It is a real check rather than
	an assertion, and it raises with the condition that failed, because a pattern
	that violates one of them produces not a wrong answer but a segmentation
	fault inside a library.

	\section{The five backends}

	\begin{table}[H]
		\centering\small
		\begin{tabular}{llll}
			\toprule
			name & factorisation & sparse & threads\\
			\midrule
			\texttt{cholmod} & supernodal $LL^{\top}$ or simplicial $LDL^{\top}$
			& yes & from the \textsc{blas}\\
			\texttt{qdldl} & simplicial $LDL^{\top}$ & yes & one\\
			\texttt{eigen\_sparse} & \texttt{SimplicialLDLT} & yes & one\\
			\texttt{lapack\_dense} & \texttt{dpotrf} or \texttt{dsytrf} & no &
			from the \textsc{blas}\\
			\texttt{eigen\_dense} & \texttt{LDLT} & no & from the \textsc{blas}\\
			\bottomrule
		\end{tabular}
	\end{table}

	Each is constructed with a flag saying whether the matrix it will receive is
	positive definite, and two of the five use it.

	\texttt{cholmod} uses it twice. It sets \texttt{supernodal} to
	\textsc{auto} when the matrix is definite and to \textsc{simplicial} when it
	is not, and it asks for a final $LL^{\top}$ in the first case and an
	$LDL^{\top}$ in the second. This is the mechanism behind the road selection
	of the previous chapter: the Schur road passes true and may therefore reach
	the supernodal algorithm of Definition~\ref{def:supernodal}, while the three
	by three road passes false and is confined to the simplicial one, its matrix
	being quasi-definite rather than definite.

	\texttt{lapack\_dense} uses it to choose between \texttt{dpotrf}, the dense
	Cholesky, and \texttt{dsytrf}, the dense symmetric indefinite factorisation
	with Bunch and Kaufman pivoting.

	The other three ignore it and always compute an $LDL^{\top}$, which is
	correct in both cases and merely leaves something on the table when the
	matrix is definite.

	\paragraph{Why five.}
	They are not alternatives of equal standing. \texttt{cholmod} is the default
	and the only one used in the benchmark. \texttt{qdldl} is vendored, has no
	dependencies, and exists so that the solver can be built and tested without
	\textsc{suitesparse} present. The two dense backends exist for small problems
	and for isolating a question: if a result differs between a sparse and a
	dense backend, the difference is in the sparse machinery and not in the
	method. \texttt{eigen\_sparse} is the reference implementation against which
	the others were checked.

	\section{Choosing one}

	\texttt{resolve\_backend} takes the requested kind and the cones, and when the
	request is \textsc{auto} it returns \texttt{cholmod}. It does not look at the
	cones; the parameter is accepted and discarded, which the code marks
	explicitly.

	That is a deliberate simplification rather than an oversight. The choice that
	matters was already made, one level up, when the road was selected: the road
	decides whether the matrix handed down is definite, and that decides whether
	\texttt{cholmod} runs supernodal or simplicial. Choosing a different backend
	on top of that would be choosing between implementations of the same
	algorithm, and on the benchmark \texttt{cholmod} was the fastest of them in
	every case.

	\section{Threads}

	Only the dense kernels are threaded, and the solver does not thread anything
	itself. \texttt{threads.cpp} declares four \textsc{blas} entry points as weak
	symbols,
	\[
	\texttt{openblas\_set\_num\_threads},\quad
	\texttt{MKL\_Set\_Num\_Threads},\quad
	\texttt{omp\_set\_num\_threads},
	\]
	together with their query counterparts, and calls whichever the linker
	actually resolved. If none did, the request is refused and the count reported
	is one.

	Two consequences. A backend that reports one thread is not necessarily
	single threaded by design; it may be that no threaded \textsc{blas} was
	linked. And the count a backend reports is what it asked for and what some
	library confirmed, not a measurement, so a build that links a threaded
	\textsc{blas} and then runs under an environment variable capping it will
	report the request rather than the cap.


	\begin{remark}[A failed factorisation is not an exception.]
		\texttt{numeric} returns false, and the caller raises the regularisation
		and tries again. The consequence is that a backend must leave itself in a
		usable state after failing, since it will be called again immediately with
		different values and the same pattern. All five do, but nothing in the
		interface says so, and a backend that freed its symbolic analysis on
		failure would work until the first ill conditioned problem.
	\end{remark}

	\chapter{Globalisation}

	\section{Where it happens}

	A globalisation rule answers one question: given a direction, how far along it
	to move. The directory is \texttt{globalisation/}, $370$ lines over four
	files, holding five rules and the merit function three of them share.

	\begin{table}[H]
		\centering\small
		\begin{tabular}{lll}
			\toprule
			rule & used by & decides from\\
			\midrule
			\texttt{unit\_step} & neither, by default & nothing\\
			\texttt{fraction\_to\_boundary} & the barrier road & the cone boundary\\
			\texttt{exact\_search} & the projection road, by default &
			$\varphi_c'$\\
			\texttt{armijo\_search} & the projection road, optional &
			$\varphi_c$ and $\varphi_c'(0)$\\
			\texttt{decrease\_search} & the projection road, optional &
			$\varphi_c$ alone\\
			\bottomrule
		\end{tabular}
	\end{table}

	Every rule returns a \texttt{StepPair}, one length for the primal pair
	$(x,s)$ and one for the dual pair $(y,z)$. Only the barrier rule returns two
	different numbers; the three searches return the same value twice, because
	the merit they work on is a function of the primal variables and has nothing
	to say about the dual. The choice of rule is made in \texttt{solve.cpp} from
	the road and, on the projection road, from a setting.

	\section{The merit along the ray}

	The three searches all work on the same function: the proximal subproblem's
	own objective $\varphi_c$, restricted to the ray. Write
	\begin{equation}
		\psi(t)=\varphi_c\bigl(x+t\Delta x,\ s+t\Delta s\bigr),
		\qquad t\geq0 .
	\end{equation}
	That choice is not free. The residual norm $\mathcal{N}_c$ is also a merit for
	the same system, and searching on it instead was measured to be worse.

	\subsection{What can be asked of it, and at what price}

	\texttt{Merit} exposes exactly two queries, and the split between them is what
	separates the three searches.

	\begin{table}[H]
		\centering\small
		\begin{tabular}{lll}
			\toprule
			call & returns & cost\\
			\midrule
			\texttt{slope}$(t)$ & $\psi'(t)$ & one projection of the slack\\
			\texttt{decrease}$(t)$ & $\psi(t)-\psi(0)$ & one projection of the
			slack\\
			\bottomrule
		\end{tabular}
	\end{table}

	Both are cheap for the same reason. Everything in $\varphi_c$ except the cone
	term is quadratic along the ray, so its contribution is a constant and a
	curvature, both computed once by \texttt{bind}. Only the cone term has to be
	recomputed at each trial:
	\begin{equation}\label{eq:meritslope}
		\psi'(t)=c+\kappa\,t
		-\bigl\langle\proj{\sigma-t\Delta s}{K^*},\ \Delta s/\mu\bigr\rangle,
		\qquad \sigma=Mz_c-s .
	\end{equation}
	A trial is therefore one projection and two inner products, against a Newton
	step that is a factorisation. This is why a few dozen trials are affordable.

	The two coefficients are
	\begin{align}
		c&=\inner{Px+q}{\Delta x}+\rho\inner{x-x_c}{\Delta x}
		+\rho_p\inner{s-s_c}{\Delta s}+\inner{\nu}{E\Delta x+\Delta s},\\
		\kappa&=\inner{\Delta x}{(P+\rho I)\Delta x}+\rho_p\norm{\Delta s}^2
		+\norm{E\Delta x+\Delta s}^2_{M_s^{-1}} .\label{eq:curvature}
	\end{align}

	\subsection{Why a bracket exists}

	The exact search looks for a zero of $\psi'$ rather than a decrease in
	$\psi$, and that is only sensible because $\psi'$ is increasing. It is, and
	the same fact supplies the bracket the search needs.

	\begin{proposition}[The restriction is strongly convex]\label{prop:meritconvex}
		Let the direction $(\Delta x,\Delta s)$ be nonzero and let $\kappa$ be given
		by \eqref{eq:curvature}. Then $\kappa>0$ and
		\begin{equation}
			\psi''(t)\ \geq\ \kappa
			\qquad\text{for every }t\geq0 ,
		\end{equation}
		so $\psi$ is strongly convex along the ray and $\psi'$ is strictly
		increasing.
	\end{proposition}

	\begin{proof}
		Every term of \eqref{eq:curvature} is nonnegative, and
		$\rho\norm{\Delta x}^2+\rho_p\norm{\Delta s}^2>0$ for a nonzero direction
		since $\rho,\rho_p>0$, so $\kappa>0$.

		Differentiating \eqref{eq:meritslope} once more, the quadratic part
		contributes $\kappa$ and the cone term contributes
		\[
		\frac{d}{dt}\Bigl(-\bigl\langle\proj{\sigma-t\Delta s}{K^*},
		\Delta s/\mu\bigr\rangle\Bigr)
		=\bigl\langle F\Delta s,\ \Delta s/\mu\bigr\rangle
		=\inner{\Delta s}{M^{-1}F\Delta s}\ \geq\ 0 ,
		\]
		using $\tfrac{d}{dt}\proj{\sigma-t\Delta s}{K^*}=-F\Delta s$ and the fact
		that $M^{-1}F$ is positive semidefinite, $F$ being so and commuting with
		$M^{-1}$ by admissibility. Hence $\psi''\geq\kappa$.
	\end{proof}

	\begin{proposition}[A bracket for the minimiser]\label{prop:bracket}
		Suppose $\psi'(0)<0$. Then $\psi$ has a unique minimiser $t^{\star}>0$ on
		the ray, and
		\begin{equation}
			0<t^{\star}\leq\frac{|\psi'(0)|}{\kappa} .
		\end{equation}
		In particular $\bigl[0,\ |\psi'(0)|/\kappa\bigr]$ brackets the zero of
		$\psi'$.
	\end{proposition}

	\begin{proof}
		By Proposition~\ref{prop:meritconvex}, integrating $\psi''\geq\kappa$ from
		$0$ to $t$ gives
		\[
		\psi'(t)\ \geq\ \psi'(0)+\kappa t .
		\]
		At $t=|\psi'(0)|/\kappa=-\psi'(0)/\kappa$ the right side is zero, so
		$\psi'$ is nonnegative there. Since $\psi'$ is strictly increasing and
		negative at $0$, it has exactly one zero, and that zero lies in
		$(0,|\psi'(0)|/\kappa]$. Strong convexity makes it the unique minimiser.
	\end{proof}

	The bound is exactly what the implementation uses. It is not a guess at a
	large enough step but the point at which a quadratic with the observed slope
	and the guaranteed curvature would turn, and the true $\psi$ turns no later
	because its curvature is at least $\kappa$ everywhere.

	\section{Fraction to the boundary}

	The rule for the barrier road, which is Definition~\ref{def:fraction}. The
	implementation adds one thing the definition does not: the fraction is not a
	single constant.
	\[
	\theta=
	\begin{cases}
		\texttt{tau} & \text{on a polyhedral block},\\
		\texttt{tau\_curved} & \text{on a curved block},
	\end{cases}
	\]
	with the second the smaller of the two by a wide margin,
	and the step is the smallest over the blocks, capped at one, with the primal
	and dual computed separately from $(s,\Delta s)$ and $(z,\Delta z)$.

	The two values differ by a great deal, and the reason is not accuracy: the
	largest feasible step is exact on the second order and semidefinite cones, by
	a quadratic root and by an eigenvalue. It is that the conditioning degrades
	much faster near a curved boundary. On the orthant, approaching it makes one
	ratio $z_i/s_i$ large; on a curved block the whole Nesterov and Todd scaling
	degenerates at once, and the next factorisation pays for it. Twenty percent of
	margin is the price of keeping the block conditioned.

	Blocks with no interior are skipped, having no boundary to respect.

	\section{The exact search}

	The default on the projection road. It finds the zero of $\psi'$ inside the
	bracket of Proposition~\ref{prop:bracket}.

	\begin{algorithm}[H]
		\caption{\texttt{exact\_search}}
		\begin{algorithmic}[1]
			\Require the direction, bound into the merit
			\State $g_0\gets\texttt{slope}(0)$
			\If{$g_0\geq0$}
			\State \Return $1$
			\Comment{not a descent direction; there is nothing to search}
			\EndIf
			\State $a\gets0$, \ $g_a\gets g_0$, \ $b\gets1$, \
			$g_b\gets\texttt{slope}(1)$
			\If{$g_b<0$}
			\Comment{the unit step does not yet reach the minimiser}
			\State $a\gets b$, \ $g_a\gets g_b$
			\State $b\gets|g_0|/\kappa$
			\Comment{Proposition~\ref{prop:bracket}; the minimiser is at or before
				this}
			\If{$b\leq a$} \State \Return $a$ \EndIf
			\State $g_b\gets\texttt{slope}(b)$
			\EndIf
			\If{$g_a\leq0\leq g_b$}
			\For{at most \texttt{ls\_max\_secant} trials}
			\State stop once $b-a$ is below $\sqrt{\varepsilon_{\mathrm{mach}}}$
			relative to $b$
			\State $t\gets$ the secant point of $(a,g_a)$ and $(b,g_b)$, or the
			midpoint if that falls outside
			\State replace whichever end shares the sign of $\texttt{slope}(t)$
			\State halve the slope kept at the other end if that end has now been
			kept twice running
			\Comment{the Illinois correction}
			\EndFor
			\EndIf
			\State \Return $\min\bigl(b,\ t_{\max}\bigr)$
			\Comment{$t_{\max}$ a fixed cap on the returned step}
		\end{algorithmic}
	\end{algorithm}

	Three points about it.

	The bracket is established rather than assumed. If the unit step still has a
	negative slope the search does not simply try larger numbers; it jumps to the
	bound of Proposition~\ref{prop:bracket}, which is guaranteed to lie at or
	beyond the minimiser, so one extra slope evaluation closes the bracket for
	good.

	The iteration is the secant method with the Illinois correction. Plain regula
	falsi keeps one endpoint fixed whenever the function is convex, which it is
	here, and then converges linearly; halving the retained slope after two
	consecutive retentions of the same end restores the superlinear rate.

	And what is returned is $b$, the end on which the slope is nonnegative, not
	the midpoint or the last trial. The step therefore never overshoots the
	minimiser, which matters because $\psi$ is only a model of the merit the outer
	loop cares about.

	\section{The Armijo search}

	Backtracking from the unit step, accepting the first $t$ with
	\begin{equation}
		\psi(t)-\psi(0)\ \leq\ \sigma_{\mathrm{ls}}\,t\,\psi'(0) ,
		\qquad \sigma_{\mathrm{ls}}=\texttt{ls\_sigma} .
	\end{equation}
	This is the standard sufficient decrease: it asks for a decrease proportional
	to what the slope at zero promised, rather than merely for a decrease, and it
	is the condition under which a convergence proof for a backtracking method is
	available. It needs $\psi'(0)$, one extra query before the loop, and returns
	the unit step immediately if that slope is not negative.

	\begin{algorithm}[H]
		\caption{\texttt{armijo\_search}}
		\begin{algorithmic}[1]
			\State $g_0\gets\texttt{slope}(0)$;\ \ \textbf{if} $g_0\geq0$
			\textbf{then} \Return $1$
			\State $t\gets1$
			\For{at most \texttt{ls\_max\_back} halvings}
			\If{$\texttt{decrease}(t)\leq\sigma_{\mathrm{ls}}\,t\,g_0$}
			\State \Return $t$
			\EndIf
			\State $t\gets t/2$
			\EndFor
			\State \Return $t$
		\end{algorithmic}
	\end{algorithm}

	\section{The decrease search}

	The same loop with the condition weakened to
	\begin{equation}
		\psi(t)-\psi(0)\ <\ 0 ,
	\end{equation}
	which proves nothing but costs one query fewer per call, needing no slope at
	all. On a strongly convex $\psi$ along a descent direction, every sufficiently
	small step satisfies it, so the loop terminates; what it does not guarantee is
	that the accepted step makes progress proportional to the slope, which is
	exactly what Armijo buys.

	Both backtracking rules give up after \texttt{ls\_max\_back} halvings and
	return the last step tried rather than failing. A step that small is not
	useful, but returning it
	lets the outer loop observe a stalled iteration and react through the
	schedule, which is a better outcome than an exception.

	\section{The unit step}

	Returns one for both pairs whatever the direction. It is the default branch of
	the dispatch and no road selects it. It exists as the null rule against which
	the others are compared: running the projection road with it is how the value
	of a line search was measured in the first place.

	\begin{remark}[The searches never look at the dual direction.]
		All three take $\Delta x$ and $\Delta s$, ignore $\Delta z$, and return the
		same length for both pairs. That is consistent, $\varphi_c$ being a function
		of the primal variables alone, but it means the dual step is chosen by a
		rule with no information about the dual residual. On the barrier road the
		two lengths are genuinely separate; on the projection road they are equal by
		omission rather than by design.
	\end{remark}

	\begin{remark}[The bracket bound depends on the penalties.]
		The bound of Proposition~\ref{prop:bracket} is $|\psi'(0)|/\kappa$, and
		$\kappa$ contains $\rho\norm{\Delta x}^2+\rho_p\norm{\Delta s}^2$. As the
		schedule drives those weights down the bound grows, so late in a solve the
		search may have to look much further out than the unit step before it
		brackets. That is not a defect of the bound, which stays valid, but it does
		mean the number of slope evaluations rises as the penalties fall.
	\end{remark}

	\chapter{The parameter schedule}

	\section{Where it happens}

	The schedule owns everything the outer loop changes between subproblems: the
	four centres, the proximal weights $\rho$ and $\rho_p$, the two penalty
	vectors, and the tolerance the inner loop stops against. It is
	\texttt{schedule/}, $407$ lines over two files, and it holds three
	implementations behind one call.

	\begin{quote}
		\texttt{update(iterate, violation, reference, primal\_residual,
			dual\_residual, inner\_converged, inner\_taken)}
	\end{quote}

	It is called once per outer iteration, at the very end, after the termination
	test and the infeasibility tests have both declined to stop. It returns
	nothing; everything it does is written into the iterate and into its own
	tolerance vector.

	\begin{table}[H]
		\centering\small
		\begin{tabular}{lll}
			\toprule
			schedule & road & penalties are\\
			\midrule
			\texttt{bcl} & projection & one value for all rows\\
			\texttt{gbcl} & projection, by default & one value per row, per block
			where curved\\
			\texttt{interior} & barrier & not used; $\rho$ and $\rho_p$ only\\
			\bottomrule
		\end{tabular}
	\end{table}

	No convergence result covers any of the three as implemented. The theory part
	gives conditions on the metrics under which the proximal iteration converges,
	summable tolerances among them, and the schedules below satisfy none of them
	exactly. What follows is what they do, and why, not a proof that it works.

    \begin{remark}[Penalty]
    	The element of $M$ are called "penalties", the term is misleading here given that when penalties get smaller constraint violation are penalize more i.e they contribute more to the objective of the smoothed problem. 
    \end{remark}

	\section{What the schedule is trying to balance}
	Every subproblem is easier when the penalties are large, because the cone term
	is then weakly weighted and the Newton system is well conditioned; and every
	subproblem is a worse approximation of the original problem for the same
	reason. The schedule moves the penalties down when the iterate is behaving
	and holds them when it is not.

	The quantity it watches is the \emph{violation}, one number per row, measuring
	how far that row is from satisfying its own cone constraint, and the
	\emph{tolerance} it is compared against, also one number per row. Those two
	vectors are the schedule's state; everything else is derived from them.

	\section{The parameters, once}

	No number appears in the algorithms below that is not a setting. The two
	projection schedules do not share a parameter block: each names its own, and
	they are tabulated where each schedule is described. The pair
	$\mu_{\mathrm{restart}}$ and $\mu_{\mathrm{th}}$ belongs to \texttt{bcl}
	alone. The defaults
	are the subject of a later chapter; what matters here is which quantity each
	one controls.

	\begin{table}[H]
		\centering\small
		\begin{tabular}{lll}
			\toprule
			symbol & setting & controls, \texttt{bcl}\\
			\midrule
			$\alpha$ & \texttt{alpha} & how fast the tolerance loosens on a reject\\
			$\beta$ & \texttt{beta} & how fast it tightens on an accept\\
			$\kappa$ & \texttt{mu\_update\_factor} & the factor a rejected penalty
			is cut by\\
			$\mu_{\min}$ & \texttt{mu\_min\_in} & the floor under every penalty\\
			$N_{\mathrm{guard}}$ & \texttt{safe\_guard} & after how many outers an
			accept is forced\\
			$\mu_{\mathrm{restart}}$, $\mu_{\mathrm{th}}$
			& \texttt{cold\_reset}$\ast$ & the restart and its trigger\\
			\bottomrule
		\end{tabular}
	\end{table}

	Two derived quantities are fixed at setup and used throughout by \texttt{bcl}:
	\begin{equation}
		\bar w_0=(0.1)^{\alpha},
		\qquad
		\varepsilon^{\min}_{\mathrm{in}}
		=\min\bigl(\varepsilon_{\mathrm{abs}},\,10^{-9}\bigr) ,
	\end{equation}
	the first the tolerance a rejected outer iteration is reset from, the second
	the floor under the inner tolerance, so that a schedule cannot ask the inner
	loop for an accuracy the outer test would not itself demand.

	\section{\texttt{bcl}, one penalty for the whole problem}

	The classical rule, and here deliberately a faithful one: it is
	\textsc{proxqp}'s \texttt{bcl\_update} written in this solver's variables,
	kept as a reference against which the generalisation is measured rather than
	tuned on its own account. A single scalar decides for every row at once, and
	the decision is taken on the \emph{primal residual} of the original problem,
	not on the violation vector.

	\begin{algorithm}[H]
		\caption{\texttt{bcl}, one outer update}
		\begin{algorithmic}[1]
			\Require the primal and dual residuals
			\State $\mu\gets\max_i(\text{penalty})_i$
			\If{$\text{primal residual}\leq\bar w$
				\textbf{ or } $\text{taken}>N_{\mathrm{guard}}$}
			\Comment{accept}
			\State move all four centres onto the current iterate
			\State $\bar w\gets\bar w\,\mu^{\beta}$,\quad
			$\varepsilon_{\mathrm{in}}\gets
			\max\bigl(\varepsilon_{\mathrm{in}}\,\mu,\
			\varepsilon^{\min}_{\mathrm{in}}\bigr)$
			\Else
			\Comment{reject}
			\State restore $y$ and $z$ to their values before the subproblem
			\State $\mu\gets\max(\mu_{\min},\ \kappa\,\mu)$ on every row
			\State $\bar w\gets\bar w_0\,\mu^{\alpha}$,\quad
			$\varepsilon_{\mathrm{in}}\gets
			\max\bigl(\mu,\ \varepsilon^{\min}_{\mathrm{in}}\bigr)$
			\EndIf
			\If{neither residual improved \textbf{and} $\mu\leq\mu_{\mathrm{th}}$}
			\State reset every penalty to $\mu_{\mathrm{restart}}$
			\EndIf
		\end{algorithmic}
	\end{algorithm}

	\section{\texttt{gbcl}, one penalty per row}

	The default, and the generalisation: \textsc{bcl}'s decision taken for each
	constraint separately rather than for the problem as a whole. It carries three
	tolerances rather than one. A bar per row, $\bar w_i$, decides which rows may
	move their centres and how hard each penalty is pushed. A scalar bar, $b$,
	decides when the proximal centre follows the iterate. A Newton tolerance,
	$\varepsilon_N$, stops the inner loop. The three are scheduled separately, and
	only the first is what \textsc{bcl} calls the bar.

	\begin{table}[H]
		\centering\small
		\begin{tabular}{lll}
			\toprule
			symbol & setting & controls\\
			\midrule
			$\mathrm{e}$ & \texttt{mu\_exp} & how hard a miss is punished\\
			$\kappa$ & \texttt{penalty\_reduction} & the hardest push allowed, and
			where it starts\\
			$a$ & \texttt{mu\_adapt} & how fast that push is retuned\\
			$\theta$ & \texttt{penalty\_reduction} & how hard a failed inner loop
			raises a penalty\\
			$\alpha$ & \texttt{alpha} & how fast a missed bar loosens\\
			$\beta$ & \texttt{beta} & how fast a met bar tightens\\
			$\mu_0$ & \texttt{mu\_in} & the first penalty, and the ceiling over
			every later one\\
			$\mu_{\mathrm{eq}}$ & \texttt{mu\_eq} & the same, for the equality
			penalties\\
			$\mu_{\min}$ & \texttt{mu\_min} & the floor under every penalty\\
			$\bar w_0$ & \texttt{eps\_outer\_init} & the bar every row starts at\\
			$\varepsilon_{N,0}$ & \texttt{eps\_newton\_init} & the Newton tolerance
			every outer starts at\\
			$\varrho_p$ & \texttt{rho\_p\_outer} & how fast the proximal slack
			falls\\
			$N$ & \texttt{safe\_guard} & how many outers the proximal centre may
			stand still for\\
			$\varepsilon$ & \texttt{eps\_abs} & the floor under the bars\\
			$\bar w_{\min}$ & \texttt{smallest\_tolerance} & the floor under the
			Newton tolerance\\
			\bottomrule
		\end{tabular}
	\end{table}

	Two floors sit under the tolerances and they are deliberately different. The
	bars stop at the accuracy the solve is asked for, since a row held tighter
	than that buys nothing. The Newton tolerance stops far below it, because the
	inner residual is not the quantity being certified: the subproblem has to be
	driven well past the target for the residual that is certified to reach it. A
	Newton tolerance floored at $\varepsilon$ instead lets the inner loop declare
	an iterate solved that the outer loop cannot improve, and the outer loop then
	repeats it unchanged.

	\begin{algorithm}[H]
		\caption{$\kappa$, retuned on the observed efficiency}
		\begin{algorithmic}[1]
			\Require the primal residual and the number of inner steps taken
			\If{an earlier outer exists \textbf{and} at least one inner step was
				taken}
			\State $\eta\gets\log_{10}\bigl(\text{previous residual}/
			\text{residual}\bigr)/\text{inner steps taken}$
			\Comment{digits per step}
			\If{$\eta$ fell short of the previous $\eta$}
			\State reverse the direction
			\EndIf
			\State $\kappa\gets\kappa\,a^{\text{direction}}$, kept clear of zero and
			of one
			\EndIf
			\State keep the residual as the previous one
		\end{algorithmic}
	\end{algorithm}

	\begin{algorithm}[H]
		\caption{\texttt{gbcl}, one outer update}
		\begin{algorithmic}[1]
			\Require the violation per row, the primal residual, whether the inner
			loop converged, the number of inner steps taken
			\For{every row $i$ with $\text{violation}_i\leq\bar w_i$}
			\State move that row's slack, dual and equality centres
			\EndFor
			\If{the inner loop did not converge}
			\State $\mu_{\mathrm{floor}}\gets\max\bigl(\mu_{\mathrm{floor}},\
			\min(\mu_{\max}/\theta,\ \mu_0)\bigr)$
			\Comment{ratcheted, never undone}
			\For{every row $i$ with $\text{violation}_i>\bar w_i$}
			\State $\mu_i\gets\min(\mu_i/\theta,\ \mu_0)$,\quad
			$\mu_{s,i}\gets\min(\mu_{s,i}/\theta,\ \mu_{\mathrm{eq}})$
			\EndFor
			\Else
			\State retune $\kappa$
			\For{every row $i$}
			\State $\displaystyle f_i\gets\min\Bigl(1,\ \max\bigl(\kappa,\
			(\text{violation}_i/\bar w_i)^{-\mathrm{e}}\bigr)\Bigr)$
			\EndFor
			\State flatten $f$ on every curved block, by the minimum
			\Comment{admissibility}
			\State $\mu\gets\max\bigl(\max(\mu_{\mathrm{floor}},\mu_{\min}),\
			\mu\odot f\bigr)$,\quad
			$\mu_s\gets\max(\mu_{\min},\ \mu_s\odot f)$
			\EndIf
			\For{every row $i$}
			\State $\bar w_i\gets
			\begin{cases}
				\max\bigl(\bar w_i\,\mu_i^{\beta},\ \varepsilon\bigr)
				 & \text{if }\text{violation}_i\leq\bar w_i,\\
				\bar w_0\,\mu_i^{\alpha} & \text{otherwise}
			\end{cases}$
			\Comment{the bar loosens on a miss}
			\EndFor
			\State $\mu_{\max}\gets\max_i\mu_i$
			\If{the primal residual met $b$ \textbf{or} $N$ outers have passed}
			\State move the proximal centre onto the iterate
			\State $\varepsilon_N\gets\max(\varepsilon_N\,\mu_{\max},\
			\bar w_{\min})$,\quad
			$b\gets\max\bigl(b\,\mu_{\max}^{\beta},\ \varepsilon\bigr)$
			\Else
			\State $\varepsilon_N\gets
			\max(\varepsilon_N\,\mu_{\max},\ \bar w_{\min})$ if no inner step was
			taken, otherwise $\varepsilon_{N,0}\,\mu_{\max}$
			\State $b\gets\bar w_0\,\mu_{\max}^{\alpha}$
			\EndIf
			\State $\rho_p\gets\max(\rho_{p,\min},\ \varrho_p\,\rho_p)$
		\end{algorithmic}
	\end{algorithm}

	Three things separate this from \textsc{bcl} beyond the per row decision. A
	failed inner loop raises the penalties of the rows that missed, rather than
	leaving them, and ratchets a floor under every penalty that no later
	reduction may cross; the floor is what stops a penalty from being driven back
	down into the region that produced the failure. The push $\kappa$ is not a
	constant: each converged outer measures how many digits of violation the
	inner loop bought per step, and moves $\kappa$ in whichever direction that
	measurement last improved. And a row that misses its bar has the bar
	loosened rather than tightened, so a row that cannot be met does not drag the
	whole schedule to a floor it will never reach.

	The proximal centre is on its own clock. It follows the iterate only when the
	primal residual meets the scalar bar $b$, or when $N$ outers have passed
	without that happening, which bounds how long the centre may stand still.
	Between those moves the Newton tolerance is reset from $\varepsilon_{N,0}$,
	except when the previous inner loop took no step at all: there the tolerance
	is tightened instead, since resetting it would ask for work already done.


	\section{\texttt{interior}, for the barrier road}

	The barrier road has no penalties to schedule, its cone term being carried by
	$\tau$ and the centring parameter, which the smoothing manages itself. What
	remains is $\rho$ and $\rho_p$, and the schedule treats them separately,
	each with its own stall counter.

	Its parameters are its own, shared with neither projection schedule.

	\begin{table}[H]
		\centering\small
		\begin{tabular}{lll}
			\toprule
			symbol & setting & controls\\
			\midrule
			$\iota$ & \texttt{improvement} & how much a residual must fall to count
			as progress\\
			$\varsigma$ & \texttt{stalled\_rate} & how much of the rate a stalled
			iteration still gets\\
			$\varphi_{\mathrm{coarse}}$ & \texttt{reg\_floor} & the floor under
			$\rho$ and $\rho_p$ at first\\
			$\varphi_{\mathrm{fine}}$ & \texttt{reg\_fine} & the floor after the
			ratchet\\
			$N_{\mathrm{stall}}$ & \texttt{stall\_before\_fine} & how many stalls
			at the floor before it drops\\
			$N_{\mathrm{early}}$ & \texttt{early\_outers} & how long a stalled
			iteration is still indulged\\
			$\beta_{\mathrm{prox}}$ & \texttt{infeasibility\_threshold} & the bar
			the proximal contribution is judged against\\
			$\varepsilon_{\mathrm{reg}}$ & \texttt{eps\_reg} & below which a
			residual counts as converged outright\\
			\bottomrule
		\end{tabular}
	\end{table}

	\begin{algorithm}[H]
		\caption{\texttt{interior}, one outer update}
		\begin{algorithmic}[1]
			\State $r\gets$ the relative decrease of the duality measure
			\State $\text{dual\_prox}\gets\rho\norm{x-x_c}_\infty$, \quad
			$\text{primal\_prox}\gets\rho_p\norm{s-s_c}_\infty$
			\Comment{how much of each residual the proximal term itself contributes}
			\If{$\text{dual residual}<\iota\cdot\text{previous}$, or it is already
				below $\varepsilon_{\mathrm{reg}}$, or the floor is fine and
				$\text{dual\_prox}<\beta_{\mathrm{prox}}$}
			\State move the $x$ centre,\quad
			$\rho\gets\max\bigl(\varphi,\ (1-r)\,\rho\bigr)$
			\Comment{$\varphi$ the current floor}
			\Else
			\State count a dual stall; reduce $\rho$ only at the damped rate
			$1-\varsigma\,r$, and only while $\text{taken}<N_{\mathrm{early}}$ or
			$\text{dual\_prox}<\beta_{\mathrm{prox}}$
			\EndIf
			\State the same, symmetrically, for the primal residual, $\rho_p$ and
			the slack and dual centres
			\If{either has stalled more than $N_{\mathrm{stall}}$ times at its floor,
				and both proximal parts are below $\beta_{\mathrm{prox}}$}
			\State lower the floor from $\varphi_{\mathrm{coarse}}$ to
			$\varphi_{\mathrm{fine}}$
			\EndIf
		\end{algorithmic}
	\end{algorithm}


	\chapter{The cone kernel}

	\section{Where it happens}

	The kernel is the largest directory in the solver, $1957$ lines over ten
	files, because it carries six independent cones behind one interface. The
	interface is a set of free functions in \texttt{cone/kernel.hpp}, dispatched
	on the cone type by a variant visit, and every cone provides the same twenty.
	Adding a cone is implementing those twenty; nothing else changes.

	\begin{table}[H]
		\centering\small
		\begin{tabular}{ll}
			\toprule
			function & what it computes\\
			\midrule
			\texttt{rows}, \texttt{degree} & the block size and the barrier
			parameter $\nu$\\
			\texttt{is\_curved} & whether admissibility forces $M=cI$ on the block\\
			\texttt{has\_interior} & whether $\intr K\neq\emptyset$\\
			\texttt{contains}, \texttt{margin} & membership, and
			Definition~\ref{def:margin}\\
			\texttt{evaluate\_projection} & $\proj{v}{K^*}$\\
			\texttt{build\_projection\_scaling} & the cache $F$ is read from, at
			$\sigma$\\
			\texttt{build\_barrier\_scaling} & the scaling point, at $(s,z)$\\
			\texttt{materialise\_operator} & the block of $G_\phi$\\
			\texttt{widest\_operator\_kind} & the widest block shape the cone will
			return\\
			\texttt{apply\_residual\_scaling} & $F$, or its barrier analogue,
			applied\\
			\texttt{complementarity\_residual} & the centrality residual\\
			\texttt{second\_order\_correction} & the Mehrotra corrector\\
			\texttt{largest\_feasible\_step} & the $\bar\alpha$ of
			Definition~\ref{def:fraction}\\
			\texttt{interior\_direction} & a point of $\intr K$, for the cold start\\
			\bottomrule
		\end{tabular}
		\caption{The kernel interface. A cone that cannot support a function raises,
			and the solver never reaches it.}
	\end{table}

	\begin{table}[H]
		\centering\small
		\begin{tabular}{lccccl}
			\toprule
			cone & $K^*$ & symmetric & $\nu$ & curved & $G$ block\\
			\midrule
			\texttt{Zero} & $\R^d$ & no & $0$ & no & diagonal\\
			\texttt{Nonneg} & itself & yes & $d$ & no & diagonal\\
			\texttt{SecondOrder} & itself & yes & $2$ & yes & diagonal, rank two\\
			\texttt{PSDTriangle} & itself & yes & $k$ & yes & dense\\
			\texttt{Exponential} & not itself & no & $3$ & yes & dense\\
			\texttt{Power} & not itself & no & $3$ & yes & dense\\
			\bottomrule
		\end{tabular}
	\end{table}

	Each cone below is given in the same order: membership, projection, the
	selected derivative, the barrier scaling, the cone operator, the boundary
	step, and cost.

	\section{The zero cone}

	$K=\{0\}\subseteq\R^d$ and $K^*=\R^d$, so the block is an equality constraint.
	It is closed and convex, not solid, and not self dual.

	The barrier is $B\equiv0$ with $\nu=0$, as fixed in the theory part: the
	relative boundary of $\{0\}$ is empty, so every condition holds vacuously. A
	zero block contributes nothing to $\tau$ and has no centrality condition.

	The projection is the identity, $F=I$ at every point, and
	\[
	G=\rho_pI+M^{-1}=\rho_p+\frac1\mu ,
	\]
	diagonal and independent of the iterate. There is no case analysis and no
	algorithm to give.

	The implementation declines the barrier road here, and the restriction is its
	own rather than the theory's: the machinery wants a strictly feasible point, a
	scaling built from it and a fraction to the boundary, none of which exists
	when $\intr K=\emptyset$. Since \texttt{has\_interior} returns false, the
	interior road routes such a block through the equality penalty, exactly as the
	projection road does.

	\section{The nonnegative orthant}

	$K=K^*=\R^d_+$, symmetric of rank $d$, with $\nu=d$; the canonical inner product
	is the Euclidean one, so the constant $C$ of the centrality discussion is one.
	The barrier is $B(s)=-\sum_i\log s_i$, giving $\nabla B=-s^{-1}$ and
	$\nabla^2B=\diag(s^{-2})$, elementwise.

	Everything on this cone is elementwise, and the formula is the algorithm:
	\[
	\proj{v}{K^*}=\max(v,0),
	\qquad
	F=\diag\bigl(\mathbf{1}_{\sigma>0}\bigr),
	\qquad
	G=\rho_p+\frac{\mathbf{1}_{\sigma>0}}{\mu} .
	\]
	The Jacobian is a $0/1$ mask, and it is idempotent, $F^2=F$, which the orthant
	shares with no other cone in the kernel. An inactive row contributes $\rho_p$
	alone, which is how the projection road identifies the active set: a row whose
	$\sigma$ is negative drops out of the system entirely.

	The barrier scaling is closed form, two square roots per row,
	\[
	w=\sqrt{s/z},
	\qquad
	\lambda=\sqrt{s\circ z},
	\qquad
	G=\rho_p+\frac{z}{s} ,
	\]
	
	The centrality residual is $\lambda\circ\lambda-\tau\mathbf{1}$, so
	$\inner{s}{z}=\sum_i\lambda_i^2=d\tau=\nu\tau$; the corrector is
	$\Delta s\circ\Delta z$, which is exactly the omitted quadratic term rather
	than an estimate of it. The step is
	$\min\{-v_i/\Delta v_i:\Delta v_i<0\}$ and the margin is $\min_iv_i$, both
	$O(d)$.

	\section{The second order cone}

	\[
	\mathcal{Q}^d=\bigl\{(t,u)\in\R\times\R^{d-1}:\norm{u}\leq t\bigr\} ,
	\]
	self dual and symmetric, with Jordan product
	$x\circ y=(\inner{x}{y},\,x_0y_{1:}+y_0x_{1:})$, identity $e=(1,0,\dots,0)$,
	determinant $\det v=t^2-\norm{u}^2$ and rank two. The barrier is
	$B=-\log(t^2-\norm{u}^2)$ with $\nabla_{\mathrm{Euc}}B=-2v^{-1}$ and $\nu=2$;
	the canonical trace inner product is twice the Euclidean one, so $C=2$ here
	against $C=1$ on the orthant and the semidefinite cone.

	\begin{proposition}[Projection onto the second order cone]\label{prop:socproj}
		Let $v=(t,u)$ and $K=\mathcal{Q}^d$. Then
		\begin{equation}
			\proj{v}{K}=
			\begin{cases}
				v & \norm{u}\leq t,\\[1mm]
				0 & \norm{u}\leq -t,\\[1mm]
				c\,\bigl(\norm{u},\,u\bigr) & \text{otherwise},
			\end{cases}
			\qquad
			c=\frac12\Bigl(1+\frac{t}{\norm{u}}\Bigr)\in(0,1) .
		\end{equation}
	\end{proposition}

	\begin{proof}
		The first two cases are $v\in K$ and $v\in K^{\circ}=-K$. Otherwise
		$|t|<\norm{u}$, so $c\in(0,1)$ and $p=c(\norm{u},u)$ lies on the boundary,
		its apex equalling its tail norm. Moreau asks three things. First $p\in K$,
		with equality. Second $v-p\in K^{\circ}$: from
		$c\norm{u}=\tfrac12(\norm{u}+t)$,
		\[
		t-c\norm{u}=-\tfrac12(\norm{u}-t),
		\qquad
		(1-c)\norm{u}=\tfrac12(\norm{u}-t),
		\]
		so the apex of $v-p$ is the negative of its tail norm. Third
		$\inner{p}{v-p}=c\norm{u}\bigl(t+\norm{u}-2c\norm{u}\bigr)=0$. Uniqueness
		completes the proof.
	\end{proof}

	\begin{algorithm}[H]
		\caption{Projection onto the second order cone}
		\begin{algorithmic}[1]
			\Require $v=(t,u)$
			\State $n\gets\norm{u}$
			\If{$n\leq t$} \State \Return $v$ \Comment{already in the cone}
			\ElsIf{$n\leq-t$} \State \Return $0$ \Comment{in the polar cone}
			\Else
			\State $c\gets\tfrac12(1+t/n)$
			\State \Return $\bigl(cn,\ cu\bigr)$
			\EndIf
		\end{algorithmic}
	\end{algorithm}

	The tests use $\leq$, so the three regions meet exactly on the boundary and no
	point falls between them.

	\begin{proposition}[The Jacobian is a rank two correction]\label{prop:socrank}
		In the boundary case, with $\widehat u=u/\norm{u}$,
		$\widehat e=(0,\widehat u)$ and $e_0$ the apex direction,
		\begin{equation}
			F=c\,I+
			\bigl[\,\widehat e\ \ e_0\,\bigr]\;T\;
			\bigl[\,\widehat e\ \ e_0\,\bigr]^{\top},
			\qquad
			T=\begin{pmatrix}
				-\tfrac12\,t/\norm{u} & \tfrac12\\[2mm]
				\tfrac12 & \tfrac12-c
			\end{pmatrix} .
		\end{equation}
	\end{proposition}

	\begin{proof}
		In that region $c$ is smooth and $\nabla_u\norm{u}=\widehat u$, so
		\[
		\nabla c=\frac{1}{2\norm{u}}
		\Bigl(1,\ -\frac{t}{\norm{u}}\,\widehat u\Bigr) .
		\]
		Differentiating the tail $cu$ gives $cI_{d-1}+u(\nabla c)^{\top}$, rank one
		along $\widehat u$ in its second term; differentiating the apex $c\norm{u}$
		gives $\norm{u}(\nabla c)^{\top}+c\widehat u^{\top}$, one further row. Both
		lie in the span of $\widehat e$ and $e_0$, and matching coefficients gives
		$T$.
	\end{proof}

	Away from the boundary the Jacobian is not that at all, and the code
	distinguishes four regions rather than three. The extra one is numerical: when
	$u$ vanishes to machine precision the boundary formula would divide by
	$\norm{u}$, so that case is decided by the sign of $t$ alone.

	\begin{algorithm}[H]
		\caption{Applying $F\in\partial\proj{\cdot}{\mathcal{Q}^d}(\sigma)$}
		\begin{algorithmic}[1]
			\Require $\sigma=(t,u)$ and a vector $v=(v_0,v_{1:})$
			\State $n\gets\norm{u}$
			\If{$n\leq\varepsilon_{\mathrm{tiny}}$}
			\State \Return $v$ if $t\geq0$, else $0$
			\Comment{$u$ vanishes; the tail carries no direction}
			\ElsIf{$n<t$} \State \Return $v$ \Comment{interior, $F=I$}
			\ElsIf{$n<-t$} \State \Return $0$ \Comment{polar interior, $F=0$}
			\Else
			\State $c\gets\tfrac12(1+t/n)$, \quad $\widehat u\gets u/n$
			\State $\pi\gets\inner{\widehat u}{v_{1:}}$, \quad
			$\kappa\gets\tfrac12\bigl(v_0-(t/n)\pi\bigr)/n$
			\State \Return
			$\bigl(\tfrac12\pi+\tfrac12v_0,\ \ c\,v_{1:}+\kappa\,n\,\widehat
			u\bigr)$
			\EndIf
		\end{algorithmic}
	\end{algorithm}

	The block $G=\rho_pI+M^{-1}F$ inherits those four regions. In three of them
	$F$ is a multiple of the identity and the block is a plain diagonal shift,
	$\rho_p+1/\mu$ inside the cone and $\rho_p$ inside the polar cone; only in the
	boundary region does it become diagonal plus rank two. This is why
	\texttt{widest\_operator\_kind} answers \texttt{DiagPlusLowRank} for this cone
	although three of its four regions produce a diagonal: the pattern must
	accommodate the widest case.

	\begin{algorithm}[H]
		\caption{Nesterov and Todd scaling for the second order cone}
		\begin{algorithmic}[1]
			\Require $s,z\in\intr\mathcal{Q}^d$
			\State $\rho_s\gets\sqrt{\det s}$, \quad $\rho_z\gets\sqrt{\det z}$
			\State $\widehat s\gets s/\rho_s$, \quad $\widehat z\gets z/\rho_z$
			\Comment{both now of unit determinant}
			\State $\gamma\gets\sqrt{(1+\inner{\widehat s}{\widehat z})/2}$
			\State $w_0\gets(\widehat z_0+\widehat s_0)/(2\gamma)$, \quad
			$w_{1:}\gets(\widehat z_{1:}-\widehat s_{1:})/(2\gamma)$
			\Comment{sum on the apex, difference on the tail}
			\State $f\gets\sqrt{\rho_z/\rho_s}$, \quad $\lambda\gets f\,W^{1/2}s$
			\Comment{applied by the arrow formula, $O(d)$}
		\end{algorithmic}
	\end{algorithm}

	Reflecting the tail of $w$ gives $W^{-1/2}$, which is why the corrector
	reflects rather than inverting anything. On the barrier road, with
	$J=2e_0e_0^{\top}-I$,
	\[
	G=\rho_pI+f^2\bigl(2ww^{\top}-J\bigr)
	=(\rho_p+f^2)I+2f^2ww^{\top}-2f^2e_0e_0^{\top} ,
	\]
	and here there is no case analysis, the barrier road visiting only the
	interior.

	The centrality residual is $\lambda\circ\lambda-2\tau e$, the factor two being
	$C=2$; the corrector is the Jordan product of the scaled steps. The step is
	the largest $\alpha$ with $\det(v+\alpha\Delta v)\geq0$ and
	$v_0+\alpha\Delta v_0\geq0$, and since the determinant is a quadratic in
	$\alpha$ this is one discriminant and a choice between two roots, $O(d)$. The
	margin is $v_0-\norm{v_{1:}}$ and the interior direction is $e_0$.

	\section{The semidefinite cone}

	The cone is $\Sp^k_+$, self dual and symmetric of rank $k$, so $\nu=k$,
	carried in a packed form of $d=k(k+1)/2$ rows. The packing is not the naive
	one.

	\begin{proposition}[\texttt{svec} is an isometry]\label{prop:svec}
		Let $\svec$ take the lower triangle in column order, multiplying every off
		diagonal entry by $\sqrt2$ and leaving the diagonal alone. Then
		$\inner{\svec X}{\svec Y}=\tr(XY)$ for all $X,Y\in\Sp^k$, and $\svec$
		carries $\Sp^k_+$ onto a closed convex self dual cone.
	\end{proposition}

	\begin{proof}
		Each off diagonal pair is counted twice in the trace, so
		$$\tr(XY)=\sum_iX_{ii}Y_{ii}+2\sum_{i>j}X_{ij}Y_{ij}$$. The packed vectors
		contribute the same diagonal sum and
		$$\sum_{i>j}(\sqrt2X_{ij})(\sqrt2Y_{ij})=2\sum_{i>j}X_{ij}Y_{ij}$$ from the
		rest. Self duality is inherited, a linear isometry carrying a self dual cone
		to a self dual cone.
	\end{proof}

	Without the $\sqrt2$ the map is still a bijection but no longer preserves
	inner products, and every pairing in the solver, the duality gap among them,
	would be wrong on semidefinite blocks by a factor depending on how much mass
	sits off the diagonal. The scaling appears in two functions only,
	\texttt{matrix\_from\_svec} and \texttt{svec\_from\_matrix}.

	The barrier is $B(X)=-\log\det X$ with $\nabla B=-X^{-1}$, and
	$\det(\lambda X)=\lambda^k\det X$ gives $\nu=k$; the canonical inner product
	is Frobenius, so $C=1$.

	\begin{proposition}[Projection onto the semidefinite cone]
		For $X=U\Lambda U^{\top}$, one has
		$\proj{X}{\Sp^k_+}=U\max(\Lambda,0)U^{\top}$.
	\end{proposition}

	\begin{proof}
		Set $P=U\max(\Lambda,0)U^{\top}\succeq0$. Then
		$X-P=U\min(\Lambda,0)U^{\top}$ lies in $-\Sp^k_+$, and their product is
		$U\max(\Lambda,0)\min(\Lambda,0)U^{\top}=0$, the two diagonal factors having
		disjoint support, so $\inner{P}{X-P}=0$. Moreau's characterisation applies.
	\end{proof}

	\begin{algorithm}[H]
		\caption{Projection onto the semidefinite cone}
		\begin{algorithmic}[1]
			\Require the packed $v$
			\State $X\gets\smat(v)$ \Comment{unpack, dividing off diagonals by
				$\sqrt2$}
			\State $(U,\Lambda)\gets$ symmetric eigendecomposition of $X$
			\Comment{$O(k^3)$}
			\State \Return $\svec\bigl(U\max(\Lambda,0)U^{\top}\bigr)$
		\end{algorithmic}
	\end{algorithm}

	There is no case analysis: the clamp handles every sign pattern uniformly, and
	the interior, boundary and polar cases of the other cones appear here only as
	the number of nonpositive eigenvalues.

	\begin{proposition}[The Jacobian is a Löwner operator]\label{prop:lowner}
		At an $X$ with eigenvalues $\lambda_i$ and eigenvectors $U$,
		\begin{equation}
			D\proj{\cdot}{\Sp^k_+}(X)[H]
			=U\Bigl(\Omega\circ\bigl(U^{\top}HU\bigr)\Bigr)U^{\top},
			\qquad
			\Omega_{ij}=
			\begin{cases}
				\dfrac{\lambda_i^+-\lambda_j^+}{\lambda_i-\lambda_j}
				& \lambda_i\neq\lambda_j,\\[3mm]
				\mathbf{1}_{\lambda_i>0} & \lambda_i=\lambda_j ,
			\end{cases}
		\end{equation}
		with $\circ$ the entrywise product and $\lambda^+=\max(\lambda,0)$.
	\end{proposition}


	The two branches of $\Omega$ are the only case analysis on this cone, and the
	test is relative rather than exact: two eigenvalues count as equal when their
	gap falls below machine epsilon times the larger of their magnitudes and one.

	\begin{algorithm}[H]
		\caption{Applying the Löwner Jacobian}
		\begin{algorithmic}[1]
			\Require the cached $U$ and $\Omega$, a packed $v$
			\State $H\gets\smat(v)$
			\State $\widetilde H\gets\Omega\circ\bigl(U^{\top}HU\bigr)$
			\Comment{into the eigenbasis, entrywise, back out}
			\State \Return $\svec\bigl(U\widetilde HU^{\top}\bigr)$
			\Comment{two congruences, $O(k^3)$}
		\end{algorithmic}
	\end{algorithm}

	\begin{algorithm}[H]
		\caption{Materialising the $G$ block}
		\begin{algorithmic}[1]
			\Require the scaling at $\sigma$, the weights $\rho_p$ and $\mu$
			\For{$j=1,\dots,d$}
			\State column $j$ of $G\gets$ the Löwner Jacobian applied to $e_j$,
			divided by $\mu$
			\Comment{one call above, $O(k^3)$}
			\EndFor
			\State add $\rho_p$ to the diagonal
		\end{algorithmic}
	\end{algorithm}

	\begin{proposition}[Cost of materialising a semidefinite block]
		\label{prop:psdcost}
		Forming the dense $G$ block of a semidefinite cone of side $k$ costs
		$O(k^5)$.
	\end{proposition}

	\begin{proof}
		The loop runs $d=k(k+1)/2$ times, and each pass applies the Löwner
		Jacobian, two matrix products and an entrywise product, $O(k^3)$. The total
		is $O(dk^3)=O(k^5)$.
	\end{proof}

	This dominates any problem with a large semidefinite block and is why
	\texttt{random\_psd} is the slowest of the generated sets. The remedy is not
	to materialise: a road that only applies $G$ to a vector needs one congruence
	per application, $O(k^3)$. Both roads materialise because both place the block
	in a sparse matrix.

	\begin{algorithm}[H]
		\caption{Nesterov and Todd scaling for the semidefinite cone}
		\begin{algorithmic}[1]
			\Require $s,z$ packed, with $X=\smat(s)\succ0$ and $Z=\smat(z)\succ0$
			\State $(X^{1/2},X^{-1/2})\gets$ symmetric square roots of $X$
			\Comment{one eigendecomposition, eigenvalues floored}
			\State $Y\gets(X^{1/2}ZX^{1/2})^{1/2}$ \Comment{a second}
			\State $W\gets X^{-1/2}YX^{-1/2}$, \quad $(Q,\Theta)\gets$ its
			eigendecomposition \Comment{a third}
			\State $W^{\pm1/2}\gets Q\Theta^{\pm1/2}Q^{\top}$
			\State $L\gets W^{1/2}XW^{1/2}$, symmetrised; $(U,\lambda)\gets$ its
			eigendecomposition \Comment{a fourth}
		\end{algorithmic}
	\end{algorithm}

	Every square root floors its eigenvalues at a relative $10^{-14}$, so a block
	that has drifted to the boundary yields a large but finite scaling rather than
	a division by zero. The symmetrisation of $L$ is not cosmetic: a product of
	two symmetric matrices is not symmetric, and passing an asymmetric matrix to a
	symmetric eigensolver silently discards half of it.

	The centrality residual is $\lambda_i^2-\tau$ on the $k$ eigenvalues,
	reassembled by congruence, so $\inner{X}{Z}=k\tau=\nu\tau$; the corrector is
	the symmetrised product of the congruence scaled steps, symmetrised for the
	reason just given. The step is $-1/\lambda_{\min}(X^{-1/2}\Delta XX^{-1/2})$
	when that eigenvalue is negative and infinite otherwise, $O(k^3)$; the margin
	is $\lambda_{\min}(X)$ and the interior direction is $\svec(I)$. Membership is
	tested by attempting a Cholesky of $X$ shifted by a relative $10^{-9}$, which
	is cheaper than an eigendecomposition and answers the same question.

	\section{The exponential and power cones}

	Neither has a closed form projection, and both reduce the dual projection to
	the primal one the same way.

	\begin{proposition}[Reduction by Moreau]\label{prop:moreaured}
		For any closed convex cone $K$ and any $v$, one has
		$\proj{v}{K^*}=v+\proj{-v}{K}$.
	\end{proposition}


	Both cones therefore implement the projection onto $K$ itself, in the single
	line \texttt{out = point + project\_onto\_cone(-point)}.

	\subsection{The exponential cone}

	\[
	\Kexp=\overline{\bigl\{(x,y,z):y>0,\ y\,e^{x/y}\leq z\bigr\}} ,
	\]
	not self dual, not symmetric, no Jordan product, block size always three,
	$\nu=3$. With the slack $\ell=y\log(z/y)-x$ the barrier is
	$B=-\log\ell-\log y-\log z$, and its gradient and Hessian are in closed form
	by the quotient rule.

	A point outside $\Kexp$ and outside its polar projects either onto a face or
	onto the smooth part of the boundary. On the smooth part the projection is
	determined by a single scalar, and the following says which.

	\begin{proposition}[The exponential boundary equation]\label{prop:expbound}
		Let $g(a,b,c)=b\,e^{a/b}-c$, so that the smooth part of $\partial\Kexp$ is
		$\{g=0,\ b>0\}$, and parametrise it by
		\begin{equation}
			p=\bigl(rb,\ b,\ b\,e^{r}\bigr),
			\qquad r=a/b\in\R,\quad b>0 .
		\end{equation}
		If $v=(x,y,z)$ projects onto that part at $p$, then
		\begin{equation}\label{eq:expnormal}
			\nabla g(p)=\bigl(e^{r},\ e^{r}(1-r),\ -1\bigr),
			\qquad
			v-p=\mu\,\nabla g(p)\ \text{ for some }\mu\geq0 ,
		\end{equation}
		and eliminating $b$ and $\mu$ leaves one equation in $r$ alone.
	\end{proposition}

	\begin{proof}
		The gradient is a direct computation:
		$\partial_ag=e^{a/b}$, $\partial_bg=e^{a/b}(1-a/b)$ and
		$\partial_cg=-1$, which at $p$ read as displayed. The projection onto a
		smooth surface has $v-p$ normal to it, hence parallel to $\nabla g(p)$,
		with a nonnegative multiplier because $v$ lies outside.

		Writing \eqref{eq:expnormal} in components,
		\[
		x-rb=\mu e^{r},
		\qquad
		y-b=\mu e^{r}(1-r),
		\qquad
		z-b\,e^{r}=-\mu .
		\]
		Dividing the second by the first removes $\mu e^{r}$ and gives
		$y-b=(1-r)(x-rb)$, that is
		\begin{equation}\label{eq:expb}
			b=\frac{y-(1-r)x}{r^{2}-r+1} ,
		\end{equation}
		the denominator being positive for every real $r$. The first and third give
		$\mu$ twice, as $(x-rb)e^{-r}$ and as $b\,e^{r}-z$, and setting those equal
		leaves an equation in $r$ and $b$; substituting \eqref{eq:expb} removes $b$.
	\end{proof}

	The implementation evaluates an algebraically equivalent rearrangement of that
	equation, chosen so that the exponential appears only as $e^{-r}$ and never as
	$e^{r}$, which cannot overflow:
	\[
	h(r)=\kappa(r)\bigl(e^{-2r}+1-r\bigr)-y+z\,e^{-r},
	\qquad
	\kappa(r)=\frac{x-r\,y}{r^2-r+1} ,
	\]
	the denominator of $\kappa$ being positive for every real $r$. Its roots are
	the roots of the equation of Proposition~\ref{prop:expbound}, and the
	projection is recovered from a root by \eqref{eq:expb} and the
	parametrisation.

	Nothing bounds $r$ in advance, which is why the search sweeps. The
	parametrisation runs over all of $\R$, $h$ is not monotone, and there is no
	analogue of the bracket the power cone enjoys below; the grid over
	$[-60,60]$ is a search for sign changes, and the interval is where $e^{-r}$
	and $e^{-2r}$ are representable.

	\begin{algorithm}[H]
		\caption{Projection onto the exponential cone}
		\begin{algorithmic}[1]
			\Require $v=(x,y,z)$
			\If{$v\in\Kexp$} \State \Return $v$
			\ElsIf{$-v\in\Kexp^*$} \State \Return $0$ \EndIf
			\State $best\gets0$, \quad $d_{best}\gets\norm{v}^2$
			\State consider $\bigl(\min(x,0),\,0,\,\max(z,0)\bigr)$
			\Comment{the axis faces}
			\If{$y>0$ and $x/y<100$}
			\State consider $\bigl(x,\,y,\,\max(z,\,y\,e^{x/y})\bigr)$
			\Comment{a vertical drop onto the smooth boundary}
			\EndIf
			\For{$1200$ values of $r$ spanning $[-60,60]$}
			\If{$h$ changes sign across the interval}
			\State bisect, at most $200$ steps, and consider the boundary point at
			the root
			\EndIf
			\EndFor
			\State \Return $best$
			\Statex
			\Statex \textbf{consider}$(p)$: if $p$ is finite and lies in $\Kexp$ to a
			tolerance, and is nearer $v$ than $best$, keep it
		\end{algorithmic}
	\end{algorithm}

	Three details of that loop are deliberate. The quantity $e^{-r}$ is carried
	multiplicatively along the sweep and recomputed from scratch every
	\texttt{kAnchorEvery} steps, costing one exponential in that many iterations
	instead of one per iteration while bounding the drift. A candidate is admitted
	only after a membership test, so a root of $h$ that is not a boundary point is
	discarded rather than returned. And the initial best is the origin at distance
	$\norm{v}$, so the search cannot return anything worse than the trivial
	answer.

	The Jacobian is not a finite difference. Once the projection is known the
	boundary is a smooth manifold near it, and the Jacobian follows from the
	implicit function theorem.

	\begin{algorithm}[H]
		\caption{The Jacobian of the projection onto the exponential cone}
		\begin{algorithmic}[1]
			\Require $\sigma\in\R^3$
			\If{$\sigma\in\Kexp$} \State \Return $I$
			\ElsIf{$-\sigma\in\Kexp^*$} \State \Return $0$ \EndIf
			\State $p\gets\proj{\sigma}{\Kexp}$
			\If{$p_y\leq\varepsilon_{\mathrm{deg}}$}
			\Comment{the projection landed on a face}
			\State \Return the diagonal with $F_{11}=1$ if $\sigma_x\leq0$ and
			$F_{33}=1$ if $\sigma_z\geq0$
			\EndIf
			\State $g\gets$ the boundary gradient at $p$, \quad
			$m\gets\inner{\sigma-p}{g}/\inner{g}{g}$, \quad
			$H\gets$ the rank one boundary Hessian
			\State \Return the leading $3\times3$ block of
			$\begin{pmatrix} I+mH & g\\ g^{\top} & 0\end{pmatrix}^{-1}$
			\Comment{solved with full pivoting}
		\end{algorithmic}
	\end{algorithm}

	The four by four system is exact up to rounding, and considerably better
	behaved than a difference quotient. The degenerate branch is the one worth
	noting: when the projection lands on a face rather than on the smooth part,
	the manifold has no tangent plane there and the implicit construction does not
	apply, so the Jacobian is taken as the diagonal indicator of the coordinates
	that survive.

	The block is dense $3\times3$ on both roads, $\rho_pI+F/\mu$ on the projection
	road and $\rho_pI+\tau_j\nabla^2B(s)$ on the barrier road, where
	$\tau_j=\inner{s_j}{z_j}/3$ is the block's own measure. The centrality
	residual is $z+\tau\nabla B(s)$, the barrier form being the only one available
	with no Jordan product to write $s\circ z$ with; Euler's identity still gives
	$\inner{s}{z}=3\tau$ at centrality. The corrector is identically zero, there
	being no product in which to form the omitted quadratic term, so these blocks
	receive pure centring. Step and margin are bisected.

	\subsection{The power cone}

	\[
	\Kpow=\bigl\{(x,y,z):x^{\alpha}y^{1-\alpha}\geq|z|,\ x\geq0,\ y\geq0\bigr\},
	\qquad\alpha\in(0,1),
	\]
	with $\alpha$ carried on the cone object, so one problem may mix blocks of
	different $\alpha$. The scheme is the same as the exponential cone's, except
	that the boundary admits an explicit parametrisation by a single multiplier,
	so the search is a bisection on a bracket known in advance rather than a sweep
	for sign changes.

	\begin{algorithm}[H]
		\caption{Projection onto the power cone}
		\begin{algorithmic}[1]
			\Require $v=(x,y,z)$ and $\alpha$
			\If{$v\in\Kpow$} \State \Return $v$
			\ElsIf{$-v\in\Kpow^*$} \State \Return $0$ \EndIf
			\State $best\gets0$, \quad $d_{best}\gets\norm{v}^2$
			\State consider $(0,\max(y,0),0)$, $(\max(x,0),0,0)$,
			$(\max(x,0),\max(y,0),0)$ \Comment{the three faces}
			\If{$|z|>0$ and the boundary value changes sign across $[0,|z|]$}
			\State bisect on the multiplier, and consider the boundary point with
			the sign of $z$ restored
			\EndIf
			\State \Return $best$
		\end{algorithmic}
	\end{algorithm}

	The parametrisation is by the normal multiplier rather than by a point of the
	boundary, and that is what makes the bracket free.

	\begin{proposition}[The power boundary equation]\label{prop:powbound}
		Take $z\geq0$, the other sign being recovered by symmetry, and let
		$g(x,y,z)=z-x^{\alpha}y^{1-\alpha}$, so that the smooth part of
		$\partial\Kpow$ is $\{g=0,\ x,y>0\}$. If $v=(X,Y,Z)$ projects onto that
		part at $p=(x,y,z)$ with multiplier $m\geq0$, then
		\begin{equation}
			\nabla g(p)=\Bigl(-\frac{\alpha z}{x},\ -\frac{(1-\alpha)z}{y},\ 1\Bigr),
			\qquad
			v-p=m\,\nabla g(p) ,
		\end{equation}
		and the three components determine $p$ from $m$ alone:
		\begin{equation}\label{eq:powfrom}
			z(m)=Z-m,
			\qquad
			x(m)^{2}-Xx(m)=m\,\alpha\,z(m),
			\qquad
			y(m)^{2}-Yy(m)=m(1-\alpha)z(m) ,
		\end{equation}
		each of the last two taken at its positive root. The projection is then a
		root of
		\begin{equation}\label{eq:powphi}
			\varphi(m)=x(m)^{\alpha}y(m)^{1-\alpha}-z(m) .
		\end{equation}
	\end{proposition}

	\begin{proof}
		On the boundary $x^{\alpha}y^{1-\alpha}=z$, so
		$\partial_xg=-\alpha x^{\alpha-1}y^{1-\alpha}=-\alpha z/x$ and likewise for
		$y$, while $\partial_zg=1$. Writing $v-p=m\nabla g(p)$ in components,
		\[
		X-x=-\frac{m\alpha z}{x},
		\qquad
		Y-y=-\frac{m(1-\alpha)z}{y},
		\qquad
		Z-z=m .
		\]
		The third is $z(m)=Z-m$. Multiplying the first by $x$ and rearranging gives
		$x^{2}-Xx=m\alpha z$, and the second likewise. Both are quadratics with a
		nonnegative constant term when $mz\geq0$, so each has exactly one
		nonnegative root. Finally $p$ lies on the boundary precisely when
		$x^{\alpha}y^{1-\alpha}=z$, which is $\varphi(m)=0$.
	\end{proof}

	\begin{proposition}[The bracket is free]\label{prop:powbracket}
		Let $v=(X,Y,Z)$ with $X,Y\geq0$ and $Z>0$ lie outside $\Kpow$. Then
		\begin{equation}
			\varphi(0)<0\ \leq\ \varphi(Z) ,
		\end{equation}
		so $[0,Z]$ brackets a root of $\varphi$ and no search for one is needed.
	\end{proposition}

	\begin{proof}
		At $m=0$ the constant terms in \eqref{eq:powfrom} vanish, so
		$x(0)=\max(X,0)=X$, $y(0)=Y$ and $z(0)=Z$, whence
		\[
		\varphi(0)=X^{\alpha}Y^{1-\alpha}-Z<0 ,
		\]
		the inequality being exactly the statement that $v\notin\Kpow$.

		At $m=Z$ one has $z(Z)=0$, so both constant terms vanish again,
		$x(Z)=X\geq0$ and $y(Z)=Y\geq0$, and
		\[
		\varphi(Z)=X^{\alpha}Y^{1-\alpha}-0\ \geq\ 0 .
		\]
		$\varphi$ is continuous on $[0,Z]$, so it has a root there.
	\end{proof}

	This is the whole difference between the two non symmetric cones. The
	exponential cone is parametrised by a ratio ranging over all of $\R$, with no
	a priori bound, so its roots must be hunted by a sweep; the power cone is
	parametrised by a multiplier that is pinned between $0$ and $Z$ by
	Proposition~\ref{prop:powbracket}, so one bisection on a known bracket
	suffices and the whole projection is $O(1)$.

	In the implementation, the two quadratics are solved by
	$\mathrm{posroot}(\ell,q)$, the positive root of $t^{2}-\ell t-q=0$, taken in
	whichever of its two algebraic forms avoids subtracting nearly equal numbers,
	and $\varphi$ is evaluated as
	$\exp\bigl(\alpha\log x+(1-\alpha)\log y\bigr)-z$ so that the power is formed
	in logarithms rather than by two exponentiations.

	Everything else is identical in form to the exponential cone: the same four
	regions with the same degenerate branch, the same implicit four by four
	construction with the power boundary function in place of the exponential one,
	a dense $3\times3$ block, the residual $z+\tau\nabla B(s)$, no corrector, and
	bisected step and margin.

	\begin{remark}[The block shape may not change between iterations.]
		The pattern is planned once, and \texttt{widest\_operator\_kind} declares
		the widest representation a cone will ever return so that the plan
		accommodates every iteration. This is why the second order cone answers
		\texttt{DiagPlusLowRank} although three of its four regions produce a
		diagonal. A cone returning a diagonal block on one iteration and a low rank
		block on the next would corrupt the factorisation silently, the numeric
		refill writing into slots that no longer mean what they meant when the
		pattern was built.
	\end{remark}

	\begin{remark}[Every scaling is built at a point and consumed at the same
		point.] \texttt{build\_projection\_scaling} caches quantities derived from
		$\sigma$, and \texttt{materialise\_operator},
		\texttt{apply\_residual\_scaling} and \texttt{complementarity\_residual}
		all read that cache. Calling any of them after the iterate has moved
		returns a stale $F$, which is a descent direction for the wrong problem and
		shows up only as a slow solve.
	\end{remark}

	\begin{remark}[The metric is uniform on a curved block, and the kernel assumes
		it.] The kernel reads \texttt{scaling.penalty(0)} for a whole curved block
		rather than one entry per row, because admissibility gives $M=cI$ there.
		The schedule and the equilibration both guarantee it by flattening. A
		change that broke the guarantee would leave the kernel reading a single
		entry of a vector whose other entries it silently ignores.
	\end{remark}

	\begin{remark}[The orthant is the wrong cone to generalise from.]
		Its $F$ is a mask, so every quantity on it is elementwise, its block is
		diagonal on both roads, and $F^2=F$ makes several identities degenerate
		that are not degenerate elsewhere. Code that assumes a diagonal $G$, or
		indexes a block one row at a time, passes every test on the orthant and
		fails on the first second order block.
	\end{remark}

	\chapter{The bindings and the Python interface}

	\section{Where it happens}

	The interface is in two pieces. \texttt{bindings/bindings.cpp}, $265$ lines,
	exposes the C++ solver to Python through \textsc{nanobind} and does nothing
	else; \texttt{python/proxgqp/}, $430$ lines over five modules, is the layer a
	caller actually uses.

	\begin{table}[H]
		\centering\small
		\begin{tabular}{lll}
			\toprule
			module & lines & what it holds\\
			\midrule
			\texttt{cones.py} & $107$ & the six cone classes and \texttt{cone\_rows}\\
			\texttt{problem.py} & $52$ & \texttt{Problem}, and the validation\\
			\texttt{settings.py} & $114$ & the knob names, the presets, \texttt{make\_settings}\\
			\texttt{solver.py} & $131$ & \texttt{solve}, \texttt{solve\_qp},
			\texttt{Solution}\\
			\texttt{\_\_init\_\_.py} & $26$ & the public names\\
			\bottomrule
		\end{tabular}
	\end{table}

	The division of labour is strict, and it is the reason the binding is as short
	as it is: \emph{nothing} in C++ knows about Python types beyond arrays and
	scalars, and nothing in Python knows about the solver beyond one function
	call.

	\section{The binding is one function}

	The whole solver is exposed as a single \texttt{solve}, taking fifteen
	arguments, all of them plain arrays and scalars:
	\[
	\begin{array}{ll}
		\text{sizes} & \texttt{columns},\ \texttt{rows}\\
		P & \texttt{objective\_start},\ \texttt{objective\_row},\
		\texttt{objective\_value}\\
		q & \texttt{linear}\\
		E & \texttt{constraint\_start},\ \texttt{constraint\_row},\
		\texttt{constraint\_value}\\
		b & \texttt{offset}\\
		K & \texttt{cone\_kind},\ \texttt{cone\_size},\ \texttt{cone\_exponent}\\
		\text{the rest} & \texttt{settings},\ \texttt{warm}
	\end{array}
	\]

	The two matrices cross as three arrays each, in compressed sparse column form,
	and are rebuilt on the C++ side; the cone list crosses as three parallel
	arrays, one integer kind, one size and one exponent per block, and is rebuilt
	by \texttt{build\_cones}. No structure is shared, and no Python object is held
	past the call.

	The warm start is the one argument with any shape to it: an optional triple
	$(x,s,z)$, which the binding unpacks into a \texttt{Results} and passes by
	pointer, or a null pointer when absent. That is the whole mechanism behind the
	sequential sets of the benchmark.

	Note what the binding does \emph{not} do. It does not validate, it does not
	convert types, and it does not name its arguments after anything in the
	solver: the constraint offset is \texttt{offset} here and $b$ inside. Whatever
	arrives is mapped into Eigen views and handed to \texttt{solve}. Everything
	that could go wrong has already been prevented one level up.

	\section{The cones}

	Each cone is a small class carrying a kind, a size and an exponent, the last
	used only by the power cone.

	\begin{table}[H]
		\centering\small
		\begin{tabular}{llll}
			\toprule
			class & kind & \texttt{size} means & rows\\
			\midrule
			\texttt{Zero}$(d)$ & $0$ & the dimension & $d$\\
			\texttt{Nonneg}$(d)$ & $1$ & the dimension & $d$\\
			\texttt{SecondOrder}$(d)$ & $2$ & the dimension & $d$\\
			\texttt{PSDTriangle}$(k)$ & $3$ & the \emph{side} $k$ & $k(k+1)/2$\\
			\texttt{Exponential}$()$ & $4$ & fixed & $3$\\
			\texttt{Power}$(\alpha)$ & $5$ & fixed & $3$\\
			\bottomrule
		\end{tabular}
		\caption{The six cone classes. Only the semidefinite one has a size that is
			not its row count, and \texttt{cone\_rows} is what converts a list of them
			into the row count of $E$.}
	\end{table}

	The constraint is written as $Ex+s=b$ with $s\in K$, and $K$ is the product of
	the cones in the order given. There is no separate way to say which rows
	belong to which block: the list determines it, and \texttt{cone\_rows}
	computes the total that $E$ and $b$ must match.

	\section{\texttt{Problem}, which is where the checking is}

	\texttt{Problem(P, q, E, b, cones)} is a validating constructor and nothing
	more. It fixes the sizes from $q$ and from the cone list, converts both
	matrices to compressed sparse column form with sorted indices, checks that
	every shape agrees, fills in zeros for a missing $P$ or $b$, and flattens the
	cone list into the three arrays the binding wants.

	Two of those are worth naming.

	\texttt{sort\_indices} is called on both matrices. Every downstream assumption
	about the pattern, and \texttt{validate\_pattern} in the backend in
	particular, requires the row indices within a column to increase, and
	\textsc{scipy} does not guarantee it.

	And the shape checks are done here rather than in C++ because this is the only
	place where a helpful message can be produced. A caller who passes an $E$ with
	the wrong number of rows learns that, with both numbers, rather than reaching a
	pattern validation error inside a backend.

	\section{Settings, and the knob names}

	The C++ settings are a nested structure: one block per schedule,
	\texttt{bcl}, \texttt{gbcl}, \texttt{semismooth} and \texttt{interior}, the
	last two inheriting a shared part, together with the tolerances, the three
	iteration limits and the road and backend choice. Python flattens that into
	dotted names and a dictionary of presets.

	\begin{quote}
		\texttt{make\_settings(preset=None, **knobs)}
	\end{quote}

	Every knob is either a top level name such as \texttt{eps\_abs}, or a dotted
	one such as \texttt{interior.mehrotra} or \texttt{gbcl.mu\_adapt}, the prefix
	naming the schedule the knob belongs to. The enumerated ones are given as
	strings, resolved through the five tables \texttt{METHOD},
	\texttt{LINE\_SEARCH}, \texttt{PENALTY}, \texttt{ROAD} and \texttt{BACKEND}.
	So \texttt{method="interior"} rather than an enumerator, and an unknown string
	raises with the list of what is accepted.

	\texttt{PRESETS} names twelve combinations, of which the useful ones are
	\texttt{"interior"} and \texttt{"semismooth"} to force a road,
	\texttt{"accurate"} and \texttt{"loose"} to move both tolerances together, and
	\texttt{"no\_equilibration"}, \texttt{"no\_mehrotra"} and \texttt{"schur"} to
	turn off one mechanism at a time. The last three exist for isolating a
	question rather than for solving faster.

	\texttt{knob\_names()} lists everything settable and \texttt{describe} returns
	the current value of each, which together are how a caller finds out what
	exists without reading the C++ headers.

	\section{Solving}

	Two entry points, and the difference between them is the only concession the
	interface makes to convenience.

	\begin{quote}
		\texttt{solve(problem, warm\_start=None, settings=None, **knobs)}
	\end{quote}

	takes a \texttt{Problem} and returns a \texttt{Solution}. The warm start may
	be a previous \texttt{Solution}, a dictionary with keys \texttt{x},
	\texttt{s}, \texttt{z}, or any triple; all three become the same tuple before
	crossing.

	\begin{quote}
		\texttt{solve\_qp(P, q, A=None, b=None, G=None, h=None, ...)}
	\end{quote}

	takes the classical quadratic programming form, with $Ax=b$ and $Gx\leq h$,
	and assembles the cone problem itself: it stacks $A$ above $G$, stacks $b$
	above $h$, and builds the cone list as a \texttt{Zero} block for the
	equalities followed by a \texttt{Nonneg} block for the inequalities. It exists
	so that a caller with an ordinary quadratic program never has to learn what a
	cone list is.

	\texttt{Solution} carries the three vectors, the four residuals of the
	termination test together with the two relative residuals, the objective, the
	two iteration counts and the elapsed time, with \texttt{solved} and \texttt{residuals} as conveniences. The status
	is a string rather than an enumerator, mapped through \texttt{STATUS}.

	\begin{remark}[The slack that comes back is not the slack that was used.]
		\texttt{Solution.s} is the recomputed $b-Ex$ of the termination chapter,
		not the iterate's own $s$. Feeding a \texttt{Solution} straight back as a
		warm start therefore hands the solver a slack consistent with $x$ by
		construction, which is the right thing, but it is not the point the previous
		solve was actually sitting at.
	\end{remark}


	\chapter{Building it}

	\section{What has to be present}

	The solver is C++17 with two external dependencies and one vendored one.

	\begin{table}[H]
		\centering\small
		\begin{tabular}{lll}
			\toprule
			dependency & what for & how obtained\\
			\midrule
			Eigen & dense linear algebra, sparse storage & found, or fetched\\
			\textsc{suitesparse} & \textsc{cholmod}, the default backend & found,
			or fetched\\
			\textsc{nanobind} & the Python bindings & found\\
			\textsc{qdldl} & the dependency free backend &
			vendored in \texttt{third\_party/}\\
			a \textsc{blas} & the dense kernels and the supernodal path & found\\
			\bottomrule
		\end{tabular}
	\end{table}

	\textsc{qdldl} is the only one carried in the tree, and it is carried so that
	the solver can be built and its tests run with no sparse library installed at
	all. 
	\section{The environment}

	\texttt{environment.yml} pins everything, including the four other solvers the
	benchmark compares against, and is meant to be used with
	\texttt{micromamba} or \texttt{conda}:

	\begin{quote}
		\texttt{micromamba env create -f environment.yml}\\
		\texttt{micromamba activate proxgqp}
	\end{quote}

	Two things about it are deliberate.

	Everything comes from \texttt{conda-forge}, including
	\textsc{suitesparse}, which is what makes \textsc{cholmod} available without
	building it. That matters for more than convenience: \textsc{cholmod} is not
	redistributed with this work, it is a dependency the environment installs, so
	nothing here ships code under its licence.

	And the competitors are pinned in the same file. Clarabel, \textsc{piqp},
	\textsc{scs}, \textsc{ecos} and \textsc{osqp} come from
	\texttt{conda-forge}; \textsc{highs} and \textsc{proxqp} from \texttt{pip},
	those not being packaged there. A benchmark whose comparison depends on which
	versions were installed is not reproducible, and this is the file that fixes
	them.

	\section{Building}

	\begin{quote}
		\texttt{cmake -S solver -B build -G Ninja}\\
		\texttt{cmake -{}-build build -j}\\
		\texttt{ctest -{}-test-dir build}
	\end{quote}

	Three targets come out of it. \texttt{\_core} is the Python extension module,
	built by \textsc{nanobind} and installed beside the pure Python package;
	\texttt{proxgqp\_core} is a static library holding the solver without the
	bindings, which is what the tests link against; and each test is an executable
	registered with \textsc{ctest}.

	The split between the two libraries is what lets the tests exercise the solver
	without Python in the loop. A test failure is then a C++ failure, and nothing
	about the bindings can produce or mask one.

	\paragraph{Threads while building.}
	The build fetches nothing by default and compiles in a few minutes, but
	\texttt{-j} with no argument will use every core. On a workstation being used
	for anything else, ask for fewer.

	\section{Building without the dependencies present}

	\begin{quote}
		\texttt{cmake -S solver -B build -DPROXGQP\_VENDOR\_DEPENDENCIES=ON}
	\end{quote}

	fetches Eigen and \textsc{suitesparse} at configure time and builds them in
	tree, pinned to the tags in the \textsc{cmake} file. It exists for the case
	where the environment cannot be used, and it is slower by a large factor,
	\textsc{suitesparse} being a substantial build of its own.

	Two settings are forced in that path and are worth knowing about.
	\textsc{fortran} is disabled, \textsc{suitesparse} otherwise requiring a
	compiler for it that a minimal environment may not have; and \textsc{cholmod}
	is linked statically, so that the resulting extension module does not depend
	on a shared library that will not be present.

	The supported path is the environment. The vendoring path is a fallback, and
	it is the one where the licence question is sharper: fetching and statically
	linking \textsc{cholmod} means shipping it, where installing it from
	\texttt{conda-forge} does not.

	\section{Installing}

	\texttt{cmake -{}-install} places the extension module and the Python package
	together, so that \texttt{import proxgqp} works from the install prefix. The
	package is pure Python except for \texttt{\_core}, and the two must stay
	beside each other; nothing looks for the extension anywhere else.

	\begin{remark}[The build has no configuration for the solver itself.]
		Nothing about the method is a compile time choice. The road, the backend,
		the smoothing, the line search and the schedule are all settings resolved at
		run time, and the same binary runs every configuration in this document. The
		only compile time option is where the dependencies come from.
	\end{remark}

	\part{Benchmark}
	
	
	\chapter{The harness}

	\section{What it is}

	The benchmark is a directory of its own, $3739$ lines over $63$ files, and it
	is a separate program from the solver: it runs every solver as a child
	process, grades all of them by one rule it implements itself, and never links
	against any of them.

	\begin{table}[H]
		\centering\small
		\begin{tabular}{lll}
			\toprule
			piece & what it holds & how many\\
			\midrule
			\texttt{utils/} & the shared machinery, the verifier, the aggregates,
			the run store & $13$ files\\
			\texttt{solvers/} & one adapter per solver & $11$ files\\
			one directory per set & how to build or read that set's problems
			& $15$ directories\\
			\texttt{prepare.py} & fetching, generating and checksumming the data
			& $232$ lines\\
			\texttt{benchmark\_report.py} & the driver and the report
			& $219$ lines\\
			\texttt{run\_set.py} & one set, from the command line & $36$ lines\\
			\bottomrule
		\end{tabular}
	\end{table}

	The configuration lives in one place, at the top of
	\texttt{benchmark\_report.py}: the five solvers, the fifteen sets, which of
	them are sequential, the tolerance $10^{-9}$, the time limit of twenty
	seconds, the worker count, which is the machine's core count less four.
	
	\section{Getting the data}

	Three of the fifteen sets are collected archives rather than generated
	problems, and one more reads market data. \texttt{prepare.py} handles all of
	it and is the first thing to run.

	\begin{table}[H]
		\centering\small
		\begin{tabular}{ll}
			\toprule
			command & what it does\\
			\midrule
			\texttt{python prepare.py generate} & builds the six planted sets\\
			\texttt{python prepare.py fetch} & downloads Maros, Netlib and
			\textsc{cblib}\\
			\texttt{python prepare.py manifest} & writes the checksums of what is
			present\\
			\texttt{python prepare.py verify} & checks every file against the
			manifest\\
			\bottomrule
		\end{tabular}
	\end{table}

	The manifest holds a \textsc{sha256} for every problem file and for every
	market data file, so a set that has been silently altered is detected rather
	than benchmarked. The generated sets are checksummed too, on the problem
	itself rather than the file, so that a change to the generator shows up as a
	mismatch even though the file is written afresh each time.

	\section{Running it}

	\begin{table}[H]
		\centering\small
		\begin{tabular}{ll}
			\toprule
			command & what it does\\
			\midrule
			\texttt{python benchmark\_report.py} & every solver on every set, then
			the report\\
			\texttt{python run\_set.py \textit{name}} & one set, same defaults\\
			\texttt{python run\_set.py \textit{name} -{}-solvers a b}
			& one set, chosen solvers\\
			\texttt{python run\_set.py \textit{name} -{}-out-of-the-box}
			& one set, every solver on its defaults\\
			\texttt{python run\_set.py \textit{name} -{}-run \textit{id}}
			& one set, appended to an existing run\\
			\texttt{python tune.py} & a small gate, for checking a change quickly\\
			\bottomrule
		\end{tabular}
	\end{table}

	Every solve is a separate process. \texttt{run\_one.py} takes a set, a
	problem, a solver and a tolerance, runs exactly one solve, and prints one
	\textsc{json} record. That is the unit of the whole harness, and it has three
	consequences worth naming.

	A solver that crashes takes down one process and produces one failed record,
	rather than ending the run. A solver that hangs is killed by the parent's
	timeout, which is how the twenty second limit is enforced regardless of
	whether the solver honours a limit of its own. And no solver can affect
	another through global state, a threading library or a memory allocator,
	because none of them is ever loaded into the same process.

	\subsection{One core per solve}

	Each child is launched under \texttt{taskset} on a core drawn from a pool, and
	its environment is set so that every threading library it might link sees a
	single thread:
	\[
	\texttt{OMP\_NUM\_THREADS},\quad
	\texttt{OPENBLAS\_NUM\_THREADS},\quad
	\texttt{MKL\_NUM\_THREADS},\quad
	\texttt{VECLIB\_MAXIMUM\_THREADS} .
	\]
	Both halves are needed and neither is sufficient. Without the environment a
	solver linked against a threaded \textsc{blas} would use every core it could
	find, and the timings of whatever ran beside it would be meaningless. Without
	the pinning the operating system would migrate the children between cores and
	the same solve would time differently on successive runs.

	\subsection{Sequential sets are different}

	Two sets measure warm starting, and they cannot be run problem by problem.
	Their problems form chains, each chain a sequence of neighbouring problems,
	and each solve is warm started from the previous one's answer.

	The driver therefore has a second path. It groups the problems into chains by
	a key the set supplies, and gives each pair of a chain and a solver to one
	worker, which runs that chain in order, carrying the answer forward in a
	scratch file. Parallelism is across chains and solvers rather than across
	problems, which is what preserves the order within a chain while still using
	the machine.

	\section{Where a run goes}

	A run owns a directory, and writes into nothing else:
	\begin{quote}
		\texttt{result/runs/<timestamp>-<uid>/}
	\end{quote}
	holding one \textsc{json} file per solver and set, named
	\texttt{<solver>\_<set>.json}, and one \texttt{meta.json} describing the run
	itself: the identifier and the uid, when it started and finished in
	\textsc{utc}, how long it took, which sets and solvers it covered, the
	settings it ran under, the versions of the rival solvers, and whether it
	finished at all. \texttt{result/runs/latest} points at the newest.


	\texttt{meta.json} records the processor, the core count, the Python version
	and the thread environment variables, and deliberately nothing else about the
	machine: no host name, no paths, no repository or branch.

	\section{Grading}

	No solver's own verdict is used. Every returned point is graded by
	\texttt{utils/verify.py}, which recomputes the four residuals from the
	original problem data:
	\[
	\begin{aligned}
		\text{primal} &=\norm{\proj{-\widehat s}{K^*}}_\infty,
		&\quad
		\text{dual} &=\norm{Px+q+E^{\top}z}_\infty,\\
		\text{cone} &=\norm{z-\proj{z}{K^*}}_\infty,
		&\quad
		\text{gap} &=\bigl|\inner{\widehat s}{z}\bigr| ,
	\end{aligned}
	\qquad \widehat s=b-Ex ,
	\]
	each against $\varepsilon_{\mathrm{abs}}+\varepsilon_{\mathrm{rel}}$ times its
	own scale, with $\varepsilon_{\mathrm{abs}}=\varepsilon_{\mathrm{rel}}=
	10^{-9}$. A solver reporting success whose point fails this is recorded as
	\texttt{solved} but not \texttt{verified}, and only \texttt{verified} counts.

	The regrading is a check on the answer, not a substitute for each solver's own
	stopping rule. Every solver is configured to stop on \emph{its own} criterion
	at the same accuracy: each adapter sets that solver's absolute tolerance to
	$10^{-9}$, its relative tolerance to $10^{-9}$, and, where the solver
	distinguishes one, its duality gap tolerance to the same. Each therefore
	decides for itself when it has finished, by its own definition of finished, at
	a target no looser than ours.

	Each is also given an unbounded iteration budget, ours included. A solver
	stopped by an iteration limit would be reported as failing on problems it
	could have solved, and every solver in the comparison does better without
	one, so none is given one; the twenty second limit is the only budget any of
	them faces.

	What the regrading adds is that a claim of success is checked rather than
	taken. A solver that reports \texttt{solved} at a point whose residuals do not
	meet the bar is recorded as unverified, and no solver's own verdict is allowed
	to stand unexamined, ours least of all.

	\subsection{What a solver is not blamed for}

	A set whose cones a solver cannot express is not counted against it. The
	adapter returns \texttt{unsupported}, those rows are dropped from that
	solver's totals rather than recorded as failures, and its score is reported
	out of what it could express. 

	\section{Statistics}

	\begin{definition}[Shifted geometric mean]\label{def:sgm}
		For values $v_1,\dots,v_N\geq0$ and a shift $\varsigma>0$,
		\begin{equation}
			\mathrm{sgm}_{\varsigma}(v)
			=\exp\Bigl(\frac1N\sum_{i=1}^{N}\log(v_i+\varsigma)\Bigr)-\varsigma .
		\end{equation}
	\end{definition}

	The geometric mean is used because a solver is judged on ratios: being twice
	as fast on a fast problem should count as much as being twice as fast on a
	slow one, which an arithmetic mean does not do. The shift keeps a single very
	small value from dominating, and the harness takes $\varsigma$ equal to one
	millisecond.

	A failed solve is scored at the cap rather than dropped, because dropping it
	would reward a solver for giving up, and excluding it would compare solvers on
	different problem sets. The cap is the twenty second limit, which is
	meaningful: a timed out solve genuinely consumed twenty seconds.



	\chapter{The benchmark sets}

	Fifteen sets, $1509$ instances, of three kinds: six generated with a planted
	solution, three collected archives, and six built from portfolio
	optimisation. This chapter says what each contains, in what form, and why
	the instance counts are what they are.

	\begin{table}[H]
		\centering\small
		\begin{tabular}{llrl}
			\toprule
			set & kind & instances & cones\\
			\midrule
			\texttt{random\_qp} & generated & $90$ & zero, orthant\\
			\texttt{random\_socp} & generated & $90$ & zero, orthant, second order\\
			\texttt{random\_psd} & generated & $60$ & zero, orthant, semidefinite\\
			\texttt{random\_exp} & generated & $60$ & zero, orthant, exponential\\
			\texttt{random\_pow} & generated & $60$ & zero, orthant, power\\
			\texttt{random\_mixed} & generated & $60$ & all six\\
			\midrule
			\texttt{maros} & collected & $138$ & zero, orthant\\
			\texttt{netlib} & collected & $109$ & zero, orthant\\
			\texttt{cblib} & collected & $134$ & exponential, power, second order\\
			\midrule
			\texttt{skfolio\_real\_qp} & portfolio & $96$ & zero, orthant\\
			\texttt{skfolio\_real\_curved} & portfolio & $192$ & second order,
			exponential\\
			\texttt{skfolio\_random\_qp} & portfolio & $60$ & zero, orthant\\
			\texttt{skfolio\_random\_curved} & portfolio & $120$ & second order,
			exponential\\
			\texttt{seq\_rebalance} & portfolio, warm & $120$ & zero, orthant\\
			\texttt{seq\_rebalance\_curved} & portfolio, warm & $120$ & second
			order, exponential\\
			\bottomrule
		\end{tabular}
	\end{table}

	\section{The generated sets}

	Six sets sweep one cone at a time over sizes from $n=30$ to $n=3000$. They
	exist because the collected archives contain almost no curved cones: without
	them the second order, semidefinite, exponential and power paths would be
	exercised only by \textsc{cblib} and by the portfolio sets, and never in
	isolation.

	Every instance is built \emph{backwards} from a solution, which is what makes
	the answer known in advance rather than agreed on by majority vote among
	solvers.

	\begin{algorithm}[H]
		\caption{Generating one instance with a planted solution}
		\begin{algorithmic}[1]
			\Require the variable count $n$, a cone specification, a seed
			\Ensure $(P,q,E,b,K)$ and the pair $(x^{\star},z^{\star})$ solving it
			\For{each cone block in the specification}
			\State draw its regime: \emph{active}, \emph{inactive} or
			\emph{degenerate}, with probabilities
			$\pi_{\mathrm{act}},\pi_{\mathrm{inact}},\pi_{\mathrm{deg}}$
			\State draw a complementary pair $(s_j,z_j)$ in that regime
			\Comment{$s_j\in K_j$, $z_j\in K_j^*$, $\inner{s_j}{z_j}=0$}
			\EndFor
			\State $s\gets$ the concatenation, $z^{\star}\gets$ the concatenation,
			$m\gets\dim s$
			\State $E\gets$ a sparse $m\times n$ matrix, ten nonzeros per row,
			normal entries
			\State $x^{\star}\gets$ a normal vector of length $n$
			\State $b\gets Ex^{\star}+s$
			\Comment{so $(x^{\star},s)$ is feasible by construction}
			\State $P\gets Q\Lambda Q^{\top}$ with $Q$ a random orthogonal matrix
			and $\Lambda$ logarithmically spaced over the requested condition
			number and rank
			\State $q\gets-Px^{\star}-E^{\top}z^{\star}$
			\Comment{so stationarity holds at $x^{\star}$}
			\State \Return the problem, together with $x^{\star}$, $z^{\star}$ and
			the objective value at $x^{\star}$
		\end{algorithmic}
	\end{algorithm}

	The last two lines are the whole construction. Feasibility is arranged by
	defining $b$ from a chosen $x^{\star}$ and a chosen cone-feasible slack;
	stationarity is arranged by defining $q$ so that $Px^{\star}+q+E^{\top}
	z^{\star}=0$; and complementarity holds because each block's pair was drawn
	complementary. The three conditions of the theory part are therefore
	satisfied by construction, and the instance carries its own optimal objective
	value.

	\paragraph{The three regimes.}
	How a complementary pair is drawn is what the regime selects, and the three
	are chosen to exercise different behaviour.

	\begin{table}[H]
		\centering\small
		\begin{tabular}{lll}
			\toprule
			regime & the pair & what it tests\\
			\midrule
			active & $s_j$ on the boundary, $z_j\neq0$ normal to it
			& the curved parts of $F$\\
			inactive & $s_j$ strictly inside, $z_j=0$
			& the identity or zero branches of $F$\\
			degenerate & $s_j=z_j=0$
			& the vertex, where $F$ is set valued\\
			\bottomrule
		\end{tabular}
	\end{table}

	On the second order cone, for instance, an active block takes
	$s_j=(\norm{u},u)$ on the boundary and $z_j$ a positive multiple of
	$(\norm{u},-u)$, which is the outward normal there; on the exponential cone
	it takes $s_j$ on the smooth boundary and $z_j$ a negative multiple of the
	gradient $\bigl(e^{r},e^{r}(1-r),-1\bigr)$ of Proposition~\ref{prop:expbound}.
	The three probabilities are $0.6$, $0.3$ and $0.1$. The degenerate regime, at
	one block in ten, is the one that puts a block
	exactly at the point where the Clarke Jacobian is a set rather than a matrix.

	\paragraph{The sweep.}
	Each set is the product of five sizes, three condition numbers, two objective
	ranks and three seeds, giving $90$, or a subset of that for the more expensive
	cones. The condition number enters through $\Lambda$, whose entries are
	logarithmically spaced from $1$ down to $10^{-c}$, so a problem with $c=6$
	has an objective spanning six orders of magnitude. The rank parameter chooses
	between a full rank $P$ and $P=0$, the second being a linear program with
	cones.

	\section{The collected archives}

	Three archives of problems from the literature, none generated here : they are graded, like everything else, by the harness
	recomputing the residuals.

	\subsection{Maros and M\'esz\'aros, $138$ instances}

	The standard repository of convex quadratic programs,
	\[
	\min_x\ \tfrac12\inner{x}{Px}+\inner{q}{x}
	\quad\text{subject to}\quad
	Ax=b,\ \ l\leq Cx\leq u,\ \ \ell\leq x\leq \upsilon ,
	\]
	converted to the split form by turning each two sided constraint into a pair
	of one sided rows and each equality into a zero cone block. It is the
	reference set for quadratic programming and the reason \textsc{proxqp} and
	\textsc{piqp} appear in the comparison at all.

	\subsection{Netlib, $109$ instances}

	The classical linear programming archive, so $P=0$ throughout, in the form
	\[
	\min_x\ \inner{q}{x}
	\quad\text{subject to}\quad
	Ax=b,\ \ x\geq0 .
	\]
	It is included because a linear program is the hardest case for a proximal
	method: with $P=0$ the only strong convexity in the subproblem is the
	proximal term itself, so the schedule has nothing else holding the system
	together as $\rho$ falls. The semismooth road's weakest result in the whole
	benchmark is on this set.

	\subsection{\textsc{cblib}, $134$ instances of $934$}

	The Conic Benchmark Library, and the only archive here with curved cones. The
	instance count needs explaining, since the library holds $934$ files and only
	$134$ are used.

	The library is assembled for conic solvers generally, and most of its
	contents are outside what this solver addresses:

	\begin{itemize}
		\item the great majority are \emph{mixed integer} conic problems, carrying
		integrality constraints on some variables. Nothing in this work addresses
		integrality, and a continuous solver run on such an instance answers a
		different question, namely the continuous relaxation;
		\item a further group are known to be primal or dual \emph{infeasible},
		and are in the library precisely to test infeasibility detection rather
		than solution;
		\item a smaller group are \emph{weakly} infeasible or unbounded, where no
		certificate of finite size exists, and where every solver in the comparison
		would be reporting a judgement rather than a result.
	\end{itemize}

	What remains, after taking the continuous, feasible and bounded members, is
	$134$ instances. That is the set used, and the filter is applied once when
	the archive is prepared rather than at run time, so every solver receives the
	same $134$.

	The cone mix is what makes it valuable: $42$ instances carry exponential
	cones, $12$ carry power cones, and the rest are second order or mixed. It is
	also by a wide margin the hardest set in the benchmark.

	\section{The portfolio sets}

	Six sets built from \texttt{skfolio}, in six formulations, on real and on
	synthetic return data. They are here because they are what the solver was
	asked to do: the exponential cone entered this work because benchmarking
	against \texttt{skfolio} showed it was needed.

	\begin{table}[H]
		\centering\small
		\begin{tabular}{lll}
			\toprule
			formulation & cones & in\\
			\midrule
			mean variance & zero, orthant & the qp sets\\
			mean \textsc{cvar} & zero, orthant & the qp sets\\
			minimum deviation & second order & the curved sets\\
			mean \textsc{evar} & exponential & the curved sets\\
			risk budgeting & exponential & the curved sets\\
			maximum diversification & second order & the curved sets\\
			\bottomrule
		\end{tabular}
		\caption{The six formulations. Two are quadratic programs and four are not,
			which is why the portfolio sets are split.}
	\end{table}

	Throughout, $R\in\R^{T\times d}$ is the matrix of returns, $T$ observations of
	$d$ assets, $\Sigma$ its covariance, $\mu$ its column means, and $w$ the
	weights. Every formulation carries the budget and long only constraints
	\[
	\mathbf{1}^{\top}w=1,
	\qquad
	w\geq0 ,
	\]
	a zero cone row and an orthant block, and differs only in what it adds.

	\paragraph{Mean variance, a quadratic program.}
	\[
	\min_{w}\ \tfrac12\inner{w}{\Sigma w}-\gamma\inner{\mu}{w} .
	\]
	Nothing beyond the budget and the sign constraints, so $P=\Sigma$ and the
	problem is an ordinary quadratic program.

	\paragraph{Mean \textsc{cvar}, a linear program.}
	With $\beta$ the confidence level, and introducing a value at risk level
	$\tau$ and one excess variable $u_t$ per observation,
	\[
	\min_{w,\tau,u}\ -\gamma\inner{\mu}{w}+\tau
	+\frac{1}{T(1-\beta)}\sum_{t=1}^{T}u_t
	\quad\text{subject to}\quad
	u_t\geq-\inner{R_t}{w}-\tau,
	\quad u\geq0 .
	\]
	The conditional value at risk is written by its Rockafellar and Uryasev
	linearisation, so $P=0$ and the two families of constraints are orthant
	blocks of $T$ rows each. It is a linear program with $d+1+T$ variables, and
	it is the larger of the two quadratic sets by a wide margin.

	\paragraph{Minimum deviation, a second order cone.}
	Writing $\Sigma=FF^{\top}$ for a factor $F$,
	\[
	\min_{w,t}\ t
	\quad\text{subject to}\quad
	\bigl(t,\ F^{\top}w\bigr)\in\mathcal{Q}^{d+1} ,
	\]
	that is $t\geq\norm{F^{\top}w}=\sqrt{\inner{w}{\Sigma w}}$. Minimising the
	standard deviation rather than the variance, which is the same optimiser but
	a conic rather than a quadratic problem, and the reason the second order cone
	appears in a portfolio set at all.

	\paragraph{Maximum diversification, a second order cone.}
	With $\sigma$ the vector of individual standard deviations,
	\[
	\min_{w,t}\ t
	\quad\text{subject to}\quad
	\inner{\sigma}{w}=1,
	\quad
	\bigl(t,\ F^{\top}w\bigr)\in\mathcal{Q}^{d+1},
	\quad w\geq0 .
	\]
	The budget row is replaced by the normalisation $\inner{\sigma}{w}=1$, and
	minimising the portfolio deviation subject to it maximises the
	diversification ratio $\inner{\sigma}{w}/\sqrt{\inner{w}{\Sigma w}}$.

	\paragraph{Mean \textsc{evar}, exponential cones.}
	The entropic value at risk is a log sum exponential, and it is written with
	one exponential cone per observation. With a scale $z$, a width $v$ and
	$u_t$ one variable per observation,
	\[
	\min_{w,z,v,u}\ -\gamma\inner{\mu}{w}+z
	+v\log\frac{1}{T(1-\beta)}
	\quad\text{subject to}\quad
	\sum_{t=1}^{T}u_t\leq v,
	\quad
	\bigl(\inner{R_t}{w}+z,\ v,\ u_t\bigr)\in\Kexp .
	\]
	Each cone membership is $v\,e^{(\inner{R_t}{w}+z)/v}\leq u_t$, which is the
	epigraph form of one term of the exponential moment, and summing the $u_t$
	under $v$ is the constraint that makes the whole a log sum exponential bound.
	This formulation, and the next, are why the exponential cone was needed.

	\paragraph{Risk budgeting, exponential cones.}
	With a prescribed budget vector $c$ summing to one, and $y$ the logarithms of
	the weights,
	\[
	\min_{w,y}\ \tfrac12\inner{w}{\Sigma w}
	\quad\text{subject to}\quad
	\inner{c}{y}\geq0,
	\quad w\geq0,
	\quad
	\bigl(y_i,\ 1,\ w_i\bigr)\in\Kexp .
	\]
	The cone membership $1\cdot e^{y_i}\leq w_i$ is $y_i\leq\log w_i$, so the
	constraint $\inner{c}{y}\geq0$ is the log barrier
	$\sum_ic_i\log w_i\geq0$ that produces a risk budgeted portfolio. There are
	$d$ exponential blocks, one per asset, which makes this the formulation with
	the most curved blocks of the six.

	\paragraph{Why the split.}
	Only two of the six formulations are quadratic programs. Before the split, a
	single \texttt{skfolio} set mixed all six, and \textsc{piqp} and
	\textsc{proxqp} could express only a third of it while being scored on the
	whole; the two halves are now separate sets so that every solver is compared
	on problems it can state.

	\paragraph{Real and synthetic.}
	The \emph{real} sets use market data, swept over datasets, estimation windows
	and eras; the \emph{random} sets generate returns instead, which removes the
	dependence on the data files and lets the same formulations be run at sizes
	the market data does not offer.

	\section{The sequential sets}

	Two sets, and the only ones that measure warm starting. Each is a family of
	chains, and within a chain each problem is a small perturbation of the last:

	\begin{table}[H]
		\centering\small
		\begin{tabular}{ll}
			\toprule
			family & what moves along the chain\\
			\midrule
			rolling & the estimation window slides forward one period at a time\\
			frontier & the risk aversion $\gamma$ is swept along the efficient
			frontier\\
			beta & a market exposure constraint is tightened by steps\\
			\bottomrule
		\end{tabular}
	\end{table}

	Three families, forty steps each, and one problem per step, giving $120$
	instances per set. \texttt{seq\_rebalance} uses the two quadratic
	formulations and \texttt{seq\_rebalance\_curved} the four curved ones.
.

	\chapter{Results}

	\section{The conditions}

	Every number below comes from one run of \texttt{benchmark\_report.py} on a
	single machine: $1509$ instances over fifteen sets, five solvers, tolerance
	$10^{-9}$ absolute and relative, a twenty second limit, one pinned core per
	solve, and every point regraded by the harness. Rows a solver cannot express
	are excluded from its totals rather than counted against it.
 It was taken on
	2026-09-05, started 12:46 UTC and finished 13:57 UTC, seventy minutes for the
	$1509$ instances, on an \texttt{x86\_64} machine with twelve cores, eight of
	them carrying the harness, one thread per solve, under Python $3.14$, against
	Clarabel $0.11.1$, \textsc{piqp} $0.6.3$ and \textsc{proxqp} $0.7.3$, the
	last as shipped in \texttt{proxsuite} $0.7.3$. Every record it produced, one
	line per instance and solver, is kept beside the code in

	\begin{center}
		\texttt{benchmark/result/runs/20260905T124650Z-db22480c/}
	\end{center}

	with a \texttt{meta.json} recording the machine, the settings and the
	duration. The identifier is what the harness assigns each run, so a later run
	lands beside this one rather than over it, and the tables here can always be
	checked against the records they came from.

	These numbers move as the solver does. They describe the build of that day and
	are meant to be replaced whenever the solver changes enough to shift them; the
	comparison is worth repeating rather than quoting, and the harness exists so
	that repeating it costs one command.



	\section{The quadratic sets}

	Six of the fifteen sets are quadratic programs, and on those all five solvers
	can be compared on identical data: \texttt{random\_qp}, \texttt{maros},
	\texttt{netlib}, \texttt{skfolio\_real\_qp}, \texttt{skfolio\_random\_qp} and
	\texttt{seq\_rebalance}, $613$ instances in all.

	\begin{table}[H]
		\centering\small
		\begin{tabular}{lrrrrr}
			\toprule
			set & gqp interior & gqp semismooth & Clarabel & PIQP & ProxQP\\
			\midrule
			\texttt{random\_qp} (90)             & \textbf{72} & \textbf{72} & \textbf{72} & 67 & 68\\
			\texttt{maros} (138)                 & \textbf{131} & 94 & 107 & 121 & 80\\
			\texttt{netlib} (109)                & \textbf{87} & 38 & 66 & 85 & 30\\
			\texttt{skfolio\_real\_qp} (96)      & \textbf{96} & 95 & 95 & 94 & 94\\
			\texttt{skfolio\_random\_qp} (60)    & \textbf{60} & 53 & 49 & 55 & 55\\
			\texttt{seq\_rebalance} (120)        & \textbf{120} & \textbf{120} & 118 & 115 & 115\\
			\midrule
			all (613)                            & \textbf{566} & 472 & 507 & 537 & 442\\
			\bottomrule
		\end{tabular}
		\caption{Quadratic sets: instances verified. Every solver can express every
			problem here, so all five are scored out of the same $613$.}
	\end{table}

	\begin{table}[H]
		\centering\small
		\begin{tabular}{lrrrrr}
			\toprule
			set & gqp interior & gqp semismooth & Clarabel & PIQP & ProxQP\\
			\midrule
			\texttt{random\_qp} (90)             & \textbf{0.146} & 0.245 & 0.176 & 0.202 & 0.341\\
			\texttt{maros} (138)                 & \textbf{0.043} & 0.266 & 0.141 & 0.056 & 0.928\\
			\texttt{netlib} (109)                & 0.225 & 2.765 & 0.629 & \textbf{0.157} & 5.590\\
			\texttt{skfolio\_real\_qp} (96)      & \textbf{0.013} & 0.019 & 0.019 & 0.017 & 0.084\\
			\texttt{skfolio\_random\_qp} (60)    & \textbf{0.041} & 0.097 & 0.112 & 0.081 & 0.139\\
			\texttt{seq\_rebalance} (120)        & \textbf{0.003} & 0.004 & 0.005 & 0.007 & 0.016\\
			\midrule
			all (613)                            & \textbf{0.035} & 0.109 & 0.072 & 0.048 & 0.288\\
			\bottomrule
		\end{tabular}
		\caption{Quadratic sets: shifted geometric mean of seconds, shift one
			millisecond, failures scored at the twenty second limit.}
	\end{table}

	

	\section{All fifteen sets}

	Only three of the five solvers carry all six cones, so the comparison over the
	whole benchmark is between those three.

	\begin{table}[H]
		\centering\small
		\begin{tabular}{lrrr}
			\toprule
			set & gqp interior & gqp semismooth & Clarabel\\
			\midrule
			\texttt{random\_qp} (90)             & \textbf{72} & \textbf{72} & \textbf{72}\\
			\texttt{random\_socp} (90)           & \textbf{71} & 69 & 69\\
			\texttt{random\_psd} (60)            & 54 & \textbf{57} & 45\\
			\texttt{random\_exp} (60)            & \textbf{45} & \textbf{45} & 44\\
			\texttt{random\_pow} (60)            & 56 & \textbf{59} & 51\\
			\texttt{random\_mixed} (60)          & 35 & \textbf{47} & 27\\
			\texttt{maros} (138)                 & \textbf{131} & 94 & 107\\
			\texttt{netlib} (109)                & \textbf{87} & 38 & 66\\
			\texttt{cblib} (134)                 & \textbf{66} & 42 & 19\\
			\texttt{skfolio\_real\_qp} (96)      & \textbf{96} & 95 & 95\\
			\texttt{skfolio\_real\_curved} (192) & \textbf{192} & 155 & 110\\
			\texttt{skfolio\_random\_qp} (60)    & \textbf{60} & 53 & 49\\
			\texttt{skfolio\_random\_curved} (120) & \textbf{105} & 72 & 82\\
			\texttt{seq\_rebalance} (120)        & \textbf{120} & \textbf{120} & 118\\
			\texttt{seq\_rebalance\_curved} (120) & \textbf{120} & 59 & 36\\
			\midrule
			all (1509)                           & \textbf{1310} & 1077 & 990\\
			\bottomrule
		\end{tabular}
		\caption{All fifteen sets: instances verified, for the three solvers
			carrying all six cones.}
	\end{table}

	\begin{table}[H]
		\centering\small
		\begin{tabular}{lrrr}
			\toprule
			set & gqp interior & gqp semismooth & Clarabel\\
			\midrule
			\texttt{random\_qp} (90)             & \textbf{0.146} & 0.245 & 0.176\\
			\texttt{random\_socp} (90)           & 0.291 & 0.288 & \textbf{0.262}\\
			\texttt{random\_psd} (60)            & 0.071 & \textbf{0.041} & 0.250\\
			\texttt{random\_exp} (60)            & 0.536 & 0.535 & \textbf{0.279}\\
			\texttt{random\_pow} (60)            & 0.211 & \textbf{0.146} & 0.154\\
			\texttt{random\_mixed} (60)          & 0.478 & \textbf{0.393} & 1.521\\
			\texttt{maros} (138)                 & \textbf{0.043} & 0.266 & 0.141\\
			\texttt{netlib} (109)                & \textbf{0.225} & 2.765 & 0.629\\
			\texttt{cblib} (134)                 & \textbf{2.435} & 5.542 & 11.248\\
			\texttt{skfolio\_real\_qp} (96)      & \textbf{0.013} & 0.019 & 0.019\\
			\texttt{skfolio\_real\_curved} (192) & \textbf{0.037} & 0.124 & 0.259\\
			\texttt{skfolio\_random\_qp} (60)    & \textbf{0.041} & 0.097 & 0.112\\
			\texttt{skfolio\_random\_curved} (120) & \textbf{0.139} & 1.174 & 0.272\\
			\texttt{seq\_rebalance} (120)        & \textbf{0.003} & 0.004 & 0.005\\
			\texttt{seq\_rebalance\_curved} (120) & \textbf{0.131} & 0.735 & 1.430\\
			\midrule
			all (1509)                           & \textbf{0.100} & 0.256 & 0.262\\
			\bottomrule
		\end{tabular}
		\caption{All fifteen sets: shifted geometric mean of seconds. Machine
			dependent, unlike the counts above.}
	\end{table}




	\chapter{Closing remark}
	
	The solver began as a study of the ProxQP paper. The first generalisation was
	implemented for the second order cone alone, on the semismooth road. A great
	deal of time then went into trying to keep the per iteration cost of the
	qp case, by changing the form of the problem; most of those attempts
	failed, and the ones that succeeded did so only in particular cases, such as a
	single second order block together with an orthant, where a pencil argument or
	a geometric condition on $C$ applies, and even there at the price of a large
	constant. Other libraries, hyhound among them, were tried in the same context.
	
	Benchmarking against skfolio showed that the exponential cone was necessary, so the codebase was restructured to support all six cones, with Claude Code serving as the coding agent. Discovering PIQP demonstrated that an interior-point version of the algorithm was feasible. This led to the development of a smoothing-point perspective, whose main advantage was that it required no changes to the underlying linear algebra. An interior-point version was subsequently planned and added.
	
	Four possible directions remain. The cost of the linear algebra is the first and most significant. Scheduling is the second. It determines the iteration count on the semismooth path, and GBCL, obtained by writing BCL for each constraint after the generalisation, already outperforms BCL in this regard. The interior-point path is no different, and its scheduling deserves the same treatment.
	
	Initialisation is another possible direction. Mixed methods are the fourth and least explored: one could apply an interior-point treatment to some blocks and a semismooth treatment to others, or switch between the two during the solve i.e using the interior-point path far from the solution, where it requires fewer iterations, and the proximal path near the solution, where warm-starting is advantageous.
	
	\begin{thebibliography}{99}
		
		\bibitem{Bhatia1997}
		R.~Bhatia. \emph{Matrix Analysis}. Graduate Texts in Mathematics 169,
		Springer, 1997.
		
		\bibitem{BolteDaniilidisLewis2009}
		J.~Bolte, A.~Daniilidis and A.~S. Lewis. Tame functions are semismooth.
		\emph{Mathematical Programming}, 117:5--19, 2009.
		
		\bibitem{BoydVandenberghe2004}
		S.~Boyd and L.~Vandenberghe. \emph{Convex Optimization}. Cambridge University
		Press, 2004.
		
		\bibitem{BunchKaufman1977}
		J.~R. Bunch and L.~Kaufman. Some stable methods for calculating inertia and
		solving symmetric linear systems. \emph{Mathematics of Computation},
		31:163--179, 1977.
		
		\bibitem{Chares2009}
		R.~Chares. \emph{Cones and Interior-Point Algorithms for Structured Convex
			Optimization involving Powers and Exponentials}. PhD thesis, Universit\'e
		catholique de Louvain, 2009.
		
		\bibitem{DavisHager1999}
		T.~A. Davis and W.~W. Hager. Modifying a sparse Cholesky factorization.
		\emph{SIAM Journal on Matrix Analysis and Applications}, 20:606--627, 1999.
		
		\bibitem{EvansGariepy2015}
		L.~C. Evans and R.~F. Gariepy. \emph{Measure Theory and Fine Properties of
			Functions}. Revised edition, CRC Press, 2015.
		
		\bibitem{FarautKoranyi1994}
		J.~Faraut and A.~Kor\'anyi. \emph{Analysis on Symmetric Cones}. Oxford
		University Press, 1994.
		
		\bibitem{GillGolubMurraySaunders1974}
		P.~E. Gill, G.~H. Golub, W.~Murray and M.~A. Saunders. Methods for modifying
		matrix factorizations. \emph{Mathematics of Computation}, 28:505--535, 1974.
		
		\bibitem{Higham2002}
		N.~J. Higham. \emph{Accuracy and Stability of Numerical Algorithms}. Second
		edition, SIAM, 2002.
		
		\bibitem{Mehrotra1992}
		S.~Mehrotra. On the implementation of a primal-dual interior point method.
		\emph{SIAM Journal on Optimization}, 2:575--601, 1992.
		
		\bibitem{Minty1962}
		G.~J. Minty. Monotone (nonlinear) operators in Hilbert space. \emph{Duke
			Mathematical Journal}, 29:341--346, 1962.
		
		\bibitem{NesterovNemirovskii1994}
		Y.~Nesterov and A.~Nemirovskii. \emph{Interior-Point Polynomial Algorithms in
			Convex Programming}. SIAM, 1994.
		
		\bibitem{NesterovTodd1997}
		Y.~E. Nesterov and M.~J. Todd. Self-scaled barriers and interior-point methods
		for convex programming. \emph{Mathematics of Operations Research}, 22:1--42,
		1997.
		
		\bibitem{Pazy1971}
		A.~Pazy. Asymptotic behavior of contractions in Hilbert space. \emph{Israel
			Journal of Mathematics}, 9:235--240, 1971.
		
		\bibitem{QiSun1993}
		L.~Qi and J.~Sun. A nonsmooth version of Newton's method. \emph{Mathematical
			Programming}, 58:353--367, 1993.
		
		\bibitem{Rockafellar1970}
		R.~T. Rockafellar. \emph{Convex Analysis}. Princeton University Press, 1970.
		
		\bibitem{Rockafellar1976a}
		R.~T. Rockafellar. Monotone operators and the proximal point algorithm.
		\emph{SIAM Journal on Control and Optimization}, 14:877--898, 1976.
		
		\bibitem{Rockafellar1976b}
		R.~T. Rockafellar. Augmented Lagrangians and applications of the proximal point
		algorithm in convex programming. \emph{Mathematics of Operations Research},
		1:97--116, 1976.
		
		\bibitem{Ruiz2001}
		D.~Ruiz. A scaling algorithm to equilibrate both rows and columns norms in
		matrices. Technical Report RAL-TR-2001-034, Rutherford Appleton Laboratory,
		2001.
		
		\bibitem{vandenDries1998}
		L.~van~den~Dries. \emph{Tame Topology and O-minimal Structures}. London
		Mathematical Society Lecture Note Series 248, Cambridge University Press, 1998.
		
		\bibitem{Wilkie1996}
		A.~J. Wilkie. Model completeness results for expansions of the ordered field of
		real numbers by restricted Pfaffian functions and the exponential function.
		\emph{Journal of the American Mathematical Society}, 9:1051--1094, 1996.
		
	\end{thebibliography}
	
\end{document}
